\documentclass[11pt,a4paper]{report}
\usepackage{fontspec,amsmath,amssymb,amsthm,microtype,xcolor}
\setmainfont{lmroman10-regular.otf}[BoldFont=lmroman10-bold.otf,ItalicFont=lmroman10-italic.otf,BoldItalicFont=lmroman10-bolditalic.otf]
\usepackage[most]{tcolorbox}
\usepackage{mathrsfs}
\providecommand{\lt}{<}
\providecommand{\gt}{>}
\setmonofont{DejaVuSansMono.ttf}[BoldFont=DejaVuSansMono-Bold.ttf,ItalicFont=DejaVuSansMono-Oblique.ttf,BoldItalicFont=DejaVuSansMono-BoldOblique.ttf,Scale=0.86]
\newfontfamily\leanfallback{STIXTwoMath-Regular.otf}[Scale=0.86]
\definecolor{leanrule}{gray}{0.62}
\definecolor{leanback}{gray}{0.962}
\definecolor{leantitle}{gray}{0.22}
\definecolor{leanmuted}{gray}{0.38}
\newtcolorbox{leanbox}[1]{enhanced,breakable,colback=leanback,colframe=leanrule,
  boxrule=0.4pt,arc=0.6mm,left=2.4mm,right=2.4mm,top=1.8mm,bottom=1.8mm,
  before skip=0.8em,after skip=1em,fonttitle=\small\bfseries,coltitle=leantitle,
  colbacktitle=leanback,titlerule=0.3pt,title={#1}}
\newlength{\leancharwidth}
\newenvironment{leancode}{\par\smallskip\begingroup\ttfamily\footnotesize\raggedright
  \settowidth{\leancharwidth}{M}\setlength{\parskip}{0pt}\setlength{\parindent}{0pt}
  \frenchspacing}{\par\endgroup\smallskip}
\newcommand{\leanline}[2]{\hangindent=\dimexpr#1\leancharwidth+4\leancharwidth\relax
  \hangafter=1\noindent\hspace*{#1\leancharwidth}#2\par}
\newcommand{\leansource}[1]{\par{\footnotesize\color{leanmuted}#1}\par}
\newcommand{\leanlabel}[1]{\par\smallskip\noindent{\small\bfseries #1}\par}
\usepackage[margin=25mm]{geometry}
\usepackage[unicode,hidelinks]{hyperref}
\hypersetup{pdfauthor={LangevinLean Project},bookmarksdepth=2}
\setlength{\emergencystretch}{3em}
\setlength{\parskip}{0.4em}
\setlength{\parindent}{0pt}
\setcounter{tocdepth}{1}
\makeatletter
\renewcommand*\l@section{\@dottedtocline{1}{1.5em}{3.8em}}
\renewcommand*\@pnumwidth{2.5em}
\renewcommand*\@tocrmarg{3.5em}
\makeatother
\setcounter{secnumdepth}{2}
\allowdisplaybreaks[2]
\raggedbottom
\title{\texorpdfstring{Particle-based Generalised Stochastic Optimisation}{Particle-based Generalised Stochastic Optimisation}\\[0.6em]\large Mathematical proofs}
\author{LangevinLean Project}
\date{Draft report LL-10\\[0.4em]\small Content revision 823682b938ce}
\begin{document}
\pagenumbering{roman}
\maketitle
\chapter*{About this document}
\addcontentsline{toc}{chapter}{About this document}
Chapter 1 states the assumptions and the main results. Chapter 2 proves them, following the argument from the hypotheses to the conclusion in order. Chapter 3 proves the supporting lemmas, each before its first use. The proofs are written for readers with an undergraduate background in analysis, linear algebra and probability; definitions are recalled where they are first used, and each inequality is justified.

\section*{Mathematics and its Lean formalization}
Every result is followed by a grey box headed \emph{Lean}. The box reproduces the exact statement of the corresponding Lean declaration, read from the source file and line range named in the box, and then explains how the Lean proof follows the mathematical argument: which hypothesis encodes which assumption, and which formal step carries out which step of the proof. The mathematics can be read without the boxes; the boxes connect each claim to the statement that Lean has checked.

\section*{Checking the proofs}
The results are verified by Lean 4 with the Mathlib library, at the versions pinned in the repository \url{https://github.com/odakyildiz/LangevinLean}. After installing the Lean toolchain manager elan and cloning the repository, run from its root:
\begin{leancode}
\leanline{0}{lake\ exe\ cache\ get}
\leanline{0}{lake\ build\ PBGSO}
\leanline{0}{lake\ env\ lean\ PBGSO/Check.lean}
\end{leancode}
The second command checks every proof in the library. The third prints the axioms on which the main theorems depend; they are only the standard axioms of propositional extensionality, choice and quotient soundness.

Prepared by the LangevinLean Project with AI assistance, these notes form part of an ongoing effort to formalize Langevin algorithms comprehensively. This is draft work and may contain mistakes.

\section*{Literature and scope}
\begin{enumerate}
\item Jiechen Jackie Zhang and O. Deniz Akyildiz (2026). \emph{Particle-based Generalised Stochastic Optimisation}. arXiv preprint arXiv:2608.02844, version 2 (9 October 2026). \url{https://arxiv.org/abs/2608.02844v2}.
\end{enumerate}
\tableofcontents
\clearpage
\pagenumbering{arabic}
\chapter{Assumptions and main results}
\section{Setting and assumptions}
\subsection{\texorpdfstring{The paper’s state spaces and notation}{The paper’s state spaces and notation}}



The paper’s objective has the integrable gradient representation





\[\nabla_\theta\ell(\theta)=\int_{\mathsf X}F(\theta,x)\,\pi_\theta(\mathrm dx),\qquad F(\theta,\cdot)\in\mathrm L^1(\pi_\theta).\]





The optimiser state is \(\vartheta\in\Theta\times\mathsf M\), with \(\theta=\Pi_\theta(\vartheta)\). The sampler state is \(Z\in\mathsf X\times\mathsf V\), with \(X=\Pi_x(Z)\). These finite-dimensional Euclidean spaces can have different dimensions. The Brownian motion takes values in \(\mathsf B\).





\[G_\nu(\vartheta)=\int_{\mathsf X\times\mathsf V} F(\Pi_\theta\vartheta,\Pi_xz)\,\nu(\mathrm dz),\qquad \nu_t=\operatorname{Law}(Z_t).\]







\[\begin{aligned}\mathrm d\vartheta_t&=\Psi_t(\vartheta_t,G_{\nu_t}(\vartheta_t))\,\mathrm dt,\\\mathrm dZ_t&=b_{\Pi_\theta\vartheta_t}(Z_t)\,\mathrm dt+\sigma_{\Pi_\theta\vartheta_t}(Z_t)\,\mathrm dB_t.\end{aligned}\]





The particle system replaces \(G_{\nu_t}\) by \(N^{-1}\sum_i F(\Pi_\theta\vartheta_t^N,\Pi_xZ_t^i)\), with \(\nu_t^N=N^{-1}\sum_i\delta_{Z_t^i}\) and independent Brownian drivers \(B^i\). The independent maps \(F,\Psi,b,\sigma\) need not come from a common potential.


\begin{leanbox}{Lean: \texttt{Model}, \texttt{meanGradient}, \texttt{meanGradient\_\allowbreak{}of\_\allowbreak{}marginal} and 1 more}
\leansource{\texttt{Model} \quad PBGSO/\allowbreak{}Framework.lean, lines 40–49}
\begin{leancode}
\leanline{0}{structure\ Model\ (S\ Θ\ Z\ X\ I\ :\ Type*)}
\leanline{4}{[NormedAddCommGroup\ S]\ [NormedSpace\ ℝ\ S]}
\leanline{4}{[NormedAddCommGroup\ Θ]\ [NormedSpace\ ℝ\ Θ]}
\leanline{4}{[NormedAddCommGroup\ Z]\ [NormedSpace\ ℝ\ Z]\ where}
\leanline{2}{parameter\ :\ S\ →\ Θ}
\leanline{2}{position\ :\ Z\ →\ X}
\leanline{2}{observation\ :\ Θ\ →\ X\ →\ Θ}
\leanline{2}{optimizer\ :\ ℝ\ →\ S\ →\ Θ\ →\ S}
\leanline{2}{drift\ :\ Θ\ →\ Z\ →\ Z}
\leanline{2}{diffusion\ :\ Θ\ →\ Z\ →\ I\ →\ Z}
\end{leancode}
\leansource{\texttt{meanGradient} \quad PBGSO/\allowbreak{}Framework.lean, lines 56–57}
\begin{leancode}
\leanline{0}{def\ meanGradient\ (ν\ :\ Measure\ Z)\ (s\ :\ S)\ :\ Θ\ :=}
\leanline{2}{∫\ z,\ M.observation\ (M.parameter\ s)\ (M.position\ z)\ ∂ν}
\end{leancode}
\leansource{\texttt{meanGradient\_\allowbreak{}of\_\allowbreak{}marginal} \quad PBGSO/\allowbreak{}Framework.lean, lines 88–95}
\begin{leancode}
\leanline{0}{theorem\ meanGradient\_of\_marginal\ [MeasurableSpace\ X]}
\leanline{4}{(s\ :\ S)\ (ν\ :\ Measure\ Z)\ (π\ :\ Measure\ X)}
\leanline{4}{(hposition\ :\ AEMeasurable\ M.position\ ν)}
\leanline{4}{(hmarginal\ :\ ν.map\ M.position\ =\ π)}
\leanline{4}{(hobs\ :\ AEStronglyMeasurable\ (M.observation\ (M.parameter\ s))\ π)\ :}
\leanline{4}{M.meanGradient\ ν\ s\ =\ ∫\ x,\ M.observation\ (M.parameter\ s)\ x\ ∂π}
\end{leancode}
\leansource{\texttt{particleGradient} \quad PBGSO/\allowbreak{}Framework.lean, lines 66–67}
\begin{leancode}
\leanline{0}{def\ particleGradient\ (N\ :\ ℕ)\ (s\ :\ S)\ (z\ :\ Fin\ N\ →\ Z)\ :\ Θ\ :=}
\leanline{2}{(N\ :\ ℝ)⁻¹\ •\ ∑\ i,\ M.observation\ (M.parameter\ s)\ (M.position\ (z\ i))}
\end{leancode}
\leanlabel{How the Lean proof follows the argument}
\textbf{Lean notation.} \texttt{S} represents \(\Theta\times\mathsf M\), \texttt{Z} represents \(\mathsf X\times\mathsf V\), \texttt{M.\allowbreak{}parameter} and \texttt{M.\allowbreak{}position} are \(\Pi_\theta,\Pi_x\), and \texttt{M.\allowbreak{}mean\allowbreak{}Gradient \ensuremath{\nu} s} is \(G_\nu(\vartheta)\). Internal source variable \texttt{s} denotes the paper’s \(\vartheta\).
\end{leanbox}
\subsection{\texorpdfstring{Distance and twisted metric}{Distance and twisted metric}}




\[\mathbf d((\vartheta,\nu),(\tilde\vartheta,\tilde\nu))=\sqrt{\|\vartheta-\tilde\vartheta\|^2+W_2^2(\nu,\tilde\nu)}.\]







\[D=\operatorname{diag}(D_\vartheta,D_z),\qquad \kappa=\frac{\lambda_{\max}(D)}{\lambda_{\min}(D)}.\]



\begin{leanbox}{Lean: \texttt{TwistedGeometry}, \texttt{conditionNumber}}
\leansource{\texttt{TwistedGeometry} \quad PBGSO/\allowbreak{}PaperTheorems.lean, lines 33–49}
\begin{leancode}
\leanline{0}{structure\ TwistedGeometry\ (S\ Z\ :\ Type*)}
\leanline{4}{[NormedAddCommGroup\ S]\ [InnerProductSpace\ ℝ\ S]}
\leanline{4}{[NormedAddCommGroup\ Z]\ [InnerProductSpace\ ℝ\ Z]\ where}
\leanline{2}{optimizerWeight\ :\ S\ →L[ℝ]\ S}
\leanline{2}{samplerWeight\ :\ Z\ →L[ℝ]\ Z}
\leanline{2}{lower\ :\ ℝ}
\leanline{2}{upper\ :\ ℝ}
\leanline{2}{lower\_pos\ :\ 0\ <\ lower}
\leanline{2}{lower\_le\_upper\ :\ lower\ ≤\ upper}
\leanline{2}{optimizer\_symmetric\ :\ ∀\ x\ y,}
\leanline{4}{{\leanfallback ⟪}x,\ optimizerWeight\ y{\leanfallback ⟫}\_ℝ\ =\ {\leanfallback ⟪}optimizerWeight\ x,\ y{\leanfallback ⟫}\_ℝ}
\leanline{2}{sampler\_symmetric\ :\ ∀\ x\ y,}
\leanline{4}{{\leanfallback ⟪}x,\ samplerWeight\ y{\leanfallback ⟫}\_ℝ\ =\ {\leanfallback ⟪}samplerWeight\ x,\ y{\leanfallback ⟫}\_ℝ}
\leanline{2}{optimizer\_lower\ :\ ∀\ s,\ lower\ *\ ‖s‖\ \textasciicircum{}\ 2\ ≤\ {\leanfallback ⟪}s,\ optimizerWeight\ s{\leanfallback ⟫}\_ℝ}
\leanline{2}{optimizer\_upper\ :\ ∀\ s,\ {\leanfallback ⟪}s,\ optimizerWeight\ s{\leanfallback ⟫}\_ℝ\ ≤\ upper\ *\ ‖s‖\ \textasciicircum{}\ 2}
\leanline{2}{sampler\_lower\ :\ ∀\ z,\ lower\ *\ ‖z‖\ \textasciicircum{}\ 2\ ≤\ {\leanfallback ⟪}z,\ samplerWeight\ z{\leanfallback ⟫}\_ℝ}
\leanline{2}{sampler\_upper\ :\ ∀\ z,\ {\leanfallback ⟪}z,\ samplerWeight\ z{\leanfallback ⟫}\_ℝ\ ≤\ upper\ *\ ‖z‖\ \textasciicircum{}\ 2}
\end{leancode}
\leansource{\texttt{conditionNumber} \quad PBGSO/\allowbreak{}PaperTheorems.lean, lines 56–56}
\begin{leancode}
\leanline{0}{def\ conditionNumber\ :\ ℝ\ :=\ D.upper\ /\ D.lower}
\end{leancode}
\leanlabel{How the Lean proof follows the argument}
The paper uses symmetric positive definite weights and sharp spectral comparisons. Lean stores certified bounds \(0<\underline\lambda\le\overline\lambda\), satisfying \(\underline\lambda\|y\|^2\le\langle y,Dy\rangle\le\overline\lambda\|y\|^2\). Taking \(\underline\lambda=\lambda_{\min}(D)\), \(\overline\lambda=\lambda_{\max}(D)\) gives the paper’s \(\kappa\). Any nonsharp certified bounds give the same statements with the corresponding larger ratio \(\overline\lambda/\underline\lambda\).
\end{leanbox}
\subsection{\texorpdfstring{A1 — Well-posedness}{A1 — Well-posedness}}



For every deterministic \(\vartheta_0\) and every \(Z_0\), independent of \(B\), with law \(\nu_0\in\mathcal P_2(\mathsf X\times\mathsf V)\), the original mean-field SDE has a pathwise unique strong solution on every finite time interval, and \(\nu_t\in\mathcal P_2\) for all finite \(t\ge0\). In addition, either (i) \(\sigma_\theta(z)\) is constant, or (ii) for one \(K\ge0\),





\[\|b_{\Pi_\theta\vartheta}(z)\|+\|\sigma_{\Pi_\theta\vartheta}(z)\|\le K(1+\|\Pi_\theta\vartheta\|+\|z\|).\]





The formal endpoints use the required A1 instances on their original and canonical Brownian spaces. These are solution existence and uniqueness premises, rather than a supplied law-comparison conclusion.

\begin{leanbox}{Lean: \texttt{Data}, \texttt{PathwiseUniqueness}, \texttt{CanonicalSolution}}
\leansource{\texttt{Data} \quad PBGSO/\allowbreak{}RawLocalStationarity.lean, lines 48–58}
\begin{leancode}
\leanline{0}{structure\ Data\ where}
\leanline{2}{optimizer\ :\ ℝ\ →\ S}
\leanline{2}{sampler\ :\ SamplerData\ Q\ D\ M\ T\ optimizer}
\leanline{2}{second\_moment\ :\ ∀\ t\ ∈\ Icc\ 0\ (T\ :\ ℝ),\ HasFiniteSecondMoment\ (Q.map\ (sampler.state\ t))}
\leanline{2}{observation\_integrable\ :\ ∀\ t\ ∈\ Icc\ 0\ (T\ :\ ℝ),\ Integrable}
\leanline{4}{(fun\ z\ =>\ M.observation\ (M.parameter\ (optimizer\ t))\ (M.position\ z))}
\leanline{6}{(Q.map\ (sampler.state\ t))}
\leanline{2}{optimizer\_field\_integrable\ :\ IntervalIntegrable}
\leanline{4}{(fun\ t\ =>\ M.optimizerField\ t\ (optimizer\ t)\ (Q.map\ (sampler.state\ t)))\ volume\ 0\ (T\ :\ ℝ)}
\leanline{2}{optimizer\_equation\ :\ ∀\ t\ ∈\ Icc\ 0\ (T\ :\ ℝ),\ optimizer\ t\ =\ optimizer\ 0\ +}
\leanline{4}{∫\ r\ in\ 0..t,\ M.optimizerField\ r\ (optimizer\ r)\ (Q.map\ (sampler.state\ r))}
\end{leancode}
\leansource{\texttt{PathwiseUniqueness} \quad PBGSO/\allowbreak{}RawLocalStationarity.lean, lines 63–66}
\begin{leancode}
\leanline{0}{def\ PathwiseUniqueness\ :\ Prop\ :=}
\leanline{2}{∀\ U\ V\ :\ Data\ Q\ D\ M\ T,\ U.optimizer\ 0\ =\ V.optimizer\ 0\ →}
\leanline{4}{U.sampler.state\ 0\ =ᵐ[Q]\ V.sampler.state\ 0\ →\ ∀\ t\ ∈\ Icc\ 0\ (T\ :\ ℝ),}
\leanline{6}{U.optimizer\ t\ =\ V.optimizer\ t\ ∧\ U.sampler.state\ t\ =ᵐ[Q]\ V.sampler.state\ t}
\end{leancode}
\leansource{\texttt{CanonicalSolution} \quad PBGSO/\allowbreak{}NaturalFiniteWorkArbitraryLaw.lean, lines 33–37}
\begin{leancode}
\leanline{0}{structure\ CanonicalSolution\ (μ\ :\ Measure\ Z)\ [IsProbabilityMeasure\ μ]\ (s\ :\ S)\ where}
\leanline{2}{data\ :\ RawLocalStationarity.Data\ (μ.prod\ (brownianProbability\ (I\ :=\ I)))}
\leanline{4}{(CanonicalBrownianDriver.driver\ (I\ :=\ I)\ μ)\ M\ T}
\leanline{2}{initial\_optimizer\ :\ data.optimizer\ 0\ =\ s}
\leanline{2}{initial\_sampler\ :\ data.sampler.state\ 0\ =ᵐ[μ.prod\ (brownianProbability\ (I\ :=\ I))]\ Prod.fst}
\end{leancode}
\end{leanbox}
\subsection{\texorpdfstring{A2 — Stationary optimiser–sampler pair}{A2 — Stationary optimiser–sampler pair}}



There are \(\vartheta_\star\), \(\theta_\star=\Pi_\theta\vartheta_\star\), and a frozen invariant law \(\nu_{\theta_\star}\in\mathcal P_2\), with





\[\nabla_\theta\ell(\theta_\star)=0,\qquad \Psi_t(\vartheta_\star,0)=0\ (t\ge0),\qquad (\Pi_x)_\#\nu_{\theta_\star}=\pi_{\theta_\star}.\]







\[G_{\nu_{\theta_\star}}(\vartheta_\star)=\int F(\theta_\star,x)\,\pi_{\theta_\star}(\mathrm dx)=\nabla_\theta\ell(\theta_\star)=0.\]





The second line uses the paper’s integrable gradient identity. The target is a critical point.

\begin{leanbox}{Lean: \texttt{FrozenSamplerData}, \texttt{stationary\_\allowbreak{}product\_\allowbreak{}law\_\allowbreak{}of\_\allowbreak{}weak\_\allowbreak{}frozen}, \texttt{stationary\_\allowbreak{}product\_\allowbreak{}law\_\allowbreak{}of\_\allowbreak{}weak\_\allowbreak{}frozen}}
\leansource{\texttt{FrozenSamplerData} \quad PBGSO/\allowbreak{}RawLocalStationarity.lean, lines 71–74}
\begin{leancode}
\leanline{0}{structure\ FrozenSamplerData\ (sStar\ :\ S)\ (νStar\ :\ Measure\ Z)\ extends}
\leanline{4}{SamplerData\ Q\ D\ M\ T\ (fun\ \_\ =>\ sStar)\ where}
\leanline{2}{initial\_square\ :\ MemLp\ (state\ 0)\ 2\ Q}
\leanline{2}{invariant\_law\ :\ ∀\ t\ ∈\ Icc\ 0\ (T\ :\ ℝ),\ Q.map\ (state\ t)\ =\ νStar}
\end{leancode}
\leansource{\texttt{stationary\_\allowbreak{}product\_\allowbreak{}law\_\allowbreak{}of\_\allowbreak{}weak\_\allowbreak{}frozen} \quad PBGSO/\allowbreak{}RawPathSolutionTransfer.lean, lines 120–140}
\begin{leancode}
\leanline{0}{theorem\ stationary\_product\_law\_of\_weak\_frozen\ [MeasurableSpace\ S]}
\leanline{4}{\{sStar\ :\ S\}\ \{νStar\ :\ Measure\ Z\}\ [IsProbabilityMeasure\ νStar]}
\leanline{4}{(F\ :\ FrozenSamplerData\ Qw\ Dw\ M\ T\ sStar\ νStar)}
\leanline{4}{(hint\ :\ Integrable\ (fun\ z\ =>\ M.observation\ (M.parameter\ sStar)\ (M.position\ z))\ νStar)}
\leanline{4}{(hg\ :\ M.meanGradient\ νStar\ sStar\ =\ 0)}
\leanline{4}{(hfix\ :\ ∀\ t\ ∈\ Icc\ 0\ (T\ :\ ℝ),\ M.optimizer\ t\ sStar\ 0\ =\ 0)}
\leanline{4}{(huniquew\ :\ PathwiseUniqueness\ Qw\ (nullAugmentedDriver\ Qw\ (continuousDriver\ Qw\ Dw))\ M\ T)}
\leanline{4}{(hexists\ :\ ∃\ V\ :\ Data\ (BrownianPathDriver.probability\ (I\ :=\ I)\ νStar)}
\leanline{8}{(BrownianPathDriver.driver\ νStar)\ M\ T,}
\leanline{6}{V.optimizer\ 0\ =\ sStar\ ∧\ V.sampler.state\ 0\ =ᵐ[}
\leanline{8}{BrownianPathDriver.probability\ (I\ :=\ I)\ νStar]\ Prod.fst)}
\leanline{4}{(U\ :\ Data\ Q\ D\ M\ T)}
\leanline{4}{(hunique\ :\ PathwiseUniqueness\ Q\ (nullAugmentedDriver\ Q\ (continuousDriver\ Q\ D))\ M\ T)}
\leanline{4}{(hs\ :\ U.optimizer\ 0\ =\ sStar)\ (hν\ :\ Q.map\ (U.sampler.state\ 0)\ =\ νStar)}
\leanline{4}{(t\ :\ ℝ)\ (ht\ :\ t\ ∈\ Icc\ 0\ (T\ :\ ℝ))\ :}
\leanline{4}{Q.map\ (fun\ ω\ =>\ (U.optimizer\ t,\ U.sampler.state\ t\ ω))\ =\ (Measure.dirac\ sStar).prod\ νStar}
\end{leancode}
\leansource{\texttt{stationary\_\allowbreak{}product\_\allowbreak{}law\_\allowbreak{}of\_\allowbreak{}weak\_\allowbreak{}frozen} \quad PBGSO/\allowbreak{}RawPathSolutionTransfer.lean, lines 120–140}
\begin{leancode}
\leanline{0}{theorem\ stationary\_product\_law\_of\_weak\_frozen\ [MeasurableSpace\ S]}
\leanline{4}{\{sStar\ :\ S\}\ \{νStar\ :\ Measure\ Z\}\ [IsProbabilityMeasure\ νStar]}
\leanline{4}{(F\ :\ FrozenSamplerData\ Qw\ Dw\ M\ T\ sStar\ νStar)}
\leanline{4}{(hint\ :\ Integrable\ (fun\ z\ =>\ M.observation\ (M.parameter\ sStar)\ (M.position\ z))\ νStar)}
\leanline{4}{(hg\ :\ M.meanGradient\ νStar\ sStar\ =\ 0)}
\leanline{4}{(hfix\ :\ ∀\ t\ ∈\ Icc\ 0\ (T\ :\ ℝ),\ M.optimizer\ t\ sStar\ 0\ =\ 0)}
\leanline{4}{(huniquew\ :\ PathwiseUniqueness\ Qw\ (nullAugmentedDriver\ Qw\ (continuousDriver\ Qw\ Dw))\ M\ T)}
\leanline{4}{(hexists\ :\ ∃\ V\ :\ Data\ (BrownianPathDriver.probability\ (I\ :=\ I)\ νStar)}
\leanline{8}{(BrownianPathDriver.driver\ νStar)\ M\ T,}
\leanline{6}{V.optimizer\ 0\ =\ sStar\ ∧\ V.sampler.state\ 0\ =ᵐ[}
\leanline{8}{BrownianPathDriver.probability\ (I\ :=\ I)\ νStar]\ Prod.fst)}
\leanline{4}{(U\ :\ Data\ Q\ D\ M\ T)}
\leanline{4}{(hunique\ :\ PathwiseUniqueness\ Q\ (nullAugmentedDriver\ Q\ (continuousDriver\ Q\ D))\ M\ T)}
\leanline{4}{(hs\ :\ U.optimizer\ 0\ =\ sStar)\ (hν\ :\ Q.map\ (U.sampler.state\ 0)\ =\ νStar)}
\leanline{4}{(t\ :\ ℝ)\ (ht\ :\ t\ ∈\ Icc\ 0\ (T\ :\ ℝ))\ :}
\leanline{4}{Q.map\ (fun\ ω\ =>\ (U.optimizer\ t,\ U.sampler.state\ t\ ω))\ =\ (Measure.dirac\ sStar).prod\ νStar}
\end{leancode}
\end{leanbox}
\subsection{\texorpdfstring{A2 — Actual SDE interpretation}{A2 — Actual SDE interpretation}}



Frozen invariance means an actual stationary weak frozen SDE on some filtered Brownian probability space. Its paths are continuous and adapted, and satisfy





\[Z_t=Z_0+\int_0^t b_{\theta_\star}(Z_r)\,\mathrm dr+\int_0^t\sigma_{\theta_\star}(Z_r)\,\mathrm dB_r,\qquad\operatorname{Law}(Z_t)=\nu_{\theta_\star}.\]





This may be a different probability space from the coupled strong solution. A generator identity alone does not supply this solution premise.

\begin{leanbox}{Lean: \texttt{SamplerData}}
\leansource{\texttt{SamplerData} \quad PBGSO/\allowbreak{}RawLocalStationarity.lean, lines 32–44}
\begin{leancode}
\leanline{0}{structure\ SamplerData\ (a\ :\ ℝ\ →\ S)\ where}
\leanline{2}{localDiffusion\ :\ RawLocalDiffusion\ (Z\ :=\ Z)\ Q\ D\ T}
\leanline{2}{state\ :\ ℝ\ →\ Ω\ →\ Z}
\leanline{2}{measurable\ :\ ∀\ t\ ∈\ Icc\ 0\ (T\ :\ ℝ),\ Measurable\ (state\ t)}
\leanline{2}{adapted\ :\ ∀\ t\ :\ ℝ≥0,\ AEStronglyMeasurable[D.filtration\ t]\ (state\ t)\ Q}
\leanline{2}{continuous\_paths\ :\ ∀ᵐ\ ω\ ∂Q,\ ContinuousOn\ (fun\ t\ =>\ state\ t\ ω)\ (Icc\ 0\ (T\ :\ ℝ))}
\leanline{2}{drift\_integrable\ :\ ∀ᵐ\ ω\ ∂Q,\ IntervalIntegrable}
\leanline{4}{(fun\ r\ =>\ M.drift\ (M.parameter\ (a\ r))\ (state\ r\ ω))\ volume\ 0\ (T\ :\ ℝ)}
\leanline{2}{diffusion\_realizes\ :\ ∀ᵐ\ ω\ ∂Q,\ ∀\ t\ ∈\ Icc\ 0\ (T\ :\ ℝ),\ ∀\ i,}
\leanline{4}{localDiffusion.sigma\ i\ t\ ω\ =\ M.diffusion\ (M.parameter\ (a\ t))\ (state\ t\ ω)\ i}
\leanline{2}{equation\ :\ ∀ᵐ\ ω\ ∂Q,\ ∀\ t\ ∈\ Icc\ 0\ (T\ :\ ℝ),\ state\ t\ ω\ =\ state\ 0\ ω\ +}
\leanline{4}{(∫\ r\ in\ 0..t,\ M.drift\ (M.parameter\ (a\ r))\ (state\ r\ ω))\ +}
\leanline{6}{localDiffusion.localNoise\ t.toNNReal\ ω}
\end{leancode}
\end{leanbox}
\subsection{\texorpdfstring{A3 — Twisted joint strong monotonicity}{A3 — Twisted joint strong monotonicity}}



There are symmetric positive definite \(D_\vartheta,D_z\) and \(\rho>0\). For every \(t\ge0\), \(\vartheta,\tilde\vartheta\), \(\nu,\tilde\nu\in\mathcal P_2\), and \(\Lambda\in\Gamma(\nu,\tilde\nu)\), put \(\delta\vartheta=\vartheta-\tilde\vartheta\), \(\delta z=z-\tilde z\), \(\theta=\Pi_\theta\vartheta\), \(\tilde\theta=\Pi_\theta\tilde\vartheta\). Then





\[\begin{aligned}&\left\langle\delta\vartheta,D_\vartheta[\Psi_t(\vartheta,G_\nu(\vartheta))-\Psi_t(\tilde\vartheta,G_{\tilde\nu}(\tilde\vartheta))]\right\rangle\\&+\int\left[\langle\delta z,D_z[b_\theta(z)-b_{\tilde\theta}(\tilde z)]\rangle+\tfrac12\|D_z^{1/2}[\sigma_\theta(z)-\sigma_{\tilde\theta}(\tilde z)]\|_{\rm HS}^2\right]\,\mathrm d\Lambda\\&\le-\rho\left(\|\delta\vartheta\|_{D_\vartheta}^2+\int\|\delta z\|_{D_z}^2\,\mathrm d\Lambda\right).\end{aligned}\]





All expectations are taken on their stated integrability domains. The ordinary finite signed-work domain needed by the constant-diffusion contraction branch is the separate orange condition R2 below.

\begin{leanbox}{Lean: \texttt{TwistedJointMonotonicity}}
\leansource{\texttt{TwistedJointMonotonicity} \quad PBGSO/\allowbreak{}PaperTheorems.lean, lines 87–101}
\begin{leancode}
\leanline{0}{def\ TwistedJointMonotonicity\ (M\ :\ Model\ S\ Θ\ Z\ X\ I)}
\leanline{4}{(D\ :\ TwistedGeometry\ S\ Z)\ (ρ\ :\ ℝ)\ :\ Prop\ :=}
\leanline{2}{∀\ t,\ 0\ ≤\ t\ →\ ∀\ s\ s'\ ν\ ν'\ (Λ\ :\ Measure\ (Z\ ×\ Z)),}
\leanline{4}{IsCoupling\ Λ\ ν\ ν'\ →}
\leanline{4}{IsProbabilityMeasure\ ν\ →\ IsProbabilityMeasure\ ν'\ →}
\leanline{4}{HasFiniteSecondMoment\ ν\ →\ HasFiniteSecondMoment\ ν'\ →}
\leanline{4}{Integrable\ (fun\ z\ =>\ M.observation\ (M.parameter\ s)\ (M.position\ z))\ ν\ →}
\leanline{4}{Integrable\ (fun\ z\ =>\ M.observation\ (M.parameter\ s')\ (M.position\ z))\ ν'\ →}
\leanline{4}{Integrable\ (samplerCouplingWork\ M\ D\ s\ s')\ Λ\ →}
\leanline{4}{Integrable\ (fun\ zz\ =>\ D.samplerSq\ (zz.1\ -\ zz.2))\ Λ\ →}
\leanline{4}{{\leanfallback ⟪}s\ -\ s',\ D.optimizerWeight}
\leanline{6}{(M.optimizerField\ t\ s\ ν\ -\ M.optimizerField\ t\ s'\ ν'){\leanfallback ⟫}\_ℝ\ +}
\leanline{6}{(∫\ zz,\ samplerCouplingWork\ M\ D\ s\ s'\ zz\ ∂Λ)\ ≤}
\leanline{6}{-ρ\ *\ (D.optimizerSq\ (s\ -\ s')\ +}
\leanline{8}{∫\ zz,\ D.samplerSq\ (zz.1\ -\ zz.2)\ ∂Λ)}
\end{leancode}
\end{leanbox}
\subsection{\texorpdfstring{A4 — Lipschitz gradient response}{A4 — Lipschitz gradient response}}


For every \(t\ge0\), \(\vartheta\) and \(g,\tilde g\in\Theta\), with \(L_\Psi\ge0\),





\[\|\Psi_t(\vartheta,g)-\Psi_t(\vartheta,\tilde g)\|\le L_\Psi\|g-\tilde g\|.\]



\begin{leanbox}{Lean: \texttt{original\_\allowbreak{}finite\_\allowbreak{}particle\_\allowbreak{}interaction\_\allowbreak{}bound}}
\leansource{\texttt{original\_\allowbreak{}finite\_\allowbreak{}particle\_\allowbreak{}interaction\_\allowbreak{}bound} \quad PBGSO/\allowbreak{}RawConstructedParticleBounds.lean, lines 109–132}
\begin{leancode}
\leanline{0}{theorem\ original\_finite\_particle\_interaction\_bound}
\leanline{4}{(hBorelF\ :\ Measurable\ (fun\ p\ :\ S\ ×\ Z\ =>\ M.observation\ (M.parameter\ p.1)\ (M.position\ p.2)))}
\leanline{4}{(hBorelb\ :\ Measurable\ (fun\ p\ :\ S\ ×\ Z\ =>\ M.drift\ (M.parameter\ p.1)\ p.2))}
\leanline{4}{(hBorelΨ\ :\ Measurable\ (fun\ p\ :\ ℝ\ ×\ (S\ ×\ Θ)\ =>\ M.optimizer\ p.1\ p.2.1\ p.2.2))}
\leanline{4}{(G\ :\ TwistedGeometry\ S\ Z)\ (ρ\ :\ ℝ)\ (hρ\ :\ 0\ <\ ρ)}
\leanline{4}{(hmono\ :\ TwistedJointMonotonicity\ M\ G\ ρ)\ (hN\ :\ 0\ <\ N)}
\leanline{4}{(LΨ\ :\ ℝ≥0)\ (V\ :\ ℝ)\ (hV\ :\ 0\ ≤\ V)}
\leanline{4}{(hΨ\ :\ ∀\ t\ ∈\ Icc\ 0\ (T\ :\ ℝ),\ ∀\ s\ g\ g',}
\leanline{6}{‖M.optimizer\ t\ s\ g-M.optimizer\ t\ s\ g'‖\ ≤\ (LΨ\ :\ ℝ)*‖g-g'‖)}
\leanline{4}{(hF\ :\ ∀\ t\ ∈\ Icc\ 0\ (T\ :\ ℝ),\ MemLp}
\leanline{6}{(fun\ z\ =>\ M.observation\ (M.parameter\ (U.optimizer\ t))\ (M.position\ z))\ 2}
\leanline{6}{(Q.map\ (U.sampler.state\ t)))}
\leanline{4}{(hVar\ :\ ∀\ t\ ∈\ Icc\ 0\ (T\ :\ ℝ),}
\leanline{6}{(∫\ z,\ ‖M.observation\ (M.parameter\ (U.optimizer\ t))\ (M.position\ z)-}
\leanline{8}{M.meanGradient\ (Q.map\ (U.sampler.state\ t))\ (U.optimizer\ t)‖\textasciicircum{}2}
\leanline{10}{∂Q.map\ (U.sampler.state\ t))\ ≤\ V)}
\leanline{4}{(hinit\ :\ ∀ᵐ\ ω\ ∂IndependentCopies.seeds\ (J\ :=\ Fin\ N)\ Q,}
\leanline{6}{oneEquiv\ N\ (Y.state\ 0\ ω)\ =\ (U.optimizer\ 0,RawIndependentCopies.cloud\ Q\ D\ N\ U\ 0\ ω))\ :}
\leanline{4}{Real.sqrt\ (∫\ ω,\ particleEmpiricalDistanceSq\ N\ (originalPair\ Q\ D\ M\ N\ U\ Y\ T\ ω)}
\leanline{6}{∂IndependentCopies.seeds\ (J\ :=\ Fin\ N)\ Q)\ ≤}
\leanline{6}{((LΨ\ :\ ℝ)/ρ)*Real.sqrt\ (G.conditionNumber*V/(N\ :\ ℝ))*(1-Real.exp\ (-ρ*T))}
\end{leancode}
\end{leanbox}
\subsection{\texorpdfstring{A5 — Uniform gradient fluctuation}{A5 — Uniform gradient fluctuation}}


Along the independent-copy mean-field path, the observation is square-integrable and, for a finite nonnegative constant,





\[\sup_{t\ge0}\mathbb E\left\|F(\Pi_\theta\bar\vartheta_t,\Pi_x\bar Z_t)-G_{\bar\nu_t}(\bar\vartheta_t)\right\|^2\le V_F<\infty.\]



\begin{leanbox}{Lean: \texttt{original\_\allowbreak{}finite\_\allowbreak{}particle\_\allowbreak{}interaction\_\allowbreak{}bound}, \texttt{original\_\allowbreak{}finite\_\allowbreak{}particle\_\allowbreak{}interaction\_\allowbreak{}bound}}
\leansource{\texttt{original\_\allowbreak{}finite\_\allowbreak{}particle\_\allowbreak{}interaction\_\allowbreak{}bound} \quad PBGSO/\allowbreak{}RawConstructedParticleBounds.lean, lines 109–132}
\begin{leancode}
\leanline{0}{theorem\ original\_finite\_particle\_interaction\_bound}
\leanline{4}{(hBorelF\ :\ Measurable\ (fun\ p\ :\ S\ ×\ Z\ =>\ M.observation\ (M.parameter\ p.1)\ (M.position\ p.2)))}
\leanline{4}{(hBorelb\ :\ Measurable\ (fun\ p\ :\ S\ ×\ Z\ =>\ M.drift\ (M.parameter\ p.1)\ p.2))}
\leanline{4}{(hBorelΨ\ :\ Measurable\ (fun\ p\ :\ ℝ\ ×\ (S\ ×\ Θ)\ =>\ M.optimizer\ p.1\ p.2.1\ p.2.2))}
\leanline{4}{(G\ :\ TwistedGeometry\ S\ Z)\ (ρ\ :\ ℝ)\ (hρ\ :\ 0\ <\ ρ)}
\leanline{4}{(hmono\ :\ TwistedJointMonotonicity\ M\ G\ ρ)\ (hN\ :\ 0\ <\ N)}
\leanline{4}{(LΨ\ :\ ℝ≥0)\ (V\ :\ ℝ)\ (hV\ :\ 0\ ≤\ V)}
\leanline{4}{(hΨ\ :\ ∀\ t\ ∈\ Icc\ 0\ (T\ :\ ℝ),\ ∀\ s\ g\ g',}
\leanline{6}{‖M.optimizer\ t\ s\ g-M.optimizer\ t\ s\ g'‖\ ≤\ (LΨ\ :\ ℝ)*‖g-g'‖)}
\leanline{4}{(hF\ :\ ∀\ t\ ∈\ Icc\ 0\ (T\ :\ ℝ),\ MemLp}
\leanline{6}{(fun\ z\ =>\ M.observation\ (M.parameter\ (U.optimizer\ t))\ (M.position\ z))\ 2}
\leanline{6}{(Q.map\ (U.sampler.state\ t)))}
\leanline{4}{(hVar\ :\ ∀\ t\ ∈\ Icc\ 0\ (T\ :\ ℝ),}
\leanline{6}{(∫\ z,\ ‖M.observation\ (M.parameter\ (U.optimizer\ t))\ (M.position\ z)-}
\leanline{8}{M.meanGradient\ (Q.map\ (U.sampler.state\ t))\ (U.optimizer\ t)‖\textasciicircum{}2}
\leanline{10}{∂Q.map\ (U.sampler.state\ t))\ ≤\ V)}
\leanline{4}{(hinit\ :\ ∀ᵐ\ ω\ ∂IndependentCopies.seeds\ (J\ :=\ Fin\ N)\ Q,}
\leanline{6}{oneEquiv\ N\ (Y.state\ 0\ ω)\ =\ (U.optimizer\ 0,RawIndependentCopies.cloud\ Q\ D\ N\ U\ 0\ ω))\ :}
\leanline{4}{Real.sqrt\ (∫\ ω,\ particleEmpiricalDistanceSq\ N\ (originalPair\ Q\ D\ M\ N\ U\ Y\ T\ ω)}
\leanline{6}{∂IndependentCopies.seeds\ (J\ :=\ Fin\ N)\ Q)\ ≤}
\leanline{6}{((LΨ\ :\ ℝ)/ρ)*Real.sqrt\ (G.conditionNumber*V/(N\ :\ ℝ))*(1-Real.exp\ (-ρ*T))}
\end{leancode}
\leansource{\texttt{original\_\allowbreak{}finite\_\allowbreak{}particle\_\allowbreak{}interaction\_\allowbreak{}bound} \quad PBGSO/\allowbreak{}RawConstructedParticleBounds.lean, lines 109–132}
\begin{leancode}
\leanline{0}{theorem\ original\_finite\_particle\_interaction\_bound}
\leanline{4}{(hBorelF\ :\ Measurable\ (fun\ p\ :\ S\ ×\ Z\ =>\ M.observation\ (M.parameter\ p.1)\ (M.position\ p.2)))}
\leanline{4}{(hBorelb\ :\ Measurable\ (fun\ p\ :\ S\ ×\ Z\ =>\ M.drift\ (M.parameter\ p.1)\ p.2))}
\leanline{4}{(hBorelΨ\ :\ Measurable\ (fun\ p\ :\ ℝ\ ×\ (S\ ×\ Θ)\ =>\ M.optimizer\ p.1\ p.2.1\ p.2.2))}
\leanline{4}{(G\ :\ TwistedGeometry\ S\ Z)\ (ρ\ :\ ℝ)\ (hρ\ :\ 0\ <\ ρ)}
\leanline{4}{(hmono\ :\ TwistedJointMonotonicity\ M\ G\ ρ)\ (hN\ :\ 0\ <\ N)}
\leanline{4}{(LΨ\ :\ ℝ≥0)\ (V\ :\ ℝ)\ (hV\ :\ 0\ ≤\ V)}
\leanline{4}{(hΨ\ :\ ∀\ t\ ∈\ Icc\ 0\ (T\ :\ ℝ),\ ∀\ s\ g\ g',}
\leanline{6}{‖M.optimizer\ t\ s\ g-M.optimizer\ t\ s\ g'‖\ ≤\ (LΨ\ :\ ℝ)*‖g-g'‖)}
\leanline{4}{(hF\ :\ ∀\ t\ ∈\ Icc\ 0\ (T\ :\ ℝ),\ MemLp}
\leanline{6}{(fun\ z\ =>\ M.observation\ (M.parameter\ (U.optimizer\ t))\ (M.position\ z))\ 2}
\leanline{6}{(Q.map\ (U.sampler.state\ t)))}
\leanline{4}{(hVar\ :\ ∀\ t\ ∈\ Icc\ 0\ (T\ :\ ℝ),}
\leanline{6}{(∫\ z,\ ‖M.observation\ (M.parameter\ (U.optimizer\ t))\ (M.position\ z)-}
\leanline{8}{M.meanGradient\ (Q.map\ (U.sampler.state\ t))\ (U.optimizer\ t)‖\textasciicircum{}2}
\leanline{10}{∂Q.map\ (U.sampler.state\ t))\ ≤\ V)}
\leanline{4}{(hinit\ :\ ∀ᵐ\ ω\ ∂IndependentCopies.seeds\ (J\ :=\ Fin\ N)\ Q,}
\leanline{6}{oneEquiv\ N\ (Y.state\ 0\ ω)\ =\ (U.optimizer\ 0,RawIndependentCopies.cloud\ Q\ D\ N\ U\ 0\ ω))\ :}
\leanline{4}{Real.sqrt\ (∫\ ω,\ particleEmpiricalDistanceSq\ N\ (originalPair\ Q\ D\ M\ N\ U\ Y\ T\ ω)}
\leanline{6}{∂IndependentCopies.seeds\ (J\ :=\ Fin\ N)\ Q)\ ≤}
\leanline{6}{((LΨ\ :\ ℝ)/ρ)*Real.sqrt\ (G.conditionNumber*V/(N\ :\ ℝ))*(1-Real.exp\ (-ρ*T))}
\end{leancode}
\end{leanbox}
\subsection{\texorpdfstring{R1 — Extra assumption: joint Borel measurability}{R1 — Extra assumption: joint Borel measurability}}



The following maps are jointly Borel:





\[(\vartheta,z)\mapsto F(\Pi_\theta\vartheta,\Pi_xz),\qquad (\vartheta,z)\mapsto b_{\Pi_\theta\vartheta}(z),\qquad(t,\vartheta,g)\mapsto\Psi_t(\vartheta,g).\]



\begin{leanbox}{Lean: \texttt{original\_\allowbreak{}finite\_\allowbreak{}particle\_\allowbreak{}interaction\_\allowbreak{}bound}, \texttt{original\_\allowbreak{}finite\_\allowbreak{}particle\_\allowbreak{}interaction\_\allowbreak{}bound}, \texttt{original\_\allowbreak{}finite\_\allowbreak{}particle\_\allowbreak{}interaction\_\allowbreak{}bound}}
\leansource{\texttt{original\_\allowbreak{}finite\_\allowbreak{}particle\_\allowbreak{}interaction\_\allowbreak{}bound} \quad PBGSO/\allowbreak{}RawConstructedParticleBounds.lean, lines 109–132}
\begin{leancode}
\leanline{0}{theorem\ original\_finite\_particle\_interaction\_bound}
\leanline{4}{(hBorelF\ :\ Measurable\ (fun\ p\ :\ S\ ×\ Z\ =>\ M.observation\ (M.parameter\ p.1)\ (M.position\ p.2)))}
\leanline{4}{(hBorelb\ :\ Measurable\ (fun\ p\ :\ S\ ×\ Z\ =>\ M.drift\ (M.parameter\ p.1)\ p.2))}
\leanline{4}{(hBorelΨ\ :\ Measurable\ (fun\ p\ :\ ℝ\ ×\ (S\ ×\ Θ)\ =>\ M.optimizer\ p.1\ p.2.1\ p.2.2))}
\leanline{4}{(G\ :\ TwistedGeometry\ S\ Z)\ (ρ\ :\ ℝ)\ (hρ\ :\ 0\ <\ ρ)}
\leanline{4}{(hmono\ :\ TwistedJointMonotonicity\ M\ G\ ρ)\ (hN\ :\ 0\ <\ N)}
\leanline{4}{(LΨ\ :\ ℝ≥0)\ (V\ :\ ℝ)\ (hV\ :\ 0\ ≤\ V)}
\leanline{4}{(hΨ\ :\ ∀\ t\ ∈\ Icc\ 0\ (T\ :\ ℝ),\ ∀\ s\ g\ g',}
\leanline{6}{‖M.optimizer\ t\ s\ g-M.optimizer\ t\ s\ g'‖\ ≤\ (LΨ\ :\ ℝ)*‖g-g'‖)}
\leanline{4}{(hF\ :\ ∀\ t\ ∈\ Icc\ 0\ (T\ :\ ℝ),\ MemLp}
\leanline{6}{(fun\ z\ =>\ M.observation\ (M.parameter\ (U.optimizer\ t))\ (M.position\ z))\ 2}
\leanline{6}{(Q.map\ (U.sampler.state\ t)))}
\leanline{4}{(hVar\ :\ ∀\ t\ ∈\ Icc\ 0\ (T\ :\ ℝ),}
\leanline{6}{(∫\ z,\ ‖M.observation\ (M.parameter\ (U.optimizer\ t))\ (M.position\ z)-}
\leanline{8}{M.meanGradient\ (Q.map\ (U.sampler.state\ t))\ (U.optimizer\ t)‖\textasciicircum{}2}
\leanline{10}{∂Q.map\ (U.sampler.state\ t))\ ≤\ V)}
\leanline{4}{(hinit\ :\ ∀ᵐ\ ω\ ∂IndependentCopies.seeds\ (J\ :=\ Fin\ N)\ Q,}
\leanline{6}{oneEquiv\ N\ (Y.state\ 0\ ω)\ =\ (U.optimizer\ 0,RawIndependentCopies.cloud\ Q\ D\ N\ U\ 0\ ω))\ :}
\leanline{4}{Real.sqrt\ (∫\ ω,\ particleEmpiricalDistanceSq\ N\ (originalPair\ Q\ D\ M\ N\ U\ Y\ T\ ω)}
\leanline{6}{∂IndependentCopies.seeds\ (J\ :=\ Fin\ N)\ Q)\ ≤}
\leanline{6}{((LΨ\ :\ ℝ)/ρ)*Real.sqrt\ (G.conditionNumber*V/(N\ :\ ℝ))*(1-Real.exp\ (-ρ*T))}
\end{leancode}
\leansource{\texttt{original\_\allowbreak{}finite\_\allowbreak{}particle\_\allowbreak{}interaction\_\allowbreak{}bound} \quad PBGSO/\allowbreak{}RawConstructedParticleBounds.lean, lines 109–132}
\begin{leancode}
\leanline{0}{theorem\ original\_finite\_particle\_interaction\_bound}
\leanline{4}{(hBorelF\ :\ Measurable\ (fun\ p\ :\ S\ ×\ Z\ =>\ M.observation\ (M.parameter\ p.1)\ (M.position\ p.2)))}
\leanline{4}{(hBorelb\ :\ Measurable\ (fun\ p\ :\ S\ ×\ Z\ =>\ M.drift\ (M.parameter\ p.1)\ p.2))}
\leanline{4}{(hBorelΨ\ :\ Measurable\ (fun\ p\ :\ ℝ\ ×\ (S\ ×\ Θ)\ =>\ M.optimizer\ p.1\ p.2.1\ p.2.2))}
\leanline{4}{(G\ :\ TwistedGeometry\ S\ Z)\ (ρ\ :\ ℝ)\ (hρ\ :\ 0\ <\ ρ)}
\leanline{4}{(hmono\ :\ TwistedJointMonotonicity\ M\ G\ ρ)\ (hN\ :\ 0\ <\ N)}
\leanline{4}{(LΨ\ :\ ℝ≥0)\ (V\ :\ ℝ)\ (hV\ :\ 0\ ≤\ V)}
\leanline{4}{(hΨ\ :\ ∀\ t\ ∈\ Icc\ 0\ (T\ :\ ℝ),\ ∀\ s\ g\ g',}
\leanline{6}{‖M.optimizer\ t\ s\ g-M.optimizer\ t\ s\ g'‖\ ≤\ (LΨ\ :\ ℝ)*‖g-g'‖)}
\leanline{4}{(hF\ :\ ∀\ t\ ∈\ Icc\ 0\ (T\ :\ ℝ),\ MemLp}
\leanline{6}{(fun\ z\ =>\ M.observation\ (M.parameter\ (U.optimizer\ t))\ (M.position\ z))\ 2}
\leanline{6}{(Q.map\ (U.sampler.state\ t)))}
\leanline{4}{(hVar\ :\ ∀\ t\ ∈\ Icc\ 0\ (T\ :\ ℝ),}
\leanline{6}{(∫\ z,\ ‖M.observation\ (M.parameter\ (U.optimizer\ t))\ (M.position\ z)-}
\leanline{8}{M.meanGradient\ (Q.map\ (U.sampler.state\ t))\ (U.optimizer\ t)‖\textasciicircum{}2}
\leanline{10}{∂Q.map\ (U.sampler.state\ t))\ ≤\ V)}
\leanline{4}{(hinit\ :\ ∀ᵐ\ ω\ ∂IndependentCopies.seeds\ (J\ :=\ Fin\ N)\ Q,}
\leanline{6}{oneEquiv\ N\ (Y.state\ 0\ ω)\ =\ (U.optimizer\ 0,RawIndependentCopies.cloud\ Q\ D\ N\ U\ 0\ ω))\ :}
\leanline{4}{Real.sqrt\ (∫\ ω,\ particleEmpiricalDistanceSq\ N\ (originalPair\ Q\ D\ M\ N\ U\ Y\ T\ ω)}
\leanline{6}{∂IndependentCopies.seeds\ (J\ :=\ Fin\ N)\ Q)\ ≤}
\leanline{6}{((LΨ\ :\ ℝ)/ρ)*Real.sqrt\ (G.conditionNumber*V/(N\ :\ ℝ))*(1-Real.exp\ (-ρ*T))}
\end{leancode}
\leansource{\texttt{original\_\allowbreak{}finite\_\allowbreak{}particle\_\allowbreak{}interaction\_\allowbreak{}bound} \quad PBGSO/\allowbreak{}RawConstructedParticleBounds.lean, lines 109–132}
\begin{leancode}
\leanline{0}{theorem\ original\_finite\_particle\_interaction\_bound}
\leanline{4}{(hBorelF\ :\ Measurable\ (fun\ p\ :\ S\ ×\ Z\ =>\ M.observation\ (M.parameter\ p.1)\ (M.position\ p.2)))}
\leanline{4}{(hBorelb\ :\ Measurable\ (fun\ p\ :\ S\ ×\ Z\ =>\ M.drift\ (M.parameter\ p.1)\ p.2))}
\leanline{4}{(hBorelΨ\ :\ Measurable\ (fun\ p\ :\ ℝ\ ×\ (S\ ×\ Θ)\ =>\ M.optimizer\ p.1\ p.2.1\ p.2.2))}
\leanline{4}{(G\ :\ TwistedGeometry\ S\ Z)\ (ρ\ :\ ℝ)\ (hρ\ :\ 0\ <\ ρ)}
\leanline{4}{(hmono\ :\ TwistedJointMonotonicity\ M\ G\ ρ)\ (hN\ :\ 0\ <\ N)}
\leanline{4}{(LΨ\ :\ ℝ≥0)\ (V\ :\ ℝ)\ (hV\ :\ 0\ ≤\ V)}
\leanline{4}{(hΨ\ :\ ∀\ t\ ∈\ Icc\ 0\ (T\ :\ ℝ),\ ∀\ s\ g\ g',}
\leanline{6}{‖M.optimizer\ t\ s\ g-M.optimizer\ t\ s\ g'‖\ ≤\ (LΨ\ :\ ℝ)*‖g-g'‖)}
\leanline{4}{(hF\ :\ ∀\ t\ ∈\ Icc\ 0\ (T\ :\ ℝ),\ MemLp}
\leanline{6}{(fun\ z\ =>\ M.observation\ (M.parameter\ (U.optimizer\ t))\ (M.position\ z))\ 2}
\leanline{6}{(Q.map\ (U.sampler.state\ t)))}
\leanline{4}{(hVar\ :\ ∀\ t\ ∈\ Icc\ 0\ (T\ :\ ℝ),}
\leanline{6}{(∫\ z,\ ‖M.observation\ (M.parameter\ (U.optimizer\ t))\ (M.position\ z)-}
\leanline{8}{M.meanGradient\ (Q.map\ (U.sampler.state\ t))\ (U.optimizer\ t)‖\textasciicircum{}2}
\leanline{10}{∂Q.map\ (U.sampler.state\ t))\ ≤\ V)}
\leanline{4}{(hinit\ :\ ∀ᵐ\ ω\ ∂IndependentCopies.seeds\ (J\ :=\ Fin\ N)\ Q,}
\leanline{6}{oneEquiv\ N\ (Y.state\ 0\ ω)\ =\ (U.optimizer\ 0,RawIndependentCopies.cloud\ Q\ D\ N\ U\ 0\ ω))\ :}
\leanline{4}{Real.sqrt\ (∫\ ω,\ particleEmpiricalDistanceSq\ N\ (originalPair\ Q\ D\ M\ N\ U\ Y\ T\ ω)}
\leanline{6}{∂IndependentCopies.seeds\ (J\ :=\ Fin\ N)\ Q)\ ≤}
\leanline{6}{((LΨ\ :\ ℝ)/ρ)*Real.sqrt\ (G.conditionNumber*V/(N\ :\ ℝ))*(1-Real.exp\ (-ρ*T))}
\end{leancode}
\leanlabel{How the Lean proof follows the argument}
Use R1 for Theorems 3 and 4 and for the constant-diffusion branch of Theorem 2. The constant-diffusion branches of Theorems 2 and 4 also assume each frozen drift is Borel, as shown in their complete Lean statements.
\end{leanbox}
\subsection{\texorpdfstring{R2 — Extra assumption: finite signed work for constant-diffusion contraction}{R2 — Extra assumption: finite signed work for constant-diffusion contraction}}



For every optimiser pair and every joint probability law with finite quadratic moment, require the ordinary finite signed-integral domain:





\[\int\left|\langle z-\tilde z,D_z[b_{\Pi_\theta\vartheta}(z)-b_{\Pi_\theta\tilde\vartheta}(\tilde z)]\rangle\right|\,\Lambda(\mathrm dz,\mathrm d\tilde z)<\infty.\]



\begin{leanbox}{Lean: \texttt{FiniteCouplingWork}}
\leansource{\texttt{FiniteCouplingWork} \quad PBGSO/\allowbreak{}DiracOptimizerGrowth.lean, lines 128–130}
\begin{leancode}
\leanline{0}{def\ FiniteCouplingWork\ :\ Prop\ :=\ ∀\ s\ s'\ (Λ\ :\ Measure\ (Z\ ×\ Z)),}
\leanline{2}{IsProbabilityMeasure\ Λ\ →\ Integrable\ (fun\ p\ :\ Z\ ×\ Z\ =>\ 1+‖p.1‖\ \textasciicircum{}\ 2+‖p.2‖\ \textasciicircum{}\ 2)\ Λ\ →}
\leanline{4}{Integrable\ (samplerCouplingWork\ M\ D\ s\ s')\ Λ}
\end{leancode}
\leanlabel{How the Lean proof follows the argument}
The Lean integrand also contains the weighted Hilbert–Schmidt diffusion term of A3; it vanishes under constant diffusion, the only branch that uses R2. This \texttt{Finite\allowbreak{}Coupling\allowbreak{}Work} condition is used for the constant-diffusion branch of Theorem 2 and its stationary uniqueness corollary. It strengthens an extended-negative-work reading of A3; it is not imposed separately on Theorems 3 or 4.
\end{leanbox}
\subsection{\texorpdfstring{R3 — Extra assumptions: measurability in the linear-growth branch}{R3 — Extra assumptions: measurability in the linear-growth branch}}



Each frozen drift and diffusion column is Borel, and \(\Pi_\theta\) is continuous. Along the canonical strong solutions, the functions





\[(r,\omega)\mapsto b_{\Pi_\theta\vartheta_r}(Z_r(\omega)),\qquad(r,\omega)\mapsto\sigma_{\Pi_\theta\vartheta_r}(Z_r(\omega))e_i\]





are jointly strongly measurable, and the drift is progressively strongly measurable with respect to the driver filtration. The linear-growth inequality is exactly A1(ii). These are extra measurability assumptions; the proof derives the moment and energy estimates.

\begin{leanbox}{Lean: \texttt{Regularity}}
\leansource{\texttt{Regularity} \quad PBGSO/\allowbreak{}RawSamplerGrowth.lean, lines 171–185}
\begin{leancode}
\leanline{0}{structure\ Regularity\ (U\ :\ RawLocalStationarity.Data\ Q\ D\ M\ T)\ where}
\leanline{2}{drift\_strong\ :\ StronglyMeasurable\ (Function.uncurry\ (originalRawDrift\ M\ U.optimizer\ U.sampler.state))}
\leanline{2}{drift\_progressive\ :\ ∀\ t\ :\ ℝ≥0,}
\leanline{4}{StronglyMeasurable[(inferInstance\ :\ MeasurableSpace\ ℝ).prod\ (D.filtration\ t)]}
\leanline{6}{(fun\ p\ :\ ℝ\ ×\ Ω\ =>\ (Ioc\ 0\ (t\ :\ ℝ)).indicator}
\leanline{8}{(fun\ r\ =>\ originalRawDrift\ M\ U.optimizer\ U.sampler.state\ r\ p.2)\ p.1)}
\leanline{2}{diffusion\_strong\ :\ ∀\ i,\ StronglyMeasurable}
\leanline{4}{(Function.uncurry\ (originalRawDiffusion\ M\ U.optimizer\ U.sampler.state\ i))}
\leanline{2}{drift\_borel\ :\ ∀\ θ,\ Measurable\ (M.drift\ θ)}
\leanline{2}{diffusion\_borel\ :\ ∀\ θ\ i,\ Measurable\ (fun\ z\ =>\ M.diffusion\ θ\ z\ i)}
\leanline{2}{projection\_continuous\ :\ Continuous\ M.parameter}
\leanline{2}{growthConstant\ :\ ℝ}
\leanline{2}{growth\_nonneg\ :\ 0\ ≤\ growthConstant}
\leanline{2}{linear\_growth\ :\ ∀\ s\ z,\ ‖M.drift\ (M.parameter\ s)\ z‖+}
\leanline{4}{‖WithLp.toLp\ 2\ (M.diffusion\ (M.parameter\ s)\ z)‖\ ≤\ growthConstant*(1+‖M.parameter\ s‖+‖z‖)}
\end{leancode}
\end{leanbox}
\section{Main theorems of the paper}
\subsection{\texorpdfstring{Theorem 1 — Invariant measure}{Theorem 1 — Invariant measure}}


\textbf{Assumptions:} A1, A2, with the actual weak-SDE meaning stated above. Initialise the optimiser at the target and the sampler independently of Brownian motion with its frozen invariant law.



\textbf{Extra assumptions: none.}





For every \(t\ge0\),





\[\operatorname{Law}(\vartheta_t,Z_t)=\mu_\star:=\delta_{\vartheta_\star}\otimes\nu_{\theta_\star}.\]







For time-dependent optimiser response this asserts a stationary solution; for autonomous response it is also ordinary invariant-measure preservation.

\begin{leanbox}{Lean: \texttt{stationary\_\allowbreak{}product\_\allowbreak{}law\_\allowbreak{}of\_\allowbreak{}weak\_\allowbreak{}frozen}}
\leansource{\texttt{stationary\_\allowbreak{}product\_\allowbreak{}law\_\allowbreak{}of\_\allowbreak{}weak\_\allowbreak{}frozen} \quad PBGSO/\allowbreak{}RawPathSolutionTransfer.lean, lines 120–140}
\begin{leancode}
\leanline{0}{theorem\ stationary\_product\_law\_of\_weak\_frozen\ [MeasurableSpace\ S]}
\leanline{4}{\{sStar\ :\ S\}\ \{νStar\ :\ Measure\ Z\}\ [IsProbabilityMeasure\ νStar]}
\leanline{4}{(F\ :\ FrozenSamplerData\ Qw\ Dw\ M\ T\ sStar\ νStar)}
\leanline{4}{(hint\ :\ Integrable\ (fun\ z\ =>\ M.observation\ (M.parameter\ sStar)\ (M.position\ z))\ νStar)}
\leanline{4}{(hg\ :\ M.meanGradient\ νStar\ sStar\ =\ 0)}
\leanline{4}{(hfix\ :\ ∀\ t\ ∈\ Icc\ 0\ (T\ :\ ℝ),\ M.optimizer\ t\ sStar\ 0\ =\ 0)}
\leanline{4}{(huniquew\ :\ PathwiseUniqueness\ Qw\ (nullAugmentedDriver\ Qw\ (continuousDriver\ Qw\ Dw))\ M\ T)}
\leanline{4}{(hexists\ :\ ∃\ V\ :\ Data\ (BrownianPathDriver.probability\ (I\ :=\ I)\ νStar)}
\leanline{8}{(BrownianPathDriver.driver\ νStar)\ M\ T,}
\leanline{6}{V.optimizer\ 0\ =\ sStar\ ∧\ V.sampler.state\ 0\ =ᵐ[}
\leanline{8}{BrownianPathDriver.probability\ (I\ :=\ I)\ νStar]\ Prod.fst)}
\leanline{4}{(U\ :\ Data\ Q\ D\ M\ T)}
\leanline{4}{(hunique\ :\ PathwiseUniqueness\ Q\ (nullAugmentedDriver\ Q\ (continuousDriver\ Q\ D))\ M\ T)}
\leanline{4}{(hs\ :\ U.optimizer\ 0\ =\ sStar)\ (hν\ :\ Q.map\ (U.sampler.state\ 0)\ =\ νStar)}
\leanline{4}{(t\ :\ ℝ)\ (ht\ :\ t\ ∈\ Icc\ 0\ (T\ :\ ℝ))\ :}
\leanline{4}{Q.map\ (fun\ ω\ =>\ (U.optimizer\ t,\ U.sampler.state\ t\ ω))\ =\ (Measure.dirac\ sStar).prod\ νStar}
\end{leancode}
\end{leanbox}
\subsection{\texorpdfstring{Theorem 2 — Contraction under joint twisted monotonicity}{Theorem 2 — Contraction under joint twisted monotonicity}}


\textbf{Assumptions:} A1, A3, deterministic optimiser initial states and initial sampler laws in the stated second-moment domain.

Extra assumptions — constant diffusion: R1 + R2 (FiniteCouplingWork); linear growth: R3.



For every \(t\ge0\),





\[\mathbf d((\vartheta_t,\nu_t),(\tilde\vartheta_t,\tilde\nu_t))\le\sqrt\kappa\,e^{-\rho t}\mathbf d((\vartheta_0,\nu_0),(\tilde\vartheta_0,\tilde\nu_0)).\]





With A2, the target is stationary and





\[W_2(\mu_t,\mu_\star)\le\sqrt\kappa\,e^{-\rho t}W_2(\mu_0,\mu_\star),\qquad\mu_t=\delta_{\vartheta_t}\otimes\nu_t.\]







The second bound follows from the first, with the stationary target of Theorem 1 as comparison solution, and the product identity \(W_2(\delta_\vartheta\otimes\nu,\delta_{\tilde\vartheta}\otimes\tilde\nu)=\mathbf d((\vartheta,\nu),(\tilde\vartheta,\tilde\nu))\) for \(\nu,\tilde\nu\in\mathcal P_2\). The stationary law is unique among finite-second-moment laws of the form prescribed by the paper, with a deterministic optimiser and the stated solution hypotheses.

\begin{leanbox}{Lean: \texttt{law\_\allowbreak{}distance\_\allowbreak{}decay}, \texttt{law\_\allowbreak{}distance\_\allowbreak{}decay}, \texttt{wasserstein\_\allowbreak{}product\_\allowbreak{}eq} and 2 more}
\leansource{\texttt{law\_\allowbreak{}distance\_\allowbreak{}decay} \quad PBGSO/\allowbreak{}NaturalFiniteWorkArbitraryLaw.lean, lines 43–88}
\begin{leancode}
\leanline{0}{theorem\ law\_distance\_decay}
\leanline{4}{(U\ :\ RawLocalStationarity.Data\ Q\ D\ M\ T)\ (V\ :\ RawLocalStationarity.Data\ R\ E\ M\ T)}
\leanline{4}{(μ\ ν\ :\ Measure\ Z)\ [IsProbabilityMeasure\ μ]\ [IsProbabilityMeasure\ ν]}
\leanline{4}{(hμ\ :\ MemLp\ (fun\ z\ :\ Z\ =>\ z)\ 2\ μ)\ (hν\ :\ MemLp\ (fun\ z\ :\ Z\ =>\ z)\ 2\ ν)}
\leanline{4}{(hU\ :\ Q.map\ (U.sampler.state\ 0)\ =\ μ)\ (hV\ :\ R.map\ (V.sampler.state\ 0)\ =\ ν)}
\leanline{4}{(A\ :\ CanonicalSolution\ (M\ :=\ M)\ (T\ :=\ T)\ μ\ (U.optimizer\ 0))}
\leanline{4}{(B\ :\ CanonicalSolution\ (M\ :=\ M)\ (T\ :=\ T)\ ν\ (V.optimizer\ 0))}
\leanline{4}{(huniqueU\ :\ PathwiseUniqueness\ Q\ (nullAugmentedDriver\ Q\ (continuousDriver\ Q\ D))\ M\ T)}
\leanline{4}{(huniqueV\ :\ PathwiseUniqueness\ R\ (nullAugmentedDriver\ R\ (continuousDriver\ R\ E))\ M\ T)}
\leanline{4}{(huniqueA\ :\ PathwiseUniqueness\ (μ.prod\ (brownianProbability\ (I\ :=\ I)))}
\leanline{6}{(nullAugmentedDriver\ \_\ (continuousDriver\ \_\ (CanonicalBrownianDriver.driver\ (I\ :=\ I)\ μ)))\ M\ T)}
\leanline{4}{(huniqueB\ :\ PathwiseUniqueness\ (ν.prod\ (brownianProbability\ (I\ :=\ I)))}
\leanline{6}{(nullAugmentedDriver\ \_\ (continuousDriver\ \_\ (CanonicalBrownianDriver.driver\ (I\ :=\ I)\ ν)))\ M\ T)}
\leanline{4}{(hpathμ\ :\ ∃\ W\ :\ RawLocalStationarity.Data\ (BrownianPathDriver.probability\ (I\ :=\ I)\ μ)}
\leanline{8}{(BrownianPathDriver.driver\ μ)\ M\ T,}
\leanline{6}{W.optimizer\ 0\ =\ U.optimizer\ 0\ ∧}
\leanline{6}{W.sampler.state\ 0\ =ᵐ[BrownianPathDriver.probability\ (I\ :=\ I)\ μ]\ Prod.fst)}
\leanline{4}{(hpathν\ :\ ∃\ W\ :\ RawLocalStationarity.Data\ (BrownianPathDriver.probability\ (I\ :=\ I)\ ν)}
\leanline{8}{(BrownianPathDriver.driver\ ν)\ M\ T,}
\leanline{6}{W.optimizer\ 0\ =\ V.optimizer\ 0\ ∧}
\leanline{6}{W.sampler.state\ 0\ =ᵐ[BrownianPathDriver.probability\ (I\ :=\ I)\ ν]\ Prod.fst)}
\leanline{4}{(G\ :\ TwistedGeometry\ S\ Z)\ (ρ\ :\ ℝ)\ (hρ\ :\ 0\ <\ ρ)}
\leanline{4}{(hmono\ :\ TwistedJointMonotonicity\ M\ G\ ρ)}
\leanline{4}{(hfull\ :\ DiracOptimizerGrowth.FiniteCouplingWork\ M\ G)}
\leanline{4}{(hσ\ :\ ∀\ s\ s'\ z\ w\ i,\ M.diffusion\ (M.parameter\ s)\ z\ i\ =\ M.diffusion\ (M.parameter\ s')\ w\ i)}
\leanline{4}{(hb\ :\ ∀\ θ,\ Measurable\ (M.drift\ θ))}
\leanline{4}{(hObs\ :\ Measurable\ (fun\ p\ :\ S\ ×\ Z\ =>\ M.observation\ (M.parameter\ p.1)\ (M.position\ p.2)))}
\leanline{4}{(hDrift\ :\ Measurable\ (fun\ p\ :\ S\ ×\ Z\ =>\ M.drift\ (M.parameter\ p.1)\ p.2))}
\leanline{4}{(hOptimizer\ :\ Measurable\ (fun\ p\ :\ ℝ\ ×\ (S\ ×\ Θ)\ =>\ M.optimizer\ p.1\ p.2.1\ p.2.2))\ :}
\leanline{4}{optimizerLawDistance\ (U.optimizer\ T)\ (V.optimizer\ T)}
\leanline{6}{(Q.map\ (U.sampler.state\ T))\ (R.map\ (V.sampler.state\ T))\ ≤}
\leanline{6}{Real.sqrt\ G.conditionNumber\ *\ Real.exp\ (-ρ\ *\ (T\ :\ ℝ))\ *}
\leanline{8}{optimizerLawDistance\ (U.optimizer\ 0)\ (V.optimizer\ 0)\ μ\ ν}
\end{leancode}
\leansource{\texttt{law\_\allowbreak{}distance\_\allowbreak{}decay} \quad PBGSO/\allowbreak{}RawSamplerGrowthArbitraryLaw.lean, lines 43–81}
\begin{leancode}
\leanline{0}{theorem\ law\_distance\_decay}
\leanline{4}{(U\ :\ RawLocalStationarity.Data\ Q\ D\ M\ T)\ (V\ :\ RawLocalStationarity.Data\ R\ E\ M\ T)}
\leanline{4}{(μ\ ν\ :\ Measure\ Z)\ [IsProbabilityMeasure\ μ]\ [IsProbabilityMeasure\ ν]}
\leanline{4}{(hμ\ :\ MemLp\ (fun\ z\ :\ Z\ =>\ z)\ 2\ μ)\ (hν\ :\ MemLp\ (fun\ z\ :\ Z\ =>\ z)\ 2\ ν)}
\leanline{4}{(hU\ :\ Q.map\ (U.sampler.state\ 0)\ =\ μ)\ (hV\ :\ R.map\ (V.sampler.state\ 0)\ =\ ν)}
\leanline{4}{(A\ :\ CanonicalSolution\ (M\ :=\ M)\ (T\ :=\ T)\ μ\ (U.optimizer\ 0))}
\leanline{4}{(B\ :\ CanonicalSolution\ (M\ :=\ M)\ (T\ :=\ T)\ ν\ (V.optimizer\ 0))}
\leanline{4}{(huniqueU\ :\ PathwiseUniqueness\ Q\ (nullAugmentedDriver\ Q\ (continuousDriver\ Q\ D))\ M\ T)}
\leanline{4}{(huniqueV\ :\ PathwiseUniqueness\ R\ (nullAugmentedDriver\ R\ (continuousDriver\ R\ E))\ M\ T)}
\leanline{4}{(huniqueA\ :\ PathwiseUniqueness\ (μ.prod\ (brownianProbability\ (I\ :=\ I)))}
\leanline{6}{(nullAugmentedDriver\ \_\ (continuousDriver\ \_\ (CanonicalBrownianDriver.driver\ (I\ :=\ I)\ μ)))\ M\ T)}
\leanline{4}{(huniqueB\ :\ PathwiseUniqueness\ (ν.prod\ (brownianProbability\ (I\ :=\ I)))}
\leanline{6}{(nullAugmentedDriver\ \_\ (continuousDriver\ \_\ (CanonicalBrownianDriver.driver\ (I\ :=\ I)\ ν)))\ M\ T)}
\leanline{4}{(hpathμ\ :\ ∃\ W\ :\ RawLocalStationarity.Data\ (BrownianPathDriver.probability\ (I\ :=\ I)\ μ)}
\leanline{8}{(BrownianPathDriver.driver\ μ)\ M\ T,}
\leanline{6}{W.optimizer\ 0\ =\ U.optimizer\ 0\ ∧}
\leanline{6}{W.sampler.state\ 0\ =ᵐ[BrownianPathDriver.probability\ (I\ :=\ I)\ μ]\ Prod.fst)}
\leanline{4}{(hpathν\ :\ ∃\ W\ :\ RawLocalStationarity.Data\ (BrownianPathDriver.probability\ (I\ :=\ I)\ ν)}
\leanline{8}{(BrownianPathDriver.driver\ ν)\ M\ T,}
\leanline{6}{W.optimizer\ 0\ =\ V.optimizer\ 0\ ∧}
\leanline{6}{W.sampler.state\ 0\ =ᵐ[BrownianPathDriver.probability\ (I\ :=\ I)\ ν]\ Prod.fst)}
\leanline{4}{(G\ :\ TwistedGeometry\ S\ Z)\ (ρ\ :\ ℝ)\ (hmono\ :\ TwistedJointMonotonicity\ M\ G\ ρ)\ :}
\leanline{4}{optimizerLawDistance\ (U.optimizer\ T)\ (V.optimizer\ T)}
\leanline{6}{(Q.map\ (U.sampler.state\ T))\ (R.map\ (V.sampler.state\ T))\ ≤}
\leanline{6}{Real.sqrt\ G.conditionNumber\ *\ Real.exp\ (-ρ\ *\ (T\ :\ ℝ))\ *}
\leanline{8}{optimizerLawDistance\ (U.optimizer\ 0)\ (V.optimizer\ 0)\ μ\ ν}
\end{leancode}
\leansource{\texttt{wasserstein\_\allowbreak{}product\_\allowbreak{}eq} \quad PBGSO/\allowbreak{}ProductTransport.lean, lines 178–207}
\begin{leancode}
\leanline{0}{theorem\ wasserstein\_product\_eq\ (s\ s'\ :\ S)\ (ν\ ν'\ :\ Measure\ Z)}
\leanline{4}{[IsProbabilityMeasure\ ν]\ [IsProbabilityMeasure\ ν']}
\leanline{4}{(hν\ :\ MemLp\ (fun\ z\ :\ Z\ =>\ z)\ 2\ ν)\ (hν'\ :\ MemLp\ (fun\ z\ :\ Z\ =>\ z)\ 2\ ν')\ :}
\leanline{4}{ULA.wasserstein2Real\ (law\ s\ ν)\ (law\ s'\ ν')\ =}
\leanline{6}{Real.sqrt\ (‖s\ -\ s'‖\ \textasciicircum{}\ 2\ +\ ULA.wasserstein2Real\ ν\ ν'\ \textasciicircum{}\ 2)}
\end{leancode}
\leansource{\texttt{stationary\_\allowbreak{}laws\_\allowbreak{}unique\_\allowbreak{}constant} \quad PBGSO/\allowbreak{}RawStationaryUniqueness.lean, lines 67–113}
\begin{leancode}
\leanline{0}{theorem\ stationary\_laws\_unique\_constant}
\leanline{4}{(F\ :\ Family\ Q\ D\ M\ s\ ν)\ (H\ :\ Family\ R\ E\ M\ s'\ ν')}
\leanline{4}{(hν\ :\ MemLp\ (fun\ z\ :\ Z\ =>\ z)\ 2\ ν)\ (hν'\ :\ MemLp\ (fun\ z\ :\ Z\ =>\ z)\ 2\ ν')}
\leanline{4}{(G\ :\ TwistedGeometry\ S\ Z)\ (ρ\ :\ ℝ)\ (hρ\ :\ 0\ <\ ρ)}
\leanline{4}{(hmono\ :\ TwistedJointMonotonicity\ M\ G\ ρ)}
\leanline{4}{(hfull\ :\ DiracOptimizerGrowth.FiniteCouplingWork\ M\ G)}
\leanline{4}{(hσ\ :\ ∀\ a\ a'\ z\ w\ i,\ M.diffusion\ (M.parameter\ a)\ z\ i\ =\ M.diffusion\ (M.parameter\ a')\ w\ i)}
\leanline{4}{(hb\ :\ ∀\ θ,\ Measurable\ (M.drift\ θ))}
\leanline{4}{(hObs\ :\ Measurable\ (fun\ p\ :\ S\ ×\ Z\ =>\ M.observation\ (M.parameter\ p.1)\ (M.position\ p.2)))}
\leanline{4}{(hDrift\ :\ Measurable\ (fun\ p\ :\ S\ ×\ Z\ =>\ M.drift\ (M.parameter\ p.1)\ p.2))}
\leanline{4}{(hOptimizer\ :\ Measurable\ (fun\ p\ :\ ℝ\ ×\ (S\ ×\ Θ)\ =>\ M.optimizer\ p.1\ p.2.1\ p.2.2))\ :}
\leanline{4}{s\ =\ s'\ ∧\ ν\ =\ ν'}
\end{leancode}
\leansource{\texttt{stationary\_\allowbreak{}laws\_\allowbreak{}unique\_\allowbreak{}linearGrowth} \quad PBGSO/\allowbreak{}RawStationaryUniqueness.lean, lines 119–161}
\begin{leancode}
\leanline{0}{theorem\ stationary\_laws\_unique\_linearGrowth}
\leanline{4}{(F\ :\ Family\ Q\ D\ M\ s\ ν)\ (H\ :\ Family\ R\ E\ M\ s'\ ν')}
\leanline{4}{(RF\ :\ ∀\ T\ :\ ℝ≥0,\ RawSamplerGrowth.Regularity\ (F.canonical\ T).data)}
\leanline{4}{(RH\ :\ ∀\ T\ :\ ℝ≥0,\ RawSamplerGrowth.Regularity\ (H.canonical\ T).data)}
\leanline{4}{(hν\ :\ MemLp\ (fun\ z\ :\ Z\ =>\ z)\ 2\ ν)\ (hν'\ :\ MemLp\ (fun\ z\ :\ Z\ =>\ z)\ 2\ ν')}
\leanline{4}{(G\ :\ TwistedGeometry\ S\ Z)\ (ρ\ :\ ℝ)\ (hρ\ :\ 0\ <\ ρ)}
\leanline{4}{(hmono\ :\ TwistedJointMonotonicity\ M\ G\ ρ)\ :\ s\ =\ s'\ ∧\ ν\ =\ ν'}
\end{leancode}
\end{leanbox}
\subsection{\texorpdfstring{Theorem 3 — Finite-particle interaction error}{Theorem 3 — Finite-particle interaction error}}


\textbf{Assumptions:} A1, A3–A5; positive population; deterministic optimiser initialisation; i.i.d. initial samplers from a common second-moment law, independent of the independent Brownian drivers; and an actual interacting solution of the original native SDE.

Extra assumption: R1.

Couple the particle and independent-copy systems synchronously: \(\vartheta_0^N=\bar\vartheta_0\), \(Z_0^i=\bar Z_0^i\), with the same driver for each pair. Put \(\bar\nu_t^{[N]}=N^{-1}\sum_i\delta_{\bar Z_t^i}\).





For every \(t\ge0\),





\[\left\|\mathbf d((\vartheta_t^N,\nu_t^N),(\bar\vartheta_t,\bar\nu_t^{[N]}))\right\|_{\mathrm L^2}\le\frac{L_\Psi}{\rho}\sqrt{\frac{\kappa V_F}{N}}(1-e^{-\rho t}).\]







The comparison is with the independent-copy empirical law. A population-law Wasserstein estimate also needs empirical sampling error.

\begin{leanbox}{Lean: \texttt{original\_\allowbreak{}finite\_\allowbreak{}particle\_\allowbreak{}interaction\_\allowbreak{}bound}}
\leansource{\texttt{original\_\allowbreak{}finite\_\allowbreak{}particle\_\allowbreak{}interaction\_\allowbreak{}bound} \quad PBGSO/\allowbreak{}RawConstructedParticleBounds.lean, lines 109–132}
\begin{leancode}
\leanline{0}{theorem\ original\_finite\_particle\_interaction\_bound}
\leanline{4}{(hBorelF\ :\ Measurable\ (fun\ p\ :\ S\ ×\ Z\ =>\ M.observation\ (M.parameter\ p.1)\ (M.position\ p.2)))}
\leanline{4}{(hBorelb\ :\ Measurable\ (fun\ p\ :\ S\ ×\ Z\ =>\ M.drift\ (M.parameter\ p.1)\ p.2))}
\leanline{4}{(hBorelΨ\ :\ Measurable\ (fun\ p\ :\ ℝ\ ×\ (S\ ×\ Θ)\ =>\ M.optimizer\ p.1\ p.2.1\ p.2.2))}
\leanline{4}{(G\ :\ TwistedGeometry\ S\ Z)\ (ρ\ :\ ℝ)\ (hρ\ :\ 0\ <\ ρ)}
\leanline{4}{(hmono\ :\ TwistedJointMonotonicity\ M\ G\ ρ)\ (hN\ :\ 0\ <\ N)}
\leanline{4}{(LΨ\ :\ ℝ≥0)\ (V\ :\ ℝ)\ (hV\ :\ 0\ ≤\ V)}
\leanline{4}{(hΨ\ :\ ∀\ t\ ∈\ Icc\ 0\ (T\ :\ ℝ),\ ∀\ s\ g\ g',}
\leanline{6}{‖M.optimizer\ t\ s\ g-M.optimizer\ t\ s\ g'‖\ ≤\ (LΨ\ :\ ℝ)*‖g-g'‖)}
\leanline{4}{(hF\ :\ ∀\ t\ ∈\ Icc\ 0\ (T\ :\ ℝ),\ MemLp}
\leanline{6}{(fun\ z\ =>\ M.observation\ (M.parameter\ (U.optimizer\ t))\ (M.position\ z))\ 2}
\leanline{6}{(Q.map\ (U.sampler.state\ t)))}
\leanline{4}{(hVar\ :\ ∀\ t\ ∈\ Icc\ 0\ (T\ :\ ℝ),}
\leanline{6}{(∫\ z,\ ‖M.observation\ (M.parameter\ (U.optimizer\ t))\ (M.position\ z)-}
\leanline{8}{M.meanGradient\ (Q.map\ (U.sampler.state\ t))\ (U.optimizer\ t)‖\textasciicircum{}2}
\leanline{10}{∂Q.map\ (U.sampler.state\ t))\ ≤\ V)}
\leanline{4}{(hinit\ :\ ∀ᵐ\ ω\ ∂IndependentCopies.seeds\ (J\ :=\ Fin\ N)\ Q,}
\leanline{6}{oneEquiv\ N\ (Y.state\ 0\ ω)\ =\ (U.optimizer\ 0,RawIndependentCopies.cloud\ Q\ D\ N\ U\ 0\ ω))\ :}
\leanline{4}{Real.sqrt\ (∫\ ω,\ particleEmpiricalDistanceSq\ N\ (originalPair\ Q\ D\ M\ N\ U\ Y\ T\ ω)}
\leanline{6}{∂IndependentCopies.seeds\ (J\ :=\ Fin\ N)\ Q)\ ≤}
\leanline{6}{((LΨ\ :\ ℝ)/ρ)*Real.sqrt\ (G.conditionNumber*V/(N\ :\ ℝ))*(1-Real.exp\ (-ρ*T))}
\end{leancode}
\end{leanbox}
\subsection{\texorpdfstring{Theorem 4 — Parameter error bound of particle system}{Theorem 4 — Parameter error bound of particle system}}


\textbf{Assumptions:} A1–A5 and the coupled initialisation/\allowbreak{}native interacting solution of Theorem 3. The parameter map is the Euclidean coordinate projection.

Extra assumptions: R1; also R3 in the linear-growth branch.



\[\mathbf d_0=\mathbf d((\vartheta_0^N,\nu_0),(\vartheta_\star,\nu_{\theta_\star})).\]







For every \(t\ge0\),





\[\|\theta_t^N-\theta_\star\|_{\mathrm L^2}\le\sqrt\kappa\,e^{-\rho t}\mathbf d_0+\frac{L_\Psi}{\rho}\sqrt{\frac{\kappa V_F}{N}}(1-e^{-\rho t}).\]





\[\lim_{N\to\infty}\limsup_{t\to\infty}\|\theta_t^N-\theta_\star\|_{\mathrm L^2}=0.\]



\begin{leanbox}{Lean: \texttt{original\_\allowbreak{}parameter\_\allowbreak{}bound\_\allowbreak{}constant}, \texttt{original\_\allowbreak{}parameter\_\allowbreak{}bound\_\allowbreak{}linearGrowth}, \texttt{parameter\_\allowbreak{}joint\_\allowbreak{}limit}}
\leansource{\texttt{original\_\allowbreak{}parameter\_\allowbreak{}bound\_\allowbreak{}constant} \quad PBGSO/\allowbreak{}RawConstructedParticleParameter.lean, lines 98–153}
\begin{leancode}
\leanline{0}{theorem\ original\_parameter\_bound\_constant}
\leanline{4}{(U\ :\ RawLocalStationarity.Data\ Q\ W\ M\ T)}
\leanline{4}{(Y\ :\ RawEmpiricalDynamics.Solution\ M\ N\ (IndependentCopies.seeds\ (J\ :=\ Fin\ N)\ Q)}
\leanline{6}{(RawIndependentCopies.driver\ Q\ W\ N)\ T)}
\leanline{4}{(μ\ νStar\ :\ Measure\ Z)\ [IsProbabilityMeasure\ μ]\ [IsProbabilityMeasure\ νStar]}
\leanline{4}{(hμ\ :\ MemLp\ (fun\ z\ :\ Z\ =>\ z)\ 2\ μ)\ (hν\ :\ MemLp\ (fun\ z\ :\ Z\ =>\ z)\ 2\ νStar)}
\leanline{4}{(hUlaw\ :\ Q.map\ (U.sampler.state\ 0)\ =\ μ)}
\leanline{4}{(A\ :\ RawLocalStationarity.Data\ (μ.prod\ (brownianProbability\ (I\ :=\ I)))}
\leanline{6}{(CanonicalBrownianDriver.driver\ (I\ :=\ I)\ μ)\ M\ T)}
\leanline{4}{(hA0\ :\ A.optimizer\ 0\ =\ U.optimizer\ 0)}
\leanline{4}{(hAz\ :\ A.sampler.state\ 0\ =ᵐ[μ.prod\ (brownianProbability\ (I\ :=\ I))]\ Prod.fst)}
\leanline{4}{(V\ :\ RawLocalStationarity.Data\ (νStar.prod\ (brownianProbability\ (I\ :=\ I)))}
\leanline{6}{(CanonicalBrownianDriver.driver\ (I\ :=\ I)\ νStar)\ M\ T)}
\leanline{4}{(hVz\ :\ V.sampler.state\ 0\ =ᵐ[νStar.prod\ (brownianProbability\ (I\ :=\ I))]\ Prod.fst)}
\leanline{4}{(huniqueU\ :\ PathwiseUniqueness\ Q\ (nullAugmentedDriver\ Q\ (continuousDriver\ Q\ W))\ M\ T)}
\leanline{4}{(huniqueA\ :\ PathwiseUniqueness\ (μ.prod\ (brownianProbability\ (I\ :=\ I)))}
\leanline{6}{(nullAugmentedDriver\ \_\ (continuousDriver\ \_\ (CanonicalBrownianDriver.driver\ (I\ :=\ I)\ μ)))\ M\ T)}
\leanline{4}{(hpath\ :\ ∃\ B\ :\ RawLocalStationarity.Data\ (BrownianPathDriver.probability\ (I\ :=\ I)\ μ)}
\leanline{8}{(BrownianPathDriver.driver\ μ)\ M\ T,}
\leanline{6}{B.optimizer\ 0\ =\ U.optimizer\ 0\ ∧}
\leanline{6}{B.sampler.state\ 0\ =ᵐ[BrownianPathDriver.probability\ (I\ :=\ I)\ μ]\ Prod.fst)}
\leanline{4}{(G\ :\ TwistedGeometry\ S\ Z)\ (ρ\ :\ ℝ)\ (hρ\ :\ 0\ <\ ρ)}
\leanline{4}{(hmono\ :\ TwistedJointMonotonicity\ M\ G\ ρ)\ (hN\ :\ 0\ <\ N)}
\leanline{4}{(hProjection\ :\ LipschitzWith\ 1\ M.parameter)}
\leanline{4}{(hσ\ :\ ∀\ s\ s'\ z\ w\ i,\ M.diffusion\ (M.parameter\ s)\ z\ i\ =\ M.diffusion\ (M.parameter\ s')\ w\ i)}
\leanline{4}{(hb\ :\ ∀\ θ,\ Measurable\ (M.drift\ θ))}
\leanline{4}{(hF\ :\ Measurable\ (fun\ p\ :\ S\ ×\ Z\ =>\ M.observation\ (M.parameter\ p.1)\ (M.position\ p.2)))}
\leanline{4}{(hDrift\ :\ Measurable\ (fun\ p\ :\ S\ ×\ Z\ =>\ M.drift\ (M.parameter\ p.1)\ p.2))}
\leanline{4}{(hOptimizer\ :\ Measurable\ (fun\ p\ :\ ℝ\ ×\ (S\ ×\ Θ)\ =>\ M.optimizer\ p.1\ p.2.1\ p.2.2))}
\leanline{4}{(LΨ\ :\ ℝ≥0)\ (VF\ :\ ℝ)}
\leanline{4}{(hΨ\ :\ ∀\ r\ ∈\ Icc\ 0\ (T\ :\ ℝ),\ ∀\ s\ g\ g',}
\leanline{6}{‖M.optimizer\ r\ s\ g-M.optimizer\ r\ s\ g'‖\ ≤\ (LΨ\ :\ ℝ)*‖g-g'‖)}
\leanline{4}{(hU2\ :\ ∀\ r\ ∈\ Icc\ 0\ (T\ :\ ℝ),\ MemLp}
\leanline{6}{(fun\ z\ =>\ M.observation\ (M.parameter\ (U.optimizer\ r))\ (M.position\ z))\ 2\ (Q.map\ (U.sampler.state\ r)))}
\leanline{4}{(hVar\ :\ ∀\ r\ ∈\ Icc\ 0\ (T\ :\ ℝ),}
\leanline{6}{(∫\ z,\ ‖M.observation\ (M.parameter\ (U.optimizer\ r))\ (M.position\ z)-}
\leanline{8}{M.meanGradient\ (Q.map\ (U.sampler.state\ r))\ (U.optimizer\ r)‖\textasciicircum{}2\ ∂Q.map\ (U.sampler.state\ r))\ ≤\ VF)}
\leanline{4}{(sStar\ :\ S)}
\leanline{4}{(hopt\ :\ ∀\ r\ ∈\ Icc\ 0\ (T\ :\ ℝ),\ V.optimizer\ r\ =\ sStar)}
\leanline{4}{(hlaw\ :\ ∀\ r\ ∈\ Icc\ 0\ (T\ :\ ℝ),\ (νStar.prod\ (brownianProbability\ (I\ :=\ I))).map\ (V.sampler.state\ r)\ =\ νStar)}
\leanline{4}{(hobs\ :\ Integrable\ (fun\ z\ =>\ M.observation\ (M.parameter\ sStar)\ (M.position\ z))\ νStar)}
\leanline{4}{(hinit\ :\ ∀ᵐ\ ω\ ∂IndependentCopies.seeds\ (J\ :=\ Fin\ N)\ Q,}
\leanline{6}{oneEquiv\ N\ (Y.state\ 0\ ω)\ =\ (U.optimizer\ 0,RawIndependentCopies.cloud\ Q\ W\ N\ U\ 0\ ω))\ :}
\leanline{4}{MemLp\ (fun\ ω\ =>\ M.parameter\ (oneEquiv\ N\ (Y.state\ T\ ω)).1-M.parameter\ sStar)}
\leanline{6}{2\ (IndependentCopies.seeds\ (J\ :=\ Fin\ N)\ Q)\ ∧}
\leanline{4}{Real.sqrt\ (∫\ ω,\ ‖M.parameter\ (oneEquiv\ N\ (Y.state\ T\ ω)).1-M.parameter\ sStar‖\textasciicircum{}2}
\leanline{6}{∂IndependentCopies.seeds\ (J\ :=\ Fin\ N)\ Q)\ ≤}
\leanline{6}{Real.sqrt\ G.conditionNumber*Real.exp\ (-ρ*(T\ :\ ℝ))*optimizerLawDistance\ (U.optimizer\ 0)\ sStar\ μ\ νStar+}
\leanline{6}{((LΨ\ :\ ℝ)/ρ)*Real.sqrt\ (G.conditionNumber*VF/(N\ :\ ℝ))*(1-Real.exp\ (-ρ*T))}
\end{leancode}
\leansource{\texttt{original\_\allowbreak{}parameter\_\allowbreak{}bound\_\allowbreak{}linearGrowth} \quad PBGSO/\allowbreak{}RawConstructedParticleParameter.lean, lines 159–226}
\begin{leancode}
\leanline{0}{theorem\ original\_parameter\_bound\_linearGrowth}
\leanline{4}{(U\ :\ RawLocalStationarity.Data\ Q\ W\ M\ T)}
\leanline{4}{(Y\ :\ RawEmpiricalDynamics.Solution\ M\ N\ (IndependentCopies.seeds\ (J\ :=\ Fin\ N)\ Q)}
\leanline{6}{(RawIndependentCopies.driver\ Q\ W\ N)\ T)}
\leanline{4}{(μ\ νStar\ :\ Measure\ Z)\ [IsProbabilityMeasure\ μ]\ [IsProbabilityMeasure\ νStar]}
\leanline{4}{(hμ\ :\ MemLp\ (fun\ z\ :\ Z\ =>\ z)\ 2\ μ)\ (hν\ :\ MemLp\ (fun\ z\ :\ Z\ =>\ z)\ 2\ νStar)}
\leanline{4}{(hUlaw\ :\ Q.map\ (U.sampler.state\ 0)\ =\ μ)}
\leanline{4}{(A\ :\ RawLocalStationarity.Data\ (μ.prod\ (brownianProbability\ (I\ :=\ I)))}
\leanline{6}{(CanonicalBrownianDriver.driver\ (I\ :=\ I)\ μ)\ M\ T)}
\leanline{4}{(RA\ :\ RawSamplerGrowth.Regularity\ A)}
\leanline{4}{(hA0\ :\ A.optimizer\ 0\ =\ U.optimizer\ 0)}
\leanline{4}{(hAz\ :\ A.sampler.state\ 0\ =ᵐ[μ.prod\ (brownianProbability\ (I\ :=\ I))]\ Prod.fst)}
\leanline{4}{(V\ :\ RawLocalStationarity.Data\ (νStar.prod\ (brownianProbability\ (I\ :=\ I)))}
\leanline{6}{(CanonicalBrownianDriver.driver\ (I\ :=\ I)\ νStar)\ M\ T)}
\leanline{4}{(RV\ :\ RawSamplerGrowth.Regularity\ V)}
\leanline{4}{(hVz\ :\ V.sampler.state\ 0\ =ᵐ[νStar.prod\ (brownianProbability\ (I\ :=\ I))]\ Prod.fst)}
\leanline{4}{(huniqueU\ :\ PathwiseUniqueness\ Q\ (nullAugmentedDriver\ Q\ (continuousDriver\ Q\ W))\ M\ T)}
\leanline{4}{(huniqueA\ :\ PathwiseUniqueness\ (μ.prod\ (brownianProbability\ (I\ :=\ I)))}
\leanline{6}{(nullAugmentedDriver\ \_\ (continuousDriver\ \_\ (CanonicalBrownianDriver.driver\ (I\ :=\ I)\ μ)))\ M\ T)}
\leanline{4}{(huniqueV\ :\ PathwiseUniqueness\ (νStar.prod\ (brownianProbability\ (I\ :=\ I)))}
\leanline{6}{(nullAugmentedDriver\ \_\ (continuousDriver\ \_\ (CanonicalBrownianDriver.driver\ (I\ :=\ I)\ νStar)))\ M\ T)}
\leanline{4}{(hpathμ\ :\ ∃\ B\ :\ RawLocalStationarity.Data\ (BrownianPathDriver.probability\ (I\ :=\ I)\ μ)}
\leanline{8}{(BrownianPathDriver.driver\ μ)\ M\ T,}
\leanline{6}{B.optimizer\ 0\ =\ U.optimizer\ 0\ ∧}
\leanline{6}{B.sampler.state\ 0\ =ᵐ[BrownianPathDriver.probability\ (I\ :=\ I)\ μ]\ Prod.fst)}
\leanline{4}{(hpathν\ :\ ∃\ B\ :\ RawLocalStationarity.Data\ (BrownianPathDriver.probability\ (I\ :=\ I)\ νStar)}
\leanline{8}{(BrownianPathDriver.driver\ νStar)\ M\ T,}
\leanline{6}{B.optimizer\ 0\ =\ V.optimizer\ 0\ ∧}
\leanline{6}{B.sampler.state\ 0\ =ᵐ[BrownianPathDriver.probability\ (I\ :=\ I)\ νStar]\ Prod.fst)}
\leanline{4}{(G\ :\ TwistedGeometry\ S\ Z)\ (ρ\ :\ ℝ)\ (hρ\ :\ 0\ <\ ρ)}
\leanline{4}{(hmono\ :\ TwistedJointMonotonicity\ M\ G\ ρ)\ (hN\ :\ 0\ <\ N)}
\leanline{4}{(hProjection\ :\ LipschitzWith\ 1\ M.parameter)}
\leanline{4}{(hF\ :\ Measurable\ (fun\ p\ :\ S\ ×\ Z\ =>\ M.observation\ (M.parameter\ p.1)\ (M.position\ p.2)))}
\leanline{4}{(hDrift\ :\ Measurable\ (fun\ p\ :\ S\ ×\ Z\ =>\ M.drift\ (M.parameter\ p.1)\ p.2))}
\leanline{4}{(hOptimizer\ :\ Measurable\ (fun\ p\ :\ ℝ\ ×\ (S\ ×\ Θ)\ =>\ M.optimizer\ p.1\ p.2.1\ p.2.2))}
\leanline{4}{(LΨ\ :\ ℝ≥0)\ (VF\ :\ ℝ)}
\leanline{4}{(hΨ\ :\ ∀\ r\ ∈\ Icc\ 0\ (T\ :\ ℝ),\ ∀\ s\ g\ g',}
\leanline{6}{‖M.optimizer\ r\ s\ g-M.optimizer\ r\ s\ g'‖\ ≤\ (LΨ\ :\ ℝ)*‖g-g'‖)}
\leanline{4}{(hU2\ :\ ∀\ r\ ∈\ Icc\ 0\ (T\ :\ ℝ),\ MemLp}
\leanline{6}{(fun\ z\ =>\ M.observation\ (M.parameter\ (U.optimizer\ r))\ (M.position\ z))\ 2\ (Q.map\ (U.sampler.state\ r)))}
\leanline{4}{(hVar\ :\ ∀\ r\ ∈\ Icc\ 0\ (T\ :\ ℝ),}
\leanline{6}{(∫\ z,\ ‖M.observation\ (M.parameter\ (U.optimizer\ r))\ (M.position\ z)-}
\leanline{8}{M.meanGradient\ (Q.map\ (U.sampler.state\ r))\ (U.optimizer\ r)‖\textasciicircum{}2\ ∂Q.map\ (U.sampler.state\ r))\ ≤\ VF)}
\leanline{4}{(sStar\ :\ S)}
\leanline{4}{(hopt\ :\ ∀\ r\ ∈\ Icc\ 0\ (T\ :\ ℝ),\ V.optimizer\ r\ =\ sStar)}
\leanline{4}{(hlaw\ :\ ∀\ r\ ∈\ Icc\ 0\ (T\ :\ ℝ),\ (νStar.prod\ (brownianProbability\ (I\ :=\ I))).map\ (V.sampler.state\ r)\ =\ νStar)}
\leanline{4}{(hinit\ :\ ∀ᵐ\ ω\ ∂IndependentCopies.seeds\ (J\ :=\ Fin\ N)\ Q,}
\leanline{6}{oneEquiv\ N\ (Y.state\ 0\ ω)\ =\ (U.optimizer\ 0,RawIndependentCopies.cloud\ Q\ W\ N\ U\ 0\ ω))\ :}
\leanline{4}{MemLp\ (fun\ ω\ =>\ M.parameter\ (oneEquiv\ N\ (Y.state\ T\ ω)).1-M.parameter\ sStar)}
\leanline{6}{2\ (IndependentCopies.seeds\ (J\ :=\ Fin\ N)\ Q)\ ∧}
\leanline{4}{Real.sqrt\ (∫\ ω,\ ‖M.parameter\ (oneEquiv\ N\ (Y.state\ T\ ω)).1-M.parameter\ sStar‖\textasciicircum{}2}
\leanline{6}{∂IndependentCopies.seeds\ (J\ :=\ Fin\ N)\ Q)\ ≤}
\leanline{6}{Real.sqrt\ G.conditionNumber*Real.exp\ (-ρ*(T\ :\ ℝ))*optimizerLawDistance\ (U.optimizer\ 0)\ sStar\ μ\ νStar+}
\leanline{6}{((LΨ\ :\ ℝ)/ρ)*Real.sqrt\ (G.conditionNumber*VF/(N\ :\ ℝ))*(1-Real.exp\ (-ρ*T))}
\end{leancode}
\leansource{\texttt{parameter\_\allowbreak{}joint\_\allowbreak{}limit} \quad PBGSO/\allowbreak{}RawConstructedParticleJointLimit.lean, lines 91–161}
\begin{leancode}
\leanline{0}{theorem\ parameter\_joint\_limit}
\leanline{4}{(hμ\ :\ MemLp\ (fun\ z\ :\ Z\ =>\ z)\ 2\ μ)\ (hν\ :\ MemLp\ (fun\ z\ :\ Z\ =>\ z)\ 2\ νStar)}
\leanline{4}{(G\ :\ TwistedGeometry\ S\ Z)\ (ρ\ :\ ℝ)\ (hρ\ :\ 0\ <\ ρ)}
\leanline{4}{(hmono\ :\ TwistedJointMonotonicity\ M\ G\ ρ)}
\leanline{4}{(hProjection\ :\ LipschitzWith\ 1\ M.parameter)}
\leanline{4}{(hb\ :\ ∀\ θ,\ Measurable\ (M.drift\ θ))}
\leanline{4}{(hF\ :\ Measurable\ (fun\ p\ :\ S\ ×\ Z\ =>\ M.observation\ (M.parameter\ p.1)\ (M.position\ p.2)))}
\leanline{4}{(hDrift\ :\ Measurable\ (fun\ p\ :\ S\ ×\ Z\ =>\ M.drift\ (M.parameter\ p.1)\ p.2))}
\leanline{4}{(hOptimizer\ :\ Measurable\ (fun\ p\ :\ ℝ\ ×\ (S\ ×\ Θ)\ =>\ M.optimizer\ p.1\ p.2.1\ p.2.2))}
\leanline{4}{(LΨ\ :\ ℝ≥0)\ (VF\ :\ ℝ)}
\leanline{4}{(hΨ\ :\ ∀\ r,\ 0\ ≤\ r\ →\ ∀\ s\ g\ g',}
\leanline{6}{‖M.optimizer\ r\ s\ g-M.optimizer\ r\ s\ g'‖\ ≤\ (LΨ\ :\ ℝ)*‖g-g'‖)}
\leanline{4}{(hU2\ :\ ∀\ T\ :\ ℝ≥0,\ ∀\ r\ ∈\ Icc\ 0\ (T\ :\ ℝ),\ MemLp}
\leanline{6}{(fun\ z\ =>\ M.observation\ (M.parameter\ ((F.reference\ T).optimizer\ r))\ (M.position\ z))}
\leanline{8}{2\ (Q.map\ ((F.reference\ T).sampler.state\ r)))}
\leanline{4}{(hVar\ :\ ∀\ T\ :\ ℝ≥0,\ ∀\ r\ ∈\ Icc\ 0\ (T\ :\ ℝ),}
\leanline{6}{(∫\ z,\ ‖M.observation\ (M.parameter\ ((F.reference\ T).optimizer\ r))\ (M.position\ z)-}
\leanline{8}{M.meanGradient\ (Q.map\ ((F.reference\ T).sampler.state\ r))\ ((F.reference\ T).optimizer\ r)‖\textasciicircum{}2}
\leanline{8}{∂Q.map\ ((F.reference\ T).sampler.state\ r))\ ≤\ VF)}
\leanline{4}{(hobs\ :\ Integrable\ (fun\ z\ =>\ M.observation\ (M.parameter\ sStar)\ (M.position\ z))\ νStar)}
\leanline{4}{(hbranch\ :\ Branch\ Q\ W\ M\ μ\ νStar\ s0\ sStar\ F)\ :}
\leanline{4}{Tendsto\ (fun\ p\ :\ ℕ\ ×\ ℝ\ =>\ parameterError\ Q\ W\ M\ μ\ νStar\ s0\ sStar\ F\ p.1\ p.2)}
\leanline{6}{(atTop\ ×ˢ\ atTop)\ ({\leanfallback 𝓝}\ 0)}
\end{leancode}
\leanlabel{How the Lean proof follows the argument}
The limit is for a family with common constants and initial law. The verified family theorem also gives a joint time/\allowbreak{}population limit. The critical-point target is the one in A2. In Lean, A2 enters as a canonical target solution whose optimiser stays at \(\vartheta_\star\) and whose sampler law stays \(\nu_{\theta_\star}\) (premises \texttt{hopt}, \texttt{hlaw}); Theorem 1 proves this stationarity from A1 and A2.
\end{leanbox}
\subsection{\texorpdfstring{Theorem New A — M-MFPO final discretisation bound}{Theorem New A — M-MFPO final discretisation bound}}


\textbf{Assumptions:} N1–N3 and the standard Euclidean cloud encoding defined above. Use the paper’s OBABO update and independent Gaussian sampler half-step innovations, with zero optimiser innovations.

Extra assumptions: N1 and N3; \(\gamma\ge0\), \(0<h\le h_{\star,M}\).

The stationary optimiser state is \(\vartheta_\star=(\theta_\star,0)\). Initialise the numerical and continuous interacting systems identically.





For every iteration \(n\in\mathbb N\), with all constants independent of time and step size,





\[\|\hat\theta_n-\theta_\star\|_{\mathrm L^2}\le D_M\sqrt h+\sqrt\kappa\,e^{-\rho nh}\mathbf d_0+\frac{2+\gamma}{\rho}\sqrt{\frac{\kappa V_F}{N}}(1-e^{-\rho nh}).\]







The full-state raw maximum-component \(L^2\) error is also at most the same discretisation term. The factor in the particle floor is the certified response constant used by this endpoint. Constants can depend on fixed population and initial law; they do not grow with iteration count.

\begin{leanbox}{Lean: \texttt{explicit\_\allowbreak{}parameter\_\allowbreak{}bound}, \texttt{explicit\_\allowbreak{}discretization\_\allowbreak{}bound}}
\leansource{\texttt{explicit\_\allowbreak{}parameter\_\allowbreak{}bound} \quad PBGSO/\allowbreak{}FullyUniformMomentum.lean, lines 256–276}
\begin{leancode}
\leanline{0}{theorem\ explicit\_parameter\_bound}
\leanline{4}{(hcloud\ :\ D.cloud\ =\ euclideanParticleEncoding\ N\ J)\ (hγ\ :\ 0\ ≤\ γ)}
\leanline{4}{(star\ :\ A)\ (νStar\ :\ Measure\ (HypocoercivePhase\ (EuclideanSpace\ ℝ\ J)))\ [IsProbabilityMeasure\ νStar]}
\leanline{4}{(hνStar\ :\ MemLp\ (fun\ z\ :\ HypocoercivePhase\ (EuclideanSpace\ ℝ\ J)\ ↦\ z)\ 2\ νStar)}
\leanline{4}{(Frozen\ :\ AdditiveStationary.FrozenTrajectory\ (ULA.CanonicalBrownian.vectorProbability\ J)}
\leanline{6}{(momentumForceDynamics\ D\ γ)\ νStar\ hνStar\ (targetState\ star)\ (input\ γ))}
\leanline{4}{(ρ\ :\ ℝ)\ (hρ\ :\ 0\ <\ ρ)\ (hmono\ :\ TwistedJointMonotonicity\ (model\ D\ γ)\ Geometry\ ρ)}
\leanline{4}{(hg\ :\ (model\ D\ γ).meanGradient\ νStar\ (targetState\ star)\ =\ 0)}
\leanline{4}{(hV\ :\ 0\ ≤\ Variance)}
\leanline{4}{(h\ :\ ℝ)\ (hh\ :\ 0\ <\ h)\ (hhmax\ :\ h\ ≤\ maximumStep\ D\ γ\ Geometry\ μ\ hγ\ ρ)\ (n\ :\ ℕ)\ :}
\leanline{8}{‖(momentumParticleIterate\ (ParticleNoiseSpace.probability\ (J\ :=\ J)\ μ\ N)}
\leanline{10}{(particleBrownian\ (J\ :=\ J)\ μ\ N)\ (particleBrownian\_isPreBrownian\ (J\ :=\ J)\ μ\ N)}
\leanline{10}{D\ γ\ h\ (initial\ D\ γ\ Geometry\ μ\ hμ\ a\ Variance\ hVariance)\ n).2.1\ -}
\leanline{10}{Lp.const\ 2\ (ParticleNoiseSpace.probability\ (J\ :=\ J)\ μ\ N)\ star‖\ ≤}
\leanline{8}{discretizationConstant\ D\ γ\ Geometry\ μ\ hμ\ a\ Variance\ hVariance\ hγ\ ρ\ *\ Real.sqrt\ h\ +}
\leanline{8}{(Real.sqrt\ Geometry.conditionNumber\ *\ Real.exp\ (-ρ\ *\ (n*h))\ *}
\leanline{10}{optimizerLawDistance\ a\ (targetState\ star)\ μ\ νStar\ +}
\leanline{10}{((2+γ)/ρ)\ *\ Real.sqrt\ (Geometry.conditionNumber\ *\ Variance/(N\ :\ ℝ))\ *}
\leanline{12}{(1-Real.exp\ (-ρ\ *\ (n*h))))}
\end{leancode}
\leansource{\texttt{explicit\_\allowbreak{}discretization\_\allowbreak{}bound} \quad PBGSO/\allowbreak{}FullyUniformMomentum.lean, lines 237–253}
\begin{leancode}
\leanline{0}{theorem\ explicit\_discretization\_bound}
\leanline{4}{(hcloud\ :\ D.cloud\ =\ euclideanParticleEncoding\ N\ J)\ (hγ\ :\ 0\ ≤\ γ)}
\leanline{4}{(star\ :\ A)\ (νStar\ :\ Measure\ (HypocoercivePhase\ (EuclideanSpace\ ℝ\ J)))\ [IsProbabilityMeasure\ νStar]}
\leanline{4}{(hνStar\ :\ MemLp\ (fun\ z\ :\ HypocoercivePhase\ (EuclideanSpace\ ℝ\ J)\ ↦\ z)\ 2\ νStar)}
\leanline{4}{(Frozen\ :\ AdditiveStationary.FrozenTrajectory\ (ULA.CanonicalBrownian.vectorProbability\ J)}
\leanline{6}{(momentumForceDynamics\ D\ γ)\ νStar\ hνStar\ (targetState\ star)\ (input\ γ))}
\leanline{4}{(ρ\ :\ ℝ)\ (hρ\ :\ 0\ <\ ρ)\ (hmono\ :\ TwistedJointMonotonicity\ (model\ D\ γ)\ Geometry\ ρ)}
\leanline{4}{(hg\ :\ (model\ D\ γ).meanGradient\ νStar\ (targetState\ star)\ =\ 0)}
\leanline{4}{(hV\ :\ 0\ ≤\ Variance)}
\leanline{4}{(h\ :\ ℝ)\ (hh\ :\ 0\ <\ h)\ (hhmax\ :\ h\ ≤\ maximumStep\ D\ γ\ Geometry\ μ\ hγ\ ρ)\ (n\ :\ ℕ)\ :}
\leanline{0}{‖momentumParticleIterate\ (ParticleNoiseSpace.probability\ (J\ :=\ J)\ μ\ N)}
\leanline{10}{(particleBrownian\ (J\ :=\ J)\ μ\ N)\ (particleBrownian\_isPreBrownian\ (J\ :=\ J)\ μ\ N)}
\leanline{10}{D\ γ\ h\ (initial\ D\ γ\ Geometry\ μ\ hμ\ a\ Variance\ hVariance)\ n\ -}
\leanline{10}{momentumUnpackLp\ (ParticleNoiseSpace.probability\ (J\ :=\ J)\ μ\ N)\ D}
\leanline{12}{((particles\ D\ γ\ Geometry\ μ\ hμ\ a\ Variance\ hVariance).state\ (n*h))‖\ ≤\ discretizationConstant\ D\ γ\ Geometry\ μ\ hμ\ a\ Variance\ hVariance\ hγ\ ρ\ *\ Real.sqrt\ h}
\end{leancode}
\end{leanbox}
\subsection{\texorpdfstring{Theorem New B — HOL-MYPGD final discretisation bound}{Theorem New B — HOL-MYPGD final discretisation bound}}


\textbf{Assumptions:} N1–N3 and the same cloud encoding. Use the paper’s parameter-first endpoint-force update, its printed coefficients, and the correlated Gaussian innovation triple, with independent time/\allowbreak{}spatial blocks.

Extra assumptions: N1 and N3; \(\alpha>0\), \(\gamma\ge0\), \(0<h\le h_{\star,H}\).

Initialise the numerical and continuous interacting systems identically.





For every iteration \(n\in\mathbb N\), with all constants independent of time and step size,





\[\|\hat\theta_n-\theta_\star\|_{\mathrm L^2}\le D_H\sqrt h+\sqrt\kappa\,e^{-\rho nh}\mathbf d_0+\frac1\rho\sqrt{\frac{\kappa V_F}{N}}(1-e^{-\rho nh}).\]







The full-state raw maximum-component \(L^2\) error has the same uniform discretisation coefficient. This is the stationary zero-observation target of the forces actually implemented; a smoothed sampler may have a different target from the original unsmoothed objective. No uniformity in particle count is claimed.

\begin{leanbox}{Lean: \texttt{explicit\_\allowbreak{}parameter\_\allowbreak{}bound}, \texttt{explicit\_\allowbreak{}discretization\_\allowbreak{}bound}}
\leansource{\texttt{explicit\_\allowbreak{}parameter\_\allowbreak{}bound} \quad PBGSO/\allowbreak{}FullyUniformHigherOrder.lean, lines 261–312}
\begin{leancode}
\leanline{0}{theorem\ explicit\_parameter\_bound\ [FiniteDimensional\ ℝ\ A]}
\leanline{4}{(hcloud\ :\ D.cloud\ =\ euclideanParticleEncoding\ N\ J)}
\leanline{4}{(hα\ :\ 0\ <\ α)\ (hγ\ :\ 0\ ≤\ γ)}
\leanline{4}{(νStar\ :\ Measure\ (HigherOrderSampler\ (EuclideanSpace\ ℝ\ J)))\ [IsProbabilityMeasure\ νStar]}
\leanline{4}{(hνStar\ :\ MemLp\ (fun\ z\ :\ HigherOrderSampler\ (EuclideanSpace\ ℝ\ J)\ ↦\ z)\ 2\ νStar)}
\leanline{4}{(aStar\ :\ A)}
\leanline{4}{(Frozen\ :\ AdditiveStationary.FrozenTrajectory\ (ULA.CanonicalBrownian.vectorProbability\ J)}
\leanline{6}{(higherOrderForceDynamics\ D\ α\ γ)\ νStar\ hνStar\ aStar\ (input\ α))}
\leanline{4}{(ρ\ :\ ℝ)\ (hρ\ :\ 0\ <\ ρ)\ (hmono\ :\ TwistedJointMonotonicity\ (sourceModel\ D\ α\ γ)\ Geo\ ρ)}
\leanline{4}{(hg\ :\ (sourceModel\ D\ α\ γ).meanGradient\ νStar\ aStar\ =\ 0)}
\leanline{4}{(hV\ :\ 0\ ≤\ Variance)}
\leanline{4}{(h\ :\ ℝ≥0)\ (n\ :\ ℕ)\ (hh\ :\ 0\ <\ (h\ :\ ℝ))}
\leanline{4}{(hhstep\ :\ (h\ :\ ℝ)\ ≤\ maximumStep\ D\ α\ γ\ Geo\ μ\ ρ)\ :}
\leanline{8}{‖(higherOrderPullbackParticleIterate\ (Fin\ N\ ×\ J)}
\leanline{10}{(ParticleNoiseSpace.probability\ (J\ :=\ J)\ μ\ N)}
\leanline{10}{(ParticleNoiseSpace.noise\_preserving\ (J\ :=\ J)\ μ\ N)\ D\ α\ γ\ h}
\leanline{10}{(initial\ D\ α\ γ\ Geo\ μ\ hμ\ a\ Variance\ hVariance)\ n).1\ -}
\leanline{10}{Lp.const\ 2\ (ParticleNoiseSpace.probability\ (J\ :=\ J)\ μ\ N)\ aStar‖\ ≤}
\leanline{8}{discretizationConstant\ D\ α\ γ\ Geo\ μ\ hμ\ a\ Variance\ hVariance\ ρ\ *\ Real.sqrt\ (h\ :\ ℝ)\ +}
\leanline{8}{(Real.sqrt\ Geo.conditionNumber\ *\ Real.exp\ (-ρ\ *\ (n\ *\ (h\ :\ ℝ)))\ *}
\leanline{10}{optimizerLawDistance\ a\ aStar\ μ\ νStar\ +}
\leanline{10}{(1\ /\ ρ)\ *\ Real.sqrt\ (Geo.conditionNumber\ *\ Variance\ /\ (N\ :\ ℝ))\ *}
\leanline{12}{(1\ -\ Real.exp\ (-ρ\ *\ (n\ *\ (h\ :\ ℝ)))))}
\end{leancode}
\leansource{\texttt{explicit\_\allowbreak{}discretization\_\allowbreak{}bound} \quad PBGSO/\allowbreak{}FullyUniformHigherOrder.lean, lines 226–240}
\begin{leancode}
\leanline{0}{theorem\ explicit\_discretization\_bound\ [FiniteDimensional\ ℝ\ A]}
\leanline{4}{(hcloud\ :\ D.cloud\ =\ euclideanParticleEncoding\ N\ J)}
\leanline{4}{(hα\ :\ 0\ <\ α)\ (hγ\ :\ 0\ ≤\ γ)}
\leanline{4}{(ρ\ :\ ℝ)\ (hρ\ :\ 0\ <\ ρ)\ (hmono\ :\ TwistedJointMonotonicity\ (sourceModel\ D\ α\ γ)\ Geo\ ρ)}
\leanline{4}{(h\ :\ ℝ≥0)\ (n\ :\ ℕ)\ (hh\ :\ 0\ <\ (h\ :\ ℝ))}
\leanline{4}{(hhmax\ :\ (h\ :\ ℝ)\ ≤\ maximumStep\ D\ α\ γ\ Geo\ μ\ ρ)\ :}
\leanline{8}{‖higherOrderPullbackParticleIterate\ (Fin\ N\ ×\ J)}
\leanline{10}{(ParticleNoiseSpace.probability\ (J\ :=\ J)\ μ\ N)}
\leanline{10}{(ParticleNoiseSpace.noise\_preserving\ (J\ :=\ J)\ μ\ N)\ D\ α\ γ\ h}
\leanline{10}{(initial\ D\ α\ γ\ Geo\ μ\ hμ\ a\ Variance\ hVariance)\ n\ -}
\leanline{10}{higherOrderUnpackLp\ (ParticleNoiseSpace.probability\ (J\ :=\ J)\ μ\ N)\ D}
\leanline{12}{((particle\ D\ α\ γ\ Geo\ μ\ hμ\ a\ Variance\ hVariance).state\ (n\ *\ (h\ :\ ℝ)))‖\ ≤}
\leanline{10}{discretizationConstant\ D\ α\ γ\ Geo\ μ\ hμ\ a\ Variance\ hVariance\ ρ\ *\ Real.sqrt\ (h\ :\ ℝ)}
\end{leancode}
\end{leanbox}
\section{Numerical setting and explicit constants}
\subsection{\texorpdfstring{N1 — Extra assumption: global force Lipschitz bounds}{N1 — Extra assumption: global force Lipschitz bounds}}



Let \(F_x(\theta,x)\) denote the sampler force: in the paper it is \(\nabla_xE(\theta,x)\) for M-MFPO and \(\nabla_xE^\lambda(\theta,x)\) for HOL-MYPGD. This additional notation avoids reusing the paper’s \(G_\nu\), which denotes the population gradient. The verified theorem also permits independent, nonconservative forces \(F,F_x\).




On the paper’s finite-dimensional optimiser and Euclidean sampler spaces, take \(\ell_F,\ell_x\ge0\), with





\[\|F(\theta,x)-F(\tilde\theta,\tilde x)\|\le\ell_F(\|\theta-\tilde\theta\|+\|x-\tilde x\|),\]







\[\|F_x(\theta,x)-F_x(\tilde\theta,\tilde x)\|\le\ell_x(\|\theta-\tilde\theta\|+\|x-\tilde x\|).\]



\begin{leanbox}{Lean: \texttt{ParticleForceData}}
\leansource{\texttt{ParticleForceData} \quad PBGSO/\allowbreak{}ParticleMethodLp.lean, lines 136–149}
\begin{leancode}
\leanline{0}{structure\ ParticleForceData\ (A\ X\ C\ :\ Type*)\ (N\ :\ ℕ)}
\leanline{4}{[NormedAddCommGroup\ A]\ [NormedSpace\ ℝ\ A]}
\leanline{4}{[NormedAddCommGroup\ X]\ [NormedSpace\ ℝ\ X]}
\leanline{4}{[NormedAddCommGroup\ C]\ [NormedSpace\ ℝ\ C]\ where}
\leanline{2}{nonempty\ :\ 0\ <\ N}
\leanline{2}{F\ :\ A\ →\ X\ →\ A}
\leanline{2}{G\ :\ A\ →\ X\ →\ X}
\leanline{2}{parameterLip\ :\ ℝ≥0}
\leanline{2}{samplerLip\ :\ ℝ≥0}
\leanline{2}{parameter\_bound\ :\ ∀\ a\ b\ x\ y,}
\leanline{4}{‖F\ a\ x\ -\ F\ b\ y‖\ ≤\ (parameterLip\ :\ ℝ)\ *\ (‖a\ -\ b‖\ +\ ‖x\ -\ y‖)}
\leanline{2}{sampler\_bound\ :\ ∀\ a\ b\ x\ y,}
\leanline{4}{‖G\ a\ x\ -\ G\ b\ y‖\ ≤\ (samplerLip\ :\ ℝ)\ *\ (‖a\ -\ b‖\ +\ ‖x\ -\ y‖)}
\leanline{2}{cloud\ :\ C\ ≃L[ℝ]\ (Fin\ N\ →\ X)}
\end{leancode}
\leanlabel{How the Lean proof follows the argument}
Force values may grow linearly. No second derivative, Hessian bound or common potential is required. In Lean, \texttt{D.\allowbreak{}F} is the observation and \texttt{D.\allowbreak{}G} is the sampler force.
\end{leanbox}
\subsection{\texorpdfstring{N2 — Initial laws, contraction and reference variance}{N2 — Initial laws, contraction and reference variance}}



Take \(N\ge1\), deterministic initial optimiser state, and i.i.d. initial sampler states with common law \(\nu_0\in\mathcal P_2\), independently of the independent Brownian drivers. Assume A3 with \(\rho>0\), A5 with \(V_F\ge0\), and target zero mean observation \(G_{\nu_{\theta_\star}}(\vartheta_\star)=0\) for the chosen complete-state model.




Use the physical noise from the paper: \(\sqrt{2\gamma}\,\mathrm dB^i\) for M-MFPO, and \(\sqrt{2\alpha}\,\mathrm dB^i\) in the final HOL auxiliary coordinate. The optimiser receives no thermal noise.

\begin{leanbox}{Lean: \texttt{explicit\_\allowbreak{}parameter\_\allowbreak{}bound}, \texttt{explicit\_\allowbreak{}parameter\_\allowbreak{}bound}, \texttt{explicit\_\allowbreak{}parameter\_\allowbreak{}bound} and 1 more}
\leansource{\texttt{explicit\_\allowbreak{}parameter\_\allowbreak{}bound} \quad PBGSO/\allowbreak{}FullyUniformMomentum.lean, lines 256–276}
\begin{leancode}
\leanline{0}{theorem\ explicit\_parameter\_bound}
\leanline{4}{(hcloud\ :\ D.cloud\ =\ euclideanParticleEncoding\ N\ J)\ (hγ\ :\ 0\ ≤\ γ)}
\leanline{4}{(star\ :\ A)\ (νStar\ :\ Measure\ (HypocoercivePhase\ (EuclideanSpace\ ℝ\ J)))\ [IsProbabilityMeasure\ νStar]}
\leanline{4}{(hνStar\ :\ MemLp\ (fun\ z\ :\ HypocoercivePhase\ (EuclideanSpace\ ℝ\ J)\ ↦\ z)\ 2\ νStar)}
\leanline{4}{(Frozen\ :\ AdditiveStationary.FrozenTrajectory\ (ULA.CanonicalBrownian.vectorProbability\ J)}
\leanline{6}{(momentumForceDynamics\ D\ γ)\ νStar\ hνStar\ (targetState\ star)\ (input\ γ))}
\leanline{4}{(ρ\ :\ ℝ)\ (hρ\ :\ 0\ <\ ρ)\ (hmono\ :\ TwistedJointMonotonicity\ (model\ D\ γ)\ Geometry\ ρ)}
\leanline{4}{(hg\ :\ (model\ D\ γ).meanGradient\ νStar\ (targetState\ star)\ =\ 0)}
\leanline{4}{(hV\ :\ 0\ ≤\ Variance)}
\leanline{4}{(h\ :\ ℝ)\ (hh\ :\ 0\ <\ h)\ (hhmax\ :\ h\ ≤\ maximumStep\ D\ γ\ Geometry\ μ\ hγ\ ρ)\ (n\ :\ ℕ)\ :}
\leanline{8}{‖(momentumParticleIterate\ (ParticleNoiseSpace.probability\ (J\ :=\ J)\ μ\ N)}
\leanline{10}{(particleBrownian\ (J\ :=\ J)\ μ\ N)\ (particleBrownian\_isPreBrownian\ (J\ :=\ J)\ μ\ N)}
\leanline{10}{D\ γ\ h\ (initial\ D\ γ\ Geometry\ μ\ hμ\ a\ Variance\ hVariance)\ n).2.1\ -}
\leanline{10}{Lp.const\ 2\ (ParticleNoiseSpace.probability\ (J\ :=\ J)\ μ\ N)\ star‖\ ≤}
\leanline{8}{discretizationConstant\ D\ γ\ Geometry\ μ\ hμ\ a\ Variance\ hVariance\ hγ\ ρ\ *\ Real.sqrt\ h\ +}
\leanline{8}{(Real.sqrt\ Geometry.conditionNumber\ *\ Real.exp\ (-ρ\ *\ (n*h))\ *}
\leanline{10}{optimizerLawDistance\ a\ (targetState\ star)\ μ\ νStar\ +}
\leanline{10}{((2+γ)/ρ)\ *\ Real.sqrt\ (Geometry.conditionNumber\ *\ Variance/(N\ :\ ℝ))\ *}
\leanline{12}{(1-Real.exp\ (-ρ\ *\ (n*h))))}
\end{leancode}
\leansource{\texttt{explicit\_\allowbreak{}parameter\_\allowbreak{}bound} \quad PBGSO/\allowbreak{}FullyUniformMomentum.lean, lines 256–276}
\begin{leancode}
\leanline{0}{theorem\ explicit\_parameter\_bound}
\leanline{4}{(hcloud\ :\ D.cloud\ =\ euclideanParticleEncoding\ N\ J)\ (hγ\ :\ 0\ ≤\ γ)}
\leanline{4}{(star\ :\ A)\ (νStar\ :\ Measure\ (HypocoercivePhase\ (EuclideanSpace\ ℝ\ J)))\ [IsProbabilityMeasure\ νStar]}
\leanline{4}{(hνStar\ :\ MemLp\ (fun\ z\ :\ HypocoercivePhase\ (EuclideanSpace\ ℝ\ J)\ ↦\ z)\ 2\ νStar)}
\leanline{4}{(Frozen\ :\ AdditiveStationary.FrozenTrajectory\ (ULA.CanonicalBrownian.vectorProbability\ J)}
\leanline{6}{(momentumForceDynamics\ D\ γ)\ νStar\ hνStar\ (targetState\ star)\ (input\ γ))}
\leanline{4}{(ρ\ :\ ℝ)\ (hρ\ :\ 0\ <\ ρ)\ (hmono\ :\ TwistedJointMonotonicity\ (model\ D\ γ)\ Geometry\ ρ)}
\leanline{4}{(hg\ :\ (model\ D\ γ).meanGradient\ νStar\ (targetState\ star)\ =\ 0)}
\leanline{4}{(hV\ :\ 0\ ≤\ Variance)}
\leanline{4}{(h\ :\ ℝ)\ (hh\ :\ 0\ <\ h)\ (hhmax\ :\ h\ ≤\ maximumStep\ D\ γ\ Geometry\ μ\ hγ\ ρ)\ (n\ :\ ℕ)\ :}
\leanline{8}{‖(momentumParticleIterate\ (ParticleNoiseSpace.probability\ (J\ :=\ J)\ μ\ N)}
\leanline{10}{(particleBrownian\ (J\ :=\ J)\ μ\ N)\ (particleBrownian\_isPreBrownian\ (J\ :=\ J)\ μ\ N)}
\leanline{10}{D\ γ\ h\ (initial\ D\ γ\ Geometry\ μ\ hμ\ a\ Variance\ hVariance)\ n).2.1\ -}
\leanline{10}{Lp.const\ 2\ (ParticleNoiseSpace.probability\ (J\ :=\ J)\ μ\ N)\ star‖\ ≤}
\leanline{8}{discretizationConstant\ D\ γ\ Geometry\ μ\ hμ\ a\ Variance\ hVariance\ hγ\ ρ\ *\ Real.sqrt\ h\ +}
\leanline{8}{(Real.sqrt\ Geometry.conditionNumber\ *\ Real.exp\ (-ρ\ *\ (n*h))\ *}
\leanline{10}{optimizerLawDistance\ a\ (targetState\ star)\ μ\ νStar\ +}
\leanline{10}{((2+γ)/ρ)\ *\ Real.sqrt\ (Geometry.conditionNumber\ *\ Variance/(N\ :\ ℝ))\ *}
\leanline{12}{(1-Real.exp\ (-ρ\ *\ (n*h))))}
\end{leancode}
\leansource{\texttt{explicit\_\allowbreak{}parameter\_\allowbreak{}bound} \quad PBGSO/\allowbreak{}FullyUniformMomentum.lean, lines 256–276}
\begin{leancode}
\leanline{0}{theorem\ explicit\_parameter\_bound}
\leanline{4}{(hcloud\ :\ D.cloud\ =\ euclideanParticleEncoding\ N\ J)\ (hγ\ :\ 0\ ≤\ γ)}
\leanline{4}{(star\ :\ A)\ (νStar\ :\ Measure\ (HypocoercivePhase\ (EuclideanSpace\ ℝ\ J)))\ [IsProbabilityMeasure\ νStar]}
\leanline{4}{(hνStar\ :\ MemLp\ (fun\ z\ :\ HypocoercivePhase\ (EuclideanSpace\ ℝ\ J)\ ↦\ z)\ 2\ νStar)}
\leanline{4}{(Frozen\ :\ AdditiveStationary.FrozenTrajectory\ (ULA.CanonicalBrownian.vectorProbability\ J)}
\leanline{6}{(momentumForceDynamics\ D\ γ)\ νStar\ hνStar\ (targetState\ star)\ (input\ γ))}
\leanline{4}{(ρ\ :\ ℝ)\ (hρ\ :\ 0\ <\ ρ)\ (hmono\ :\ TwistedJointMonotonicity\ (model\ D\ γ)\ Geometry\ ρ)}
\leanline{4}{(hg\ :\ (model\ D\ γ).meanGradient\ νStar\ (targetState\ star)\ =\ 0)}
\leanline{4}{(hV\ :\ 0\ ≤\ Variance)}
\leanline{4}{(h\ :\ ℝ)\ (hh\ :\ 0\ <\ h)\ (hhmax\ :\ h\ ≤\ maximumStep\ D\ γ\ Geometry\ μ\ hγ\ ρ)\ (n\ :\ ℕ)\ :}
\leanline{8}{‖(momentumParticleIterate\ (ParticleNoiseSpace.probability\ (J\ :=\ J)\ μ\ N)}
\leanline{10}{(particleBrownian\ (J\ :=\ J)\ μ\ N)\ (particleBrownian\_isPreBrownian\ (J\ :=\ J)\ μ\ N)}
\leanline{10}{D\ γ\ h\ (initial\ D\ γ\ Geometry\ μ\ hμ\ a\ Variance\ hVariance)\ n).2.1\ -}
\leanline{10}{Lp.const\ 2\ (ParticleNoiseSpace.probability\ (J\ :=\ J)\ μ\ N)\ star‖\ ≤}
\leanline{8}{discretizationConstant\ D\ γ\ Geometry\ μ\ hμ\ a\ Variance\ hVariance\ hγ\ ρ\ *\ Real.sqrt\ h\ +}
\leanline{8}{(Real.sqrt\ Geometry.conditionNumber\ *\ Real.exp\ (-ρ\ *\ (n*h))\ *}
\leanline{10}{optimizerLawDistance\ a\ (targetState\ star)\ μ\ νStar\ +}
\leanline{10}{((2+γ)/ρ)\ *\ Real.sqrt\ (Geometry.conditionNumber\ *\ Variance/(N\ :\ ℝ))\ *}
\leanline{12}{(1-Real.exp\ (-ρ\ *\ (n*h))))}
\end{leancode}
\leansource{\texttt{explicit\_\allowbreak{}parameter\_\allowbreak{}bound} \quad PBGSO/\allowbreak{}FullyUniformMomentum.lean, lines 256–276}
\begin{leancode}
\leanline{0}{theorem\ explicit\_parameter\_bound}
\leanline{4}{(hcloud\ :\ D.cloud\ =\ euclideanParticleEncoding\ N\ J)\ (hγ\ :\ 0\ ≤\ γ)}
\leanline{4}{(star\ :\ A)\ (νStar\ :\ Measure\ (HypocoercivePhase\ (EuclideanSpace\ ℝ\ J)))\ [IsProbabilityMeasure\ νStar]}
\leanline{4}{(hνStar\ :\ MemLp\ (fun\ z\ :\ HypocoercivePhase\ (EuclideanSpace\ ℝ\ J)\ ↦\ z)\ 2\ νStar)}
\leanline{4}{(Frozen\ :\ AdditiveStationary.FrozenTrajectory\ (ULA.CanonicalBrownian.vectorProbability\ J)}
\leanline{6}{(momentumForceDynamics\ D\ γ)\ νStar\ hνStar\ (targetState\ star)\ (input\ γ))}
\leanline{4}{(ρ\ :\ ℝ)\ (hρ\ :\ 0\ <\ ρ)\ (hmono\ :\ TwistedJointMonotonicity\ (model\ D\ γ)\ Geometry\ ρ)}
\leanline{4}{(hg\ :\ (model\ D\ γ).meanGradient\ νStar\ (targetState\ star)\ =\ 0)}
\leanline{4}{(hV\ :\ 0\ ≤\ Variance)}
\leanline{4}{(h\ :\ ℝ)\ (hh\ :\ 0\ <\ h)\ (hhmax\ :\ h\ ≤\ maximumStep\ D\ γ\ Geometry\ μ\ hγ\ ρ)\ (n\ :\ ℕ)\ :}
\leanline{8}{‖(momentumParticleIterate\ (ParticleNoiseSpace.probability\ (J\ :=\ J)\ μ\ N)}
\leanline{10}{(particleBrownian\ (J\ :=\ J)\ μ\ N)\ (particleBrownian\_isPreBrownian\ (J\ :=\ J)\ μ\ N)}
\leanline{10}{D\ γ\ h\ (initial\ D\ γ\ Geometry\ μ\ hμ\ a\ Variance\ hVariance)\ n).2.1\ -}
\leanline{10}{Lp.const\ 2\ (ParticleNoiseSpace.probability\ (J\ :=\ J)\ μ\ N)\ star‖\ ≤}
\leanline{8}{discretizationConstant\ D\ γ\ Geometry\ μ\ hμ\ a\ Variance\ hVariance\ hγ\ ρ\ *\ Real.sqrt\ h\ +}
\leanline{8}{(Real.sqrt\ Geometry.conditionNumber\ *\ Real.exp\ (-ρ\ *\ (n*h))\ *}
\leanline{10}{optimizerLawDistance\ a\ (targetState\ star)\ μ\ νStar\ +}
\leanline{10}{((2+γ)/ρ)\ *\ Real.sqrt\ (Geometry.conditionNumber\ *\ Variance/(N\ :\ ℝ))\ *}
\leanline{12}{(1-Real.exp\ (-ρ\ *\ (n*h))))}
\end{leancode}
\end{leanbox}
\subsection{\texorpdfstring{N3 — Extra assumption: canonical frozen-target realization}{N3 — Extra assumption: canonical frozen-target realization}}



For the target pair, require \(\nu_{\theta_\star}\in\mathcal P_2\). On \(\nu_{\theta_\star}\) times the canonical Brownian probability, supply a norm-continuous \(\mathrm L^2\) frozen sampler trajectory, initialized by the initial-coordinate projection, which satisfies





\[Z_t^\star-Z_0^\star=\int_0^t b_{\theta_\star}(Z_r^\star)\,\mathrm dr+\sigma_\star B_t,\qquad\operatorname{Law}(Z_t^\star)=\nu_{\theta_\star}\quad(t\ge0).\]



\begin{leanbox}{Lean: \texttt{FrozenTrajectory}}
\leansource{\texttt{FrozenTrajectory} \quad PBGSO/\allowbreak{}AdditiveStationaryTarget.lean, lines 35–42}
\begin{leancode}
\leanline{0}{structure\ FrozenTrajectory\ where}
\leanline{2}{sampler\ :\ ℝ\ →\ Lp\ Z\ 2\ (νStar.prod\ R)}
\leanline{2}{continuous\ :\ Continuous\ sampler}
\leanline{2}{initial\ :\ sampler\ 0\ =\ AdditiveCanonical.initialSampler\ R\ νStar\ hνStar}
\leanline{2}{equation\ :\ ∀\ t,\ 0\ ≤\ t\ →\ sampler\ t\ -\ sampler\ 0\ =}
\leanline{4}{(∫\ r\ in\ 0..t,\ liftMap\ (νStar.prod\ R)\ (L.sampler\ aStar)\ (L.sampler\_state\ aStar)}
\leanline{6}{(sampler\ r))\ +\ AdditiveCanonical.input\ R\ νStar\ B\ t}
\leanline{2}{stationary\_law\ :\ ∀\ t,\ 0\ ≤\ t\ →\ (νStar.prod\ R).map\ (sampler\ t)\ =\ νStar}
\end{leancode}
\leanlabel{How the Lean proof follows the argument}
The constant physical diffusion is the one in N2. This explicit \texttt{Frozen\allowbreak{}Trajectory} input is stronger than the weak frozen-SDE realization used in Theorem 1. The numerical endpoint retains it as a visible assumption.
\end{leanbox}
\subsection{\texorpdfstring{The two method models and their raw coordinates}{The two method models and their raw coordinates}}



As in the paper’s algorithm section, suppress the population superscript in the method coordinates. For M-MFPO:





\[\begin{aligned}\dot\theta_t&=m_t,\qquad\dot m_t=-N^{-1}\sum_i F(\theta_t,X_t^i)-\gamma m_t,\\\mathrm dX_t^i&=V_t^i\,\mathrm dt,\\\mathrm dV_t^i&=(-F_x(\theta_t,X_t^i)-\gamma V_t^i)\,\mathrm dt+\sqrt{2\gamma}\,\mathrm dB_t^i.\end{aligned}\]





For HOL-MYPGD:





\[\begin{aligned}\dot\theta_t&=-N^{-1}\sum_iF(\theta_t,X_t^i),\qquad\mathrm dX_t^i=U_t^i\,\mathrm dt,\\\mathrm dU_t^i&=(-F_x(\theta_t,X_t^i)+\gamma V_t^i)\,\mathrm dt,\\\mathrm dV_t^i&=(-\gamma U_t^i-\alpha V_t^i)\,\mathrm dt+\sqrt{2\alpha}\,\mathrm dB_t^i.\end{aligned}\]





Write \(y\) for the implementation phase, an additional symbol distinct from the paper’s sampler \(Z\). Its four coordinates are \(((m,V^{[N]}),(\theta,X^{[N]}))\) for M-MFPO and \((\theta,X^{[N]},U^{[N]},V^{[N]})\) for HOL. The exact raw norm is \(\|y\|_{\max,2}=\max_{1\le k\le4}\|y_k\|_{\mathrm L^2}\). The numerical and continuous initial phases coincide.

\begin{leanbox}{Lean: \texttt{heavyBallModel}, \texttt{higherOrderModel}}
\leansource{\texttt{heavyBallModel} \quad PBGSO/\allowbreak{}HypocoerciveCriterion.lean, lines 289–297}
\begin{leancode}
\leanline{0}{def\ heavyBallModel\ (γ\ :\ ℝ)\ (F\ :\ A\ →\ X\ →\ A)\ (H\ :\ A\ →\ X\ →\ X)}
\leanline{4}{(σ\ :\ I\ →\ HypocoercivePhase\ X)\ :}
\leanline{4}{Model\ (HypocoercivePhase\ A)\ A\ (HypocoercivePhase\ X)\ X\ I\ where}
\leanline{2}{parameter\ :=\ WithLp.fst}
\leanline{2}{position\ :=\ WithLp.fst}
\leanline{2}{observation\ :=\ F}
\leanline{2}{optimizer\ :=\ fun\ \_\ s\ g\ =>\ WithLp.toLp\ 2\ (s.snd,\ -g\ -\ γ\ •\ s.snd)}
\leanline{2}{drift\ :=\ fun\ θ\ z\ =>\ WithLp.toLp\ 2\ (z.snd,\ -H\ θ\ z.fst\ -\ γ\ •\ z.snd)}
\leanline{2}{diffusion\ :=\ fun\ \_\ \_\ =>\ σ}
\end{leancode}
\leansource{\texttt{higherOrderModel} \quad PBGSO/\allowbreak{}ImplementationDynamics.lean, lines 174–180}
\begin{leancode}
\leanline{0}{def\ higherOrderModel\ (σ\ :\ I\ →\ X)\ :\ Model\ A\ A\ (HigherOrderSampler\ X)\ X\ I\ where}
\leanline{2}{parameter\ :=\ id}
\leanline{2}{position\ :=\ WithLp.fst}
\leanline{2}{observation\ :=\ F}
\leanline{2}{optimizer\ :=\ fun\ \_\ \_\ g\ =>\ -g}
\leanline{2}{drift\ :=\ higherOrderSamplerField\ G\ α\ γ}
\leanline{2}{diffusion\ :=\ fun\ \_\ \_\ i\ =>\ WithLp.toLp\ 2\ (0,\ WithLp.toLp\ 2\ (0,\ σ\ i))}
\end{leancode}
\end{leanbox}
\subsection{\texorpdfstring{Explicit force and cloud constants}{Explicit force and cloud constants}}



Let \(d=Nd_x\), and let the standard cloud map be the concrete reindexing \(\mathcal E_N:\mathbb R^{Nd_x}\to(\mathbb R^{d_x})^N\), \((\mathcal E_Nx)_i=(x_{i,j})_{j=1}^{d_x}\). Its domain has the Euclidean norm and its target the maximum of particle Euclidean norms. Define the actual finite-cloud maps





\[\bar F_N(\theta,x)=\frac1N\sum_{i=1}^NF(\theta,(\mathcal E_Nx)_i),\qquad\mathcal F_{x,N}(\theta,x)=\mathcal E_N^{-1}\bigl((F_x(\theta,(\mathcal E_Nx)_i))_{i=1}^N\bigr).\]





For the paper’s positive position dimension \(d_x\ge1\), the exact norm identities \(\|\mathcal E_N\|=1\), \(\|\mathcal E_N^{-1}\|=\sqrt N\) give the explicit certified force constants



\[L_F=2\ell_F,\qquad L_x=2\sqrt N\,\ell_x.\]



The origin identities on the right prove that these are exactly the \(L^2\) norms of the lifted mean and sampler forces. In either raw phase, the full drift at zero has norm \(b_\circ\). These formulas specify the coordinate norms, the force values and every input used below; \(G_\nu\) retains its meaning from the paper.

\begin{leanbox}{Lean: \texttt{euclideanParticleEncoding}, \texttt{meanLip\_\allowbreak{}euclidean}, \texttt{samplerForceLip\_\allowbreak{}euclidean} and 6 more}
\leansource{\texttt{euclideanParticleEncoding} \quad PBGSO/\allowbreak{}ParticleMethodLp.lean, lines 395–400}
\begin{leancode}
\leanline{0}{def\ euclideanParticleEncoding\ (N\ :\ ℕ)\ (I\ :\ Type*)\ [Fintype\ I]\ :}
\leanline{4}{EuclideanSpace\ ℝ\ (Fin\ N\ ×\ I)\ ≃L[ℝ]\ (Fin\ N\ →\ EuclideanSpace\ ℝ\ I)\ :=}
\leanline{2}{((WithLp.linearEquiv\ 2\ ℝ\ (Fin\ N\ ×\ I\ →\ ℝ)).trans}
\leanline{4}{((LinearEquiv.curry\ ℝ\ ℝ\ (Fin\ N)\ I).trans}
\leanline{6}{(LinearEquiv.piCongrRight\ (fun\ \_\ :\ Fin\ N\ ↦}
\leanline{8}{(WithLp.linearEquiv\ 2\ ℝ\ (I\ →\ ℝ)).symm)))).toContinuousLinearEquiv}
\end{leancode}
\leansource{\texttt{meanLip\_\allowbreak{}euclidean} \quad PBGSO/\allowbreak{}CloudEncodingConstants.lean, lines 113–118}
\begin{leancode}
\leanline{0}{theorem\ meanLip\_euclidean\ [Nonempty\ J]}
\leanline{4}{(hcloud\ :\ D.cloud\ =\ euclideanParticleEncoding\ N\ J)\ :}
\leanline{4}{(D.meanLip\ :\ ℝ)\ =\ 2*(D.parameterLip\ :\ ℝ)}
\end{leancode}
\leansource{\texttt{samplerForceLip\_\allowbreak{}euclidean} \quad PBGSO/\allowbreak{}CloudEncodingConstants.lean, lines 121–128}
\begin{leancode}
\leanline{0}{theorem\ samplerForceLip\_euclidean\ [Nonempty\ J]}
\leanline{4}{(hcloud\ :\ D.cloud\ =\ euclideanParticleEncoding\ N\ J)\ :}
\leanline{4}{(D.samplerForceLip\ :\ ℝ)\ =\ 2*Real.sqrt\ (N:ℝ)*(D.samplerLip\ :\ ℝ)}
\end{leancode}
\leansource{\texttt{meanLip\_\allowbreak{}euclidean\_\allowbreak{}zero} \quad PBGSO/\allowbreak{}CloudEncodingConstants.lean, lines 131–136}
\begin{leancode}
\leanline{0}{theorem\ meanLip\_euclidean\_zero\ [IsEmpty\ J]}
\leanline{4}{(hcloud\ :\ D.cloud\ =\ euclideanParticleEncoding\ N\ J)\ :}
\leanline{4}{(D.meanLip\ :\ ℝ)\ =\ (D.parameterLip\ :\ ℝ)}
\end{leancode}
\leansource{\texttt{samplerForceLip\_\allowbreak{}euclidean\_\allowbreak{}zero} \quad PBGSO/\allowbreak{}CloudEncodingConstants.lean, lines 139–145}
\begin{leancode}
\leanline{0}{theorem\ samplerForceLip\_euclidean\_zero\ [IsEmpty\ J]}
\leanline{4}{(hcloud\ :\ D.cloud\ =\ euclideanParticleEncoding\ N\ J)\ :}
\leanline{4}{(D.samplerForceLip\ :\ ℝ)\ =\ 0}
\end{leancode}
\leansource{\texttt{meanForceLp\_\allowbreak{}zero\_\allowbreak{}norm} \quad PBGSO/\allowbreak{}UniformOriginConstants.lean, lines 55–57}
\begin{leancode}
\leanline{0}{theorem\ meanForceLp\_zero\_norm\ :\ ‖D.meanForceLp\ Q\ 0\ 0‖\ =\ ‖D.F\ 0\ 0‖}
\end{leancode}
\leansource{\texttt{samplerForceLp\_\allowbreak{}zero\_\allowbreak{}norm} \quad PBGSO/\allowbreak{}UniformOriginConstants.lean, lines 60–63}
\begin{leancode}
\leanline{0}{theorem\ samplerForceLp\_zero\_norm\ (hcloud\ :\ D.cloud\ =\ euclideanParticleEncoding\ N\ J)\ :}
\leanline{4}{‖D.samplerForceLp\ Q\ 0\ 0‖\ =\ Real.sqrt\ (N\ :\ ℝ)\ *\ ‖D.G\ 0\ 0‖}
\end{leancode}
\leansource{\texttt{coupledForce\_\allowbreak{}zero\_\allowbreak{}norm} \quad PBGSO/\allowbreak{}UniformOriginConstants.lean, lines 66–70}
\begin{leancode}
\leanline{0}{theorem\ coupledForce\_zero\_norm\ (hcloud\ :\ D.cloud\ =\ euclideanParticleEncoding\ N\ J)\ :}
\leanline{4}{‖coupledForce\ (D.meanForceLp\ Q)\ (D.samplerForceLp\ Q)\ 0‖\ =}
\leanline{6}{max\ ‖D.F\ 0\ 0‖\ (Real.sqrt\ (N\ :\ ℝ)\ *\ ‖D.G\ 0\ 0‖)}
\end{leancode}
\leansource{\texttt{holDrift\_\allowbreak{}zero\_\allowbreak{}norm} \quad PBGSO/\allowbreak{}UniformOriginConstants.lean, lines 73–80}
\begin{leancode}
\leanline{0}{theorem\ holDrift\_zero\_norm\ (hcloud\ :\ D.cloud\ =\ euclideanParticleEncoding\ N\ J)\ (α\ γ\ :\ ℝ)\ :}
\leanline{4}{‖holDrift\ (D.meanForceLp\ Q)\ (D.samplerForceLp\ Q)\ α\ γ\ 0‖\ =}
\leanline{6}{max\ ‖D.F\ 0\ 0‖\ (Real.sqrt\ (N\ :\ ℝ)\ *\ ‖D.G\ 0\ 0‖)}
\end{leancode}
\leanlabel{How the Lean proof follows the argument}
The formal development also permits \(d_x=0\); then use \(L_F=\ell_F\), \(L_x=0\). Both cases are proved in Lean, so no positive-dimension assumption is hidden.





\[f_0=\|F(0,0)\|,\qquad g_0=\sqrt N\,\|F_x(0,0)\|,\qquad b_\circ=\max\{f_0,g_0\}.\]
\end{leanbox}
\subsection{\texorpdfstring{New A — Fully expanded local coefficients}{New A — Fully expanded local coefficients}}




\[L_M=2\max\{L_F,L_x\},\qquad K_M=L_M+\gamma+1,\qquad\Xi_M=\sqrt d,\qquad\beta_M=\sqrt\gamma\sqrt{2d}.\]







\[T_M=b_\circ+K_Me^{K_M+1}(b_\circ+\beta_M)+\beta_M.\]







\[c_{0,M}=\left(\frac{\gamma+1}{2}+\frac{L_M}{4}\right)b_\circ+\left(\frac{L_M}{2}+\sqrt\gamma+\frac32\gamma\sqrt\gamma\right)\sqrt d+\frac23K_MT_M.\]







\[c_{1,M}=\frac{\gamma^2+\gamma+L_M}{2}+\left(\frac{\gamma+1}{2}+\frac{L_M}{4}\right)L_M+\frac23K_M^2e^{K_M+1}.\]



\begin{leanbox}{Lean: \texttt{momentumExpandedConstant0}, \texttt{momentumExpandedConstant1}, \texttt{momentumUniformDefect\_\allowbreak{}expanded}}
\leansource{\texttt{momentumExpandedConstant0} \quad PBGSO/\allowbreak{}UniformConstantExpansions.lean, lines 15–19}
\begin{leancode}
\leanline{0}{def\ momentumExpandedConstant0\ (L\ γ\ fzero\ Ξ\ β\ :\ ℝ)\ :\ ℝ\ :=}
\leanline{2}{((γ\ +\ 1)\ /\ 2\ +\ L\ /\ 4)\ *\ fzero\ +}
\leanline{4}{(L\ /\ 2\ +\ Real.sqrt\ γ\ +\ (3\ /\ 2\ :\ ℝ)\ *\ γ\ *\ Real.sqrt\ γ)\ *\ Ξ\ +}
\leanline{4}{(2\ /\ 3\ :\ ℝ)\ *\ (L\ +\ γ\ +\ 1)\ *}
\leanline{6}{(fzero\ +\ (L\ +\ γ\ +\ 1)\ *\ Real.exp\ (L\ +\ γ\ +\ 2)\ *\ (fzero\ +\ β)\ +\ β)}
\end{leancode}
\leansource{\texttt{momentumExpandedConstant1} \quad PBGSO/\allowbreak{}UniformConstantExpansions.lean, lines 22–24}
\begin{leancode}
\leanline{0}{def\ momentumExpandedConstant1\ (L\ γ\ :\ ℝ)\ :\ ℝ\ :=}
\leanline{2}{(γ\ \textasciicircum{}\ 2\ +\ γ\ +\ L)\ /\ 2\ +\ ((γ\ +\ 1)\ /\ 2\ +\ L\ /\ 4)\ *\ L\ +}
\leanline{4}{(2\ /\ 3\ :\ ℝ)\ *\ (L\ +\ γ\ +\ 1)\ \textasciicircum{}\ 2\ *\ Real.exp\ (L\ +\ γ\ +\ 2)}
\end{leancode}
\leansource{\texttt{momentumUniformDefect\_\allowbreak{}expanded} \quad PBGSO/\allowbreak{}UniformConstantExpansions.lean, lines 26–34}
\begin{leancode}
\leanline{0}{theorem\ momentumUniformDefect\_expanded\ (L\ γ\ :\ ℝ≥0)\ (fzero\ Ξ\ β\ :\ ℝ)\ :}
\leanline{4}{momentumUniformDefect\ L\ γ\ fzero\ Ξ\ β\ =}
\leanline{6}{(momentumExpandedConstant0\ L\ γ\ fzero\ Ξ\ β,}
\leanline{7}{momentumExpandedConstant1\ L\ γ)}
\end{leancode}
\leanlabel{How the Lean proof follows the argument}
These are nonnegative coefficients of the proved local error, with its dependence on the numerical input norm explicit. The Lean equality on the right verifies the arithmetic expansion.
\end{leanbox}
\subsection{\texorpdfstring{New B — Fully expanded local coefficients}{New B — Fully expanded local coefficients}}




\[K_H=4(L_F+L_x+1+2\gamma+\alpha),\qquad\beta_H=\sqrt\alpha\sqrt{2d}.\]





The coefficient-remainder constants below are distinct from the paper’s step-dependent \(\mu_{ij}\):





\[r_{12}=\frac{\gamma^2}{2\alpha},\quad r_{13}=\frac\gamma2,\quad r_{22}=\frac{\gamma^2}{2},\quad r_{23}=\frac{\gamma\alpha}{2},\quad r_{31}=\frac\gamma2,\]







\[r_{32}=\frac{2\gamma^3}{\alpha}+\frac{\gamma^3}{2}+\frac{\gamma\alpha}{2},\qquad r_{33}=\frac{3\gamma^2}{2}+\frac{\alpha^2}{2}+\frac{\alpha\gamma^2}{2}.\]







\[A_H=\frac{1+\gamma}{2}\bigl[g_0+2L_x+L_x(f_0+2L_F)\bigr]+(r_{12}+r_{22}+r_{32})+(r_{13}+r_{23}+r_{33})+L_x(f_0+2L_F),\]







\[N_H=\sqrt{2\alpha}\left(\frac{3\gamma}{2}+\alpha+r_{33}\right),\qquad T_H=b_\circ+K_He^{K_H+1}(b_\circ+\beta_H)+\beta_H,\]







\[c_{0,H}=A_H+\sqrt d\,N_H+\frac83K_HT_H,\qquad c_{1,H}=A_H+\frac83K_H^2e^{K_H+1}.\]





Every intermediate coefficient is defined here. The dimension factor \(\sqrt d\) occurs once in the Gaussian term. The identities on the right certify these expansions for \(\alpha>0\), \(\gamma\ge0\).

\begin{leanbox}{Lean: \texttt{holExpandedRemainder32}, \texttt{holExpandedRemainder33}, \texttt{holExpandedGrowthConstant} and 5 more}
\leansource{\texttt{holExpandedRemainder32} \quad PBGSO/\allowbreak{}UniformConstantExpansions.lean, lines 37–38}
\begin{leancode}
\leanline{0}{def\ holExpandedRemainder32\ (α\ γ\ :\ ℝ)\ :\ ℝ\ :=}
\leanline{2}{2\ *\ γ\ \textasciicircum{}\ 3\ /\ α\ +\ γ\ \textasciicircum{}\ 3\ /\ 2\ +\ γ\ *\ α\ /\ 2}
\end{leancode}
\leansource{\texttt{holExpandedRemainder33} \quad PBGSO/\allowbreak{}UniformConstantExpansions.lean, lines 41–42}
\begin{leancode}
\leanline{0}{def\ holExpandedRemainder33\ (α\ γ\ :\ ℝ)\ :\ ℝ\ :=}
\leanline{2}{3\ *\ γ\ \textasciicircum{}\ 2\ /\ 2\ +\ α\ \textasciicircum{}\ 2\ /\ 2\ +\ α\ *\ γ\ \textasciicircum{}\ 2\ /\ 2}
\end{leancode}
\leansource{\texttt{holExpandedGrowthConstant} \quad PBGSO/\allowbreak{}UniformConstantExpansions.lean, lines 57–62}
\begin{leancode}
\leanline{0}{def\ holExpandedGrowthConstant\ (α\ γ\ LF\ LG\ fzero\ gzero\ :\ ℝ)\ :\ ℝ\ :=}
\leanline{2}{((1\ +\ γ)\ /\ 2)\ *\ (gzero\ +\ 2\ *\ LG\ +\ LG\ *\ (fzero\ +\ 2\ *\ LF))\ +}
\leanline{4}{(γ\ \textasciicircum{}\ 2\ /\ (2\ *\ α)\ +\ γ\ \textasciicircum{}\ 2\ /\ 2\ +}
\leanline{6}{2\ *\ γ\ \textasciicircum{}\ 3\ /\ α\ +\ γ\ \textasciicircum{}\ 3\ /\ 2\ +\ γ\ *\ α\ /\ 2)\ +}
\leanline{4}{(γ\ /\ 2\ +\ γ\ *\ α\ /\ 2\ +\ 3\ *\ γ\ \textasciicircum{}\ 2\ /\ 2\ +\ α\ \textasciicircum{}\ 2\ /\ 2\ +\ α\ *\ γ\ \textasciicircum{}\ 2\ /\ 2)\ +}
\leanline{4}{LG\ *\ (fzero\ +\ 2\ *\ LF)}
\end{leancode}
\leansource{\texttt{holExpandedNoiseConstant} \quad PBGSO/\allowbreak{}UniformConstantExpansions.lean, lines 65–67}
\begin{leancode}
\leanline{0}{def\ holExpandedNoiseConstant\ (α\ γ\ :\ ℝ)\ :\ ℝ\ :=}
\leanline{2}{Real.sqrt\ (2\ *\ α)\ *}
\leanline{4}{(3\ *\ γ\ /\ 2\ +\ α\ +\ 3\ *\ γ\ \textasciicircum{}\ 2\ /\ 2\ +\ α\ \textasciicircum{}\ 2\ /\ 2\ +\ α\ *\ γ\ \textasciicircum{}\ 2\ /\ 2)}
\end{leancode}
\leansource{\texttt{holExpandedConstant0} \quad PBGSO/\allowbreak{}UniformConstantExpansions.lean, lines 97–102}
\begin{leancode}
\leanline{0}{def\ holExpandedConstant0\ (d\ α\ γ\ LF\ LG\ fzero\ gzero\ bzero\ :\ ℝ)\ :\ ℝ\ :=}
\leanline{2}{let\ K\ :=\ 4\ *\ (LF\ +\ LG\ +\ 1\ +\ 2\ *\ γ\ +\ α)}
\leanline{2}{let\ β\ :=\ Real.sqrt\ α\ *\ Real.sqrt\ (2\ *\ d)}
\leanline{2}{holExpandedGrowthConstant\ α\ γ\ LF\ LG\ fzero\ gzero\ +}
\leanline{4}{Real.sqrt\ d\ *\ holExpandedNoiseConstant\ α\ γ\ +}
\leanline{4}{(8\ /\ 3\ :\ ℝ)\ *\ K\ *\ (bzero\ +\ K\ *\ Real.exp\ (K\ +\ 1)\ *\ (bzero\ +\ β)\ +\ β)}
\end{leancode}
\leansource{\texttt{holExpandedConstant1} \quad PBGSO/\allowbreak{}UniformConstantExpansions.lean, lines 105–108}
\begin{leancode}
\leanline{0}{def\ holExpandedConstant1\ (α\ γ\ LF\ LG\ fzero\ gzero\ :\ ℝ)\ :\ ℝ\ :=}
\leanline{2}{let\ K\ :=\ 4\ *\ (LF\ +\ LG\ +\ 1\ +\ 2\ *\ γ\ +\ α)}
\leanline{2}{holExpandedGrowthConstant\ α\ γ\ LF\ LG\ fzero\ gzero\ +}
\leanline{4}{(8\ /\ 3\ :\ ℝ)\ *\ K\ \textasciicircum{}\ 2\ *\ Real.exp\ (K\ +\ 1)}
\end{leancode}
\leansource{\texttt{holUniformLocalConstant0\_\allowbreak{}expanded} \quad PBGSO/\allowbreak{}UniformConstantExpansions.lean, lines 110–119}
\begin{leancode}
\leanline{0}{theorem\ holUniformLocalConstant0\_expanded\ (I\ :\ Type*)\ [Fintype\ I]}
\leanline{4}{(α\ γ\ LF\ LG\ fzero\ gzero\ bzero\ :\ ℝ)}
\leanline{4}{(hα\ :\ 0\ <\ α)\ (hγ\ :\ 0\ ≤\ γ)\ (hLF\ :\ 0\ ≤\ LF)\ (hLG\ :\ 0\ ≤\ LG)\ :}
\leanline{4}{holUniformLocalConstant0\ I\ α\ γ\ LF\ LG\ fzero\ gzero\ bzero\ =}
\leanline{6}{holExpandedConstant0\ (Fintype.card\ I)\ α\ γ\ LF\ LG\ fzero\ gzero\ bzero}
\end{leancode}
\leansource{\texttt{holUniformLocalConstant1\_\allowbreak{}expanded} \quad PBGSO/\allowbreak{}UniformConstantExpansions.lean, lines 121–128}
\begin{leancode}
\leanline{0}{theorem\ holUniformLocalConstant1\_expanded\ (α\ γ\ LF\ LG\ fzero\ gzero\ :\ ℝ)}
\leanline{4}{(hα\ :\ 0\ <\ α)\ (hγ\ :\ 0\ ≤\ γ)\ (hLF\ :\ 0\ ≤\ LF)\ (hLG\ :\ 0\ ≤\ LG)\ :}
\leanline{4}{holUniformLocalConstant1\ α\ γ\ LF\ LG\ fzero\ gzero\ =}
\leanline{6}{holExpandedConstant1\ α\ γ\ LF\ LG\ fzero\ gzero}
\end{leancode}
\end{leanbox}
\subsection{\texorpdfstring{Both methods — Explicit metric comparisons}{Both methods — Explicit metric comparisons}}



For either chosen method, use its source A3 weights and certified \(\underline\lambda,\overline\lambda\). Define





\[\underline d_N=\underline\lambda\min\{1,N^{-1}\},\qquad\overline d_N=\overline\lambda\max\{1,N^{-1}\},\qquad m_N=\frac{\sqrt{\underline d_N}}2,\qquad U_N=2\sqrt{\overline d_N}.\]





Since \(N\ge1\), these are \(m_N=\tfrac12\sqrt{\underline\lambda/N}\) and \(U_N=2\sqrt{\overline\lambda}\). The coordinate comparisons and source A3 construct a linear metric factor \(C\) with \(m_N\|y\|_{\max,2}\le\|Cy\|\le U_N\|y\|_{\max,2}\), together with exact-flow contraction at rate \(\rho\). No unknown coordinate-map norm enters these metric constants.

\begin{leanbox}{Lean: \texttt{metricCoercivity}, \texttt{metricUpper}, \texttt{momentum\_\allowbreak{}metric\_\allowbreak{}explicit} and 1 more}
\leansource{\texttt{metricCoercivity} \quad PBGSO/\allowbreak{}FullyUniformMomentum.lean, lines 45–45}
\begin{leancode}
\leanline{0}{def\ metricCoercivity\ :\ ℝ\ :=\ Real.sqrt\ (EmpiricalMetric.lower\ N\ Geometry)\ /\ 2}
\end{leancode}
\leansource{\texttt{metricUpper} \quad PBGSO/\allowbreak{}FullyUniformMomentum.lean, lines 48–48}
\begin{leancode}
\leanline{0}{def\ metricUpper\ :\ ℝ\ :=\ 2\ *\ Real.sqrt\ (EmpiricalMetric.upper\ N\ Geometry)}
\end{leancode}
\leansource{\texttt{momentum\_\allowbreak{}metric\_\allowbreak{}explicit} \quad PBGSO/\allowbreak{}MethodMetricComparisons.lean, lines 209–236}
\begin{leancode}
\leanline{0}{theorem\ momentum\_metric\_explicit\ (hcloud\ :\ D.cloud\ =\ euclideanParticleEncoding\ N\ J)}
\leanline{4}{(γ\ :\ ℝ)\ (σ\ :\ I\ →\ HypocoercivePhase\ (EuclideanSpace\ ℝ\ J))}
\leanline{4}{(G\ :\ TwistedGeometry\ (HypocoercivePhase\ A)\ (HypocoercivePhase\ (EuclideanSpace\ ℝ\ J)))}
\leanline{4}{(ρ\ :\ ℝ)\ (hmono\ :\ TwistedJointMonotonicity\ (heavyBallModel\ γ\ D.F\ D.G\ σ)\ G\ ρ)\ :}
\leanline{4}{∃\ n\ :\ ℕ,\ ∃\ T\ :\ ((Lp\ A\ 2\ Q\ ×\ Lp\ (EuclideanSpace\ ℝ\ (Fin\ N\ ×\ J))\ 2\ Q)\ ×}
\leanline{6}{(Lp\ A\ 2\ Q\ ×\ Lp\ (EuclideanSpace\ ℝ\ (Fin\ N\ ×\ J))\ 2\ Q))\ →L[ℝ]}
\leanline{8}{Lp\ (EuclideanSpace\ ℝ\ (Fin\ n))\ 2\ Q,}
\leanline{6}{(∀\ x,\ (Real.sqrt\ (EmpiricalMetric.lower\ N\ G)/2)*‖x‖≤‖T\ x‖)\ ∧}
\leanline{6}{(∀\ x,\ ‖T\ x‖≤(2*Real.sqrt\ (EmpiricalMetric.upper\ N\ G))*‖x‖)\ ∧}
\leanline{6}{(∀\ x\ y,\ {\leanfallback ⟪}T\ (x-y),T\ (}
\leanline{8}{(-coupledForce\ (D.meanForceLp\ Q)\ (D.samplerForceLp\ Q)\ x.2-γ\ •\ x.1,x.1)-}
\leanline{8}{(-coupledForce\ (D.meanForceLp\ Q)\ (D.samplerForceLp\ Q)\ y.2-γ\ •\ y.1,y.1)){\leanfallback ⟫}\_ℝ}
\leanline{10}{≤\ -ρ*‖T\ (x-y)‖\textasciicircum{}2)}
\end{leancode}
\leansource{\texttt{higherOrder\_\allowbreak{}metric\_\allowbreak{}explicit} \quad PBGSO/\allowbreak{}MethodMetricComparisons.lean, lines 240–264}
\begin{leancode}
\leanline{0}{theorem\ higherOrder\_metric\_explicit\ (hcloud\ :\ D.cloud\ =\ euclideanParticleEncoding\ N\ J)}
\leanline{4}{(α\ γ\ :\ ℝ)\ (σ\ :\ I\ →\ EuclideanSpace\ ℝ\ J)}
\leanline{4}{(G\ :\ TwistedGeometry\ A\ (HigherOrderSampler\ (EuclideanSpace\ ℝ\ J)))\ (ρ\ :\ ℝ)}
\leanline{4}{(hmono\ :\ TwistedJointMonotonicity\ (higherOrderModel\ D.F\ D.G\ α\ γ\ σ)\ G\ ρ)\ :}
\leanline{4}{∃\ n\ :\ ℕ,\ ∃\ T\ :\ HOLState\ (Lp\ A\ 2\ Q)\ (Lp\ (EuclideanSpace\ ℝ\ (Fin\ N\ ×\ J))\ 2\ Q)\ →L[ℝ]}
\leanline{8}{Lp\ (EuclideanSpace\ ℝ\ (Fin\ n))\ 2\ Q,}
\leanline{6}{(∀\ x,\ (Real.sqrt\ (EmpiricalMetric.lower\ N\ G)/2)*‖x‖≤‖T\ x‖)\ ∧}
\leanline{6}{(∀\ x,\ ‖T\ x‖≤(2*Real.sqrt\ (EmpiricalMetric.upper\ N\ G))*‖x‖)\ ∧}
\leanline{6}{(∀\ x\ y,\ {\leanfallback ⟪}T\ (x-y),T\ (holDrift\ (D.meanForceLp\ Q)\ (D.samplerForceLp\ Q)\ α\ γ\ x-}
\leanline{8}{holDrift\ (D.meanForceLp\ Q)\ (D.samplerForceLp\ Q)\ α\ γ\ y){\leanfallback ⟫}\_ℝ\ ≤\ -ρ*‖T\ (x-y)‖\textasciicircum{}2)}
\end{leancode}
\end{leanbox}
\subsection{\texorpdfstring{Both methods — Explicit uniform constants and allowed steps}{Both methods — Explicit uniform constants and allowed steps}}



For \(j=M,H\), use the corresponding \(K_j,\beta_j,c_{0,j},c_{1,j}\) above and the same initial raw phase \(y_0\). All remaining constants are





\[S_j=U_Ne^{K_j+1}(b_\circ+\beta_j),\qquad R_j=U_N\|y_0\|_{\max,2}+\frac{S_j}{1-e^{-\rho}}+S_j,\]







\[a_j=U_Nc_{0,j},\qquad b_j=\frac{U_Nc_{1,j}}{m_N},\qquad D_j=\frac{4(a_j+b_jR_j)}{\rho m_N},\]







\[h_{\star,j}=\min\left\{1,\frac1\rho,\left(\frac{\rho}{4(b_j+1)}\right)^2\right\}>0.\]





Thus \(D_j\) and \(h_{\star,j}\) depend on the fixed population, initial second moments, force values and Lipschitz constants, friction parameters and A3 metric/\allowbreak{}rate. They are independent of \(h,n\) and time horizon. The proof derives the reference radius and absorbs the local errors; they are not additional moment or numerical-stability assumptions.

\begin{leanbox}{Lean: \texttt{zeroStartMoment}, \texttt{referenceRadius}, \texttt{radiusGrowth} and 7 more}
\leansource{\texttt{zeroStartMoment} \quad PBGSO/\allowbreak{}FullyUniformMomentum.lean, lines 77–78}
\begin{leancode}
\leanline{0}{def\ zeroStartMoment\ (hγ\ :\ 0\ ≤\ γ)\ :\ ℝ\ :=}
\leanline{2}{metricUpper\ (N\ :=\ N)\ Geometry\ *\ zeroStartRadius\ D\ γ\ μ\ hγ}
\end{leancode}
\leansource{\texttt{referenceRadius} \quad PBGSO/\allowbreak{}FullyUniformMomentum.lean, lines 81–84}
\begin{leancode}
\leanline{0}{def\ referenceRadius\ (hγ\ :\ 0\ ≤\ γ)\ (ρ\ :\ ℝ)\ :\ ℝ\ :=}
\leanline{2}{let\ S\ :=\ zeroStartMoment\ D\ γ\ Geometry\ μ\ hγ}
\leanline{2}{metricUpper\ (N\ :=\ N)\ Geometry\ *\ ‖initial\ D\ γ\ Geometry\ μ\ hμ\ a\ Variance\ hVariance‖\ +}
\leanline{4}{S\ /\ (1\ -\ Real.exp\ (-ρ))\ +\ S}
\end{leancode}
\leansource{\texttt{radiusGrowth} \quad PBGSO/\allowbreak{}FullyUniformMomentum.lean, lines 87–88}
\begin{leancode}
\leanline{0}{def\ radiusGrowth\ (hγ\ :\ 0\ ≤\ γ)\ :\ ℝ\ :=}
\leanline{2}{metricUpper\ (N\ :=\ N)\ Geometry\ *\ localConstant1\ D\ γ\ μ\ hγ\ /\ metricCoercivity\ (N\ :=\ N)\ Geometry}
\end{leancode}
\leansource{\texttt{discretizationConstant} \quad PBGSO/\allowbreak{}FullyUniformMomentum.lean, lines 91–94}
\begin{leancode}
\leanline{0}{def\ discretizationConstant\ (hγ\ :\ 0\ ≤\ γ)\ (ρ\ :\ ℝ)\ :\ ℝ\ :=}
\leanline{2}{4\ *\ (metricUpper\ (N\ :=\ N)\ Geometry\ *\ localConstant0\ D\ γ\ μ\ hγ\ +}
\leanline{4}{radiusGrowth\ D\ γ\ Geometry\ μ\ hγ\ *\ referenceRadius\ D\ γ\ Geometry\ μ\ hμ\ a\ Variance\ hVariance\ hγ\ ρ)\ /}
\leanline{6}{(ρ\ *\ metricCoercivity\ (N\ :=\ N)\ Geometry)}
\end{leancode}
\leansource{\texttt{maximumStep} \quad PBGSO/\allowbreak{}FullyUniformMomentum.lean, lines 97–98}
\begin{leancode}
\leanline{0}{def\ maximumStep\ (hγ\ :\ 0\ ≤\ γ)\ (ρ\ :\ ℝ)\ :\ ℝ\ :=}
\leanline{2}{UniformStep.maximumStep\ ρ\ (radiusGrowth\ D\ γ\ Geometry\ μ\ hγ)}
\end{leancode}
\leansource{\texttt{zeroStartMoment} \quad PBGSO/\allowbreak{}FullyUniformHigherOrder.lean, lines 86–86}
\begin{leancode}
\leanline{0}{def\ zeroStartMoment\ :\ ℝ\ :=\ metricUpper\ (N\ :=\ N)\ Geo\ *\ zeroStartRadius\ D\ α\ γ\ μ}
\end{leancode}
\leansource{\texttt{referenceRadius} \quad PBGSO/\allowbreak{}FullyUniformHigherOrder.lean, lines 89–92}
\begin{leancode}
\leanline{0}{def\ referenceRadius\ (ρ\ :\ ℝ)\ :\ ℝ\ :=}
\leanline{2}{let\ S\ :=\ zeroStartMoment\ D\ α\ γ\ Geo\ μ}
\leanline{2}{metricUpper\ (N\ :=\ N)\ Geo\ *\ ‖initial\ D\ α\ γ\ Geo\ μ\ hμ\ a\ Variance\ hVariance‖\ +}
\leanline{4}{S\ /\ (1\ -\ Real.exp\ (-ρ))\ +\ S}
\end{leancode}
\leansource{\texttt{radiusGrowth} \quad PBGSO/\allowbreak{}FullyUniformHigherOrder.lean, lines 95–96}
\begin{leancode}
\leanline{0}{def\ radiusGrowth\ :\ ℝ\ :=}
\leanline{2}{metricUpper\ (N\ :=\ N)\ Geo\ *\ localConstant1\ D\ α\ γ\ μ\ /\ metricCoercivity\ (N\ :=\ N)\ Geo}
\end{leancode}
\leansource{\texttt{discretizationConstant} \quad PBGSO/\allowbreak{}FullyUniformHigherOrder.lean, lines 99–102}
\begin{leancode}
\leanline{0}{def\ discretizationConstant\ (ρ\ :\ ℝ)\ :\ ℝ\ :=}
\leanline{2}{(4\ *\ (metricUpper\ (N\ :=\ N)\ Geo\ *\ localConstant0\ D\ α\ γ\ μ\ +}
\leanline{4}{radiusGrowth\ D\ α\ γ\ Geo\ μ\ *\ referenceRadius\ D\ α\ γ\ Geo\ μ\ hμ\ a\ Variance\ hVariance\ ρ)\ /\ ρ)\ /}
\leanline{6}{metricCoercivity\ (N\ :=\ N)\ Geo}
\end{leancode}
\leansource{\texttt{maximumStep} \quad PBGSO/\allowbreak{}FullyUniformHigherOrder.lean, lines 105–105}
\begin{leancode}
\leanline{0}{def\ maximumStep\ (ρ\ :\ ℝ)\ :\ ℝ\ :=\ UniformStep.maximumStep\ ρ\ (radiusGrowth\ D\ α\ γ\ Geo\ μ)}
\end{leancode}
\end{leanbox}
\chapter{Mathematical argument}
\section{\texorpdfstring{Define the full state, expectations and metric}{Define the full state, expectations and metric}}


Fix a complete optimiser state \(s\), sampler state \(z\), parameter projection \(P\), position projection \(R\), observation \(F\), optimiser response \(\Psi\), and sampler coefficients \(b,\sigma_i\). The mean observation is the actual integral \(G_\nu(s)=\int F(P(s),R(z))d\nu\). Its empirical version is integration against \(N^{-1}\sum_i\delta_{z_i}\), which equals the finite observation average for positive particle count. In particular the empirical optimiser uses the same response as the law-level optimiser.



Choose symmetric positive quadratic weights with common bounds \(0<\ell\le u\), and write \(\kappa=u/\ell\). For a coupling of two sampler laws, add the weighted optimiser square to its expected sampler square to obtain an energy. The joint dissipativity assumption bounds half its derivative by \(-\rho E\), including the finite sum of diffusion-difference squares. Every observation and signed work expression used in this assumption must be integrable.

\begin{leanbox}{Lean: \texttt{Model}, \texttt{meanGradient\_\allowbreak{}empirical}, \texttt{TwistedJointMonotonicity}}
\leansource{\texttt{Model} \quad PBGSO/\allowbreak{}Framework.lean, lines 40–49}
\begin{leancode}
\leanline{0}{structure\ Model\ (S\ Θ\ Z\ X\ I\ :\ Type*)}
\leanline{4}{[NormedAddCommGroup\ S]\ [NormedSpace\ ℝ\ S]}
\leanline{4}{[NormedAddCommGroup\ Θ]\ [NormedSpace\ ℝ\ Θ]}
\leanline{4}{[NormedAddCommGroup\ Z]\ [NormedSpace\ ℝ\ Z]\ where}
\leanline{2}{parameter\ :\ S\ →\ Θ}
\leanline{2}{position\ :\ Z\ →\ X}
\leanline{2}{observation\ :\ Θ\ →\ X\ →\ Θ}
\leanline{2}{optimizer\ :\ ℝ\ →\ S\ →\ Θ\ →\ S}
\leanline{2}{drift\ :\ Θ\ →\ Z\ →\ Z}
\leanline{2}{diffusion\ :\ Θ\ →\ Z\ →\ I\ →\ Z}
\end{leancode}
\leansource{\texttt{meanGradient\_\allowbreak{}empirical} \quad PBGSO/\allowbreak{}Framework.lean, lines 76–84}
\begin{leancode}
\leanline{0}{theorem\ meanGradient\_empirical\ [MeasurableSingletonClass\ Z]\ [CompleteSpace\ Θ]}
\leanline{4}{(N\ :\ ℕ)\ (s\ :\ S)\ (z\ :\ Fin\ N\ →\ Z)\ :}
\leanline{4}{M.meanGradient\ (empiricalLaw\ N\ z)\ s\ =\ M.particleGradient\ N\ s\ z}
\end{leancode}
\leansource{\texttt{TwistedJointMonotonicity} \quad PBGSO/\allowbreak{}PaperTheorems.lean, lines 87–101}
\begin{leancode}
\leanline{0}{def\ TwistedJointMonotonicity\ (M\ :\ Model\ S\ Θ\ Z\ X\ I)}
\leanline{4}{(D\ :\ TwistedGeometry\ S\ Z)\ (ρ\ :\ ℝ)\ :\ Prop\ :=}
\leanline{2}{∀\ t,\ 0\ ≤\ t\ →\ ∀\ s\ s'\ ν\ ν'\ (Λ\ :\ Measure\ (Z\ ×\ Z)),}
\leanline{4}{IsCoupling\ Λ\ ν\ ν'\ →}
\leanline{4}{IsProbabilityMeasure\ ν\ →\ IsProbabilityMeasure\ ν'\ →}
\leanline{4}{HasFiniteSecondMoment\ ν\ →\ HasFiniteSecondMoment\ ν'\ →}
\leanline{4}{Integrable\ (fun\ z\ =>\ M.observation\ (M.parameter\ s)\ (M.position\ z))\ ν\ →}
\leanline{4}{Integrable\ (fun\ z\ =>\ M.observation\ (M.parameter\ s')\ (M.position\ z))\ ν'\ →}
\leanline{4}{Integrable\ (samplerCouplingWork\ M\ D\ s\ s')\ Λ\ →}
\leanline{4}{Integrable\ (fun\ zz\ =>\ D.samplerSq\ (zz.1\ -\ zz.2))\ Λ\ →}
\leanline{4}{{\leanfallback ⟪}s\ -\ s',\ D.optimizerWeight}
\leanline{6}{(M.optimizerField\ t\ s\ ν\ -\ M.optimizerField\ t\ s'\ ν'){\leanfallback ⟫}\_ℝ\ +}
\leanline{6}{(∫\ zz,\ samplerCouplingWork\ M\ D\ s\ s'\ zz\ ∂Λ)\ ≤}
\leanline{6}{-ρ\ *\ (D.optimizerSq\ (s\ -\ s')\ +}
\leanline{8}{∫\ zz,\ D.samplerSq\ (zz.1\ -\ zz.2)\ ∂Λ)}
\end{leancode}
\leanlabel{How the Lean proof follows the argument}
\texttt{Model} retains independent maps and full auxiliary coordinates. \texttt{mean\allowbreak{}Gradient\_\allowbreak{}empirical} identifies the actual empirical expectation. \texttt{Twisted\allowbreak{}Joint\allowbreak{}Monotonicity} is the complete joint assumption with its finite-integral domain. These definitions fix the objects used in every subsequent theorem.
\end{leanbox}
\section{\texorpdfstring{Construct and transfer the stationary reference}{Construct and transfer the stationary reference}}


Let the frozen sampler at parameter \(P(s_\star)\) have invariant law \(\nu_\star\). Require integrability of its observation, \(G_{\nu_\star}(s_\star)=0\), and \(\Psi(t,s_\star,0)=0\). Then adjoining the constant optimiser to this sampler solves the coupled equation. Strong existence on initial-law times Wiener-path space and same-driver uniqueness transfer its law to any original strong solution with the same initial optimiser and sampler law. Thus its complete state has law \(\delta_{s_\star}\otimes\nu_\star\). The stationarity claim retains the frozen invariant sampler and strong existence/\allowbreak{}uniqueness assumptions; it is not inferred from the quadratic inequality alone.

\begin{leanbox}{Lean: \texttt{stationary\_\allowbreak{}product\_\allowbreak{}law\_\allowbreak{}of\_\allowbreak{}weak\_\allowbreak{}frozen}}
\leansource{\texttt{stationary\_\allowbreak{}product\_\allowbreak{}law\_\allowbreak{}of\_\allowbreak{}weak\_\allowbreak{}frozen} \quad PBGSO/\allowbreak{}RawPathSolutionTransfer.lean, lines 120–140}
\begin{leancode}
\leanline{0}{theorem\ stationary\_product\_law\_of\_weak\_frozen\ [MeasurableSpace\ S]}
\leanline{4}{\{sStar\ :\ S\}\ \{νStar\ :\ Measure\ Z\}\ [IsProbabilityMeasure\ νStar]}
\leanline{4}{(F\ :\ FrozenSamplerData\ Qw\ Dw\ M\ T\ sStar\ νStar)}
\leanline{4}{(hint\ :\ Integrable\ (fun\ z\ =>\ M.observation\ (M.parameter\ sStar)\ (M.position\ z))\ νStar)}
\leanline{4}{(hg\ :\ M.meanGradient\ νStar\ sStar\ =\ 0)}
\leanline{4}{(hfix\ :\ ∀\ t\ ∈\ Icc\ 0\ (T\ :\ ℝ),\ M.optimizer\ t\ sStar\ 0\ =\ 0)}
\leanline{4}{(huniquew\ :\ PathwiseUniqueness\ Qw\ (nullAugmentedDriver\ Qw\ (continuousDriver\ Qw\ Dw))\ M\ T)}
\leanline{4}{(hexists\ :\ ∃\ V\ :\ Data\ (BrownianPathDriver.probability\ (I\ :=\ I)\ νStar)}
\leanline{8}{(BrownianPathDriver.driver\ νStar)\ M\ T,}
\leanline{6}{V.optimizer\ 0\ =\ sStar\ ∧\ V.sampler.state\ 0\ =ᵐ[}
\leanline{8}{BrownianPathDriver.probability\ (I\ :=\ I)\ νStar]\ Prod.fst)}
\leanline{4}{(U\ :\ Data\ Q\ D\ M\ T)}
\leanline{4}{(hunique\ :\ PathwiseUniqueness\ Q\ (nullAugmentedDriver\ Q\ (continuousDriver\ Q\ D))\ M\ T)}
\leanline{4}{(hs\ :\ U.optimizer\ 0\ =\ sStar)\ (hν\ :\ Q.map\ (U.sampler.state\ 0)\ =\ νStar)}
\leanline{4}{(t\ :\ ℝ)\ (ht\ :\ t\ ∈\ Icc\ 0\ (T\ :\ ℝ))\ :}
\leanline{4}{Q.map\ (fun\ ω\ =>\ (U.optimizer\ t,\ U.sampler.state\ t\ ω))\ =\ (Measure.dirac\ sStar).prod\ νStar}
\end{leancode}
\leanlabel{How the Lean proof follows the argument}
\texttt{stationary\_\allowbreak{}product\_\allowbreak{}law\_\allowbreak{}of\_\allowbreak{}weak\_\allowbreak{}frozen} constructs the constant-optimiser solution, identifies both it and the original solution with the genuine path-space solution by uniqueness, and converts the resulting marginal identities into the full product-law identity.
\end{leanbox}
\section{\texorpdfstring{Derive the law contraction from actual stochastic equations}{Derive the law contraction from actual stochastic equations}}


For each initial coupling, place the two initial samplers with common Brownian coordinates on a canonical source. Localise the raw diffusion by its pathwise squared coefficient energy. Stopped coefficients have finite expected energy, so the actual Brownian integral and its quadratic identity apply. Nested stopping times glue these integrals consistently to the original local equation. Joint dissipativity then gives \(E'\le-2\rho E\), and scalar comparison yields \(E(T)\le e^{-2\rho T}E(0)\).



Convert the weighted energies to Euclidean ones using \(\ell,u\), take square roots, and minimise over initial couplings. With \(d((s,\mu),(s',\nu))^2=\|s-s'\|^2+W_2(\mu,\nu)^2\), this gives \(d_T\le\sqrt\kappa e^{-\rho T}d_0\). The passage from arbitrary original probability spaces uses path-source existence and same-driver uniqueness to identify their laws with the canonical ones. There are two explicitly separate domains: primitive coefficient linear growth with measurable/\allowbreak{}progressive realizations, or constant diffusion with universal finite signed coupling work and Borel fields.

\begin{leanbox}{Lean: \texttt{RawLocalDiffusion}, \texttt{law\_\allowbreak{}distance\_\allowbreak{}decay}, \texttt{law\_\allowbreak{}distance\_\allowbreak{}decay}}
\leansource{\texttt{RawLocalDiffusion} \quad PBGSO/\allowbreak{}RawLocalIntegral.lean, lines 26–37}
\begin{leancode}
\leanline{0}{structure\ RawLocalDiffusion\ (T\ :\ ℝ≥0)\ where}
\leanline{2}{sigma\ :\ I\ →\ ℝ\ →\ Ω\ →\ Z}
\leanline{2}{sigma\_strong\ :\ ∀\ i,\ StronglyMeasurable\ (Function.uncurry\ (sigma\ i))}
\leanline{2}{sigma\_adapted\ :\ ∀\ t\ :\ ℝ≥0,\ ∀\ i,\ AEStronglyMeasurable[D.filtration\ t]\ (sigma\ i\ t)\ Q}
\leanline{2}{q\ :\ ℝ\ →\ Ω\ →\ ℝ}
\leanline{2}{q\_strong\ :\ StronglyMeasurable\ (Function.uncurry\ q)}
\leanline{2}{q\_nonneg\ :\ ∀\ r\ ω,\ 0\ ≤\ q\ r\ ω}
\leanline{2}{q\_local\ :\ ∀\ ω,\ IntegrableOn\ (fun\ r\ =>\ q\ r\ ω)\ (Ioc\ 0\ (T\ :\ ℝ))}
\leanline{2}{q\_progressive\ :\ ∀\ t\ :\ ℝ≥0,}
\leanline{4}{StronglyMeasurable[(inferInstance\ :\ MeasurableSpace\ ℝ).prod\ (D.filtration\ t)]}
\leanline{6}{(fun\ p\ :\ ℝ\ ×\ Ω\ =>\ (Ioc\ 0\ (t\ :\ ℝ)).indicator\ (fun\ r\ =>\ q\ r\ p.2)\ p.1)}
\leanline{2}{sigma\_control\ :\ ∀\ r\ ω,\ ∑\ i,\ ‖sigma\ i\ r\ ω‖\ \textasciicircum{}\ 2\ ≤\ q\ r\ ω}
\end{leancode}
\leansource{\texttt{law\_\allowbreak{}distance\_\allowbreak{}decay} \quad PBGSO/\allowbreak{}RawSamplerGrowthArbitraryLaw.lean, lines 43–81}
\begin{leancode}
\leanline{0}{theorem\ law\_distance\_decay}
\leanline{4}{(U\ :\ RawLocalStationarity.Data\ Q\ D\ M\ T)\ (V\ :\ RawLocalStationarity.Data\ R\ E\ M\ T)}
\leanline{4}{(μ\ ν\ :\ Measure\ Z)\ [IsProbabilityMeasure\ μ]\ [IsProbabilityMeasure\ ν]}
\leanline{4}{(hμ\ :\ MemLp\ (fun\ z\ :\ Z\ =>\ z)\ 2\ μ)\ (hν\ :\ MemLp\ (fun\ z\ :\ Z\ =>\ z)\ 2\ ν)}
\leanline{4}{(hU\ :\ Q.map\ (U.sampler.state\ 0)\ =\ μ)\ (hV\ :\ R.map\ (V.sampler.state\ 0)\ =\ ν)}
\leanline{4}{(A\ :\ CanonicalSolution\ (M\ :=\ M)\ (T\ :=\ T)\ μ\ (U.optimizer\ 0))}
\leanline{4}{(B\ :\ CanonicalSolution\ (M\ :=\ M)\ (T\ :=\ T)\ ν\ (V.optimizer\ 0))}
\leanline{4}{(huniqueU\ :\ PathwiseUniqueness\ Q\ (nullAugmentedDriver\ Q\ (continuousDriver\ Q\ D))\ M\ T)}
\leanline{4}{(huniqueV\ :\ PathwiseUniqueness\ R\ (nullAugmentedDriver\ R\ (continuousDriver\ R\ E))\ M\ T)}
\leanline{4}{(huniqueA\ :\ PathwiseUniqueness\ (μ.prod\ (brownianProbability\ (I\ :=\ I)))}
\leanline{6}{(nullAugmentedDriver\ \_\ (continuousDriver\ \_\ (CanonicalBrownianDriver.driver\ (I\ :=\ I)\ μ)))\ M\ T)}
\leanline{4}{(huniqueB\ :\ PathwiseUniqueness\ (ν.prod\ (brownianProbability\ (I\ :=\ I)))}
\leanline{6}{(nullAugmentedDriver\ \_\ (continuousDriver\ \_\ (CanonicalBrownianDriver.driver\ (I\ :=\ I)\ ν)))\ M\ T)}
\leanline{4}{(hpathμ\ :\ ∃\ W\ :\ RawLocalStationarity.Data\ (BrownianPathDriver.probability\ (I\ :=\ I)\ μ)}
\leanline{8}{(BrownianPathDriver.driver\ μ)\ M\ T,}
\leanline{6}{W.optimizer\ 0\ =\ U.optimizer\ 0\ ∧}
\leanline{6}{W.sampler.state\ 0\ =ᵐ[BrownianPathDriver.probability\ (I\ :=\ I)\ μ]\ Prod.fst)}
\leanline{4}{(hpathν\ :\ ∃\ W\ :\ RawLocalStationarity.Data\ (BrownianPathDriver.probability\ (I\ :=\ I)\ ν)}
\leanline{8}{(BrownianPathDriver.driver\ ν)\ M\ T,}
\leanline{6}{W.optimizer\ 0\ =\ V.optimizer\ 0\ ∧}
\leanline{6}{W.sampler.state\ 0\ =ᵐ[BrownianPathDriver.probability\ (I\ :=\ I)\ ν]\ Prod.fst)}
\leanline{4}{(G\ :\ TwistedGeometry\ S\ Z)\ (ρ\ :\ ℝ)\ (hmono\ :\ TwistedJointMonotonicity\ M\ G\ ρ)\ :}
\leanline{4}{optimizerLawDistance\ (U.optimizer\ T)\ (V.optimizer\ T)}
\leanline{6}{(Q.map\ (U.sampler.state\ T))\ (R.map\ (V.sampler.state\ T))\ ≤}
\leanline{6}{Real.sqrt\ G.conditionNumber\ *\ Real.exp\ (-ρ\ *\ (T\ :\ ℝ))\ *}
\leanline{8}{optimizerLawDistance\ (U.optimizer\ 0)\ (V.optimizer\ 0)\ μ\ ν}
\end{leancode}
\leansource{\texttt{law\_\allowbreak{}distance\_\allowbreak{}decay} \quad PBGSO/\allowbreak{}NaturalFiniteWorkArbitraryLaw.lean, lines 43–88}
\begin{leancode}
\leanline{0}{theorem\ law\_distance\_decay}
\leanline{4}{(U\ :\ RawLocalStationarity.Data\ Q\ D\ M\ T)\ (V\ :\ RawLocalStationarity.Data\ R\ E\ M\ T)}
\leanline{4}{(μ\ ν\ :\ Measure\ Z)\ [IsProbabilityMeasure\ μ]\ [IsProbabilityMeasure\ ν]}
\leanline{4}{(hμ\ :\ MemLp\ (fun\ z\ :\ Z\ =>\ z)\ 2\ μ)\ (hν\ :\ MemLp\ (fun\ z\ :\ Z\ =>\ z)\ 2\ ν)}
\leanline{4}{(hU\ :\ Q.map\ (U.sampler.state\ 0)\ =\ μ)\ (hV\ :\ R.map\ (V.sampler.state\ 0)\ =\ ν)}
\leanline{4}{(A\ :\ CanonicalSolution\ (M\ :=\ M)\ (T\ :=\ T)\ μ\ (U.optimizer\ 0))}
\leanline{4}{(B\ :\ CanonicalSolution\ (M\ :=\ M)\ (T\ :=\ T)\ ν\ (V.optimizer\ 0))}
\leanline{4}{(huniqueU\ :\ PathwiseUniqueness\ Q\ (nullAugmentedDriver\ Q\ (continuousDriver\ Q\ D))\ M\ T)}
\leanline{4}{(huniqueV\ :\ PathwiseUniqueness\ R\ (nullAugmentedDriver\ R\ (continuousDriver\ R\ E))\ M\ T)}
\leanline{4}{(huniqueA\ :\ PathwiseUniqueness\ (μ.prod\ (brownianProbability\ (I\ :=\ I)))}
\leanline{6}{(nullAugmentedDriver\ \_\ (continuousDriver\ \_\ (CanonicalBrownianDriver.driver\ (I\ :=\ I)\ μ)))\ M\ T)}
\leanline{4}{(huniqueB\ :\ PathwiseUniqueness\ (ν.prod\ (brownianProbability\ (I\ :=\ I)))}
\leanline{6}{(nullAugmentedDriver\ \_\ (continuousDriver\ \_\ (CanonicalBrownianDriver.driver\ (I\ :=\ I)\ ν)))\ M\ T)}
\leanline{4}{(hpathμ\ :\ ∃\ W\ :\ RawLocalStationarity.Data\ (BrownianPathDriver.probability\ (I\ :=\ I)\ μ)}
\leanline{8}{(BrownianPathDriver.driver\ μ)\ M\ T,}
\leanline{6}{W.optimizer\ 0\ =\ U.optimizer\ 0\ ∧}
\leanline{6}{W.sampler.state\ 0\ =ᵐ[BrownianPathDriver.probability\ (I\ :=\ I)\ μ]\ Prod.fst)}
\leanline{4}{(hpathν\ :\ ∃\ W\ :\ RawLocalStationarity.Data\ (BrownianPathDriver.probability\ (I\ :=\ I)\ ν)}
\leanline{8}{(BrownianPathDriver.driver\ ν)\ M\ T,}
\leanline{6}{W.optimizer\ 0\ =\ V.optimizer\ 0\ ∧}
\leanline{6}{W.sampler.state\ 0\ =ᵐ[BrownianPathDriver.probability\ (I\ :=\ I)\ ν]\ Prod.fst)}
\leanline{4}{(G\ :\ TwistedGeometry\ S\ Z)\ (ρ\ :\ ℝ)\ (hρ\ :\ 0\ <\ ρ)}
\leanline{4}{(hmono\ :\ TwistedJointMonotonicity\ M\ G\ ρ)}
\leanline{4}{(hfull\ :\ DiracOptimizerGrowth.FiniteCouplingWork\ M\ G)}
\leanline{4}{(hσ\ :\ ∀\ s\ s'\ z\ w\ i,\ M.diffusion\ (M.parameter\ s)\ z\ i\ =\ M.diffusion\ (M.parameter\ s')\ w\ i)}
\leanline{4}{(hb\ :\ ∀\ θ,\ Measurable\ (M.drift\ θ))}
\leanline{4}{(hObs\ :\ Measurable\ (fun\ p\ :\ S\ ×\ Z\ =>\ M.observation\ (M.parameter\ p.1)\ (M.position\ p.2)))}
\leanline{4}{(hDrift\ :\ Measurable\ (fun\ p\ :\ S\ ×\ Z\ =>\ M.drift\ (M.parameter\ p.1)\ p.2))}
\leanline{4}{(hOptimizer\ :\ Measurable\ (fun\ p\ :\ ℝ\ ×\ (S\ ×\ Θ)\ =>\ M.optimizer\ p.1\ p.2.1\ p.2.2))\ :}
\leanline{4}{optimizerLawDistance\ (U.optimizer\ T)\ (V.optimizer\ T)}
\leanline{6}{(Q.map\ (U.sampler.state\ T))\ (R.map\ (V.sampler.state\ T))\ ≤}
\leanline{6}{Real.sqrt\ G.conditionNumber\ *\ Real.exp\ (-ρ\ *\ (T\ :\ ℝ))\ *}
\leanline{8}{optimizerLawDistance\ (U.optimizer\ 0)\ (V.optimizer\ 0)\ μ\ ν}
\end{leancode}
\leanlabel{How the Lean proof follows the argument}
\texttt{Raw\allowbreak{}Local\allowbreak{}Diffusion} supplies only primitive local coefficient data. The two \texttt{law\_\allowbreak{}distance\_\allowbreak{}decay} theorems identify canonical and original laws first, apply their established canonical contraction theorem, and rewrite its conclusion. The linear-growth and finite-work hypotheses remain distinct.
\end{leanbox}
\section{\texorpdfstring{Control the empirical interaction by independent copies}{Control the empirical interaction by independent copies}}


Form \(N>0\) independent mean-field references, driven by the same individual Brownian coordinates as the interacting particles. Start with identical optimiser and particle states. If the centred observation variance is at most \(V\), independence cancels every cross term in the squared empirical average, leaving variance \(V/N\). An optimiser response bound \(L_\Psi\) turns this into the only forcing term in the coupling energy.



Cauchy–Schwarz gives a differential inequality \(E'\le-2\rho E+2c\sqrt E\) with zero initial energy. To handle zeros rigorously, apply the derivative comparison to \(\sqrt{E+\varepsilon^2}\) and then let \(\varepsilon\downarrow0\). Conversion from the weighted energy and the labelled-particle empirical coupling gives interaction error at most \(\frac{L_\Psi}{\rho}\sqrt{\kappa V/N}(1-e^{-\rho T})\). The references' product law and the matching with the original native empirical state are proved by construction.

\begin{leanbox}{Lean: \texttt{forced\_\allowbreak{}energy\_\allowbreak{}bound}, \texttt{original\_\allowbreak{}finite\_\allowbreak{}particle\_\allowbreak{}interaction\_\allowbreak{}bound}}
\leansource{\texttt{forced\_\allowbreak{}energy\_\allowbreak{}bound} \quad PBGSO/\allowbreak{}PaperTheorems.lean, lines 123–180}
\begin{leancode}
\leanline{0}{theorem\ forced\_energy\_bound\ \{E\ E'\ :\ ℝ\ →\ ℝ\}\ \{ρ\ c\ T\ :\ ℝ\}}
\leanline{4}{(hρ\ :\ 0\ <\ ρ)\ (hc\ :\ 0\ ≤\ c)}
\leanline{4}{(hcont\ :\ ContinuousOn\ E\ (Icc\ 0\ T))}
\leanline{4}{(hder\ :\ ∀\ t\ ∈\ Ico\ 0\ T,\ HasDerivWithinAt\ E\ (E'\ t)\ (Ici\ t)\ t)}
\leanline{4}{(hE\ :\ ∀\ t\ ∈\ Icc\ 0\ T,\ 0\ ≤\ E\ t)\ (hzero\ :\ E\ 0\ =\ 0)}
\leanline{4}{(hdiss\ :\ ∀\ t\ ∈\ Ico\ 0\ T,}
\leanline{6}{E'\ t\ ≤\ -2\ *\ ρ\ *\ E\ t\ +\ 2\ *\ c\ *\ Real.sqrt\ (E\ t))}
\leanline{4}{(t\ :\ ℝ)\ (ht\ :\ t\ ∈\ Icc\ 0\ T)\ :}
\leanline{4}{Real.sqrt\ (E\ t)\ ≤\ c\ /\ ρ\ *\ (1\ -\ Real.exp\ (-ρ\ *\ t))}
\end{leancode}
\leansource{\texttt{original\_\allowbreak{}finite\_\allowbreak{}particle\_\allowbreak{}interaction\_\allowbreak{}bound} \quad PBGSO/\allowbreak{}RawConstructedParticleBounds.lean, lines 109–132}
\begin{leancode}
\leanline{0}{theorem\ original\_finite\_particle\_interaction\_bound}
\leanline{4}{(hBorelF\ :\ Measurable\ (fun\ p\ :\ S\ ×\ Z\ =>\ M.observation\ (M.parameter\ p.1)\ (M.position\ p.2)))}
\leanline{4}{(hBorelb\ :\ Measurable\ (fun\ p\ :\ S\ ×\ Z\ =>\ M.drift\ (M.parameter\ p.1)\ p.2))}
\leanline{4}{(hBorelΨ\ :\ Measurable\ (fun\ p\ :\ ℝ\ ×\ (S\ ×\ Θ)\ =>\ M.optimizer\ p.1\ p.2.1\ p.2.2))}
\leanline{4}{(G\ :\ TwistedGeometry\ S\ Z)\ (ρ\ :\ ℝ)\ (hρ\ :\ 0\ <\ ρ)}
\leanline{4}{(hmono\ :\ TwistedJointMonotonicity\ M\ G\ ρ)\ (hN\ :\ 0\ <\ N)}
\leanline{4}{(LΨ\ :\ ℝ≥0)\ (V\ :\ ℝ)\ (hV\ :\ 0\ ≤\ V)}
\leanline{4}{(hΨ\ :\ ∀\ t\ ∈\ Icc\ 0\ (T\ :\ ℝ),\ ∀\ s\ g\ g',}
\leanline{6}{‖M.optimizer\ t\ s\ g-M.optimizer\ t\ s\ g'‖\ ≤\ (LΨ\ :\ ℝ)*‖g-g'‖)}
\leanline{4}{(hF\ :\ ∀\ t\ ∈\ Icc\ 0\ (T\ :\ ℝ),\ MemLp}
\leanline{6}{(fun\ z\ =>\ M.observation\ (M.parameter\ (U.optimizer\ t))\ (M.position\ z))\ 2}
\leanline{6}{(Q.map\ (U.sampler.state\ t)))}
\leanline{4}{(hVar\ :\ ∀\ t\ ∈\ Icc\ 0\ (T\ :\ ℝ),}
\leanline{6}{(∫\ z,\ ‖M.observation\ (M.parameter\ (U.optimizer\ t))\ (M.position\ z)-}
\leanline{8}{M.meanGradient\ (Q.map\ (U.sampler.state\ t))\ (U.optimizer\ t)‖\textasciicircum{}2}
\leanline{10}{∂Q.map\ (U.sampler.state\ t))\ ≤\ V)}
\leanline{4}{(hinit\ :\ ∀ᵐ\ ω\ ∂IndependentCopies.seeds\ (J\ :=\ Fin\ N)\ Q,}
\leanline{6}{oneEquiv\ N\ (Y.state\ 0\ ω)\ =\ (U.optimizer\ 0,RawIndependentCopies.cloud\ Q\ D\ N\ U\ 0\ ω))\ :}
\leanline{4}{Real.sqrt\ (∫\ ω,\ particleEmpiricalDistanceSq\ N\ (originalPair\ Q\ D\ M\ N\ U\ Y\ T\ ω)}
\leanline{6}{∂IndependentCopies.seeds\ (J\ :=\ Fin\ N)\ Q)\ ≤}
\leanline{6}{((LΨ\ :\ ℝ)/ρ)*Real.sqrt\ (G.conditionNumber*V/(N\ :\ ℝ))*(1-Real.exp\ (-ρ*T))}
\end{leancode}
\leanlabel{How the Lean proof follows the argument}
\texttt{forced\_\allowbreak{}energy\_\allowbreak{}bound} performs the positive square-root regularisation. \texttt{original\_\allowbreak{}finite\_\allowbreak{}particle\_\allowbreak{}interaction\_\allowbreak{}bound} builds the native pair internally from the Borel model maps and the actual empirical solution, applies the paired estimate, and rewrites the integral using the proved almost-sure state identity.
\end{leanbox}
\section{\texorpdfstring{Add stationary relaxation and empirical error}{Add stationary relaxation and empirical error}}


Assume the parameter projection is \(1\)-Lipschitz. The \(L^2\) triangle inequality splits its distance from \(P(s_\star)\) into interaction error and the mean-field optimiser distance from \(s_\star\). The latter is bounded by the optimiser–law distance, whose transport-square summand is nonnegative. Combining with stationary law contraction yields

\[\|P(S^N_T)-P(s_\star)\|_{L^2}\le\sqrt\kappa e^{-\rho T}d_0+\frac{L_\Psi}{\rho}\sqrt{\frac{\kappa V}{N}}(1-e^{-\rho T}).\]

The initial prefactor is the true initial optimiser–law distance. The constant-diffusion theorem here derives the required comparison using the stated response and variance assumptions and actual stationary strong solution; it does not add a universal finite-work premise.

\begin{leanbox}{Lean: \texttt{original\_\allowbreak{}parameter\_\allowbreak{}bound\_\allowbreak{}constant}}
\leansource{\texttt{original\_\allowbreak{}parameter\_\allowbreak{}bound\_\allowbreak{}constant} \quad PBGSO/\allowbreak{}RawConstructedParticleParameter.lean, lines 98–153}
\begin{leancode}
\leanline{0}{theorem\ original\_parameter\_bound\_constant}
\leanline{4}{(U\ :\ RawLocalStationarity.Data\ Q\ W\ M\ T)}
\leanline{4}{(Y\ :\ RawEmpiricalDynamics.Solution\ M\ N\ (IndependentCopies.seeds\ (J\ :=\ Fin\ N)\ Q)}
\leanline{6}{(RawIndependentCopies.driver\ Q\ W\ N)\ T)}
\leanline{4}{(μ\ νStar\ :\ Measure\ Z)\ [IsProbabilityMeasure\ μ]\ [IsProbabilityMeasure\ νStar]}
\leanline{4}{(hμ\ :\ MemLp\ (fun\ z\ :\ Z\ =>\ z)\ 2\ μ)\ (hν\ :\ MemLp\ (fun\ z\ :\ Z\ =>\ z)\ 2\ νStar)}
\leanline{4}{(hUlaw\ :\ Q.map\ (U.sampler.state\ 0)\ =\ μ)}
\leanline{4}{(A\ :\ RawLocalStationarity.Data\ (μ.prod\ (brownianProbability\ (I\ :=\ I)))}
\leanline{6}{(CanonicalBrownianDriver.driver\ (I\ :=\ I)\ μ)\ M\ T)}
\leanline{4}{(hA0\ :\ A.optimizer\ 0\ =\ U.optimizer\ 0)}
\leanline{4}{(hAz\ :\ A.sampler.state\ 0\ =ᵐ[μ.prod\ (brownianProbability\ (I\ :=\ I))]\ Prod.fst)}
\leanline{4}{(V\ :\ RawLocalStationarity.Data\ (νStar.prod\ (brownianProbability\ (I\ :=\ I)))}
\leanline{6}{(CanonicalBrownianDriver.driver\ (I\ :=\ I)\ νStar)\ M\ T)}
\leanline{4}{(hVz\ :\ V.sampler.state\ 0\ =ᵐ[νStar.prod\ (brownianProbability\ (I\ :=\ I))]\ Prod.fst)}
\leanline{4}{(huniqueU\ :\ PathwiseUniqueness\ Q\ (nullAugmentedDriver\ Q\ (continuousDriver\ Q\ W))\ M\ T)}
\leanline{4}{(huniqueA\ :\ PathwiseUniqueness\ (μ.prod\ (brownianProbability\ (I\ :=\ I)))}
\leanline{6}{(nullAugmentedDriver\ \_\ (continuousDriver\ \_\ (CanonicalBrownianDriver.driver\ (I\ :=\ I)\ μ)))\ M\ T)}
\leanline{4}{(hpath\ :\ ∃\ B\ :\ RawLocalStationarity.Data\ (BrownianPathDriver.probability\ (I\ :=\ I)\ μ)}
\leanline{8}{(BrownianPathDriver.driver\ μ)\ M\ T,}
\leanline{6}{B.optimizer\ 0\ =\ U.optimizer\ 0\ ∧}
\leanline{6}{B.sampler.state\ 0\ =ᵐ[BrownianPathDriver.probability\ (I\ :=\ I)\ μ]\ Prod.fst)}
\leanline{4}{(G\ :\ TwistedGeometry\ S\ Z)\ (ρ\ :\ ℝ)\ (hρ\ :\ 0\ <\ ρ)}
\leanline{4}{(hmono\ :\ TwistedJointMonotonicity\ M\ G\ ρ)\ (hN\ :\ 0\ <\ N)}
\leanline{4}{(hProjection\ :\ LipschitzWith\ 1\ M.parameter)}
\leanline{4}{(hσ\ :\ ∀\ s\ s'\ z\ w\ i,\ M.diffusion\ (M.parameter\ s)\ z\ i\ =\ M.diffusion\ (M.parameter\ s')\ w\ i)}
\leanline{4}{(hb\ :\ ∀\ θ,\ Measurable\ (M.drift\ θ))}
\leanline{4}{(hF\ :\ Measurable\ (fun\ p\ :\ S\ ×\ Z\ =>\ M.observation\ (M.parameter\ p.1)\ (M.position\ p.2)))}
\leanline{4}{(hDrift\ :\ Measurable\ (fun\ p\ :\ S\ ×\ Z\ =>\ M.drift\ (M.parameter\ p.1)\ p.2))}
\leanline{4}{(hOptimizer\ :\ Measurable\ (fun\ p\ :\ ℝ\ ×\ (S\ ×\ Θ)\ =>\ M.optimizer\ p.1\ p.2.1\ p.2.2))}
\leanline{4}{(LΨ\ :\ ℝ≥0)\ (VF\ :\ ℝ)}
\leanline{4}{(hΨ\ :\ ∀\ r\ ∈\ Icc\ 0\ (T\ :\ ℝ),\ ∀\ s\ g\ g',}
\leanline{6}{‖M.optimizer\ r\ s\ g-M.optimizer\ r\ s\ g'‖\ ≤\ (LΨ\ :\ ℝ)*‖g-g'‖)}
\leanline{4}{(hU2\ :\ ∀\ r\ ∈\ Icc\ 0\ (T\ :\ ℝ),\ MemLp}
\leanline{6}{(fun\ z\ =>\ M.observation\ (M.parameter\ (U.optimizer\ r))\ (M.position\ z))\ 2\ (Q.map\ (U.sampler.state\ r)))}
\leanline{4}{(hVar\ :\ ∀\ r\ ∈\ Icc\ 0\ (T\ :\ ℝ),}
\leanline{6}{(∫\ z,\ ‖M.observation\ (M.parameter\ (U.optimizer\ r))\ (M.position\ z)-}
\leanline{8}{M.meanGradient\ (Q.map\ (U.sampler.state\ r))\ (U.optimizer\ r)‖\textasciicircum{}2\ ∂Q.map\ (U.sampler.state\ r))\ ≤\ VF)}
\leanline{4}{(sStar\ :\ S)}
\leanline{4}{(hopt\ :\ ∀\ r\ ∈\ Icc\ 0\ (T\ :\ ℝ),\ V.optimizer\ r\ =\ sStar)}
\leanline{4}{(hlaw\ :\ ∀\ r\ ∈\ Icc\ 0\ (T\ :\ ℝ),\ (νStar.prod\ (brownianProbability\ (I\ :=\ I))).map\ (V.sampler.state\ r)\ =\ νStar)}
\leanline{4}{(hobs\ :\ Integrable\ (fun\ z\ =>\ M.observation\ (M.parameter\ sStar)\ (M.position\ z))\ νStar)}
\leanline{4}{(hinit\ :\ ∀ᵐ\ ω\ ∂IndependentCopies.seeds\ (J\ :=\ Fin\ N)\ Q,}
\leanline{6}{oneEquiv\ N\ (Y.state\ 0\ ω)\ =\ (U.optimizer\ 0,RawIndependentCopies.cloud\ Q\ W\ N\ U\ 0\ ω))\ :}
\leanline{4}{MemLp\ (fun\ ω\ =>\ M.parameter\ (oneEquiv\ N\ (Y.state\ T\ ω)).1-M.parameter\ sStar)}
\leanline{6}{2\ (IndependentCopies.seeds\ (J\ :=\ Fin\ N)\ Q)\ ∧}
\leanline{4}{Real.sqrt\ (∫\ ω,\ ‖M.parameter\ (oneEquiv\ N\ (Y.state\ T\ ω)).1-M.parameter\ sStar‖\textasciicircum{}2}
\leanline{6}{∂IndependentCopies.seeds\ (J\ :=\ Fin\ N)\ Q)\ ≤}
\leanline{6}{Real.sqrt\ G.conditionNumber*Real.exp\ (-ρ*(T\ :\ ℝ))*optimizerLawDistance\ (U.optimizer\ 0)\ sStar\ μ\ νStar+}
\leanline{6}{((LΨ\ :\ ℝ)/ρ)*Real.sqrt\ (G.conditionNumber*VF/(N\ :\ ℝ))*(1-Real.exp\ (-ρ*T))}
\end{leancode}
\leanlabel{How the Lean proof follows the argument}
\texttt{original\_\allowbreak{}parameter\_\allowbreak{}bound\_\allowbreak{}constant} combines the native parameter interaction bound, \texttt{Raw\allowbreak{}Natural\allowbreak{}Constant\allowbreak{}Arbitrary\allowbreak{}Law.\allowbreak{}law\_\allowbreak{}decay}, and the optimiser projection inequality. It also proves square integrability of the actual parameter discrepancy.
\end{leanbox}
\section{\texorpdfstring{Verify joint dissipativity from the appendix assumptions}{Verify joint dissipativity from the appendix assumptions}}


Three criteria produce the joint inequality used above. For separated monotonicity, the optimiser and sampler coupling radii reduce the work to a symmetric two-dimensional quadratic form. Positive diagonal and determinant conditions make its smallest eigenvalue positive; its exact value is \((a+b-\sqrt{(a-b)^2+c^2})/2\).



For projected monotonicity, projected force dissipation and auxiliary work give the block matrices \(K_o=\operatorname{diag}(cI,D_m)\), \(K_s=\operatorname{diag}(cI,D_v)\), with \(c=\beta r-A_o-A_s>0\). Whitening each by its metric matrix gives the positive generalized-eigenvalue rate. Constant diffusion cancels, but signed-work integrability must still be asserted on its finite domain. A finite second moment alone does not suffice: the strongly dissipative drift \(-x-x^3\) has nonintegrable signed work under a probability law with finite second but infinite fourth moment.



For heavy-ball dynamics, use the quadratic form \(\gamma\|q\|^2+2\langle q,p\rangle+2\|p\|^2/\gamma\). The threshold \(\gamma>\sqrt2L/\sqrt r\) gives an explicit Young parameter between zero and one, leaving positive position and momentum dissipation and therefore a positive contraction rate. These criteria allow independent forces without a common potential.

\begin{leanbox}{Lean: \texttt{separated\_\allowbreak{}source\_\allowbreak{}criterion}, \texttt{canonical\_\allowbreak{}projected\_\allowbreak{}source\_\allowbreak{}criterion}, \texttt{canonical\_\allowbreak{}projected\_\allowbreak{}source\_\allowbreak{}criterion\_\allowbreak{}finite} and 2 more}
\leansource{\texttt{separated\_\allowbreak{}source\_\allowbreak{}criterion} \quad PBGSO/\allowbreak{}SmallGainSpectrum.lean, lines 105–126}
\begin{leancode}
\leanline{0}{theorem\ separated\_source\_criterion\ (M\ :\ Model\ S\ Θ\ Z\ X\ I)\ (D\ :\ TwistedGeometry\ S\ Z)}
\leanline{4}{(rO\ rS\ :\ ℝ)\ (LΨ\ LF\ Lb\ Lσ\ :\ ℝ≥0)}
\leanline{4}{(H\ :\ SeparatedAssumptions\ M\ D\ rO\ rS\ LΨ\ LF\ Lb\ Lσ)}
\leanline{4}{(hO\ :\ 0\ <\ 2\ *\ rO\ -\ D.conditionNumber\ *\ (Lσ\ :\ ℝ)\ \textasciicircum{}\ 2)}
\leanline{4}{(hS\ :\ 0\ <\ 2\ *\ rS\ -\ D.conditionNumber\ *\ (Lσ\ :\ ℝ)\ \textasciicircum{}\ 2)}
\leanline{4}{(hdet\ :\ D.conditionNumber\ *\ ((LΨ\ :\ ℝ)\ *\ (LF\ :\ ℝ)\ +\ (Lb\ :\ ℝ))\ \textasciicircum{}\ 2\ <}
\leanline{6}{(2\ *\ rO\ -\ D.conditionNumber\ *\ (Lσ\ :\ ℝ)\ \textasciicircum{}\ 2)\ *}
\leanline{8}{(2\ *\ rS\ -\ D.conditionNumber\ *\ (Lσ\ :\ ℝ)\ \textasciicircum{}\ 2))\ :}
\leanline{4}{(separatedSmallGainMatrix\ D\ rO\ rS\ LΨ\ LF\ Lb\ Lσ).PosDef\ ∧}
\leanline{6}{0\ <\ separatedMinimumEigenvalue\ D\ rO\ rS\ LΨ\ LF\ Lb\ Lσ\ ∧}
\leanline{6}{TwistedJointMonotonicity\ M\ D\ (separatedMinimumEigenvalue\ D\ rO\ rS\ LΨ\ LF\ Lb\ Lσ)}
\end{leancode}
\leansource{\texttt{canonical\_\allowbreak{}projected\_\allowbreak{}source\_\allowbreak{}criterion} \quad PBGSO/\allowbreak{}CanonicalProjectedCriterion.lean, lines 264–295}
\begin{leancode}
\leanline{0}{theorem\ canonical\_projected\_source\_criterion}
\leanline{4}{(M\ :\ Model\ (BlockSpace\ a\ u)\ (EuclideanSpace\ ℝ\ a)\ (BlockSpace\ x\ v)\ (EuclideanSpace\ ℝ\ x)\ I)}
\leanline{4}{(hparameter\ :\ M.parameter\ =\ WithLp.fst)\ (hposition\ :\ M.position\ =\ WithLp.fst)}
\leanline{4}{(H\ :\ EuclideanSpace\ ℝ\ a\ →\ EuclideanSpace\ ℝ\ x\ →\ EuclideanSpace\ ℝ\ x)}
\leanline{4}{(DOpt\ :\ Matrix\ (a\ ⊕\ u)\ (a\ ⊕\ u)\ ℝ)\ (DSam\ :\ Matrix\ (x\ ⊕\ v)\ (x\ ⊕\ v)\ ℝ)}
\leanline{4}{(Dm\ :\ Matrix\ u\ u\ ℝ)\ (Dv\ :\ Matrix\ v\ v\ ℝ)}
\leanline{4}{(hDOpt\ :\ DOpt.PosDef)\ (hDSam\ :\ DSam.PosDef)\ (hDm\ :\ Dm.PosDef)\ (hDv\ :\ Dv.PosDef)}
\leanline{4}{(β\ r\ AOpt\ ASam\ :\ ℝ)\ (hβ\ :\ 0\ ≤\ β)\ (hc\ :\ 0\ <\ β\ *\ r\ -\ AOpt\ -\ ASam)}
\leanline{4}{(hmono\ :\ (canonicalData\ M\ hparameter\ hposition\ H\ Dm\ Dv).StrongMonotonicity\ r)}
\leanline{4}{(haux\ :\ ProjectedAuxiliaryDissipation\ M\ (canonicalGeometry\ DOpt\ DSam\ hDOpt\ hDSam)}
\leanline{6}{(canonicalData\ M\ hparameter\ hposition\ H\ Dm\ Dv)\ β\ AOpt\ ASam)}
\leanline{4}{(σ\ :\ I\ →\ BlockSpace\ x\ v)\ (hσ\ :\ ∀\ θ\ z,\ M.diffusion\ θ\ z\ =\ σ)}
\leanline{4}{(hInt\ :\ ∀\ s\ s'\ ν\ ν'\ Λ,}
\leanline{6}{ProjectedAdmissibleCoupling\ M\ (canonicalGeometry\ DOpt\ DSam\ hDOpt\ hDSam)\ s\ s'\ ν\ ν'\ Λ\ →}
\leanline{6}{(canonicalData\ M\ hparameter\ hposition\ H\ Dm\ Dv).Integrability\ s\ s'\ Λ)\ :}
\leanline{4}{let\ ρ\ :=\ min}
\leanline{6}{(generalizedMinimum\ hDOpt\ (blockK\_posDef\ (a\ :=\ a)\ \_\ hc\ Dm\ hDm))}
\leanline{6}{(generalizedMinimum\ hDSam\ (blockK\_posDef\ (a\ :=\ x)\ \_\ hc\ Dv\ hDv))}
\leanline{4}{0\ <\ ρ\ ∧\ TwistedJointMonotonicity\ M\ (canonicalGeometry\ DOpt\ DSam\ hDOpt\ hDSam)\ ρ}
\end{leancode}
\leansource{\texttt{canonical\_\allowbreak{}projected\_\allowbreak{}source\_\allowbreak{}criterion\_\allowbreak{}finite} \quad PBGSO/\allowbreak{}ProjectedFiniteWork.lean, lines 45–68}
\begin{leancode}
\leanline{0}{theorem\ canonical\_projected\_source\_criterion\_finite}
\leanline{4}{(M\ :\ Model\ (BlockSpace\ a\ u)\ (EuclideanSpace\ ℝ\ a)\ (BlockSpace\ x\ v)\ (EuclideanSpace\ ℝ\ x)\ I)}
\leanline{4}{(hparameter\ :\ M.parameter\ =\ WithLp.fst)\ (hposition\ :\ M.position\ =\ WithLp.fst)}
\leanline{4}{(H\ :\ EuclideanSpace\ ℝ\ a\ →\ EuclideanSpace\ ℝ\ x\ →\ EuclideanSpace\ ℝ\ x)}
\leanline{4}{(DOpt\ :\ Matrix\ (a\ ⊕\ u)\ (a\ ⊕\ u)\ ℝ)\ (DSam\ :\ Matrix\ (x\ ⊕\ v)\ (x\ ⊕\ v)\ ℝ)}
\leanline{4}{(Dm\ :\ Matrix\ u\ u\ ℝ)\ (Dv\ :\ Matrix\ v\ v\ ℝ)}
\leanline{4}{(hDOpt\ :\ DOpt.PosDef)\ (hDSam\ :\ DSam.PosDef)\ (hDm\ :\ Dm.PosDef)\ (hDv\ :\ Dv.PosDef)}
\leanline{4}{(β\ r\ AOpt\ ASam\ :\ ℝ)\ (hβ\ :\ 0\ ≤\ β)\ (hc\ :\ 0\ <\ β\ *\ r\ -\ AOpt\ -\ ASam)}
\leanline{4}{(hmono\ :\ (canonicalData\ M\ hparameter\ hposition\ H\ Dm\ Dv).StrongMonotonicity\ r)}
\leanline{4}{(haux\ :\ ProjectedAuxiliaryDissipation\ M\ (canonicalGeometry\ DOpt\ DSam\ hDOpt\ hDSam)}
\leanline{6}{(canonicalData\ M\ hparameter\ hposition\ H\ Dm\ Dv)\ β\ AOpt\ ASam)}
\leanline{4}{(σ\ :\ I\ →\ BlockSpace\ x\ v)\ (hσ\ :\ ∀\ θ\ z,\ M.diffusion\ θ\ z\ =\ σ)}
\leanline{4}{(hInt\ :\ ∀\ s\ s'\ ν\ ν'\ Λ,}
\leanline{6}{ProjectedAdmissibleCoupling\ M\ (canonicalGeometry\ DOpt\ DSam\ hDOpt\ hDSam)\ s\ s'\ ν\ ν'\ Λ\ →}
\leanline{6}{(canonicalData\ M\ hparameter\ hposition\ H\ Dm\ Dv).Integrability\ s\ s'\ Λ)}
\leanline{4}{(hfinite\ :\ FiniteSamplerDriftDomain\ M\ (canonicalGeometry\ DOpt\ DSam\ hDOpt\ hDSam))\ :}
\leanline{4}{let\ ρ\ :=\ min}
\leanline{6}{(generalizedMinimum\ hDOpt\ (blockK\_posDef\ (a\ :=\ a)\ \_\ hc\ Dm\ hDm))}
\leanline{6}{(generalizedMinimum\ hDSam\ (blockK\_posDef\ (a\ :=\ x)\ \_\ hc\ Dv\ hDv))}
\leanline{4}{0\ <\ ρ\ ∧\ TwistedJointMonotonicity\ M\ (canonicalGeometry\ DOpt\ DSam\ hDOpt\ hDSam)\ ρ\ ∧}
\leanline{6}{DiracOptimizerGrowth.FiniteCouplingWork\ M\ (canonicalGeometry\ DOpt\ DSam\ hDOpt\ hDSam)}
\end{leancode}
\leansource{\texttt{heavyBall\_\allowbreak{}source\_\allowbreak{}criterion} \quad PBGSO/\allowbreak{}HypocoerciveCriterion.lean, lines 644–657}
\begin{leancode}
\leanline{0}{theorem\ heavyBall\_source\_criterion\ (γ\ r\ :\ ℝ)\ (L\ :\ ℝ≥0)}
\leanline{4}{(hγ\ :\ 0\ <\ γ)\ (hr\ :\ 0\ <\ r)}
\leanline{4}{(hsource\ :\ Real.sqrt\ 2\ *\ (L\ :\ ℝ)\ /\ Real.sqrt\ r\ <\ γ)}
\leanline{4}{(F\ :\ A\ →\ X\ →\ A)\ (H\ :\ A\ →\ X\ →\ X)\ (σ\ :\ I\ →\ HypocoercivePhase\ X)}
\leanline{4}{(hmono\ :\ ∀\ θ\ θ'\ x\ x',\ r\ *\ (‖θ\ -\ θ'‖\ \textasciicircum{}\ 2\ +\ ‖x\ -\ x'‖\ \textasciicircum{}\ 2)\ ≤}
\leanline{6}{{\leanfallback ⟪}θ\ -\ θ',\ F\ θ\ x\ -\ F\ θ'\ x'{\leanfallback ⟫}\_ℝ\ +\ {\leanfallback ⟪}x\ -\ x',\ H\ θ\ x\ -\ H\ θ'\ x'{\leanfallback ⟫}\_ℝ)}
\leanline{4}{(hLip\ :\ LipschitzWith\ L\ (heavyBallProjectedField\ F\ H))\ :}
\leanline{4}{0\ <\ hypocoerciveRate\ γ\ r\ L\ ∧}
\leanline{6}{TwistedJointMonotonicity\ (heavyBallModel\ γ\ F\ H\ σ)\ (hypocoerciveGeometry\ γ\ hγ)}
\leanline{8}{(hypocoerciveRate\ γ\ r\ L)}
\end{leancode}
\leansource{\texttt{exists\_\allowbreak{}p2\_\allowbreak{}nonintegrable\_\allowbreak{}work} \quad PBGSO/\allowbreak{}FiniteWorkScopeAudit.lean, lines 49–66}
\begin{leancode}
\leanline{0}{theorem\ exists\_p2\_nonintegrable\_work\ :}
\leanline{4}{∃\ μ\ :\ Measure\ ℝ,\ IsProbabilityMeasure\ μ\ ∧\ MemLp\ (fun\ x\ :\ ℝ\ =>\ x)\ 2\ μ\ ∧}
\leanline{6}{¬\ Integrable\ (fun\ x\ :\ ℝ\ =>\ x*(drift\ x-drift\ 0))\ μ}
\end{leancode}
\leanlabel{How the Lean proof follows the argument}
The cited criteria derive their matrix and metric comparisons from the primitive assumptions. The finite-domain theorem adds its work-integrability component explicitly. The scope audit proves the quadratic-versus-quartic obstruction, preventing a silent strengthening of the law theorem's domain.
\end{leanbox}
\section{\texorpdfstring{Construct the actual particle forces and contracting metric}{Construct the actual particle forces and contracting metric}}


The force data provide a nonempty cloud encoding and independent globally Lipschitz parameter and sampler forces. Finite averaging gives the actual empirical mean-force Lipschitz bound; pointwise application and integration lift both forces to square-integrable random variables. Their raw numerical updates agree almost surely with the physical particle formulas.



Source joint dissipativity is then applied to empirical measures and transported through the physical/\allowbreak{}raw coordinate identification. Factoring the resulting positive quadratic form constructs a bounded linear map \(C\) and positive \(m\) such that \(m\|z\|\le\|Cz\|\), with drift pairing \(\langle C(x-y),C(f(x)-f(y))\rangle\le-\rho\|C(x-y)\|^2\). Both the momentum and four-coordinate higher-order schemes use this construction. Numerical contractivity is not a premise.

\begin{leanbox}{Lean: \texttt{meanForce\_\allowbreak{}lipschitz}, \texttt{momentum\_\allowbreak{}metric}, \texttt{higherOrder\_\allowbreak{}metric}}
\leansource{\texttt{meanForce\_\allowbreak{}lipschitz} \quad PBGSO/\allowbreak{}ParticleMethodLp.lean, lines 174–192}
\begin{leancode}
\leanline{0}{theorem\ meanForce\_lipschitz\ :\ LipschitzWith\ D.meanLip\ D.meanForce}
\end{leancode}
\leansource{\texttt{momentum\_\allowbreak{}metric} \quad PBGSO/\allowbreak{}ImplementationMetric.lean, lines 96–111}
\begin{leancode}
\leanline{0}{theorem\ momentum\_metric\ [FiniteDimensional\ ℝ\ A]\ [FiniteDimensional\ ℝ\ X]}
\leanline{4}{(γ\ :\ ℝ)\ (σ\ :\ I\ →\ HypocoercivePhase\ X)}
\leanline{4}{(G\ :\ TwistedGeometry\ (HypocoercivePhase\ A)\ (HypocoercivePhase\ X))\ (ρ\ :\ ℝ)}
\leanline{4}{(hmono\ :\ TwistedJointMonotonicity\ (heavyBallModel\ γ\ D.F\ D.G\ σ)\ G\ ρ)\ :}
\leanline{4}{∃\ n\ :\ ℕ,\ ∃\ T\ :\ ((Lp\ A\ 2\ Q\ ×\ Lp\ C\ 2\ Q)\ ×\ (Lp\ A\ 2\ Q\ ×\ Lp\ C\ 2\ Q))\ →L[ℝ]}
\leanline{8}{Lp\ (EuclideanSpace\ ℝ\ (Fin\ n))\ 2\ Q,}
\leanline{6}{(∀\ x,\ (Real.sqrt\ (EmpiricalMetric.lower\ N\ G)/}
\leanline{8}{(‖(momentumPhysicalLpEquiv\ Q\ D).symm.toContinuousLinearMap‖+1))*‖x‖≤‖T\ x‖)\ ∧}
\leanline{6}{(∀\ x,\ ‖T\ x‖≤(Real.sqrt\ (EmpiricalMetric.upper\ N\ G)*}
\leanline{8}{‖(momentumPhysicalLpEquiv\ Q\ D).toContinuousLinearMap‖)*‖x‖)\ ∧}
\leanline{6}{(∀\ x\ y,\ {\leanfallback ⟪}T\ (x-y),T\ ((-coupledForce\ (D.meanForceLp\ Q)\ (D.samplerForceLp\ Q)\ x.2-γ\ •\ x.1,x.1)-}
\leanline{8}{(-coupledForce\ (D.meanForceLp\ Q)\ (D.samplerForceLp\ Q)\ y.2-γ\ •\ y.1,y.1)){\leanfallback ⟫}\_ℝ\ ≤\ -ρ*‖T\ (x-y)‖\textasciicircum{}2)}
\end{leancode}
\leansource{\texttt{higherOrder\_\allowbreak{}metric} \quad PBGSO/\allowbreak{}ImplementationMetric.lean, lines 114–129}
\begin{leancode}
\leanline{0}{theorem\ higherOrder\_metric\ [FiniteDimensional\ ℝ\ A]\ [FiniteDimensional\ ℝ\ X]}
\leanline{4}{(α\ γ\ :\ ℝ)\ (σ\ :\ I\ →\ X)}
\leanline{4}{(G\ :\ TwistedGeometry\ A\ (HigherOrderSampler\ X))\ (ρ\ :\ ℝ)}
\leanline{4}{(hmono\ :\ TwistedJointMonotonicity\ (higherOrderModel\ D.F\ D.G\ α\ γ\ σ)\ G\ ρ)\ :}
\leanline{4}{∃\ n\ :\ ℕ,\ ∃\ T\ :\ HOLState\ (Lp\ A\ 2\ Q)\ (Lp\ C\ 2\ Q)\ →L[ℝ]}
\leanline{8}{Lp\ (EuclideanSpace\ ℝ\ (Fin\ n))\ 2\ Q,}
\leanline{6}{(∀\ x,\ (Real.sqrt\ (EmpiricalMetric.lower\ N\ G)/}
\leanline{8}{(‖(higherOrderPhysicalLpEquiv\ Q\ D).symm.toContinuousLinearMap‖+1))*‖x‖≤‖T\ x‖)\ ∧}
\leanline{6}{(∀\ x,\ ‖T\ x‖≤(Real.sqrt\ (EmpiricalMetric.upper\ N\ G)*}
\leanline{8}{‖(higherOrderPhysicalLpEquiv\ Q\ D).toContinuousLinearMap‖)*‖x‖)\ ∧}
\leanline{6}{(∀\ x\ y,\ {\leanfallback ⟪}T\ (x-y),T\ (holDrift\ (D.meanForceLp\ Q)\ (D.samplerForceLp\ Q)\ α\ γ\ x-}
\leanline{8}{holDrift\ (D.meanForceLp\ Q)\ (D.samplerForceLp\ Q)\ α\ γ\ y){\leanfallback ⟫}\_ℝ\ ≤\ -ρ*‖T\ (x-y)‖\textasciicircum{}2)}
\end{leancode}
\leanlabel{How the Lean proof follows the argument}
\texttt{mean\allowbreak{}Force\_\allowbreak{}lipschitz} derives the empirical bound from the actual force data. \texttt{momentum\_\allowbreak{}metric} and \texttt{higher\allowbreak{}Order\_\allowbreak{}metric} turn empirical source monotonicity into the implemented drift pairing through their coordinate-compatible dynamics. Their returned metric factors supply both coercivity and the dissipation used by the restart estimates.
\end{leanbox}
\section{\texorpdfstring{Prove moments and absorb accumulated local defects}{Prove moments and absorb accumulated local defects}}


For an additive integral trajectory with globally Lipschitz drift, input increments bounded by \(\Xi\sqrt t\) on intervals of length at most one, and drift value \(b_0\) at zero, the integral equation gives radius \(e^{L+1}(\|z\|+b_0+\Xi)\). Compare each unit-time exact trajectory with a zero-start exact trajectory using the same input. Exact drift contraction and a geometric sum give a radius \(R\) uniform over all time.



For a numerical step \(A_k\) and its exact restart \(\Phi_k\), suppose the actual local error is bounded in the transformed norm by \(h\sqrt h(A+B\|Cx\|)\). Then, writing \(e_k=\|C(X_k-Y(kh))\|\), the triangle inequality and \(\|CX_k\|\le e_k+R\) give \(e_{k+1}\le(e^{-\rho h}+Bh\sqrt h)e_k+(A+BR)h\sqrt h\). If \(\rho h\le1\) and \(B\sqrt h\le\rho/4\), the coefficient is at most \(1-\rho h/4\). Summing yields \(e_n\le4(A+BR)\sqrt h/\rho\) for all \(n\), when \(e_0=0\).

\begin{leanbox}{Lean: \texttt{uniform\_\allowbreak{}moment}, \texttt{one\_\allowbreak{}step\_\allowbreak{}same\_\allowbreak{}initial\_\allowbreak{}uniform\_\allowbreak{}norm}}
\leansource{\texttt{uniform\_\allowbreak{}moment} \quad PBGSO/\allowbreak{}UniformIntegralMoments.lean, lines 166–238}
\begin{leancode}
\leanline{0}{theorem\ uniform\_moment\ (f\ :\ E\ →\ E)\ (L\ :\ ℝ≥0)\ (hf\ :\ LipschitzWith\ L\ f)}
\leanline{4}{(C\ :\ E\ →L[ℝ]\ H)\ (ρ\ :\ ℝ)\ (hρ\ :\ 0\ <\ ρ)}
\leanline{4}{(hmono\ :\ ∀\ x\ y,\ {\leanfallback ⟪}C\ (x-y),C\ (f\ x-f\ y){\leanfallback ⟫}\_ℝ\ ≤\ -ρ*‖C\ (x-y)‖\textasciicircum{}2)}
\leanline{4}{(Y\ B\ :\ ℝ\ →\ E)\ (hY\ :\ Continuous\ Y)\ (hB\ :\ Continuous\ B)}
\leanline{4}{(heq\ :\ ∀\ t,\ 0\ ≤\ t\ →\ Y\ t-Y\ 0\ =\ (∫\ r\ in\ 0..t,\ f\ (Y\ r))+B\ t)}
\leanline{4}{(Ξ\ :\ ℝ)\ (hΞ\ :\ 0\ ≤\ Ξ)}
\leanline{4}{(hincrement\ :\ ∀\ a\ s,\ 0\ ≤\ a\ →\ s\ ∈\ Icc\ 0\ 1\ →}
\leanline{6}{‖B\ (a+s)-B\ a‖\ ≤\ Ξ*Real.sqrt\ s)}
\leanline{4}{(t\ :\ ℝ)\ (ht\ :\ 0\ ≤\ t)\ :}
\leanline{4}{‖C\ (Y\ t)‖\ ≤\ ‖C\ (Y\ 0)‖\ +}
\leanline{6}{(‖C‖*radius\ L\ ‖f\ 0‖\ Ξ\ 0)/(1-Real.exp\ (-ρ))\ +}
\leanline{6}{‖C‖*radius\ L\ ‖f\ 0‖\ Ξ\ 0}
\end{leancode}
\leansource{\texttt{one\_\allowbreak{}step\_\allowbreak{}same\_\allowbreak{}initial\_\allowbreak{}uniform\_\allowbreak{}norm} \quad PBGSO/\allowbreak{}UniformDiscrete.lean, lines 170–182}
\begin{leancode}
\leanline{0}{theorem\ one\_step\_same\_initial\_uniform\_norm\ (C\ :\ H\ →L[ℝ]\ K)\ (A\ Φ\ :\ ℕ\ →\ H\ →\ H)\ (X\ Y\ :\ ℕ\ →\ H)}
\leanline{4}{(ρ\ h\ a\ b\ R\ m\ :\ ℝ)\ (hρ\ :\ 0\ <\ ρ)\ (hh\ :\ 0\ <\ h)\ (ha\ :\ 0\ ≤\ a)\ (hb\ :\ 0\ ≤\ b)}
\leanline{4}{(hR\ :\ 0\ ≤\ R)\ (hm\ :\ 0\ <\ m)\ (hcoerc\ :\ ∀\ z,\ m*‖z‖≤‖C\ z‖)}
\leanline{4}{(hρh\ :\ ρ*h≤1)\ (hsmall\ :\ b*Real.sqrt\ h≤ρ/4)}
\leanline{4}{(hX\ :\ ∀\ n,\ X\ (n+1)=A\ n\ (X\ n))\ (hY\ :\ ∀\ n,\ Y\ (n+1)=Φ\ n\ (Y\ n))}
\leanline{4}{(hcontract\ :\ ∀\ n\ x\ y,\ ‖C\ (Φ\ n\ x-Φ\ n\ y)‖\ ≤\ Real.exp\ (-ρ*h)*‖C\ (x-y)‖)}
\leanline{4}{(hdefect\ :\ ∀\ n\ x,\ ‖C\ (A\ n\ x-Φ\ n\ x)‖\ ≤\ (h*Real.sqrt\ h)*(a+b*‖C\ x‖))}
\leanline{4}{(hreference\ :\ ∀\ n,\ ‖C\ (Y\ n)‖\ ≤\ R)\ (hinit\ :\ X\ 0=Y\ 0)\ (n\ :\ ℕ)\ :}
\leanline{4}{‖X\ n-Y\ n‖\ ≤\ ((4*(a+b*R)/ρ)*Real.sqrt\ h)/m}
\end{leancode}
\leanlabel{How the Lean proof follows the argument}
\texttt{uniform\_\allowbreak{}moment} derives the reference radius from the original integral equation and exact-flow contraction. \texttt{one\_\allowbreak{}step\_\allowbreak{}same\_\allowbreak{}initial\_\allowbreak{}uniform\_\allowbreak{}norm} proves the recurrence absorption and divides the final transformed estimate by coercivity. Neither moment bounds nor a final numerical recurrence are assumed in the particle application.
\end{leanbox}
\section{\texorpdfstring{The time-uniform momentum parameter bound}{The time-uniform momentum parameter bound}}


The actual OBABO step is compared with its exact integral restart using its implemented force kicks, position drifts and Gaussian thermostats. Primitive Lipschitz force bounds, Brownian second moments and the derived short-time trajectory modulus give a local defect affine in the initial radius and of order \(h\sqrt h\). The contracting empirical metric and uniform reference radius supply the preceding absorption argument.



Consequently, on a strictly positive step interval computed from the model and \(\rho\), the numerical-to-exact error is at most \(K\sqrt h\) for every iteration. Project to the parameter coordinate and add the exact interacting-particle bound: the total error is numerical bias \(K\sqrt h\), initial-law relaxation \(\sqrt\kappa e^{-\rho nh}d_0\), and particle error \((L_\Psi/\rho)\sqrt{\kappa V/N}(1-e^{-\rho nh})\), with the response factor fixed by this model. All constants are defined from the actual force, metric, initial state and Brownian data.

\begin{leanbox}{Lean: \texttt{obabo\_\allowbreak{}affine\_\allowbreak{}local\_\allowbreak{}defect}, \texttt{explicit\_\allowbreak{}parameter\_\allowbreak{}bound}}
\leansource{\texttt{obabo\_\allowbreak{}affine\_\allowbreak{}local\_\allowbreak{}defect} \quad PBGSO/\allowbreak{}MomentumUniformDefect.lean, lines 66–119}
\begin{leancode}
\leanline{0}{theorem\ obabo\_affine\_local\_defect\ (G\ :\ E\ →\ E)\ (L\ γ\ :\ ℝ≥0)}
\leanline{4}{(hLip\ :\ LipschitzWith\ L\ G)\ (Y\ :\ ℝ\ →\ E\ ×\ E)\ (b\ :\ ℝ\ →\ E)}
\leanline{4}{(hY\ :\ Continuous\ Y)\ (hb\ :\ Continuous\ b)}
\leanline{4}{(Ξ\ β\ h\ :\ ℝ)\ (hΞ\ :\ 0\ ≤\ Ξ)\ (hβ\ :\ 0\ ≤\ β)\ (hh\ :\ 0\ <\ h)\ (hh1\ :\ h\ ≤\ 1)}
\leanline{4}{(heq\ :\ ∀\ t\ ∈\ Icc\ 0\ h,\ Y\ t\ -\ Y\ 0\ =}
\leanline{6}{(∫\ r\ in\ 0..t,\ uniformMomentumDrift\ G\ γ\ (Y\ r))\ +\ (b\ t,\ 0))}
\leanline{4}{(hinput\ :\ ∀\ t\ ∈\ Icc\ 0\ h,\ ‖b\ t‖\ ≤\ β\ *\ Real.sqrt\ t)}
\leanline{4}{(ξ\ ξ'\ :\ E)\ (hξ\ :\ ‖ξ‖\ ≤\ Ξ)\ (hξ'\ :\ ‖ξ'‖\ ≤\ Ξ)}
\leanline{4}{(hbrown\ :\ b\ h\ =\ Real.sqrt\ ((γ\ :\ ℝ)\ *\ h)\ •\ (ξ\ +\ ξ'))\ :}
\leanline{4}{phaseError}
\leanline{6}{(KineticLangevin.obaboStep\ G\ h\ (Real.exp\ (-(γ\ :\ ℝ)\ *\ h\ /\ 2))}
\leanline{8}{(Real.sqrt\ (1\ -\ Real.exp\ (-(γ\ :\ ℝ)\ *\ h)))\ (Y\ 0)\ ξ\ ξ')}
\leanline{6}{(Y\ h)\ ≤}
\leanline{6}{((momentumUniformDefect\ L\ γ\ ‖G\ 0‖\ Ξ\ β).1\ +}
\leanline{8}{(momentumUniformDefect\ L\ γ\ ‖G\ 0‖\ Ξ\ β).2\ *\ ‖Y\ 0‖)\ *\ h\ *\ Real.sqrt\ h}
\end{leancode}
\leansource{\texttt{explicit\_\allowbreak{}parameter\_\allowbreak{}bound} \quad PBGSO/\allowbreak{}FullyUniformMomentum.lean, lines 256–276}
\begin{leancode}
\leanline{0}{theorem\ explicit\_parameter\_bound}
\leanline{4}{(hcloud\ :\ D.cloud\ =\ euclideanParticleEncoding\ N\ J)\ (hγ\ :\ 0\ ≤\ γ)}
\leanline{4}{(star\ :\ A)\ (νStar\ :\ Measure\ (HypocoercivePhase\ (EuclideanSpace\ ℝ\ J)))\ [IsProbabilityMeasure\ νStar]}
\leanline{4}{(hνStar\ :\ MemLp\ (fun\ z\ :\ HypocoercivePhase\ (EuclideanSpace\ ℝ\ J)\ ↦\ z)\ 2\ νStar)}
\leanline{4}{(Frozen\ :\ AdditiveStationary.FrozenTrajectory\ (ULA.CanonicalBrownian.vectorProbability\ J)}
\leanline{6}{(momentumForceDynamics\ D\ γ)\ νStar\ hνStar\ (targetState\ star)\ (input\ γ))}
\leanline{4}{(ρ\ :\ ℝ)\ (hρ\ :\ 0\ <\ ρ)\ (hmono\ :\ TwistedJointMonotonicity\ (model\ D\ γ)\ Geometry\ ρ)}
\leanline{4}{(hg\ :\ (model\ D\ γ).meanGradient\ νStar\ (targetState\ star)\ =\ 0)}
\leanline{4}{(hV\ :\ 0\ ≤\ Variance)}
\leanline{4}{(h\ :\ ℝ)\ (hh\ :\ 0\ <\ h)\ (hhmax\ :\ h\ ≤\ maximumStep\ D\ γ\ Geometry\ μ\ hγ\ ρ)\ (n\ :\ ℕ)\ :}
\leanline{8}{‖(momentumParticleIterate\ (ParticleNoiseSpace.probability\ (J\ :=\ J)\ μ\ N)}
\leanline{10}{(particleBrownian\ (J\ :=\ J)\ μ\ N)\ (particleBrownian\_isPreBrownian\ (J\ :=\ J)\ μ\ N)}
\leanline{10}{D\ γ\ h\ (initial\ D\ γ\ Geometry\ μ\ hμ\ a\ Variance\ hVariance)\ n).2.1\ -}
\leanline{10}{Lp.const\ 2\ (ParticleNoiseSpace.probability\ (J\ :=\ J)\ μ\ N)\ star‖\ ≤}
\leanline{8}{discretizationConstant\ D\ γ\ Geometry\ μ\ hμ\ a\ Variance\ hVariance\ hγ\ ρ\ *\ Real.sqrt\ h\ +}
\leanline{8}{(Real.sqrt\ Geometry.conditionNumber\ *\ Real.exp\ (-ρ\ *\ (n*h))\ *}
\leanline{10}{optimizerLawDistance\ a\ (targetState\ star)\ μ\ νStar\ +}
\leanline{10}{((2+γ)/ρ)\ *\ Real.sqrt\ (Geometry.conditionNumber\ *\ Variance/(N\ :\ ℝ))\ *}
\leanline{12}{(1-Real.exp\ (-ρ\ *\ (n*h))))}
\end{leancode}
\leanlabel{How the Lean proof follows the argument}
\texttt{obabo\_\allowbreak{}affine\_\allowbreak{}local\_\allowbreak{}defect} derives the actual step-versus-integral defect. \texttt{Fully\allowbreak{}Uniform\allowbreak{}Momentum.\allowbreak{}explicit\_\allowbreak{}parameter\_\allowbreak{}bound} supplies the computed admissible step and numerical constant, projects the uniform strong error and adds the constructed exact particle parameter bound by the triangle inequality.
\end{leanbox}
\section{\texorpdfstring{The time-uniform higher-order parameter bound}{The time-uniform higher-order parameter bound}}


For the seven printed coefficients, substitute \(r=e^{-\alpha h}-1+\alpha h\), with \(0\le r\le\alpha^2h^2/2\). This gives all quadratic coefficient remainders explicitly. The exact restart is constructed from the global force bounds and the same physical Brownian path as the Gaussian update. Its short-time radius and drift modulus, together with the actual kernel mismatch, yield local error \((c_0+c_1\|z\|)h\sqrt h\).



Exact restart agreement and drift contraction give the uniform comparison in the preceding step. With coercivity \(m\), data-only metric upper bound \(U\), radius \(R\) and \(B=Uc_1/m\), one may use \(h_{\max}=\min\{1,1/\rho,(\rho/(4(B+1)))^2\}>0\) and \(K=4(Uc_0+BR)/(m\rho)\). Projection and the exact particle estimate give

\[\|\theta_n-a_\star\|_{L^2}\le K\sqrt h+\sqrt\kappa e^{-\rho nh}d_0+\frac1\rho\sqrt{\kappa V/N}(1-e^{-\rho nh}).\]

The rate applies to the printed numerical map. Its coefficients are not treated as an exact fundamental matrix: at unit parameters and unit time the necessary identity \(\mu_{22}'=\mu_{32}\) fails by \(1-3e^{-1}\ne0\). The proven local-error budget explicitly retains the actual coefficient and innovation differences.

\begin{leanbox}{Lean: \texttt{holCoefficient\_\allowbreak{}bounds}, \texttt{higherOrderRestart\_\allowbreak{}contracts}, \texttt{higherOrderRestart\_\allowbreak{}local\_\allowbreak{}error} and 3 more}
\leansource{\texttt{holCoefficient\_\allowbreak{}bounds} \quad PBGSO/\allowbreak{}HigherOrderRates.lean, lines 37–131}
\begin{leancode}
\leanline{0}{theorem\ holCoefficient\_bounds\ (α\ γ\ h\ :\ ℝ)\ (hα\ :\ 0\ <\ α)\ (hγ\ :\ 0\ ≤\ γ)}
\leanline{4}{(hh\ :\ 0\ <\ h)\ (hh1\ :\ h\ ≤\ 1)\ :}
\leanline{4}{HOLCoefficientBounds\ (holCoefficients\ α\ γ\ h)\ (holRemainderConstants\ α\ γ)\ α\ γ\ h}
\end{leancode}
\leansource{\texttt{higherOrderRestart\_\allowbreak{}contracts} \quad PBGSO/\allowbreak{}HOLRestart.lean, lines 183–216}
\begin{leancode}
\leanline{0}{theorem\ higherOrderRestart\_contracts}
\leanline{4}{\{H\ :\ Type*\}\ [NormedAddCommGroup\ H]\ [InnerProductSpace\ ℝ\ H]}
\leanline{4}{(C\ :\ HOLState\ Θ\ (Lp\ (EuclideanSpace\ ℝ\ I)\ 2\ Q)\ →L[ℝ]\ H)\ (ρ\ :\ ℝ)}
\leanline{4}{(hmono\ :\ ∀\ x\ y,}
\leanline{6}{inner\ ℝ\ (C\ (x\ -\ y))\ (C\ (holDrift\ Fbar\ G\ α\ γ\ x\ -\ holDrift\ Fbar\ G\ α\ γ\ y))\ ≤}
\leanline{8}{-ρ\ *\ ‖C\ (x\ -\ y)‖\ \textasciicircum{}\ 2)}
\leanline{4}{(z\ w\ :\ HOLState\ Θ\ (Lp\ (EuclideanSpace\ ℝ\ I)\ 2\ Q))\ (t\ :\ ℝ)\ (ht\ :\ 0\ ≤\ t)\ :}
\leanline{4}{‖C\ (higherOrderRestart\ Q\ hf\ Fbar\ G\ LF\ LG\ α\ γ\ hLF\ hLG\ hF\ hG\ z\ a\ t\ -}
\leanline{8}{higherOrderRestart\ Q\ hf\ Fbar\ G\ LF\ LG\ α\ γ\ hLF\ hLG\ hF\ hG\ w\ a\ t)‖\ ≤}
\leanline{6}{Real.exp\ (-ρ\ *\ t)\ *\ ‖C\ (z\ -\ w)‖}
\end{leancode}
\leansource{\texttt{higherOrderRestart\_\allowbreak{}local\_\allowbreak{}error} \quad PBGSO/\allowbreak{}HOLRestartLocal.lean, lines 60–121}
\begin{leancode}
\leanline{0}{theorem\ higherOrderRestart\_local\_error}
\leanline{4}{(hα\ :\ 0\ <\ α)\ (hγ\ :\ 0\ ≤\ γ)\ (a\ h\ :\ ℝ≥0)}
\leanline{4}{(hh\ :\ 0\ <\ (h\ :\ ℝ))\ (hh1\ :\ (h\ :\ ℝ)\ ≤\ 1)}
\leanline{4}{(z\ :\ HOLState\ Θ\ (Lp\ (EuclideanSpace\ ℝ\ I)\ 2\ Q))\ :}
\leanline{4}{holError\ (higherOrderStep\ Fbar\ G\ α\ γ\ h\ z}
\leanline{8}{(holPullbackVectorNoise\ I\ hf\ α\ γ\ a\ h\ 0,}
\leanline{9}{holPullbackVectorNoise\ I\ hf\ α\ γ\ a\ h\ 1,}
\leanline{9}{holPullbackVectorNoise\ I\ hf\ α\ γ\ a\ h\ 2))}
\leanline{6}{(higherOrderRestart\ Q\ hf\ Fbar\ G\ LF\ LG\ α\ γ\ hLF\ hLG\ hF\ hG\ z\ a\ h)\ ≤}
\leanline{6}{(holUniformLocalConstant0\ I\ α\ γ\ LF\ LG\ ‖Fbar\ 0\ 0‖\ ‖G\ 0\ 0‖}
\leanline{10}{‖holDrift\ Fbar\ G\ α\ γ\ 0‖\ +}
\leanline{8}{holUniformLocalConstant1\ α\ γ\ LF\ LG\ ‖Fbar\ 0\ 0‖\ ‖G\ 0\ 0‖\ *\ ‖z‖)\ *}
\leanline{10}{(h\ :\ ℝ)\ *\ Real.sqrt\ (h\ :\ ℝ)}
\end{leancode}
\leansource{\texttt{higherOrder\_\allowbreak{}time\_\allowbreak{}uniform\_\allowbreak{}of\_\allowbreak{}integral} \quad PBGSO/\allowbreak{}HOLUniform.lean, lines 43–124}
\begin{leancode}
\leanline{0}{theorem\ higherOrder\_time\_uniform\_of\_integral}
\leanline{4}{(hα\ :\ 0\ <\ α)\ (hγ\ :\ 0\ ≤\ γ)}
\leanline{4}{(C\ :\ HOLState\ Θ\ (Lp\ (EuclideanSpace\ ℝ\ I)\ 2\ Q)\ →L[ℝ]\ H)}
\leanline{4}{(m\ ρ\ :\ ℝ)\ (hm\ :\ 0\ <\ m)\ (hρ\ :\ 0\ <\ ρ)}
\leanline{4}{(hcoerc\ :\ ∀\ x,\ m\ *\ ‖x‖\ ≤\ ‖C\ x‖)}
\leanline{4}{(hmono\ :\ ∀\ x\ y,}
\leanline{6}{inner\ ℝ\ (C\ (x\ -\ y))\ (C\ (holDrift\ Fbar\ G\ α\ γ\ x\ -\ holDrift\ Fbar\ G\ α\ γ\ y))\ ≤}
\leanline{8}{-ρ\ *\ ‖C\ (x\ -\ y)‖\ \textasciicircum{}\ 2)}
\leanline{4}{(Y\ :\ ℝ\ →\ HOLState\ Θ\ (Lp\ (EuclideanSpace\ ℝ\ I)\ 2\ Q))\ (hY\ :\ Continuous\ Y)}
\leanline{4}{(heq\ :\ ∀\ t,\ 0\ ≤\ t\ →\ Y\ t\ -\ Y\ 0\ =}
\leanline{6}{(∫\ s\ in\ 0..t,\ holDrift\ Fbar\ G\ α\ γ\ (Y\ s))\ +\ higherOrderRestartInput\ Q\ hf\ α\ 0\ t)}
\leanline{4}{(h\ :\ ℝ≥0)\ (hh\ :\ 0\ <\ (h\ :\ ℝ))\ (hh1\ :\ (h\ :\ ℝ)\ ≤\ 1)}
\leanline{4}{(hρh\ :\ ρ\ *\ (h\ :\ ℝ)\ ≤\ 1)}
\leanline{4}{(hsmall\ :\ (‖C‖\ *\ holUniformLocalConstant1\ α\ γ\ LF\ LG\ ‖Fbar\ 0\ 0‖\ ‖G\ 0\ 0‖\ /\ m)\ *}
\leanline{6}{Real.sqrt\ (h\ :\ ℝ)\ ≤\ ρ\ /\ 4)\ (n\ :\ ℕ)\ :}
\leanline{4}{let\ c₀\ :=\ holUniformLocalConstant0\ I\ α\ γ\ LF\ LG\ ‖Fbar\ 0\ 0‖\ ‖G\ 0\ 0‖\ ‖holDrift\ Fbar\ G\ α\ γ\ 0‖}
\leanline{4}{let\ c₁\ :=\ holUniformLocalConstant1\ α\ γ\ LF\ LG\ ‖Fbar\ 0\ 0‖\ ‖G\ 0\ 0‖}
\leanline{4}{let\ R\ :=\ holUniformReferenceRadius\ (holDrift\ Fbar\ G\ α\ γ)\ (holDriftGlobalLip\ LF\ LG\ α\ γ)}
\leanline{6}{C\ ρ\ (holUniformNoiseRadius\ I\ α)\ (Y\ 0)}
\leanline{4}{‖trajectory\ (higherOrderStep\ Fbar\ G\ α\ γ\ h)}
\leanline{8}{(higherOrderPullbackGaussianNoise\ I\ hf\ α\ γ\ h)\ (Y\ 0)\ n\ -\ Y\ ((n\ :\ ℝ)\ *\ (h\ :\ ℝ))‖\ ≤}
\leanline{6}{((4\ *\ (‖C‖\ *\ c₀\ +\ (‖C‖\ *\ c₁\ /\ m)\ *\ R)\ /\ ρ)\ *\ Real.sqrt\ (h\ :\ ℝ))\ /\ m}
\end{leancode}
\leansource{\texttt{explicit\_\allowbreak{}parameter\_\allowbreak{}bound} \quad PBGSO/\allowbreak{}FullyUniformHigherOrder.lean, lines 261–312}
\begin{leancode}
\leanline{0}{theorem\ explicit\_parameter\_bound\ [FiniteDimensional\ ℝ\ A]}
\leanline{4}{(hcloud\ :\ D.cloud\ =\ euclideanParticleEncoding\ N\ J)}
\leanline{4}{(hα\ :\ 0\ <\ α)\ (hγ\ :\ 0\ ≤\ γ)}
\leanline{4}{(νStar\ :\ Measure\ (HigherOrderSampler\ (EuclideanSpace\ ℝ\ J)))\ [IsProbabilityMeasure\ νStar]}
\leanline{4}{(hνStar\ :\ MemLp\ (fun\ z\ :\ HigherOrderSampler\ (EuclideanSpace\ ℝ\ J)\ ↦\ z)\ 2\ νStar)}
\leanline{4}{(aStar\ :\ A)}
\leanline{4}{(Frozen\ :\ AdditiveStationary.FrozenTrajectory\ (ULA.CanonicalBrownian.vectorProbability\ J)}
\leanline{6}{(higherOrderForceDynamics\ D\ α\ γ)\ νStar\ hνStar\ aStar\ (input\ α))}
\leanline{4}{(ρ\ :\ ℝ)\ (hρ\ :\ 0\ <\ ρ)\ (hmono\ :\ TwistedJointMonotonicity\ (sourceModel\ D\ α\ γ)\ Geo\ ρ)}
\leanline{4}{(hg\ :\ (sourceModel\ D\ α\ γ).meanGradient\ νStar\ aStar\ =\ 0)}
\leanline{4}{(hV\ :\ 0\ ≤\ Variance)}
\leanline{4}{(h\ :\ ℝ≥0)\ (n\ :\ ℕ)\ (hh\ :\ 0\ <\ (h\ :\ ℝ))}
\leanline{4}{(hhstep\ :\ (h\ :\ ℝ)\ ≤\ maximumStep\ D\ α\ γ\ Geo\ μ\ ρ)\ :}
\leanline{8}{‖(higherOrderPullbackParticleIterate\ (Fin\ N\ ×\ J)}
\leanline{10}{(ParticleNoiseSpace.probability\ (J\ :=\ J)\ μ\ N)}
\leanline{10}{(ParticleNoiseSpace.noise\_preserving\ (J\ :=\ J)\ μ\ N)\ D\ α\ γ\ h}
\leanline{10}{(initial\ D\ α\ γ\ Geo\ μ\ hμ\ a\ Variance\ hVariance)\ n).1\ -}
\leanline{10}{Lp.const\ 2\ (ParticleNoiseSpace.probability\ (J\ :=\ J)\ μ\ N)\ aStar‖\ ≤}
\leanline{8}{discretizationConstant\ D\ α\ γ\ Geo\ μ\ hμ\ a\ Variance\ hVariance\ ρ\ *\ Real.sqrt\ (h\ :\ ℝ)\ +}
\leanline{8}{(Real.sqrt\ Geo.conditionNumber\ *\ Real.exp\ (-ρ\ *\ (n\ *\ (h\ :\ ℝ)))\ *}
\leanline{10}{optimizerLawDistance\ a\ aStar\ μ\ νStar\ +}
\leanline{10}{(1\ /\ ρ)\ *\ Real.sqrt\ (Geo.conditionNumber\ *\ Variance\ /\ (N\ :\ ℝ))\ *}
\leanline{12}{(1\ -\ Real.exp\ (-ρ\ *\ (n\ *\ (h\ :\ ℝ)))))}
\end{leancode}
\leansource{\texttt{hol\_\allowbreak{}linear\_\allowbreak{}coefficient\_\allowbreak{}not\_\allowbreak{}exact} \quad PBGSO/\allowbreak{}HOLCoefficientAudit.lean, lines 39–48}
\begin{leancode}
\leanline{0}{theorem\ hol\_linear\_coefficient\_not\_exact\ :}
\leanline{4}{deriv\ (fun\ h\ :\ ℝ\ =>\ (holCoefficients\ 1\ 1\ h).μ22)\ 1\ ≠}
\leanline{6}{(holCoefficients\ 1\ 1\ 1).μ32}
\end{leancode}
\leanlabel{How the Lean proof follows the argument}
The coefficient, restart, local-error and uniform-error theorems form this chain in order. \texttt{Fully\allowbreak{}Uniform\allowbreak{}Higher\allowbreak{}Order.\allowbreak{}explicit\_\allowbreak{}parameter\_\allowbreak{}bound} instantiates it with the physical particle source and adds the exact-particle error. \texttt{hol\_\allowbreak{}linear\_\allowbreak{}coefficient\_\allowbreak{}not\_\allowbreak{}exact} verifies the scope correction by exact real arithmetic.
\end{leanbox}
\section{\texorpdfstring{Verify both continuous and discrete volume statements}{Verify both continuous and discrete volume statements}}


The frictionless phase vector field \((p,q)\mapsto(-F(q),p)\) has zero divergence. A \(C^1\) force and globally defined actual ODE trajectories give joint \(C^1\) dependence by compact-orbit localization and Picard construction. The variational equation and zero divergence imply Jacobian determinant one. Uniqueness gives the inverse negative-time flow, and change of variables then proves Lebesgue preservation.



The discrete B–A–B map is a composition of two force kicks and one position drift. Each is a translation shear, so a measurable force suffices for product-Haar preservation by Tonelli and translation invariance. If the force is differentiable in finite dimension, its actual derivatives are triangular block matrices with identity diagonals; the determinant chain rule gives signed determinant one for the full discrete step. These statements concern the frictionless maps.

\begin{leanbox}{Lean: \texttt{phase\_\allowbreak{}flow\_\allowbreak{}fderiv\_\allowbreak{}det\_\allowbreak{}one\_\allowbreak{}source}, \texttt{phase\_\allowbreak{}flow\_\allowbreak{}measurePreserving\_\allowbreak{}source}, \texttt{interwovenBAB\_\allowbreak{}preserves} and 1 more}
\leansource{\texttt{phase\_\allowbreak{}flow\_\allowbreak{}fderiv\_\allowbreak{}det\_\allowbreak{}one\_\allowbreak{}source} \quad PBGSO/\allowbreak{}FlowSourceRegularity.lean, lines 150–158}
\begin{leancode}
\leanline{0}{theorem\ phase\_flow\_fderiv\_det\_one\_source}
\leanline{4}{(F\ :\ (ι\ →\ ℝ)\ →\ (ι\ →\ ℝ))\ (hF\ :\ ContDiff\ ℝ\ 1\ F)}
\leanline{4}{(Φ\ :\ ℝ\ ×\ Phase\ ι\ →\ Phase\ ι)}
\leanline{4}{(hODE\ :\ ∀\ t\ x,\ HasDerivAt\ (fun\ s\ =>\ Φ\ (s,\ x))\ (phaseField\ F\ (Φ\ (t,\ x)))\ t)}
\leanline{4}{(hinit\ :\ ∀\ x,\ Φ\ (0,\ x)\ =\ x)\ (t\ :\ ℝ)\ (x\ :\ Phase\ ι)\ :}
\leanline{4}{(fderiv\ ℝ\ (fun\ y\ =>\ Φ\ (t,\ y))\ x).det\ =\ 1}
\end{leancode}
\leansource{\texttt{phase\_\allowbreak{}flow\_\allowbreak{}measurePreserving\_\allowbreak{}source} \quad PBGSO/\allowbreak{}ContinuousFlowMeasure.lean, lines 90–100}
\begin{leancode}
\leanline{0}{theorem\ phase\_flow\_measurePreserving\_source}
\leanline{4}{(F\ :\ (ι\ →\ ℝ)\ →\ (ι\ →\ ℝ))\ (hF\ :\ ContDiff\ ℝ\ 1\ F)}
\leanline{4}{(Φ\ :\ ℝ\ ×\ ContinuousVolume.Phase\ ι\ →\ ContinuousVolume.Phase\ ι)}
\leanline{4}{(hODE\ :\ ∀\ t\ x,\ HasDerivAt\ (fun\ s\ =>\ Φ\ (s,\ x))}
\leanline{6}{(ContinuousVolume.phaseField\ F\ (Φ\ (t,\ x)))\ t)}
\leanline{4}{(hinit\ :\ ∀\ x,\ Φ\ (0,\ x)\ =\ x)\ (t\ :\ ℝ)\ :}
\leanline{4}{MeasurePreserving\ (fun\ x\ =>\ Φ\ (t,\ x))\ volume\ volume}
\end{leancode}
\leansource{\texttt{interwovenBAB\_\allowbreak{}preserves} \quad PBGSO/\allowbreak{}Volume.lean, lines 46–49}
\begin{leancode}
\leanline{0}{theorem\ interwovenBAB\_preserves\ (μ\ :\ Measure\ Q)\ [SFinite\ μ]}
\leanline{4}{[Measure.IsAddLeftInvariant\ μ]\ (F\ :\ Q\ →\ Q)\ (hF\ :\ Measurable\ F)\ (h\ :\ ℝ)\ :}
\leanline{4}{MeasurePreserving\ (interwovenBAB\ F\ h)\ (μ.prod\ μ)\ (μ.prod\ μ)}
\end{leancode}
\leansource{\texttt{interwovenBAB\_\allowbreak{}fderiv\_\allowbreak{}det\_\allowbreak{}one} \quad PBGSO/\allowbreak{}DiscreteVolume.lean, lines 130–136}
\begin{leancode}
\leanline{0}{theorem\ interwovenBAB\_fderiv\_det\_one}
\leanline{4}{\{Q\ :\ Type*\}\ [NormedAddCommGroup\ Q]\ [InnerProductSpace\ ℝ\ Q]}
\leanline{4}{[CompleteSpace\ Q]\ [MeasurableSpace\ Q]\ [BorelSpace\ Q]\ [SecondCountableTopology\ Q]}
\leanline{4}{[FiniteDimensional\ ℝ\ Q]\ (F\ :\ Q\ →\ Q)\ (hF\ :\ Differentiable\ ℝ\ F)}
\leanline{4}{(h\ :\ ℝ)\ (z\ :\ Q\ ×\ Q)\ :}
\leanline{4}{(fderiv\ ℝ\ (interwovenBAB\ F\ h)\ z).det\ =\ 1}
\end{leancode}
\leanlabel{How the Lean proof follows the argument}
The continuous results derive the solution-map regularity, variational identity, inverse and measure pushforward in sequence. The discrete measure and determinant statements have separate regularity hypotheses and apply directly to \texttt{interwoven\allowbreak{}BAB}.
\end{leanbox}
\chapter{Supporting proofs}
The following sections prove the results behind the boxes of the proof roadmap, in roadmap order.

\section{\texorpdfstring{Independent optimiser and sampler equations}{Independent optimiser and sampler equations}}


The optimiser state \(s\in S\), reported parameter \(\theta=P(s)\in\Theta\), sampler state \(z\in Z\), and reported position \(x=R(z)\in X\) can have different dimensions. Auxiliary velocities belong to \(S\) or \(Z\). The observation \(F(\theta,x)\in\Theta\), optimiser response \(\Psi(t,s,g)\in S\), sampler drift \(b(\theta,z)\in Z\), and finitely many diffusion columns \(\sigma_i(\theta,z)\in Z\) are independent choices.



For a probability law \(\nu\), define \(G_\nu(s)=\int F(P(s),R(z))\,d\nu(z)\), whenever the integrand is Bochner integrable. The equations are \(\dot s=\Psi(t,s,G_\nu(s))\) and \(dZ=b(P(s),Z)\,dt+\sum_i\sigma_i(P(s),Z)\,dB^i\). The probability and finite-second-moment conditions are separate assumptions.

\subsection{\texorpdfstring{Model and empirical expectation}{Model and empirical expectation}}


\textbf{Claim.} The six maps just specified define the model without a differentiability, convexity or potential assumption. For \(N\) points \(z_1,\ldots,z_N\), put \(\widehat\nu=N^{-1}\sum_i\delta_{z_i}\) and \(\widehat G=N^{-1}\sum_iF(P(s),R(z_i))\). In a complete parameter space with measurable singletons, integration against this finite atomic measure gives \(G_{\widehat\nu}(s)=\widehat G\). If \(N>0\), \(\widehat\nu\) is a probability measure.



\textbf{Proof.} The integral is linear in a finite sum of measures, and \(\int f\,d\delta_z=f(z)\); every function on a singleton is integrable. Multiplication of the measure by the nonnegative scalar \(1/N\) multiplies the integral by that scalar. Thus the integral is the displayed average. Its total mass is \(N^{-1}\sum_{i=1}^N1=N^{-1}N=1\), using \(N>0\). At \(N=0\) the totalised inverse and empty sum both give zero; the identity remains true, but the resulting measure is not a probability law. \(\square\)

\begin{leanbox}{Lean: \texttt{HasFiniteSecondMoment}, \texttt{Model}, \texttt{meanGradient} and 5 more}
\leansource{\texttt{HasFiniteSecondMoment} \quad PBGSO/\allowbreak{}Framework.lean, lines 35–36}
\begin{leancode}
\leanline{0}{def\ HasFiniteSecondMoment\ (ν\ :\ Measure\ Z)\ :\ Prop\ :=}
\leanline{2}{Integrable\ (fun\ z\ =>\ ‖z‖\ \textasciicircum{}\ 2)\ ν}
\end{leancode}
\leansource{\texttt{Model} \quad PBGSO/\allowbreak{}Framework.lean, lines 40–49}
\begin{leancode}
\leanline{0}{structure\ Model\ (S\ Θ\ Z\ X\ I\ :\ Type*)}
\leanline{4}{[NormedAddCommGroup\ S]\ [NormedSpace\ ℝ\ S]}
\leanline{4}{[NormedAddCommGroup\ Θ]\ [NormedSpace\ ℝ\ Θ]}
\leanline{4}{[NormedAddCommGroup\ Z]\ [NormedSpace\ ℝ\ Z]\ where}
\leanline{2}{parameter\ :\ S\ →\ Θ}
\leanline{2}{position\ :\ Z\ →\ X}
\leanline{2}{observation\ :\ Θ\ →\ X\ →\ Θ}
\leanline{2}{optimizer\ :\ ℝ\ →\ S\ →\ Θ\ →\ S}
\leanline{2}{drift\ :\ Θ\ →\ Z\ →\ Z}
\leanline{2}{diffusion\ :\ Θ\ →\ Z\ →\ I\ →\ Z}
\end{leancode}
\leansource{\texttt{meanGradient} \quad PBGSO/\allowbreak{}Framework.lean, lines 56–57}
\begin{leancode}
\leanline{0}{def\ meanGradient\ (ν\ :\ Measure\ Z)\ (s\ :\ S)\ :\ Θ\ :=}
\leanline{2}{∫\ z,\ M.observation\ (M.parameter\ s)\ (M.position\ z)\ ∂ν}
\end{leancode}
\leansource{\texttt{optimizerField} \quad PBGSO/\allowbreak{}Framework.lean, lines 60–61}
\begin{leancode}
\leanline{0}{def\ optimizerField\ (t\ :\ ℝ)\ (s\ :\ S)\ (ν\ :\ Measure\ Z)\ :\ S\ :=}
\leanline{2}{M.optimizer\ t\ s\ (M.meanGradient\ ν\ s)}
\end{leancode}
\leansource{\texttt{particleGradient} \quad PBGSO/\allowbreak{}Framework.lean, lines 66–67}
\begin{leancode}
\leanline{0}{def\ particleGradient\ (N\ :\ ℕ)\ (s\ :\ S)\ (z\ :\ Fin\ N\ →\ Z)\ :\ Θ\ :=}
\leanline{2}{(N\ :\ ℝ)⁻¹\ •\ ∑\ i,\ M.observation\ (M.parameter\ s)\ (M.position\ (z\ i))}
\end{leancode}
\leansource{\texttt{empiricalLaw} \quad PBGSO/\allowbreak{}Framework.lean, lines 71–72}
\begin{leancode}
\leanline{0}{def\ empiricalLaw\ (N\ :\ ℕ)\ (z\ :\ Fin\ N\ →\ Z)\ :\ Measure\ Z\ :=}
\leanline{2}{(N\ :\ ℝ≥0∞)⁻¹\ •\ ∑\ i,\ Measure.dirac\ (z\ i)}
\end{leancode}
\leansource{\texttt{meanGradient\_\allowbreak{}empirical} \quad PBGSO/\allowbreak{}Framework.lean, lines 76–84}
\begin{leancode}
\leanline{0}{theorem\ meanGradient\_empirical\ [MeasurableSingletonClass\ Z]\ [CompleteSpace\ Θ]}
\leanline{4}{(N\ :\ ℕ)\ (s\ :\ S)\ (z\ :\ Fin\ N\ →\ Z)\ :}
\leanline{4}{M.meanGradient\ (empiricalLaw\ N\ z)\ s\ =\ M.particleGradient\ N\ s\ z}
\end{leancode}
\leansource{\texttt{empiricalLaw\_\allowbreak{}probability} \quad PBGSO/\allowbreak{}Framework.lean, lines 98–102}
\begin{leancode}
\leanline{0}{theorem\ empiricalLaw\_probability\ (N\ :\ ℕ)\ (hN\ :\ 0\ <\ N)\ (z\ :\ Fin\ N\ →\ Z)\ :}
\leanline{4}{IsProbabilityMeasure\ (empiricalLaw\ N\ z)}
\end{leancode}
\leanlabel{How the Lean proof follows the argument}
\texttt{Model} records the maps as independent fields. \texttt{Has\allowbreak{}Finite\allowbreak{}Second\allowbreak{}Moment} means integrability of the squared norm and does not include probability normalisation. \texttt{mean\allowbreak{}Gradient} is the Bochner integral; \texttt{optimizer\allowbreak{}Field} substitutes it into the optimiser. \texttt{mean\allowbreak{}Gradient\_\allowbreak{}empirical} unfolds the definitions, uses \texttt{integral\_\allowbreak{}smul\_\allowbreak{}measure} and \texttt{integral\_\allowbreak{}finset\allowbreak{}Sum\_\allowbreak{}measure}, then simplifies the Dirac integrals. \texttt{empirical\allowbreak{}Law\_\allowbreak{}probability} proves total mass one using \texttt{h\allowbreak{}N} to justify cancellation of the particle count. Extended nonnegative scalars in a measure are written \texttt{\ensuremath{\mathbb R}\ensuremath{\ge}0\ensuremath{\infty}}; ordinary vector scalars are real numbers.
\end{leanbox}
\subsection{\texorpdfstring{Stationarity at the weak equation level}{Stationarity at the weak equation level}}


\textbf{Claim.} Choose a family of twice differentiable test functions for which all integrals below are finite. Write \(L_\theta\varphi=D\varphi[b(\theta,\cdot)]+\tfrac12\sum_iD^2\varphi[\sigma_i,\sigma_i]\). Suppose \(\nu_\star\) is a probability law with finite second moment, all observations and tests and generator images are integrable, and \(\int L_{P(s_\star)}\varphi\,d\nu_\star=0\). If \(G_{\nu_\star}(s_\star)=0\) and \(\Psi(t,s_\star,0)=0\) for every \(t\ge0\), then the constant pair \((s_\star,\nu_\star)\) satisfies the optimiser ODE and the weak forward equation. If solutions of precisely these equations with the same initial data are unique, every such solution started at this pair remains equal to it.



\textbf{Proof.} The derivative of the constant optimiser is zero, equal to \(\Psi(t,s_\star,G_{\nu_\star}(s_\star))\) by the two vanishing assumptions. For each test, \(t\mapsto\int\varphi\,d\nu_\star\) is constant, hence has derivative zero; frozen stationarity identifies the right-hand side of its forward equation with zero. All probability, moment and integrability requirements are inherited from the assumed frozen law. Apply the assumed uniqueness to this constant solution and the given solution. Their initial optimiser states and laws agree, so they agree at every nonnegative time. \(\square\)

\begin{leanbox}{Lean: \texttt{generator}, \texttt{IsWeakSolution}, \texttt{FrozenStationary} and 3 more}
\leansource{\texttt{generator} \quad PBGSO/\allowbreak{}Framework.lean, lines 118–121}
\begin{leancode}
\leanline{0}{def\ generator\ (θ\ :\ Θ)\ (φ\ :\ Z\ →\ ℝ)\ (z\ :\ Z)\ :\ ℝ\ :=}
\leanline{2}{fderiv\ ℝ\ φ\ z\ (M.drift\ θ\ z)\ +}
\leanline{4}{(1\ /\ 2\ :\ ℝ)\ *\ ∑\ i,}
\leanline{6}{(fderiv\ ℝ\ (fderiv\ ℝ\ φ)\ z\ (M.diffusion\ θ\ z\ i))\ (M.diffusion\ θ\ z\ i)}
\end{leancode}
\leansource{\texttt{IsWeakSolution} \quad PBGSO/\allowbreak{}Framework.lean, lines 126–139}
\begin{leancode}
\leanline{0}{structure\ IsWeakSolution\ (tests\ :\ Set\ (Z\ →\ ℝ))\ (s\ :\ ℝ\ →\ S)}
\leanline{4}{(ν\ :\ ℝ\ →\ Measure\ Z)\ :\ Prop\ where}
\leanline{2}{probability\ :\ ∀\ t,\ 0\ ≤\ t\ →\ IsProbabilityMeasure\ (ν\ t)}
\leanline{2}{second\_moment\ :\ ∀\ t,\ 0\ ≤\ t\ →\ HasFiniteSecondMoment\ (ν\ t)}
\leanline{2}{gradient\_integrable\ :\ ∀\ t,\ 0\ ≤\ t\ →}
\leanline{4}{Integrable\ (fun\ z\ =>\ M.observation\ (M.parameter\ (s\ t))\ (M.position\ z))\ (ν\ t)}
\leanline{2}{optimizer\_ode\ :\ ∀\ t,\ 0\ ≤\ t\ →}
\leanline{4}{HasDerivWithinAt\ s\ (M.optimizerField\ t\ (s\ t)\ (ν\ t))\ (Ici\ 0)\ t}
\leanline{2}{test\_integrable\ :\ ∀\ φ\ ∈\ tests,\ ∀\ t,\ 0\ ≤\ t\ →\ Integrable\ φ\ (ν\ t)}
\leanline{2}{generator\_integrable\ :\ ∀\ φ\ ∈\ tests,\ ∀\ t,\ 0\ ≤\ t\ →}
\leanline{4}{Integrable\ (M.generator\ (M.parameter\ (s\ t))\ φ)\ (ν\ t)}
\leanline{2}{forward\_equation\ :\ ∀\ φ\ ∈\ tests,\ ∀\ t,\ 0\ ≤\ t\ →}
\leanline{4}{HasDerivWithinAt\ (fun\ r\ =>\ ∫\ z,\ φ\ z\ ∂ν\ r)}
\leanline{6}{(∫\ z,\ M.generator\ (M.parameter\ (s\ t))\ φ\ z\ ∂ν\ t)\ (Ici\ 0)\ t}
\end{leancode}
\leansource{\texttt{FrozenStationary} \quad PBGSO/\allowbreak{}Framework.lean, lines 144–149}
\begin{leancode}
\leanline{0}{structure\ FrozenStationary\ (tests\ :\ Set\ (Z\ →\ ℝ))\ (θ\ :\ Θ)\ (ν\ :\ Measure\ Z)\ :\ Prop\ where}
\leanline{2}{probability\ :\ IsProbabilityMeasure\ ν}
\leanline{2}{second\_moment\ :\ HasFiniteSecondMoment\ ν}
\leanline{2}{test\_integrable\ :\ ∀\ φ\ ∈\ tests,\ Integrable\ φ\ ν}
\leanline{2}{generator\_integrable\ :\ ∀\ φ\ ∈\ tests,\ Integrable\ (M.generator\ θ\ φ)\ ν}
\leanline{2}{generator\_zero\ :\ ∀\ φ\ ∈\ tests,\ (∫\ z,\ M.generator\ θ\ φ\ z\ ∂ν)\ =\ 0}
\end{leancode}
\leansource{\texttt{WeakUniqueness} \quad PBGSO/\allowbreak{}Framework.lean, lines 153–155}
\begin{leancode}
\leanline{0}{def\ WeakUniqueness\ (tests\ :\ Set\ (Z\ →\ ℝ))\ :\ Prop\ :=}
\leanline{2}{∀\ s\ ν\ s'\ ν',\ M.IsWeakSolution\ tests\ s\ ν\ →\ M.IsWeakSolution\ tests\ s'\ ν'\ →}
\leanline{4}{s\ 0\ =\ s'\ 0\ →\ ν\ 0\ =\ ν'\ 0\ →\ ∀\ t,\ 0\ ≤\ t\ →\ s\ t\ =\ s'\ t\ ∧\ ν\ t\ =\ ν'\ t}
\end{leancode}
\leansource{\texttt{constant\_\allowbreak{}isWeakSolution} \quad PBGSO/\allowbreak{}Framework.lean, lines 160–179}
\begin{leancode}
\leanline{0}{theorem\ constant\_isWeakSolution\ (tests\ :\ Set\ (Z\ →\ ℝ))\ (sStar\ :\ S)\ (νStar\ :\ Measure\ Z)}
\leanline{4}{(hν\ :\ M.FrozenStationary\ tests\ (M.parameter\ sStar)\ νStar)}
\leanline{4}{(hint\ :\ Integrable}
\leanline{6}{(fun\ z\ =>\ M.observation\ (M.parameter\ sStar)\ (M.position\ z))\ νStar)}
\leanline{4}{(hg\ :\ M.meanGradient\ νStar\ sStar\ =\ 0)}
\leanline{4}{(hfix\ :\ ∀\ t,\ 0\ ≤\ t\ →\ M.optimizer\ t\ sStar\ 0\ =\ 0)\ :}
\leanline{4}{M.IsWeakSolution\ tests\ (fun\ \_\ =>\ sStar)\ (fun\ \_\ =>\ νStar)}
\end{leancode}
\leansource{\texttt{stationary\_\allowbreak{}of\_\allowbreak{}frozen} \quad PBGSO/\allowbreak{}Framework.lean, lines 183–194}
\begin{leancode}
\leanline{0}{theorem\ stationary\_of\_frozen\ (tests\ :\ Set\ (Z\ →\ ℝ))}
\leanline{4}{(hunique\ :\ M.WeakUniqueness\ tests)\ (sStar\ :\ S)\ (νStar\ :\ Measure\ Z)}
\leanline{4}{(hν\ :\ M.FrozenStationary\ tests\ (M.parameter\ sStar)\ νStar)}
\leanline{4}{(hint\ :\ Integrable}
\leanline{6}{(fun\ z\ =>\ M.observation\ (M.parameter\ sStar)\ (M.position\ z))\ νStar)}
\leanline{4}{(hg\ :\ M.meanGradient\ νStar\ sStar\ =\ 0)}
\leanline{4}{(hfix\ :\ ∀\ t,\ 0\ ≤\ t\ →\ M.optimizer\ t\ sStar\ 0\ =\ 0)}
\leanline{4}{(s\ :\ ℝ\ →\ S)\ (ν\ :\ ℝ\ →\ Measure\ Z)\ (hsol\ :\ M.IsWeakSolution\ tests\ s\ ν)}
\leanline{4}{(hs\ :\ s\ 0\ =\ sStar)\ (hn\ :\ ν\ 0\ =\ νStar)\ (t\ :\ ℝ)\ (ht\ :\ 0\ ≤\ t)\ :}
\leanline{4}{s\ t\ =\ sStar\ ∧\ ν\ t\ =\ νStar}
\end{leancode}
\leanlabel{How the Lean proof follows the argument}
\texttt{generator} uses first and second Fréchet derivatives. \texttt{Is\allowbreak{}Weak\allowbreak{}Solution} lists probability, moment, integrability and right derivative conditions explicitly; \texttt{Has\allowbreak{}Deriv\allowbreak{}Within\allowbreak{}At} on the nonnegative half-line permits the right derivative at zero. \texttt{constant\_\allowbreak{}is\allowbreak{}Weak\allowbreak{}Solution} fills every field, with both derivative proofs obtained from \texttt{has\allowbreak{}Deriv\allowbreak{}Within\allowbreak{}At\_\allowbreak{}const}. \texttt{stationary\_\allowbreak{}of\_\allowbreak{}frozen} applies \texttt{hunique} to the constructed constant solution. This argument concerns the stated weak equation; existence of strong stochastic solutions is not a field of \texttt{Model}.
\end{leanbox}
\section{\texorpdfstring{The exact two-dimensional small-gain rate}{The exact two-dimensional small-gain rate}}


Separate optimiser and sampler dissipation produces a two-dimensional quadratic form in the square roots of the optimiser and sampler coupling energies. Its least eigenvalue is the contraction rate. The result identifies that eigenvalue exactly, including the cross-dependence of both components and the diffusion correction.

\subsection{\texorpdfstring{Smaller eigenvalue and the source criterion}{Smaller eigenvalue and the source criterion}}


\textbf{Claim.} Suppose the optimiser contracts at rate \(r_o\) with the law fixed, responds to a change of mean observation with constant \(L_\Psi\), the observation is sampler-Lipschitz with constant \(L_F\), the frozen sampler contracts at rate \(r_s\), and its parameter and diffusion dependence satisfy the separate bounds with constants \(L_b,L_\sigma\ge0\). All these bounds are in the quadratic geometry above and hold on the probability laws with finite second moment. Set

\[a=r_o-\kappa L_\sigma^2/2,\quad b=r_s-\kappa L_\sigma^2/2,\quad c=\sqrt\kappa(L_\Psi L_F+L_b),\qquad H=\begin{pmatrix}a&-c/2\\-c/2&b\end{pmatrix}.\]

If \(2r_o-\kappa L_\sigma^2>0\), \(2r_s-\kappa L_\sigma^2>0\), and \(\kappa(L_\Psi L_F+L_b)^2<(2r_o-\kappa L_\sigma^2)(2r_s-\kappa L_\sigma^2)\), then \(H\) is positive definite and the joint dissipativity inequality holds with its positive smallest eigenvalue

\[\rho=\frac{a+b-\sqrt{(a-b)^2+c^2}}2.\]

The separate bounds used here mean precisely: fixed-law optimiser work at most \(-r_o|s-s'|_o^2\); optimiser law-response norm at most \(L_\Psi\|G_\nu(s)-G_{\nu'}(s)\|\); observation Lipschitz constant \(L_F\) in the full sampler variable; fixed-parameter weighted sampler work at most \(-r_s|z-z'|_s^2\); drift parameter-change norm at most \(L_b\|s-s'\|\); and \(\sum_i\|\sigma_i(P(s),z)-\sigma_i(P(s'),z')\|^2\le L_\sigma^2(\|s-s'\|^2+\|z-z'\|^2)\). The optimiser bounds hold at every nonnegative time and on every probability law of finite second moment.



\textbf{Proof.} A real symmetric matrix has two real eigenvalues. Its trace and determinant give their sum \(a+b\) and product \(ab-c^2/4\); the quadratic formula therefore gives the stated smaller eigenvalue. The assumptions say \(a,b>0\) and \(c^2<4ab\). Thus \((a-b)^2+c^2<(a+b)^2\), and the positivity of \(a+b\) shows \(\rho>0\). Both eigenvalues are at least \(\rho\), proving positive definiteness.



For an admissible coupling write \(u=|s-s'|_o\) and \(v=(\int|z-z'|_s^2d\Lambda)^{1/2}\). The separate coefficient bounds, the coupling marginal identities and Cauchy–Schwarz give the total work bound \(-au^2-bv^2+cuv\). The matrix spectral bound gives \(au^2+bv^2-cuv\ge\rho(u^2+v^2)\). Negating proves exactly the joint inequality. Observation integrability follows from its Lipschitz bound and finite second moment; sampler signed-work integrability remains among the admissible coupling premises. \(\square\)

\begin{leanbox}{Lean: \texttt{smallGainMatrix}, \texttt{smallGainMatrix\_\allowbreak{}min\_\allowbreak{}eigenvalue}, \texttt{separatedSmallGainMatrix} and 2 more}
\leansource{\texttt{smallGainMatrix} \quad PBGSO/\allowbreak{}SmallGainSpectrum.lean, lines 16–17}
\begin{leancode}
\leanline{0}{def\ smallGainMatrix\ (a\ b\ c\ :\ ℝ)\ :\ Matrix\ (Fin\ 2)\ (Fin\ 2)\ ℝ\ :=}
\leanline{2}{!![a,\ -c\ /\ 2;\ -c\ /\ 2,\ b]}
\end{leancode}
\leansource{\texttt{smallGainMatrix\_\allowbreak{}min\_\allowbreak{}eigenvalue} \quad PBGSO/\allowbreak{}SmallGainSpectrum.lean, lines 25–57}
\begin{leancode}
\leanline{0}{theorem\ smallGainMatrix\_min\_eigenvalue\ (a\ b\ c\ :\ ℝ)\ :}
\leanline{4}{min\ ((smallGainMatrix\_isHermitian\ a\ b\ c).eigenvalues\ 0)}
\leanline{6}{((smallGainMatrix\_isHermitian\ a\ b\ c).eigenvalues\ 1)\ =\ smallGainSpectralRate\ a\ b\ c}
\end{leancode}
\leansource{\texttt{separatedSmallGainMatrix} \quad PBGSO/\allowbreak{}SmallGainSpectrum.lean, lines 78–83}
\begin{leancode}
\leanline{0}{def\ separatedSmallGainMatrix\ (D\ :\ TwistedGeometry\ S\ Z)\ (rO\ rS\ :\ ℝ)}
\leanline{4}{(LΨ\ LF\ Lb\ Lσ\ :\ ℝ≥0)\ :\ Matrix\ (Fin\ 2)\ (Fin\ 2)\ ℝ\ :=}
\leanline{2}{smallGainMatrix}
\leanline{4}{(rO\ -\ D.conditionNumber\ *\ (Lσ\ :\ ℝ)\ \textasciicircum{}\ 2\ /\ 2)}
\leanline{4}{(rS\ -\ D.conditionNumber\ *\ (Lσ\ :\ ℝ)\ \textasciicircum{}\ 2\ /\ 2)}
\leanline{4}{(Real.sqrt\ D.conditionNumber\ *\ ((LΨ\ :\ ℝ)\ *\ (LF\ :\ ℝ)\ +\ (Lb\ :\ ℝ)))}
\end{leancode}
\leansource{\texttt{separatedMinimumEigenvalue} \quad PBGSO/\allowbreak{}SmallGainSpectrum.lean, lines 91–94}
\begin{leancode}
\leanline{0}{def\ separatedMinimumEigenvalue\ (D\ :\ TwistedGeometry\ S\ Z)\ (rO\ rS\ :\ ℝ)}
\leanline{4}{(LΨ\ LF\ Lb\ Lσ\ :\ ℝ≥0)\ :\ ℝ\ :=}
\leanline{2}{min\ ((separatedSmallGainMatrix\_isHermitian\ D\ rO\ rS\ LΨ\ LF\ Lb\ Lσ).eigenvalues\ 0)}
\leanline{4}{((separatedSmallGainMatrix\_isHermitian\ D\ rO\ rS\ LΨ\ LF\ Lb\ Lσ).eigenvalues\ 1)}
\end{leancode}
\leansource{\texttt{separated\_\allowbreak{}source\_\allowbreak{}criterion} \quad PBGSO/\allowbreak{}SmallGainSpectrum.lean, lines 105–126}
\begin{leancode}
\leanline{0}{theorem\ separated\_source\_criterion\ (M\ :\ Model\ S\ Θ\ Z\ X\ I)\ (D\ :\ TwistedGeometry\ S\ Z)}
\leanline{4}{(rO\ rS\ :\ ℝ)\ (LΨ\ LF\ Lb\ Lσ\ :\ ℝ≥0)}
\leanline{4}{(H\ :\ SeparatedAssumptions\ M\ D\ rO\ rS\ LΨ\ LF\ Lb\ Lσ)}
\leanline{4}{(hO\ :\ 0\ <\ 2\ *\ rO\ -\ D.conditionNumber\ *\ (Lσ\ :\ ℝ)\ \textasciicircum{}\ 2)}
\leanline{4}{(hS\ :\ 0\ <\ 2\ *\ rS\ -\ D.conditionNumber\ *\ (Lσ\ :\ ℝ)\ \textasciicircum{}\ 2)}
\leanline{4}{(hdet\ :\ D.conditionNumber\ *\ ((LΨ\ :\ ℝ)\ *\ (LF\ :\ ℝ)\ +\ (Lb\ :\ ℝ))\ \textasciicircum{}\ 2\ <}
\leanline{6}{(2\ *\ rO\ -\ D.conditionNumber\ *\ (Lσ\ :\ ℝ)\ \textasciicircum{}\ 2)\ *}
\leanline{8}{(2\ *\ rS\ -\ D.conditionNumber\ *\ (Lσ\ :\ ℝ)\ \textasciicircum{}\ 2))\ :}
\leanline{4}{(separatedSmallGainMatrix\ D\ rO\ rS\ LΨ\ LF\ Lb\ Lσ).PosDef\ ∧}
\leanline{6}{0\ <\ separatedMinimumEigenvalue\ D\ rO\ rS\ LΨ\ LF\ Lb\ Lσ\ ∧}
\leanline{6}{TwistedJointMonotonicity\ M\ D\ (separatedMinimumEigenvalue\ D\ rO\ rS\ LΨ\ LF\ Lb\ Lσ)}
\end{leancode}
\leanlabel{How the Lean proof follows the argument}
\texttt{small\allowbreak{}Gain\allowbreak{}Matrix\_\allowbreak{}min\_\allowbreak{}eigenvalue} uses the spectral theorem trace and determinant identities, then splits on the order of the two indexed eigenvalues. It identifies the nonnegative square root with their nonnegative difference, so the arbitrary eigenvalue indexing disappears. \texttt{separated\_\allowbreak{}source\_\allowbreak{}criterion} first rewrites the smallest eigenvalue using \texttt{separated\allowbreak{}Minimum\allowbreak{}Eigenvalue\_\allowbreak{}eq\_\allowbreak{}rate}, obtains positivity from \texttt{separated\allowbreak{}Spectral\allowbreak{}Rate\_\allowbreak{}pos}, and uses \texttt{pos\allowbreak{}Def\_\allowbreak{}iff\_\allowbreak{}eigenvalues\_\allowbreak{}pos} for positive definiteness. Its last component is \texttt{separated\_\allowbreak{}joint\_\allowbreak{}monotonicity\_\allowbreak{}spectral}, the established coupling-work estimate from the primitive \texttt{Separated\allowbreak{}Assumptions}.
\end{leanbox}
\section{\texorpdfstring{Quadratic coupling geometry and scalar comparison}{Quadratic coupling geometry and scalar comparison}}


Let \(D_o,D_s\) be symmetric bounded quadratic weights. Assume \(0<\ell\le u\) and \(\ell\|v\|^2\le\langle v,Dv\rangle\le u\|v\|^2\) for each of the two state spaces. Write \(|s|_o^2=\langle s,D_os\rangle\), \(|z|_s^2=\langle z,D_sz\rangle\), and \(\kappa=u/\ell\). For a coupling \(\Lambda\) of sampler laws, the joint energy is \(E=|s-s'|_o^2+\int|z-z'|_s^2\,d\Lambda\).



The joint dissipativity assumption includes both drift work and the diffusion correction arising from a synchronous stochastic coupling. The scalar comparison results apply once a continuous energy and its right derivative have been justified.

\subsection{\texorpdfstring{The complete joint dissipativity assumption}{The complete joint dissipativity assumption}}


\textbf{Claim.} A joint dissipativity rate \(\rho\) means the following inequality for every \(t\ge0\), every optimiser pair, and every probability coupling \(\Lambda\) of finite-second-moment laws \(\nu,\nu'\), whenever both observations, the weighted coupling cost and the signed sampler work are integrable:

\[\begin{aligned} &\langle s-s',D_o(\Psi(t,s,G_\nu(s))-\Psi(t,s',G_{\nu'}(s')))\rangle\\&\quad+\int\left(\langle z-z',D_s(b(P(s),z)-b(P(s'),z'))\rangle+\frac12\sum_i|\sigma_i(P(s),z)-\sigma_i(P(s'),z')|_s^2\right)d\Lambda\le-\rho E.\end{aligned}\]

This is a hypothesis on the actual coefficients and actual expectations, rather than a claim that follows merely from positive weights. Probability couplings have pushforwards \(\Lambda\circ\operatorname{pr}_1^{-1}=\nu\), \(\Lambda\circ\operatorname{pr}_2^{-1}=\nu'\). The finite work assumptions make the signed expression an ordinary finite integral. \(\square\)

\begin{leanbox}{Lean: \texttt{TwistedGeometry}, \texttt{IsCoupling}, \texttt{samplerCouplingWork} and 1 more}
\leansource{\texttt{TwistedGeometry} \quad PBGSO/\allowbreak{}PaperTheorems.lean, lines 33–49}
\begin{leancode}
\leanline{0}{structure\ TwistedGeometry\ (S\ Z\ :\ Type*)}
\leanline{4}{[NormedAddCommGroup\ S]\ [InnerProductSpace\ ℝ\ S]}
\leanline{4}{[NormedAddCommGroup\ Z]\ [InnerProductSpace\ ℝ\ Z]\ where}
\leanline{2}{optimizerWeight\ :\ S\ →L[ℝ]\ S}
\leanline{2}{samplerWeight\ :\ Z\ →L[ℝ]\ Z}
\leanline{2}{lower\ :\ ℝ}
\leanline{2}{upper\ :\ ℝ}
\leanline{2}{lower\_pos\ :\ 0\ <\ lower}
\leanline{2}{lower\_le\_upper\ :\ lower\ ≤\ upper}
\leanline{2}{optimizer\_symmetric\ :\ ∀\ x\ y,}
\leanline{4}{{\leanfallback ⟪}x,\ optimizerWeight\ y{\leanfallback ⟫}\_ℝ\ =\ {\leanfallback ⟪}optimizerWeight\ x,\ y{\leanfallback ⟫}\_ℝ}
\leanline{2}{sampler\_symmetric\ :\ ∀\ x\ y,}
\leanline{4}{{\leanfallback ⟪}x,\ samplerWeight\ y{\leanfallback ⟫}\_ℝ\ =\ {\leanfallback ⟪}samplerWeight\ x,\ y{\leanfallback ⟫}\_ℝ}
\leanline{2}{optimizer\_lower\ :\ ∀\ s,\ lower\ *\ ‖s‖\ \textasciicircum{}\ 2\ ≤\ {\leanfallback ⟪}s,\ optimizerWeight\ s{\leanfallback ⟫}\_ℝ}
\leanline{2}{optimizer\_upper\ :\ ∀\ s,\ {\leanfallback ⟪}s,\ optimizerWeight\ s{\leanfallback ⟫}\_ℝ\ ≤\ upper\ *\ ‖s‖\ \textasciicircum{}\ 2}
\leanline{2}{sampler\_lower\ :\ ∀\ z,\ lower\ *\ ‖z‖\ \textasciicircum{}\ 2\ ≤\ {\leanfallback ⟪}z,\ samplerWeight\ z{\leanfallback ⟫}\_ℝ}
\leanline{2}{sampler\_upper\ :\ ∀\ z,\ {\leanfallback ⟪}z,\ samplerWeight\ z{\leanfallback ⟫}\_ℝ\ ≤\ upper\ *\ ‖z‖\ \textasciicircum{}\ 2}
\end{leancode}
\leansource{\texttt{IsCoupling} \quad PBGSO/\allowbreak{}PaperTheorems.lean, lines 70–71}
\begin{leancode}
\leanline{0}{def\ IsCoupling\ (Λ\ :\ Measure\ (Z\ ×\ Z))\ (ν\ ν'\ :\ Measure\ Z)\ :\ Prop\ :=}
\leanline{2}{Λ.map\ Prod.fst\ =\ ν\ ∧\ Λ.map\ Prod.snd\ =\ ν'}
\end{leancode}
\leansource{\texttt{samplerCouplingWork} \quad PBGSO/\allowbreak{}PaperTheorems.lean, lines 75–81}
\begin{leancode}
\leanline{0}{def\ samplerCouplingWork\ (M\ :\ Model\ S\ Θ\ Z\ X\ I)\ (D\ :\ TwistedGeometry\ S\ Z)}
\leanline{4}{(s\ s'\ :\ S)\ (zz\ :\ Z\ ×\ Z)\ :\ ℝ\ :=}
\leanline{2}{{\leanfallback ⟪}zz.1\ -\ zz.2,\ D.samplerWeight}
\leanline{6}{(M.drift\ (M.parameter\ s)\ zz.1\ -\ M.drift\ (M.parameter\ s')\ zz.2){\leanfallback ⟫}\_ℝ\ +}
\leanline{4}{(1\ /\ 2\ :\ ℝ)\ *\ ∑\ i,}
\leanline{6}{D.samplerSq\ (M.diffusion\ (M.parameter\ s)\ zz.1\ i\ -}
\leanline{8}{M.diffusion\ (M.parameter\ s')\ zz.2\ i)}
\end{leancode}
\leansource{\texttt{TwistedJointMonotonicity} \quad PBGSO/\allowbreak{}PaperTheorems.lean, lines 87–101}
\begin{leancode}
\leanline{0}{def\ TwistedJointMonotonicity\ (M\ :\ Model\ S\ Θ\ Z\ X\ I)}
\leanline{4}{(D\ :\ TwistedGeometry\ S\ Z)\ (ρ\ :\ ℝ)\ :\ Prop\ :=}
\leanline{2}{∀\ t,\ 0\ ≤\ t\ →\ ∀\ s\ s'\ ν\ ν'\ (Λ\ :\ Measure\ (Z\ ×\ Z)),}
\leanline{4}{IsCoupling\ Λ\ ν\ ν'\ →}
\leanline{4}{IsProbabilityMeasure\ ν\ →\ IsProbabilityMeasure\ ν'\ →}
\leanline{4}{HasFiniteSecondMoment\ ν\ →\ HasFiniteSecondMoment\ ν'\ →}
\leanline{4}{Integrable\ (fun\ z\ =>\ M.observation\ (M.parameter\ s)\ (M.position\ z))\ ν\ →}
\leanline{4}{Integrable\ (fun\ z\ =>\ M.observation\ (M.parameter\ s')\ (M.position\ z))\ ν'\ →}
\leanline{4}{Integrable\ (samplerCouplingWork\ M\ D\ s\ s')\ Λ\ →}
\leanline{4}{Integrable\ (fun\ zz\ =>\ D.samplerSq\ (zz.1\ -\ zz.2))\ Λ\ →}
\leanline{4}{{\leanfallback ⟪}s\ -\ s',\ D.optimizerWeight}
\leanline{6}{(M.optimizerField\ t\ s\ ν\ -\ M.optimizerField\ t\ s'\ ν'){\leanfallback ⟫}\_ℝ\ +}
\leanline{6}{(∫\ zz,\ samplerCouplingWork\ M\ D\ s\ s'\ zz\ ∂Λ)\ ≤}
\leanline{6}{-ρ\ *\ (D.optimizerSq\ (s\ -\ s')\ +}
\leanline{8}{∫\ zz,\ D.samplerSq\ (zz.1\ -\ zz.2)\ ∂Λ)}
\end{leancode}
\leanlabel{How the Lean proof follows the argument}
\texttt{Twisted\allowbreak{}Geometry} stores symmetry and both spectral bounds, including \texttt{lower\_\allowbreak{}pos}. \texttt{Is\allowbreak{}Coupling} specifies both marginals by measure pushforward. \texttt{sampler\allowbreak{}Coupling\allowbreak{}Work} adds the weighted drift pairing to one half of the finite diffusion-column square sum. \texttt{Twisted\allowbreak{}Joint\allowbreak{}Monotonicity} quantifies over all the times, states, laws and couplings above. Its integrability premises occur before its inequality; they are not conclusions inferred from that inequality.
\end{leanbox}
\subsection{\texorpdfstring{Exponential comparison of squared energies}{Exponential comparison of squared energies}}


\textbf{Claim.} Suppose \(E\) is continuous on \([0,T]\), has right derivative \(E'(t)\) on \([0,T)\), and \(E'(t)\le-2\rho E(t)\). Then \(E(t)\le E(0)e^{-2\rho t}\) for \(0\le t\le T\); the real number \(\rho\) need not be positive.



\textbf{Proof.} Multiplication by the integrating factor gives \((e^{2\rho t}E(t))'\le0\) at each right differentiation point. The right derivative comparison principle, valid for continuous functions, makes this product nonincreasing. Evaluating at zero yields \(e^{2\rho t}E(t)\le E(0)\). The exponential is strictly positive, so division gives the stated estimate, including the endpoints by continuity. \(\square\)

\begin{leanbox}{Lean: \texttt{energy\_\allowbreak{}decay}}
\leansource{\texttt{energy\_\allowbreak{}decay} \quad PBGSO/\allowbreak{}PaperTheorems.lean, lines 108–117}
\begin{leancode}
\leanline{0}{theorem\ energy\_decay\ \{E\ E'\ :\ ℝ\ →\ ℝ\}\ \{ρ\ T\ :\ ℝ\}}
\leanline{4}{(hcont\ :\ ContinuousOn\ E\ (Icc\ 0\ T))}
\leanline{4}{(hder\ :\ ∀\ t\ ∈\ Ico\ 0\ T,\ HasDerivWithinAt\ E\ (E'\ t)\ (Ici\ t)\ t)}
\leanline{4}{(hdiss\ :\ ∀\ t\ ∈\ Ico\ 0\ T,\ E'\ t\ ≤\ -2\ *\ ρ\ *\ E\ t)}
\leanline{4}{(t\ :\ ℝ)\ (ht\ :\ t\ ∈\ Icc\ 0\ T)\ :}
\leanline{4}{E\ t\ ≤\ E\ 0\ *\ Real.exp\ (-2\ *\ ρ\ *\ t)}
\end{leancode}
\leanlabel{How the Lean proof follows the argument}
\texttt{energy\_\allowbreak{}decay} calls \texttt{le\_\allowbreak{}gronwall\allowbreak{}Bound\_\allowbreak{}of\_\allowbreak{}liminf\_\allowbreak{}deriv\_\allowbreak{}right\_\allowbreak{}le} with coefficient \texttt{-\allowbreak{}2 * \ensuremath{\rho}} and zero forcing. The derivative hypothesis supplies the required lower limit of right difference quotients through \texttt{liminf\_\allowbreak{}right\_\allowbreak{}slope\_\allowbreak{}le}. The final simplification uses \texttt{gronwall\allowbreak{}Bound\_\allowbreak{}\ensuremath{\varepsilon}0}. No differentiability at the right endpoint or nonnegativity of the energy is required by this particular scalar lemma.
\end{leanbox}
\subsection{\texorpdfstring{Forced energy and the square root at zero}{Forced energy and the square root at zero}}


\textbf{Claim.} Suppose \(\rho>0\), \(c\ge0\), \(E\ge0\) is continuous on \([0,T]\), \(E(0)=0\), and its right derivative satisfies \(E'\le-2\rho E+2c\sqrt E\) on \([0,T)\). Then \(\sqrt{E(t)}\le(c/\rho)(1-e^{-\rho t})\).



\textbf{Proof.} Fix \(\varepsilon>0\) and set \(v=\sqrt{E+\varepsilon^2}\). Then \(v\ge\varepsilon\), \(\sqrt E\le v\), and the chain rule is valid because its argument is positive. Hence

\[v'=\frac{E'}{2v}\le-\rho v+\rho\frac{\varepsilon^2}{v}+c\frac{\sqrt E}{v}\le-\rho v+\rho\varepsilon+c.\]

Scalar Grönwall with \(v(0)=\varepsilon\) gives \(v(t)\le\varepsilon+(c/\rho)(1-e^{-\rho t})\). Since \(\sqrt E\le v\), the same upper bound holds for \(\sqrt E\). It holds for every positive \(\varepsilon\), so taking their infimum removes \(\varepsilon\). This proves the estimate without differentiating \(\sqrt E\) at a zero. \(\square\)

\begin{leanbox}{Lean: \texttt{forced\_\allowbreak{}energy\_\allowbreak{}bound}}
\leansource{\texttt{forced\_\allowbreak{}energy\_\allowbreak{}bound} \quad PBGSO/\allowbreak{}PaperTheorems.lean, lines 123–180}
\begin{leancode}
\leanline{0}{theorem\ forced\_energy\_bound\ \{E\ E'\ :\ ℝ\ →\ ℝ\}\ \{ρ\ c\ T\ :\ ℝ\}}
\leanline{4}{(hρ\ :\ 0\ <\ ρ)\ (hc\ :\ 0\ ≤\ c)}
\leanline{4}{(hcont\ :\ ContinuousOn\ E\ (Icc\ 0\ T))}
\leanline{4}{(hder\ :\ ∀\ t\ ∈\ Ico\ 0\ T,\ HasDerivWithinAt\ E\ (E'\ t)\ (Ici\ t)\ t)}
\leanline{4}{(hE\ :\ ∀\ t\ ∈\ Icc\ 0\ T,\ 0\ ≤\ E\ t)\ (hzero\ :\ E\ 0\ =\ 0)}
\leanline{4}{(hdiss\ :\ ∀\ t\ ∈\ Ico\ 0\ T,}
\leanline{6}{E'\ t\ ≤\ -2\ *\ ρ\ *\ E\ t\ +\ 2\ *\ c\ *\ Real.sqrt\ (E\ t))}
\leanline{4}{(t\ :\ ℝ)\ (ht\ :\ t\ ∈\ Icc\ 0\ T)\ :}
\leanline{4}{Real.sqrt\ (E\ t)\ ≤\ c\ /\ ρ\ *\ (1\ -\ Real.exp\ (-ρ\ *\ t))}
\end{leancode}
\leanlabel{How the Lean proof follows the argument}
\texttt{forced\_\allowbreak{}energy\_\allowbreak{}bound} begins with \texttt{le\_\allowbreak{}of\_\allowbreak{}forall\_\allowbreak{}pos\_\allowbreak{}le\_\allowbreak{}add}, so it suffices to prove the conclusion with an arbitrary positive error. It defines \texttt{u} as the regularised square root, proves its positivity, squared identity and lower bound by the error, and applies right derivative Grönwall to that function. The order argument at the start then removes the regularisation. The hypotheses \texttt{h\allowbreak{}E} and \texttt{hzero} are used to justify square roots and their initial value.
\end{leanbox}
\section{\texorpdfstring{Projected monotonicity in canonical block coordinates}{Projected monotonicity in canonical block coordinates}}


Split the optimiser into a parameter and auxiliary variable, \(s=(\theta,m)\), and the sampler into a position and auxiliary variable, \(z=(x,v)\). The parameter and position spaces are nonempty finite-dimensional Euclidean spaces; either auxiliary dimension may be zero. Positive definite matrices \(D_o,D_s\) define the complete-state metric, while positive definite auxiliary matrices \(D_m,D_v\) describe auxiliary dissipation.

\subsection{\texorpdfstring{Block dissipation and generalized eigenvalues}{Block dissipation and generalized eigenvalues}}


\textbf{Claim.} Assume the parameter and position maps are the first block projections. Suppose the projected joint force is strongly monotone with constant \(r\), and the complete drift satisfies the projected auxiliary work inequality with constants \(\beta\ge0,A_o,A_s\). Suppose \(c=\beta r-A_o-A_s>0\), all four matrices above are positive definite, sampler diffusion is constant, and each admissible coupling makes every projected work expression integrable. Define \(K_o=\operatorname{diag}(cI,D_m)\), \(K_s=\operatorname{diag}(cI,D_v)\), and

\[\rho=\min\{\lambda_{\min}(D_o^{-1/2}K_oD_o^{-1/2}),\lambda_{\min}(D_s^{-1/2}K_sD_s^{-1/2})\}.\]

Then \(\rho>0\) and the complete model satisfies joint dissipativity at rate \(\rho\).



More explicitly, the projected monotonicity assumption is \(\langle\theta-\theta',F(\theta,x)-F(\theta',x')\rangle+\langle x-x',H(\theta,x)-H(\theta',x')\rangle\ge r(\|\theta-\theta'\|^2+\|x-x'\|^2)\) pointwise. For an admissible coupling write \(Q=\|\theta-\theta'\|^2+\int\|x-x'\|^2d\Lambda\), \(A_m=\langle m-m',D_m(m-m')\rangle\), \(A_v=\int\langle v-v',D_v(v-v')\rangle d\Lambda\), \(P_o=\langle\theta-\theta',G_\nu(s)-G_{\nu'}(s')\rangle\), and \(P_s=\int\langle x-x',H(\theta,x)-H(\theta',x')\rangle d\Lambda\). The two auxiliary inequalities say that the actual weighted optimiser work is at most \(-\beta P_o-A_m+A_oQ\), and the integrated weighted sampler drift work is at most \(-\beta P_s-A_v+A_sQ\). Admissibility includes the probability and marginal conditions, finite second moments, integrable observations, signed sampler work and metric cost. The additional projected-integrability hypothesis makes the position square, score work defining \(P_s\), and auxiliary square defining \(A_v\) integrable for every such coupling.



\textbf{Proof.} For \((q,p)\ne0\), \(c\|q\|^2+p^TD_mp>0\), since either \(q\ne0\) or \(p\ne0\). This also covers a zero auxiliary dimension. Thus both block matrices are positive definite. Their congruences by the invertible positive square roots of the metric matrices remain positive definite, so their least eigenvalues and their minimum are positive.



The projected monotonicity and auxiliary work assumptions sum to minus the two block forms. The constant diffusion columns cancel between the two synchronously coupled samplers, so the diffusion-difference square term vanishes. The spectral inequality \(D^{-1/2}KD^{-1/2}\ge\rho I\), multiplied on both sides by \(D^{1/2}\), gives \(K\ge\rho D\). Applying it to the optimiser difference and integrating its sampler version bounds the work by minus \(\rho\) times the complete coupling energy. The prescribed integrability justifies the sum and these integral inequalities. \(\square\)

\begin{leanbox}{Lean: \texttt{blockK}, \texttt{canonical\_\allowbreak{}projected\_\allowbreak{}source\_\allowbreak{}criterion}}
\leansource{\texttt{blockK} \quad PBGSO/\allowbreak{}CanonicalProjectedCriterion.lean, lines 163–164}
\begin{leancode}
\leanline{0}{def\ blockK\ (c\ :\ ℝ)\ (A\ :\ Matrix\ u\ u\ ℝ)\ :\ Matrix\ (a\ ⊕\ u)\ (a\ ⊕\ u)\ ℝ\ :=}
\leanline{2}{Matrix.fromBlocks\ (c\ •\ (1\ :\ Matrix\ a\ a\ ℝ))\ 0\ 0\ A}
\end{leancode}
\leansource{\texttt{canonical\_\allowbreak{}projected\_\allowbreak{}source\_\allowbreak{}criterion} \quad PBGSO/\allowbreak{}CanonicalProjectedCriterion.lean, lines 264–295}
\begin{leancode}
\leanline{0}{theorem\ canonical\_projected\_source\_criterion}
\leanline{4}{(M\ :\ Model\ (BlockSpace\ a\ u)\ (EuclideanSpace\ ℝ\ a)\ (BlockSpace\ x\ v)\ (EuclideanSpace\ ℝ\ x)\ I)}
\leanline{4}{(hparameter\ :\ M.parameter\ =\ WithLp.fst)\ (hposition\ :\ M.position\ =\ WithLp.fst)}
\leanline{4}{(H\ :\ EuclideanSpace\ ℝ\ a\ →\ EuclideanSpace\ ℝ\ x\ →\ EuclideanSpace\ ℝ\ x)}
\leanline{4}{(DOpt\ :\ Matrix\ (a\ ⊕\ u)\ (a\ ⊕\ u)\ ℝ)\ (DSam\ :\ Matrix\ (x\ ⊕\ v)\ (x\ ⊕\ v)\ ℝ)}
\leanline{4}{(Dm\ :\ Matrix\ u\ u\ ℝ)\ (Dv\ :\ Matrix\ v\ v\ ℝ)}
\leanline{4}{(hDOpt\ :\ DOpt.PosDef)\ (hDSam\ :\ DSam.PosDef)\ (hDm\ :\ Dm.PosDef)\ (hDv\ :\ Dv.PosDef)}
\leanline{4}{(β\ r\ AOpt\ ASam\ :\ ℝ)\ (hβ\ :\ 0\ ≤\ β)\ (hc\ :\ 0\ <\ β\ *\ r\ -\ AOpt\ -\ ASam)}
\leanline{4}{(hmono\ :\ (canonicalData\ M\ hparameter\ hposition\ H\ Dm\ Dv).StrongMonotonicity\ r)}
\leanline{4}{(haux\ :\ ProjectedAuxiliaryDissipation\ M\ (canonicalGeometry\ DOpt\ DSam\ hDOpt\ hDSam)}
\leanline{6}{(canonicalData\ M\ hparameter\ hposition\ H\ Dm\ Dv)\ β\ AOpt\ ASam)}
\leanline{4}{(σ\ :\ I\ →\ BlockSpace\ x\ v)\ (hσ\ :\ ∀\ θ\ z,\ M.diffusion\ θ\ z\ =\ σ)}
\leanline{4}{(hInt\ :\ ∀\ s\ s'\ ν\ ν'\ Λ,}
\leanline{6}{ProjectedAdmissibleCoupling\ M\ (canonicalGeometry\ DOpt\ DSam\ hDOpt\ hDSam)\ s\ s'\ ν\ ν'\ Λ\ →}
\leanline{6}{(canonicalData\ M\ hparameter\ hposition\ H\ Dm\ Dv).Integrability\ s\ s'\ Λ)\ :}
\leanline{4}{let\ ρ\ :=\ min}
\leanline{6}{(generalizedMinimum\ hDOpt\ (blockK\_posDef\ (a\ :=\ a)\ \_\ hc\ Dm\ hDm))}
\leanline{6}{(generalizedMinimum\ hDSam\ (blockK\_posDef\ (a\ :=\ x)\ \_\ hc\ Dv\ hDv))}
\leanline{4}{0\ <\ ρ\ ∧\ TwistedJointMonotonicity\ M\ (canonicalGeometry\ DOpt\ DSam\ hDOpt\ hDSam)\ ρ}
\end{leancode}
\leanlabel{How the Lean proof follows the argument}
\texttt{block\allowbreak{}K} is the literal block matrix. \texttt{canonical\_\allowbreak{}projected\_\allowbreak{}source\_\allowbreak{}criterion} applies \texttt{projected\_\allowbreak{}twisted\allowbreak{}Joint\allowbreak{}Monotonicity\_\allowbreak{}spectral} to the constructed \texttt{canonical\allowbreak{}Geometry} and \texttt{canonical\allowbreak{}Data}. The four final goals identify their quadratic forms with the supplied matrices and with \texttt{block\allowbreak{}K\_\allowbreak{}coordinate\_\allowbreak{}form}. The hypotheses \texttt{h\allowbreak{}Int}, \texttt{hmono} and \texttt{haux} retain the substantive integrability and work assumptions. The coordinate maps and metric comparisons are derived; they are not additional assumptions.
\end{leanbox}
\section{\texorpdfstring{Remainders of the seven printed coefficients}{Remainders of the seven printed coefficients}}


The seven coefficients in the higher-order numerical map are explicit real expressions in \(\alpha,\gamma,h\) and \(e^{-\alpha h}\). The coefficient estimates below follow from one nonnegative exponential remainder. They concern the printed numerical map itself, without identifying it with an exact linear flow.

\subsection{\texorpdfstring{Quadratic remainders with explicit constants}{Quadratic remainders with explicit constants}}


\textbf{Claim.} Suppose \(\alpha>0\), \(\gamma\ge0\), and \(0<h\le1\). The printed coefficients satisfy

\[\begin{aligned}|\mu_{12}-h|&\le c_{12}h^2,& |\mu_{13}|&\le c_{13}h^2,&|\mu_{22}-1|&\le c_{22}h^2,\\|\mu_{23}-\gamma h|&\le c_{23}h^2,&|\mu_{31}h|&\le c_{31}h^2,&|\mu_{31}|&\le c_{31}h,\\|\mu_{32}+\gamma h|&\le c_{32}h^2,&|\mu_{33}-1+\alpha h|&\le c_{33}h^2,\end{aligned}\]

where \(c_{12}=\gamma^2/(2\alpha)\), \(c_{13}=c_{31}=\gamma/2\), \(c_{22}=\gamma^2/2\), \(c_{23}=\gamma\alpha/2\), and

\[\begin{aligned}c_{32}&=\gamma^3/\alpha+(\gamma^3/\alpha^2+2\gamma^3/\alpha^3+\gamma/\alpha)\alpha^2/2,\\c_{33}&=\gamma^2+(1+\gamma^2/\alpha+\gamma^2/\alpha^2)\alpha^2/2.\end{aligned}\]

\textbf{Proof.} Put \(r=e^{-\alpha h}-1+\alpha h\). The elementary exponential bounds \(e^{-x}\ge1-x\) and \(e^{-x}\le1-x+x^2/2\), for \(x\ge0\), give \(0\le r\le\alpha^2h^2/2\). Substitution of \(e^{-\alpha h}=1-\alpha h+r\) gives exactly

\[\begin{aligned}\mu_{12}-h&=-\gamma^2h^2/(2\alpha)+\gamma^2r/\alpha^3,&\mu_{13}&=\gamma r/\alpha^2,\\\mu_{22}-1&=-\gamma^2r/\alpha^2,&\mu_{23}-\gamma h&=-\gamma r/\alpha,\\\mu_{31}h&=\gamma r/\alpha^2,\\\mu_{32}+\gamma h&=-\gamma^3h^2/\alpha+(\gamma^3h/\alpha^2+2\gamma^3/\alpha^3+\gamma/\alpha)r,\\\mu_{33}-1+\alpha h&=-\gamma^2h^2+(1+\gamma^2h/\alpha+\gamma^2/\alpha^2)r.\end{aligned}\]

Every multiplier of \(r\) is nonnegative. Its upper bound proves the first six estimates; the bound on \(\mu_{31}\) follows by dividing the bound on \(\mu_{31}h\) by positive \(h\). For the last two estimates use \(h\le1\) in the multipliers, and bound the absolute value of a difference of two nonnegative terms by the sum of their upper bounds. This yields exactly the displayed conservative constants. \(\square\)

\begin{leanbox}{Lean: \texttt{holRemainderConstants}, \texttt{HOLCoefficientBounds}, \texttt{holCoefficient\_\allowbreak{}bounds}}
\leansource{\texttt{holRemainderConstants} \quad PBGSO/\allowbreak{}HigherOrderRates.lean, lines 14–22}
\begin{leancode}
\leanline{0}{def\ holRemainderConstants\ (α\ γ\ :\ ℝ)\ :\ HOLCoefficients\ where}
\leanline{2}{μ12\ :=\ γ\ \textasciicircum{}\ 2\ /\ (2\ *\ α)}
\leanline{2}{μ13\ :=\ γ\ /\ 2}
\leanline{2}{μ22\ :=\ γ\ \textasciicircum{}\ 2\ /\ 2}
\leanline{2}{μ23\ :=\ γ\ *\ α\ /\ 2}
\leanline{2}{μ31\ :=\ γ\ /\ 2}
\leanline{2}{μ32\ :=\ γ\ \textasciicircum{}\ 3\ /\ α\ +}
\leanline{4}{(γ\ \textasciicircum{}\ 3\ /\ α\ \textasciicircum{}\ 2\ +\ 2\ *\ γ\ \textasciicircum{}\ 3\ /\ α\ \textasciicircum{}\ 3\ +\ γ\ /\ α)\ *\ α\ \textasciicircum{}\ 2\ /\ 2}
\leanline{2}{μ33\ :=\ γ\ \textasciicircum{}\ 2\ +\ (1\ +\ γ\ \textasciicircum{}\ 2\ /\ α\ +\ γ\ \textasciicircum{}\ 2\ /\ α\ \textasciicircum{}\ 2)\ *\ α\ \textasciicircum{}\ 2\ /\ 2}
\end{leancode}
\leansource{\texttt{HOLCoefficientBounds} \quad PBGSO/\allowbreak{}HigherOrderRates.lean, lines 26–34}
\begin{leancode}
\leanline{0}{structure\ HOLCoefficientBounds\ (μ\ c\ :\ HOLCoefficients)\ (α\ γ\ h\ :\ ℝ)\ :\ Prop\ where}
\leanline{2}{bound12\ :\ |μ.μ12\ -\ h|\ ≤\ c.μ12\ *\ h\ \textasciicircum{}\ 2}
\leanline{2}{bound13\ :\ |μ.μ13|\ ≤\ c.μ13\ *\ h\ \textasciicircum{}\ 2}
\leanline{2}{bound22\ :\ |μ.μ22\ -\ 1|\ ≤\ c.μ22\ *\ h\ \textasciicircum{}\ 2}
\leanline{2}{bound23\ :\ |μ.μ23\ -\ γ\ *\ h|\ ≤\ c.μ23\ *\ h\ \textasciicircum{}\ 2}
\leanline{2}{bound31\ :\ |μ.μ31\ *\ h|\ ≤\ c.μ31\ *\ h\ \textasciicircum{}\ 2}
\leanline{2}{bound31'\ :\ |μ.μ31|\ ≤\ c.μ31\ *\ h}
\leanline{2}{bound32\ :\ |μ.μ32\ +\ γ\ *\ h|\ ≤\ c.μ32\ *\ h\ \textasciicircum{}\ 2}
\leanline{2}{bound33\ :\ |μ.μ33\ -\ 1\ +\ α\ *\ h|\ ≤\ c.μ33\ *\ h\ \textasciicircum{}\ 2}
\end{leancode}
\leansource{\texttt{holCoefficient\_\allowbreak{}bounds} \quad PBGSO/\allowbreak{}HigherOrderRates.lean, lines 37–131}
\begin{leancode}
\leanline{0}{theorem\ holCoefficient\_bounds\ (α\ γ\ h\ :\ ℝ)\ (hα\ :\ 0\ <\ α)\ (hγ\ :\ 0\ ≤\ γ)}
\leanline{4}{(hh\ :\ 0\ <\ h)\ (hh1\ :\ h\ ≤\ 1)\ :}
\leanline{4}{HOLCoefficientBounds\ (holCoefficients\ α\ γ\ h)\ (holRemainderConstants\ α\ γ)\ α\ γ\ h}
\end{leancode}
\leanlabel{How the Lean proof follows the argument}
\texttt{hol\allowbreak{}Coefficient\_\allowbreak{}bounds} defines \texttt{r}, proves \texttt{hr} and \texttt{hr'} from the elementary exponential estimates, and records the seven exact substitutions as \texttt{e12} through \texttt{e33}. The subsequent \texttt{b12} through \texttt{b33} apply the nonnegative remainder bound. For \texttt{bound31'} it cancels the positive step. The final tuple fills the eight fields of \texttt{HOLCoefficient\allowbreak{}Bounds} using \texttt{abs\_\allowbreak{}le} or known signs. The coefficient identities are polynomial/\allowbreak{}rational simplifications after the substitution, not an assumed matrix-exponential identity.
\end{leanbox}
\section{\texorpdfstring{Measure preservation of the frictionless splitting}{Measure preservation of the frictionless splitting}}


Use velocity-first coordinates \((p,q)\). A force kick \(K_h(p,q)=(p-hF(q),q)\) and position drift \(A_h(p,q)=(p,q+hp)\) are translations of one coordinate while the other is fixed. Their composition \(K_{h/2}\circ A_h\circ K_{h/2}\) is the frictionless B–A–B part of the algorithm.

\subsection{\texorpdfstring{Translation invariance of each shear}{Translation invariance of each shear}}


\textbf{Claim.} Let \(Q\) be a complete real Hilbert space with its Borel measurable structure and a second-countable topology. Let \(\mu\) be an S-finite translation-invariant measure on \(Q\). For any measurable \(F:Q\to Q\) and any real step \(h\), both shears and their B–A–B composition preserve \(\mu\otimes\mu\).



\textbf{Proof.} For a nonnegative measurable test \(g\), Tonelli gives

\[\int g(p-hF(q),q)\,d\mu(p)d\mu(q)=\int\left[\int g(p,q)\,d\mu(p)\right]d\mu(q),\]

because for each fixed \(q\) the substitution \(p\mapsto p-hF(q)\) is a translation. This proves the kick assertion. Reverse the integration order and translate \(q\mapsto q+hp\) to prove the drift assertion. The assumptions on \(F\) ensure joint measurability of the kick, while the drift is continuous. A composition of measure-preserving maps is measure preserving, so the three-stage map preserves product measure. S-finiteness is the measure-theoretic hypothesis used for product integration. Neither differentiability nor a gradient representation of \(F\) is needed. Thermostats and friction are not among these three maps. \(\square\)

\begin{leanbox}{Lean: \texttt{forceKick}, \texttt{positionDrift}, \texttt{interwovenBAB} and 3 more}
\leansource{\texttt{forceKick} \quad PBGSO/\allowbreak{}Volume.lean, lines 20–21}
\begin{leancode}
\leanline{0}{def\ forceKick\ (F\ :\ Q\ →\ Q)\ (h\ :\ ℝ)\ :\ Q\ ×\ Q\ →\ Q\ ×\ Q\ :=}
\leanline{2}{KineticLangevin.kickStep\ F\ h}
\end{leancode}
\leansource{\texttt{positionDrift} \quad PBGSO/\allowbreak{}Volume.lean, lines 24–25}
\begin{leancode}
\leanline{0}{def\ positionDrift\ (h\ :\ ℝ)\ :\ Q\ ×\ Q\ →\ Q\ ×\ Q\ :=}
\leanline{2}{KineticLangevin.driftStep\ h}
\end{leancode}
\leansource{\texttt{interwovenBAB} \quad PBGSO/\allowbreak{}Volume.lean, lines 29–30}
\begin{leancode}
\leanline{0}{def\ interwovenBAB\ (F\ :\ Q\ →\ Q)\ (h\ :\ ℝ)\ :\ Q\ ×\ Q\ →\ Q\ ×\ Q\ :=}
\leanline{2}{KineticLangevin.verletStep\ F\ h}
\end{leancode}
\leansource{\texttt{forceKick\_\allowbreak{}preserves} \quad PBGSO/\allowbreak{}Volume.lean, lines 33–36}
\begin{leancode}
\leanline{0}{theorem\ forceKick\_preserves\ (μ\ :\ Measure\ Q)\ [SFinite\ μ]\ [Measure.IsAddLeftInvariant\ μ]}
\leanline{4}{(F\ :\ Q\ →\ Q)\ (hF\ :\ Measurable\ F)\ (h\ :\ ℝ)\ :}
\leanline{4}{MeasurePreserving\ (forceKick\ F\ h)\ (μ.prod\ μ)\ (μ.prod\ μ)}
\end{leancode}
\leansource{\texttt{positionDrift\_\allowbreak{}preserves} \quad PBGSO/\allowbreak{}Volume.lean, lines 39–42}
\begin{leancode}
\leanline{0}{theorem\ positionDrift\_preserves\ (μ\ :\ Measure\ Q)\ [SFinite\ μ]}
\leanline{4}{[Measure.IsAddLeftInvariant\ μ]\ (h\ :\ ℝ)\ :}
\leanline{4}{MeasurePreserving\ (positionDrift\ h)\ (μ.prod\ μ)\ (μ.prod\ μ)}
\end{leancode}
\leansource{\texttt{interwovenBAB\_\allowbreak{}preserves} \quad PBGSO/\allowbreak{}Volume.lean, lines 46–49}
\begin{leancode}
\leanline{0}{theorem\ interwovenBAB\_preserves\ (μ\ :\ Measure\ Q)\ [SFinite\ μ]}
\leanline{4}{[Measure.IsAddLeftInvariant\ μ]\ (F\ :\ Q\ →\ Q)\ (hF\ :\ Measurable\ F)\ (h\ :\ ℝ)\ :}
\leanline{4}{MeasurePreserving\ (interwovenBAB\ F\ h)\ (μ.prod\ μ)\ (μ.prod\ μ)}
\end{leancode}
\leanlabel{How the Lean proof follows the argument}
\texttt{force\allowbreak{}Kick}, \texttt{position\allowbreak{}Drift} and \texttt{interwoven\allowbreak{}BAB} are the exact kinetic kick, drift and Verlet maps. The three preservation theorems apply the previously proved \texttt{Kinetic\allowbreak{}Langevin} shear and composition results. \texttt{Measurable F} is the only force regularity premise in this measure-theoretic argument; \texttt{Measure.\allowbreak{}Is\allowbreak{}Add\allowbreak{}Left\allowbreak{}Invariant} supplies translation invariance, and \texttt{SFinite} is retained explicitly.
\end{leanbox}
\section{\texorpdfstring{Construction of the local Brownian integral}{Construction of the local Brownian integral}}


Work on a probability space with a specified finite family of filtered Brownian coordinates. Let \(\sigma_i(t,\omega)\) be jointly strongly measurable, adapted diffusion columns with pathwise locally integrable squared size. A nonnegative progressive energy bound \(q\) satisfies \(\sum_i\|\sigma_i\|^2\le q\) and \(\int_0^Tq(t,\omega)dt<\infty\) for every path. Expected energy need not be finite before stopping.

\subsection{\texorpdfstring{First-hit localization from primitive coefficients}{First-hit localization from primitive coefficients}}


\textbf{Claim.} Under the preceding hypotheses, put \(A_t=\int_0^{\min(t,T)}q(r)dr\) and let \(\tau_n\) be its first hit of level \(n\), capped at \(T\). Then \(A\) is continuous, nondecreasing and adapted, \(\tau_n\) is a stopping time, \(\int_0^{\tau_n}q\le n\), and each stopped column \(1_{r\le\tau_n}\sigma_i(r)\) is square integrable on \([0,T]\times\Omega\). Its Brownian integral \(M^n\) has a continuous martingale version with \(\mathbb E\|M_T^n\|^2\le n\).



\textbf{Proof.} Pathwise integrability and nonnegativity make the indefinite integral continuous and monotone. Progressive measurability makes \(A_t\) measurable in the time-\(t\) past, so the event that it has reached a fixed level by \(t\) is measurable there. Continuity at its first hit gives \(A_{\tau_n}\le n\). The stopped square sum is bounded pathwise by the stopped integral of \(q\), hence by \(n\). Tonelli, which applies to these nonnegative measurable functions, gives \(\sum_i\int_0^T\mathbb E\|1_{r\le\tau_n}\sigma_i(r)\|^2dr\le n\); each summand is finite. The square-integrable Brownian integral can now be formed by completion of elementary adapted integrands. Its isometry identifies its terminal squared norm with this coefficient-energy integral. A continuous martingale version has the same value almost surely at each fixed time. \(\square\)

\begin{leanbox}{Lean: \texttt{RawLocalDiffusion}, \texttt{pathEnergy}, \texttt{stop} and 4 more}
\leansource{\texttt{RawLocalDiffusion} \quad PBGSO/\allowbreak{}RawLocalIntegral.lean, lines 26–37}
\begin{leancode}
\leanline{0}{structure\ RawLocalDiffusion\ (T\ :\ ℝ≥0)\ where}
\leanline{2}{sigma\ :\ I\ →\ ℝ\ →\ Ω\ →\ Z}
\leanline{2}{sigma\_strong\ :\ ∀\ i,\ StronglyMeasurable\ (Function.uncurry\ (sigma\ i))}
\leanline{2}{sigma\_adapted\ :\ ∀\ t\ :\ ℝ≥0,\ ∀\ i,\ AEStronglyMeasurable[D.filtration\ t]\ (sigma\ i\ t)\ Q}
\leanline{2}{q\ :\ ℝ\ →\ Ω\ →\ ℝ}
\leanline{2}{q\_strong\ :\ StronglyMeasurable\ (Function.uncurry\ q)}
\leanline{2}{q\_nonneg\ :\ ∀\ r\ ω,\ 0\ ≤\ q\ r\ ω}
\leanline{2}{q\_local\ :\ ∀\ ω,\ IntegrableOn\ (fun\ r\ =>\ q\ r\ ω)\ (Ioc\ 0\ (T\ :\ ℝ))}
\leanline{2}{q\_progressive\ :\ ∀\ t\ :\ ℝ≥0,}
\leanline{4}{StronglyMeasurable[(inferInstance\ :\ MeasurableSpace\ ℝ).prod\ (D.filtration\ t)]}
\leanline{6}{(fun\ p\ :\ ℝ\ ×\ Ω\ =>\ (Ioc\ 0\ (t\ :\ ℝ)).indicator\ (fun\ r\ =>\ q\ r\ p.2)\ p.1)}
\leanline{2}{sigma\_control\ :\ ∀\ r\ ω,\ ∑\ i,\ ‖sigma\ i\ r\ ω‖\ \textasciicircum{}\ 2\ ≤\ q\ r\ ω}
\end{leancode}
\leansource{\texttt{pathEnergy} \quad PBGSO/\allowbreak{}RawLocalIntegral.lean, lines 42–42}
\begin{leancode}
\leanline{0}{def\ pathEnergy\ :=\ PathwiseLocalEnergy.energy\ T\ L.q}
\end{leancode}
\leansource{\texttt{stop} \quad PBGSO/\allowbreak{}RawLocalIntegral.lean, lines 44–44}
\begin{leancode}
\leanline{0}{def\ stop\ (n\ :\ ℕ)\ :\ Ω\ →\ ℝ≥0\ :=\ PathEnergyStopping.firstHit\ T\ L.pathEnergy\ n}
\end{leancode}
\leansource{\texttt{stop\_\allowbreak{}energy\_\allowbreak{}bound} \quad PBGSO/\allowbreak{}RawLocalIntegral.lean, lines 62–67}
\begin{leancode}
\leanline{0}{theorem\ stop\_energy\_bound\ (n\ :\ ℕ)\ (ω\ :\ Ω)\ :}
\leanline{4}{(∫\ r\ in\ Ioc\ 0\ (L.stop\ n\ ω\ :\ ℝ),\ L.q\ r\ ω)\ ≤\ n}
\end{leancode}
\leansource{\texttt{stoppedSigma\_\allowbreak{}square} \quad PBGSO/\allowbreak{}RawLocalIntegral.lean, lines 90–99}
\begin{leancode}
\leanline{0}{theorem\ stoppedSigma\_square\ (n\ :\ ℕ)\ (i\ :\ I)\ :}
\leanline{4}{Integrable\ (fun\ p\ :\ ℝ\ ×\ Ω\ =>\ ‖L.stoppedSigma\ n\ i\ p.1\ p.2‖\ \textasciicircum{}\ 2)}
\leanline{6}{((volume.restrict\ (Ioc\ 0\ (T\ :\ ℝ))).prod\ Q)}
\end{leancode}
\leansource{\texttt{noise} \quad PBGSO/\allowbreak{}RawLocalIntegral.lean, lines 102–104}
\begin{leancode}
\leanline{0}{def\ noise\ (n\ :\ ℕ)\ :\ ℝ≥0\ →\ Lp\ Z\ 2\ Q\ :=}
\leanline{2}{rawDiffusionIntegral\ Q\ D\ T\ (L.stoppedSigma\ n)\ (L.stoppedSigma\_strong\ n)}
\leanline{4}{(L.stoppedSigma\_square\ n)\ (L.stoppedSigma\_adapted\ n)}
\end{leancode}
\leansource{\texttt{noise\_\allowbreak{}terminal\_\allowbreak{}sq\_\allowbreak{}le} \quad PBGSO/\allowbreak{}RawLocalIntegral.lean, lines 128–150}
\begin{leancode}
\leanline{0}{theorem\ noise\_terminal\_sq\_le\ (n\ :\ ℕ)\ :\ ‖L.noise\ n\ T‖\ \textasciicircum{}\ 2\ ≤\ n}
\end{leancode}
\leanlabel{How the Lean proof follows the argument}
\texttt{Raw\allowbreak{}Local\allowbreak{}Diffusion} records exactly the joint, adapted and progressive measurability and pathwise bounds. \texttt{path\allowbreak{}Energy} and \texttt{stop} construct the energy and first-hit time. \texttt{stopped\allowbreak{}Sigma\_\allowbreak{}square} invokes the proved first-hit column-integrability result; it does not assume a finite global expected energy. \texttt{noise} is the actual \texttt{raw\allowbreak{}Diffusion\allowbreak{}Integral}. \texttt{noise\_\allowbreak{}terminal\_\allowbreak{}sq\_\allowbreak{}le} rewrites its norm by \texttt{raw\allowbreak{}Diffusion\allowbreak{}Integral\_\allowbreak{}norm\_\allowbreak{}sq}, uses product integration and the finite column sum, and applies the stopped-energy bound. Here \texttt{Lp Z 2 Q} is the space of square-integrable random variables modulo almost-sure equality.
\end{leanbox}
\subsection{\texorpdfstring{Gluing one continuous adapted local process}{Gluing one continuous adapted local process}}


\textbf{Claim.} There is one almost surely continuous local integral \(M\), initially zero, whose stopped version at \(\tau_n\) is \(M^n\). At each fixed \(t\le T\), it is measurable in the completed time-\(t\) past.



\textbf{Proof.} The thresholds are nested. For \(n\le m\), stopping the \(m\)-th coefficient at \(\tau_n\) gives the \(n\)-th coefficient exactly. Consistency of the completed Brownian integral and continuity of its versions imply \(M^n_t=M^m_{t\wedge\tau_n}\) simultaneously at all times, outside one null set for each pair. The union of the countably many exceptional sets is still null.



For each path the finite energy \(A_T\) lies below some integer threshold, so eventually \(\tau_n=T\). Choose the first such integer \(k(\omega)\) and set \(M_t(\omega)=M^{k(\omega)}_t(\omega)\). Pairwise consistency shows that this value agrees with any sufficiently high threshold, and gives \(M^n_t=M_{t\wedge\tau_n}\). On the common full-measure set, each chosen path is continuous and starts at zero. For fixed \(t\le T\), \(M_t^n\) is eventually equal to \(M_t\). Its limit is measurable in the completed past, since every \(M_t^n\) is adapted. This proves adaptation without requiring the future-dependent selector \(k\) itself to be past-measurable. \(\square\)

\begin{leanbox}{Lean: \texttt{version\_\allowbreak{}consistent}, \texttt{horizonIndex}, \texttt{localNoise} and 5 more}
\leansource{\texttt{version\_\allowbreak{}consistent} \quad PBGSO/\allowbreak{}RawLocalIntegral.lean, lines 161–174}
\begin{leancode}
\leanline{0}{theorem\ version\_consistent\ \{n\ m\ :\ ℕ\}\ (h\ :\ n\ ≤\ m)\ :}
\leanline{4}{∀ᵐ\ ω\ ∂Q,\ ∀\ t,\ L.version\ n\ t\ ω\ =\ L.version\ m\ (min\ t\ (L.stop\ n\ ω))\ ω}
\end{leancode}
\leansource{\texttt{horizonIndex} \quad PBGSO/\allowbreak{}RawLocalIntegral.lean, lines 178–178}
\begin{leancode}
\leanline{0}{def\ horizonIndex\ (ω\ :\ Ω)\ :\ ℕ\ :=\ Nat.find\ ((L.stop\_eventually\_horizon\ ω).exists)}
\end{leancode}
\leansource{\texttt{localNoise} \quad PBGSO/\allowbreak{}RawLocalIntegral.lean, lines 185–185}
\begin{leancode}
\leanline{0}{def\ localNoise\ (t\ :\ ℝ≥0)\ (ω\ :\ Ω)\ :\ Z\ :=\ L.version\ (L.horizonIndex\ ω)\ t\ ω}
\end{leancode}
\leansource{\texttt{localNoise\_\allowbreak{}stopped} \quad PBGSO/\allowbreak{}RawLocalIntegral.lean, lines 201–209}
\begin{leancode}
\leanline{0}{theorem\ localNoise\_stopped\ :\ ∀ᵐ\ ω\ ∂Q,\ ∀\ n\ :\ ℕ,\ ∀\ t\ :\ ℝ≥0,\ t\ ≤\ T\ →}
\leanline{4}{L.version\ n\ t\ ω\ =\ L.localNoise\ (min\ t\ (L.stop\ n\ ω))\ ω}
\end{leancode}
\leansource{\texttt{localNoise\_\allowbreak{}continuous} \quad PBGSO/\allowbreak{}RawLocalIntegral.lean, lines 211–213}
\begin{leancode}
\leanline{0}{theorem\ localNoise\_continuous\ :\ ∀ᵐ\ ω\ ∂Q,\ Continuous\ (fun\ t\ =>\ L.localNoise\ t\ ω)}
\end{leancode}
\leansource{\texttt{localNoise\_\allowbreak{}zero} \quad PBGSO/\allowbreak{}RawLocalIntegral.lean, lines 221–223}
\begin{leancode}
\leanline{0}{theorem\ localNoise\_zero\ :\ L.localNoise\ 0\ =ᵐ[Q]\ (fun\ \_\ =>\ (0\ :\ Z))}
\end{leancode}
\leansource{\texttt{version\_\allowbreak{}eventually\_\allowbreak{}localNoise} \quad PBGSO/\allowbreak{}RawLocalIntegral.lean, lines 227–232}
\begin{leancode}
\leanline{0}{theorem\ version\_eventually\_localNoise\ :\ ∀ᵐ\ ω\ ∂Q,\ ∀\ t\ :\ ℝ≥0,\ t\ ≤\ T\ →}
\leanline{4}{∀ᶠ\ n\ :\ ℕ\ in\ atTop,\ L.version\ n\ t\ ω\ =\ L.localNoise\ t\ ω}
\end{leancode}
\leansource{\texttt{localNoise\_\allowbreak{}adapted} \quad PBGSO/\allowbreak{}RawLocalIntegral.lean, lines 237–249}
\begin{leancode}
\leanline{0}{theorem\ localNoise\_adapted\ (t\ :\ ℝ≥0)\ (ht\ :\ t\ ≤\ T)\ :}
\leanline{4}{AEStronglyMeasurable[D.filtration\ t]\ (L.localNoise\ t)\ Q}
\end{leancode}
\leanlabel{How the Lean proof follows the argument}
\texttt{version\_\allowbreak{}consistent} applies the actual stopped-integral theorem after checking equality of the two nested coefficient masks. \texttt{horizon\allowbreak{}Index} is the least integer provided by eventual attainment of the horizon, and \texttt{local\allowbreak{}Noise} selects its continuous martingale version. \texttt{local\allowbreak{}Noise\_\allowbreak{}stopped} compares both indices with their maximum. \texttt{local\allowbreak{}Noise\_\allowbreak{}continuous} and \texttt{local\allowbreak{}Noise\_\allowbreak{}zero} take a countable intersection of the corresponding almost-sure properties. \texttt{local\allowbreak{}Noise\_\allowbreak{}adapted} represents the local noise at a fixed time as the almost-sure limit of the adapted versions; \texttt{version\_\allowbreak{}eventually\_\allowbreak{}local\allowbreak{}Noise} identifies that limit.
\end{leanbox}
\section{\texorpdfstring{The finite signed-work domain of the projected criterion}{The finite signed-work domain of the projected criterion}}


A dissipativity inequality stated only when its integral exists does not establish integrability for every finite-second-moment coupling. The projected criterion therefore records its finite signed drift-work domain separately. For constant diffusion, that domain is exactly the sampler coupling-work domain.

\subsection{\texorpdfstring{Constant noise cancels but does not create integrability}{Constant noise cancels but does not create integrability}}


\textbf{Claim.} Assume all hypotheses of the canonical projected criterion: literal block projections, positive definite optimiser and sampler metric and auxiliary matrices, \(\beta\ge0\), \(\beta r-A_o-A_s>0\), projected strong monotonicity, auxiliary dissipation, constant diffusion and integrability of its projected expressions. Assume additionally that for every pair of optimiser states and every probability coupling \(\Lambda\) with \(\int(1+\|z\|^2+\|z'\|^2)d\Lambda<\infty\), the signed sampler drift work is integrable. Then the same positive generalized-eigenvalue rate gives joint dissipativity, and the complete sampler coupling work is integrable for all these couplings.



\textbf{Proof.} The preceding projected criterion supplies the positive rate and its joint inequality. In the sampler coupling work, every column difference is \(\sigma_i-\sigma_i=0\). Hence the finite sum of their quadratic squares is zero, and the sampler work equals the signed drift work pointwise. The added domain hypothesis makes precisely that remaining expression integrable. Combine this integrability assertion with the two conclusions of the projected criterion. No inequality about a totalised integral is used to infer finiteness. \(\square\)

\begin{leanbox}{Lean: \texttt{FiniteSamplerDriftDomain}, \texttt{finiteCouplingWork\_\allowbreak{}of\_\allowbreak{}U5\_\allowbreak{}domain}, \texttt{canonical\_\allowbreak{}projected\_\allowbreak{}source\_\allowbreak{}criterion\_\allowbreak{}finite}}
\leansource{\texttt{FiniteSamplerDriftDomain} \quad PBGSO/\allowbreak{}ProjectedFiniteWork.lean, lines 24–29}
\begin{leancode}
\leanline{0}{def\ FiniteSamplerDriftDomain}
\leanline{4}{(M\ :\ Model\ (BlockSpace\ a\ u)\ (EuclideanSpace\ ℝ\ a)\ (BlockSpace\ x\ v)\ (EuclideanSpace\ ℝ\ x)\ I)}
\leanline{4}{(D\ :\ TwistedGeometry\ (BlockSpace\ a\ u)\ (BlockSpace\ x\ v))\ :\ Prop\ :=}
\leanline{2}{∀\ s\ s'\ (Λ\ :\ Measure\ (BlockSpace\ x\ v\ ×\ BlockSpace\ x\ v)),}
\leanline{4}{IsProbabilityMeasure\ Λ\ →\ Integrable\ (fun\ p\ =>\ 1+‖p.1‖\ \textasciicircum{}\ 2+‖p.2‖\ \textasciicircum{}\ 2)\ Λ\ →}
\leanline{6}{Integrable\ (samplerDriftWork\ M\ D\ s\ s')\ Λ}
\end{leancode}
\leansource{\texttt{finiteCouplingWork\_\allowbreak{}of\_\allowbreak{}U5\_\allowbreak{}domain} \quad PBGSO/\allowbreak{}ProjectedFiniteWork.lean, lines 33–40}
\begin{leancode}
\leanline{0}{theorem\ finiteCouplingWork\_of\_U5\_domain}
\leanline{4}{(M\ :\ Model\ (BlockSpace\ a\ u)\ (EuclideanSpace\ ℝ\ a)\ (BlockSpace\ x\ v)\ (EuclideanSpace\ ℝ\ x)\ I)}
\leanline{4}{(D\ :\ TwistedGeometry\ (BlockSpace\ a\ u)\ (BlockSpace\ x\ v))}
\leanline{4}{(σ\ :\ I\ →\ BlockSpace\ x\ v)\ (hσ\ :\ ∀\ θ\ z,\ M.diffusion\ θ\ z\ =\ σ)}
\leanline{4}{(hfinite\ :\ FiniteSamplerDriftDomain\ M\ D)\ :\ DiracOptimizerGrowth.FiniteCouplingWork\ M\ D}
\end{leancode}
\leansource{\texttt{canonical\_\allowbreak{}projected\_\allowbreak{}source\_\allowbreak{}criterion\_\allowbreak{}finite} \quad PBGSO/\allowbreak{}ProjectedFiniteWork.lean, lines 45–68}
\begin{leancode}
\leanline{0}{theorem\ canonical\_projected\_source\_criterion\_finite}
\leanline{4}{(M\ :\ Model\ (BlockSpace\ a\ u)\ (EuclideanSpace\ ℝ\ a)\ (BlockSpace\ x\ v)\ (EuclideanSpace\ ℝ\ x)\ I)}
\leanline{4}{(hparameter\ :\ M.parameter\ =\ WithLp.fst)\ (hposition\ :\ M.position\ =\ WithLp.fst)}
\leanline{4}{(H\ :\ EuclideanSpace\ ℝ\ a\ →\ EuclideanSpace\ ℝ\ x\ →\ EuclideanSpace\ ℝ\ x)}
\leanline{4}{(DOpt\ :\ Matrix\ (a\ ⊕\ u)\ (a\ ⊕\ u)\ ℝ)\ (DSam\ :\ Matrix\ (x\ ⊕\ v)\ (x\ ⊕\ v)\ ℝ)}
\leanline{4}{(Dm\ :\ Matrix\ u\ u\ ℝ)\ (Dv\ :\ Matrix\ v\ v\ ℝ)}
\leanline{4}{(hDOpt\ :\ DOpt.PosDef)\ (hDSam\ :\ DSam.PosDef)\ (hDm\ :\ Dm.PosDef)\ (hDv\ :\ Dv.PosDef)}
\leanline{4}{(β\ r\ AOpt\ ASam\ :\ ℝ)\ (hβ\ :\ 0\ ≤\ β)\ (hc\ :\ 0\ <\ β\ *\ r\ -\ AOpt\ -\ ASam)}
\leanline{4}{(hmono\ :\ (canonicalData\ M\ hparameter\ hposition\ H\ Dm\ Dv).StrongMonotonicity\ r)}
\leanline{4}{(haux\ :\ ProjectedAuxiliaryDissipation\ M\ (canonicalGeometry\ DOpt\ DSam\ hDOpt\ hDSam)}
\leanline{6}{(canonicalData\ M\ hparameter\ hposition\ H\ Dm\ Dv)\ β\ AOpt\ ASam)}
\leanline{4}{(σ\ :\ I\ →\ BlockSpace\ x\ v)\ (hσ\ :\ ∀\ θ\ z,\ M.diffusion\ θ\ z\ =\ σ)}
\leanline{4}{(hInt\ :\ ∀\ s\ s'\ ν\ ν'\ Λ,}
\leanline{6}{ProjectedAdmissibleCoupling\ M\ (canonicalGeometry\ DOpt\ DSam\ hDOpt\ hDSam)\ s\ s'\ ν\ ν'\ Λ\ →}
\leanline{6}{(canonicalData\ M\ hparameter\ hposition\ H\ Dm\ Dv).Integrability\ s\ s'\ Λ)}
\leanline{4}{(hfinite\ :\ FiniteSamplerDriftDomain\ M\ (canonicalGeometry\ DOpt\ DSam\ hDOpt\ hDSam))\ :}
\leanline{4}{let\ ρ\ :=\ min}
\leanline{6}{(generalizedMinimum\ hDOpt\ (blockK\_posDef\ (a\ :=\ a)\ \_\ hc\ Dm\ hDm))}
\leanline{6}{(generalizedMinimum\ hDSam\ (blockK\_posDef\ (a\ :=\ x)\ \_\ hc\ Dv\ hDv))}
\leanline{4}{0\ <\ ρ\ ∧\ TwistedJointMonotonicity\ M\ (canonicalGeometry\ DOpt\ DSam\ hDOpt\ hDSam)\ ρ\ ∧}
\leanline{6}{DiracOptimizerGrowth.FiniteCouplingWork\ M\ (canonicalGeometry\ DOpt\ DSam\ hDOpt\ hDSam)}
\end{leancode}
\leanlabel{How the Lean proof follows the argument}
\texttt{Finite\allowbreak{}Sampler\allowbreak{}Drift\allowbreak{}Domain} quantifies over all probability couplings with the displayed quadratic weight integrable. \texttt{finite\allowbreak{}Coupling\allowbreak{}Work\_\allowbreak{}of\_\allowbreak{}U5\_\allowbreak{}domain} rewrites the work by \texttt{sampler\allowbreak{}Coupling\allowbreak{}Work\_\allowbreak{}eq\_\allowbreak{}drift\_\allowbreak{}of\_\allowbreak{}constant} and applies \texttt{hfinite}. \texttt{canonical\_\allowbreak{}projected\_\allowbreak{}source\_\allowbreak{}criterion\_\allowbreak{}finite} obtains the rate and monotonicity from \texttt{canonical\_\allowbreak{}projected\_\allowbreak{}source\_\allowbreak{}criterion}, then packages this separate finite-work result as its third component.
\end{leanbox}
\section{\texorpdfstring{An explicit heavy-ball metric and positive rate}{An explicit heavy-ball metric and positive rate}}


For a position–momentum difference \((q,p)\), the quadratic energy is \(\gamma\|q\|^2+2\langle q,p\rangle+2\|p\|^2/\gamma\), where \(\gamma>0\). Its mixed term transfers position monotonicity to a contraction estimate for the full heavy-ball dynamics. The optimiser and sampler forces may be independent and need not arise from a common potential.

\subsection{\texorpdfstring{Heavy-ball source criterion}{Heavy-ball source criterion}}


\textbf{Claim.} Let \(F(\theta,x)\) and \(H(\theta,x)\) be forces on complete real Hilbert parameter and position spaces. Suppose their paired field is \(L\)-Lipschitz, \(L\ge0\), and

\[\langle\theta-\theta',F(\theta,x)-F(\theta',x')\rangle+\langle x-x',H(\theta,x)-H(\theta',x')\rangle\ge r(\|\theta-\theta'\|^2+\|x-x'\|^2),\]

with \(r>0\). Assume \(\gamma>\sqrt2L/\sqrt r\), and take arbitrary constant sampler diffusion columns. Put \(\eta=(1+2L^2/(r\gamma^2))/2\), \(a=1-\eta\), \(b=L^2/(\eta\gamma^2)\), \(c=r-2b\), and \(u=\gamma+1+2/\gamma\). Then \(\rho=\min(c,a)/u>0\), and the heavy-ball model satisfies the joint dissipativity condition in the displayed quadratic metric.



\textbf{Proof.} Squaring the friction condition is legitimate because both sides are nonnegative, and gives \(2L^2<r\gamma^2\). Thus \(0<\eta<1\), \(a>0\), and \(2L^2<r\eta\gamma^2\), so \(c>0\). For the heavy-ball difference field \((p,-\gamma p-f)\), its weighted pairing is

\[-\langle q,f\rangle-\|p\|^2-(2/\gamma)\langle p,f\rangle.\]

Cauchy–Schwarz and Young's inequality give \((2/\gamma)|\langle p,f\rangle|\le\eta\|p\|^2+\|f\|^2/(\eta\gamma^2)\). Applied to the optimiser mean force and sampler force, their squared bounds are controlled by the paired Lipschitz condition and Jensen's inequality for the mean. Together with joint monotonicity the total work is at most \(-c\|q\|^2-a\|p\|^2\). The energy is at most \(u(\|q\|^2+\|p\|^2)\), by \(2|\langle q,p\rangle|\le\|q\|^2+\|p\|^2\). Its lower bound is \(\min(\gamma/4,1/(2\gamma))(\|q\|^2+\|p\|^2)\), obtained by Young's inequality with position coefficient \(3\gamma/4\). Therefore it is positive definite, and the work bound is at most minus \(\rho\) times the energy. Lipschitz growth and finite second moments ensure the projected expressions are integrable. Constant noise cancels in the coupling. \(\square\)

\begin{leanbox}{Lean: \texttt{hypocoerciveYoung}, \texttt{hypocoerciveRate}, \texttt{heavyBall\_\allowbreak{}source\_\allowbreak{}criterion}}
\leansource{\texttt{hypocoerciveYoung} \quad PBGSO/\allowbreak{}HypocoerciveCriterion.lean, lines 18–18}
\begin{leancode}
\leanline{0}{def\ hypocoerciveYoung\ (γ\ r\ L\ :\ ℝ)\ :\ ℝ\ :=\ (1\ +\ 2\ *\ L\ \textasciicircum{}\ 2\ /\ (r\ *\ γ\ \textasciicircum{}\ 2))\ /\ 2}
\end{leancode}
\leansource{\texttt{hypocoerciveRate} \quad PBGSO/\allowbreak{}HypocoerciveCriterion.lean, lines 24–26}
\begin{leancode}
\leanline{0}{def\ hypocoerciveRate\ (γ\ r\ L\ :\ ℝ)\ :\ ℝ\ :=}
\leanline{2}{min\ (hypocoerciveRemainder\ γ\ r\ L)\ (hypocoerciveAuxiliary\ γ\ r\ L)\ /}
\leanline{4}{hypocoerciveUpper\ γ}
\end{leancode}
\leansource{\texttt{heavyBall\_\allowbreak{}source\_\allowbreak{}criterion} \quad PBGSO/\allowbreak{}HypocoerciveCriterion.lean, lines 644–657}
\begin{leancode}
\leanline{0}{theorem\ heavyBall\_source\_criterion\ (γ\ r\ :\ ℝ)\ (L\ :\ ℝ≥0)}
\leanline{4}{(hγ\ :\ 0\ <\ γ)\ (hr\ :\ 0\ <\ r)}
\leanline{4}{(hsource\ :\ Real.sqrt\ 2\ *\ (L\ :\ ℝ)\ /\ Real.sqrt\ r\ <\ γ)}
\leanline{4}{(F\ :\ A\ →\ X\ →\ A)\ (H\ :\ A\ →\ X\ →\ X)\ (σ\ :\ I\ →\ HypocoercivePhase\ X)}
\leanline{4}{(hmono\ :\ ∀\ θ\ θ'\ x\ x',\ r\ *\ (‖θ\ -\ θ'‖\ \textasciicircum{}\ 2\ +\ ‖x\ -\ x'‖\ \textasciicircum{}\ 2)\ ≤}
\leanline{6}{{\leanfallback ⟪}θ\ -\ θ',\ F\ θ\ x\ -\ F\ θ'\ x'{\leanfallback ⟫}\_ℝ\ +\ {\leanfallback ⟪}x\ -\ x',\ H\ θ\ x\ -\ H\ θ'\ x'{\leanfallback ⟫}\_ℝ)}
\leanline{4}{(hLip\ :\ LipschitzWith\ L\ (heavyBallProjectedField\ F\ H))\ :}
\leanline{4}{0\ <\ hypocoerciveRate\ γ\ r\ L\ ∧}
\leanline{6}{TwistedJointMonotonicity\ (heavyBallModel\ γ\ F\ H\ σ)\ (hypocoerciveGeometry\ γ\ hγ)}
\leanline{8}{(hypocoerciveRate\ γ\ r\ L)}
\end{leancode}
\leanlabel{How the Lean proof follows the argument}
\texttt{heavy\allowbreak{}Ball\_\allowbreak{}source\_\allowbreak{}criterion} obtains the two component force estimates from \texttt{heavy\allowbreak{}Ball\allowbreak{}Projected\allowbreak{}Field\_\allowbreak{}component\_\allowbreak{}bounds} and passes them to \texttt{heavy\allowbreak{}Ball\_\allowbreak{}twisted\allowbreak{}Joint\allowbreak{}Monotonicity\_\allowbreak{}of\_\allowbreak{}source}. The latter proves projected integrability and applies the explicit hypocoercive criterion. The definitions \texttt{hypocoercive\allowbreak{}Young} and \texttt{hypocoercive\allowbreak{}Rate} are the constants above. The earlier lemmas in this file establish their positivity, the energy bounds, and the weighted drift pairing. The force maps remain independent arguments throughout.
\end{leanbox}
\section{\texorpdfstring{Smooth dependence and the continuous Jacobian}{Smooth dependence and the continuous Jacobian}}


Consider an actual globally defined flow \(\Phi_t\) of a finite-dimensional autonomous \(C^1\) vector field. Joint differentiability in time and initial state is derived from these hypotheses. For the frictionless phase field \((p,q)\mapsto(-F(q),p)\), the divergence is zero even when \(F\) is not a gradient.

\subsection{\texorpdfstring{From a primitive differential equation to determinant one}{From a primitive differential equation to determinant one}}


\textbf{Claim.} Suppose \(F:\mathbb R^d\to\mathbb R^d\) is \(C^1\), and for every initial phase point \(x\) there is a global trajectory \(\Phi(t,x)\) satisfying \(\partial_t\Phi(t,x)=(-F(q_t),p_t)\) and \(\Phi(0,x)=x\). Then \(\Phi\) is jointly \(C^1\), and \(\det D_x\Phi(t,x)=1\) for every real \(t\) and every \(x\), including zero dimension.



\textbf{Proof.} A \(C^1\) field is Lipschitz on a compact ball because its continuous derivative is bounded there. On any compact time interval, both continuous trajectories lie in one such ball, so the usual difference estimate and Grönwall give uniqueness. To obtain differentiability in the initial value near a fixed trajectory, choose a smooth compact cutoff equal to one on an open neighbourhood of its compact orbit. The cut-off field is globally Lipschitz and \(C^1\), and global Picard iteration constructs its flow with \(C^1\) dependence. Its continuous dependence keeps nearby trajectories inside that open neighbourhood for the whole compact time interval. Uniqueness then identifies the cut-off and original trajectories locally. Thus their solution maps agree in a neighbourhood of each \((t,x)\), proving joint \(C^1\) regularity.



Let \(J(t)=D_x\Phi(t,x)\). Differentiating the ODE in the initial condition gives \(J'=Df(\Phi_t(x))J\), \(J(0)=I\). The derivative of the phase field has blocks \(\begin{pmatrix}0&-DF(q)\\I&0\end{pmatrix}\), whose trace is zero. The determinant differentiation formula for this variational equation gives \((\det J)'=\operatorname{tr}(Df)\det J=0\). Hence \(\det J(t)=\det I=1\), for positive or negative time. \(\square\)

\begin{leanbox}{Lean: \texttt{flow\_\allowbreak{}contDiff}, \texttt{phase\_\allowbreak{}flow\_\allowbreak{}fderiv\_\allowbreak{}det\_\allowbreak{}one\_\allowbreak{}source}}
\leansource{\texttt{flow\_\allowbreak{}contDiff} \quad PBGSO/\allowbreak{}FlowSourceRegularity.lean, lines 114–120}
\begin{leancode}
\leanline{0}{theorem\ flow\_contDiff\ (Φ\ :\ ℝ\ ×\ E\ →\ E)\ (f\ :\ E\ →\ E)\ (hf\ :\ ContDiff\ ℝ\ 1\ f)}
\leanline{4}{(hODE\ :\ ∀\ t\ x,\ HasDerivAt\ (fun\ s\ =>\ Φ\ (s,\ x))\ (f\ (Φ\ (t,\ x)))\ t)}
\leanline{4}{(hinit\ :\ ∀\ x,\ Φ\ (0,\ x)\ =\ x)\ :\ ContDiff\ ℝ\ 1\ Φ}
\end{leancode}
\leansource{\texttt{phase\_\allowbreak{}flow\_\allowbreak{}fderiv\_\allowbreak{}det\_\allowbreak{}one\_\allowbreak{}source} \quad PBGSO/\allowbreak{}FlowSourceRegularity.lean, lines 150–158}
\begin{leancode}
\leanline{0}{theorem\ phase\_flow\_fderiv\_det\_one\_source}
\leanline{4}{(F\ :\ (ι\ →\ ℝ)\ →\ (ι\ →\ ℝ))\ (hF\ :\ ContDiff\ ℝ\ 1\ F)}
\leanline{4}{(Φ\ :\ ℝ\ ×\ Phase\ ι\ →\ Phase\ ι)}
\leanline{4}{(hODE\ :\ ∀\ t\ x,\ HasDerivAt\ (fun\ s\ =>\ Φ\ (s,\ x))\ (phaseField\ F\ (Φ\ (t,\ x)))\ t)}
\leanline{4}{(hinit\ :\ ∀\ x,\ Φ\ (0,\ x)\ =\ x)\ (t\ :\ ℝ)\ (x\ :\ Phase\ ι)\ :}
\leanline{4}{(fderiv\ ℝ\ (fun\ y\ =>\ Φ\ (t,\ y))\ x).det\ =\ 1}
\end{leancode}
\leanlabel{How the Lean proof follows the argument}
\texttt{flow\_\allowbreak{}cont\allowbreak{}Diff} discharges the cutoff-solution existence interface using \texttt{exists\_\allowbreak{}global\_\allowbreak{}ode\_\allowbreak{}flow}. The preceding \texttt{flow\_\allowbreak{}cont\allowbreak{}Diff\_\allowbreak{}of\_\allowbreak{}cutoff\_\allowbreak{}solutions} constructs a compact orbit, applies \texttt{exists\_\allowbreak{}global\_\allowbreak{}lipschitz\_\allowbreak{}C1\_\allowbreak{}cutoff}, proves the orbit equality by uniqueness, obtains a neighbourhood of nearby initial values through the compact-open topology, and transfers differentiability by eventual equality. \texttt{phase\_\allowbreak{}flow\_\allowbreak{}fderiv\_\allowbreak{}det\_\allowbreak{}one\_\allowbreak{}source} supplies this derived regularity to \texttt{phase\_\allowbreak{}flow\_\allowbreak{}fderiv\_\allowbreak{}det\_\allowbreak{}one\_\allowbreak{}C1}. Thus neither a Jacobian identity nor differentiability of a chosen solution map is an input.
\end{leanbox}
\section{\texorpdfstring{Force data of the particle methods and the Lipschitz bounds of the cloud forces}{Force data of the particle methods and the Lipschitz bounds of the cloud forces}}


This section fixes the force data shared by the two numerical methods (assumption N1) and proves the Lipschitz estimates for the forces acting on a whole particle cloud. Let \(\Theta\), \(\mathcal X\) and \(\mathcal C\) be real normed spaces: \(\Theta\) contains the optimiser parameter, \(\mathcal X\) the position of a single particle, and \(\mathcal C\) a coordinate vector describing a cloud of \(N\ge1\) particles. A continuous linear bijection \(\mathcal E:\mathcal C\to\mathcal X^N\) with continuous inverse reads off the particles; \(\mathcal X^N\) carries the norm \(\|x\|_\infty=\max_{1\le i\le N}\|x_i\|\), and \(\Theta\times\mathcal C\) carries the norm \(\|(\theta,c)\|=\max\{\|\theta\|,\|c\|\}\). The observation \(F:\Theta\times\mathcal X\to\Theta\) and the sampler force \(F_x:\Theta\times\mathcal X\to\mathcal X\) are two unrelated maps with joint Lipschitz constants \(\ell_F,\ell_x\ge0\). From them we build the empirical mean force \(\bar F(\theta,c)=N^{-1}\sum_iF(\theta,(\mathcal Ec)_i)\), which drives the parameter, and the stacked sampler force \(\mathcal F_x(\theta,c)=\mathcal E^{-1}\bigl((F_x(\theta,(\mathcal Ec)_i))_{i=1}^N\bigr)\), which drives the cloud.




The main result is that \(\bar F\) is globally Lipschitz with constant \(L_F=\ell_F(1+\|\mathcal E\|)\), and \(\mathcal F_x\) with constant \(L_x=\|\mathcal E^{-1}\|\,\ell_x(1+\|\mathcal E\|)\). Both forces are then evaluated at square-integrable random arguments \((\theta,X)\) defined on one probability space, so that the parameter and all particles may be dependent. The lifted forces act sample by sample and satisfy the same Lipschitz bounds in \(L^2\). We also show that one step of either method, computed on square-integrable random variables, coincides almost surely with the update applied to the sampled values. For the standard Euclidean encoding of \(\mathbb R^{Nd_x}\) with \(d_x\ge1\) one has \(\|\mathcal E\|=1\) and \(\|\mathcal E^{-1}\|=\sqrt N\), hence \(L_F=2\ell_F\) and \(L_x=2\sqrt N\,\ell_x\); these are the force constants in Theorems New A and New B.

\subsection{\texorpdfstring{Force data (assumption N1)}{Force data (assumption N1)}}


\textbf{Definition.} Let \(\Theta\), \(\mathcal X\), \(\mathcal C\) be real normed spaces and \(N\in\mathbb N\). Force data for \(N\) particles consist of



\begin{enumerate}


\item the condition \(N\ge1\);

\item two maps \(F:\Theta\times\mathcal X\to\Theta\) (the observation) and \(F_x:\Theta\times\mathcal X\to\mathcal X\) (the sampler force);

\item two numbers \(\ell_F,\ell_x\in[0,\infty)\) such that for all \(\theta,\theta'\in\Theta\) and \(x,x'\in\mathcal X\)
\[\begin{aligned}\|F(\theta,x)-F(\theta',x')\|&\le\ell_F\bigl(\|\theta-\theta'\|+\|x-x'\|\bigr),\qquad(1)\\ \|F_x(\theta,x)-F_x(\theta',x')\|&\le\ell_x\bigl(\|\theta-\theta'\|+\|x-x'\|\bigr);\qquad(2)\end{aligned}\]

\item a linear bijection \(\mathcal E:\mathcal C\to\mathcal X^N\) such that \(\mathcal E\) and \(\mathcal E^{-1}\) are bounded, with operator norms \(\|\mathcal E\|\) and \(\|\mathcal E^{-1}\|\) taken with respect to \(\|x\|_\infty=\max_i\|x_i\|\) on \(\mathcal X^N\).

\end{enumerate}



Conditions (1) and (2) constitute assumption N1. They are global, so the forces may grow linearly: taking \(\theta'=0\), \(x'=0\) in (1) and using the triangle inequality gives \(\|F(\theta,x)\|\le\|F(0,0)\|+\ell_F(\|\theta\|+\|x\|)\). No differentiability, no common potential and no relation between \(F\) and \(F_x\) is assumed. In the paper, \(F\) is the observation whose population mean is the gradient estimate \(G_\nu\), and \(F_x\) is \(\nabla_xE\) for M-MFPO and the gradient of the Moreau–Yosida envelope \(\nabla_xE^\lambda\) for HOL-MYPGD. The vector \(c\in\mathcal C\) is a coordinate description of the cloud and \((\mathcal Ec)_i\in\mathcal X\) is the position of particle \(i\).

\begin{leanbox}{Lean: \texttt{ParticleForceData}}
\leansource{\texttt{ParticleForceData} \quad PBGSO/\allowbreak{}ParticleMethodLp.lean, lines 136–149}
\begin{leancode}
\leanline{0}{structure\ ParticleForceData\ (A\ X\ C\ :\ Type*)\ (N\ :\ ℕ)}
\leanline{4}{[NormedAddCommGroup\ A]\ [NormedSpace\ ℝ\ A]}
\leanline{4}{[NormedAddCommGroup\ X]\ [NormedSpace\ ℝ\ X]}
\leanline{4}{[NormedAddCommGroup\ C]\ [NormedSpace\ ℝ\ C]\ where}
\leanline{2}{nonempty\ :\ 0\ <\ N}
\leanline{2}{F\ :\ A\ →\ X\ →\ A}
\leanline{2}{G\ :\ A\ →\ X\ →\ X}
\leanline{2}{parameterLip\ :\ ℝ≥0}
\leanline{2}{samplerLip\ :\ ℝ≥0}
\leanline{2}{parameter\_bound\ :\ ∀\ a\ b\ x\ y,}
\leanline{4}{‖F\ a\ x\ -\ F\ b\ y‖\ ≤\ (parameterLip\ :\ ℝ)\ *\ (‖a\ -\ b‖\ +\ ‖x\ -\ y‖)}
\leanline{2}{sampler\_bound\ :\ ∀\ a\ b\ x\ y,}
\leanline{4}{‖G\ a\ x\ -\ G\ b\ y‖\ ≤\ (samplerLip\ :\ ℝ)\ *\ (‖a\ -\ b‖\ +\ ‖x\ -\ y‖)}
\leanline{2}{cloud\ :\ C\ ≃L[ℝ]\ (Fin\ N\ →\ X)}
\end{leancode}
\leanlabel{How the Lean proof follows the argument}
\texttt{Particle\allowbreak{}Force\allowbreak{}Data A X C N} is a structure; its fields are the data of the definition. \texttt{nonempty : 0 < N} is item 1. \texttt{F : A \ensuremath{\to} X \ensuremath{\to} A} and \texttt{G : A \ensuremath{\to} X \ensuremath{\to} X} are the two forces in curried form, so \texttt{G} is the sampler force \(F_x\). The constants \texttt{parameter\allowbreak{}Lip} and \texttt{sampler\allowbreak{}Lip} have type \texttt{\ensuremath{\mathbb R}\ensuremath{\ge}0}, the nonnegative reals; this type encodes \(\ell_F,\ell_x\ge0\), and the coercion \texttt{(parameter\allowbreak{}Lip : \ensuremath{\mathbb R})} turns them into real numbers in the inequalities \texttt{parameter\_\allowbreak{}bound} (1) and \texttt{sampler\_\allowbreak{}bound} (2). The field \texttt{cloud : C ≃L[\ensuremath{\mathbb R}] (Fin N \ensuremath{\to} X)} is a continuous linear equivalence, that is, a linear bijection with both directions continuous; in a normed space continuity of a linear map is boundedness. The function space \texttt{Fin N \ensuremath{\to} X} carries Mathlib's supremum norm, which is \(\|x\|_\infty\).
\end{leanbox}
\subsection{\texorpdfstring{The empirical mean force and the stacked sampler force}{The empirical mean force and the stacked sampler force}}


\textbf{Definition.} For \((\theta,c)\in\Theta\times\mathcal C\) put


\[\bar F(\theta,c)=\frac1N\sum_{i=1}^NF\bigl(\theta,(\mathcal Ec)_i\bigr)\in\Theta,\qquad \mathcal F_x(\theta,c)=\mathcal E^{-1}\Bigl(F_x\bigl(\theta,(\mathcal Ec)_1\bigr),\dots,F_x\bigl(\theta,(\mathcal Ec)_N\bigr)\Bigr)\in\mathcal C.\]


All particles see the same parameter \(\theta\). The vector \(\bar F(\theta,c)\) is the integral of \(F(\theta,\cdot)\) against the empirical law \(N^{-1}\sum_i\delta_{(\mathcal Ec)_i}\) of the cloud, that is, the particle estimate of the population observation \(G_\nu(\theta)=\int F(\theta,x)\,\nu(dx)\). The sampler force applies \(F_x\) to every particle separately and writes the result back in cloud coordinates. The space \(\Theta\times\mathcal C\) is normed by \(\|(\theta,c)\|=\max\{\|\theta\|,\|c\|\}\).

\begin{leanbox}{Lean: \texttt{meanForce}, \texttt{samplerForce}}
\leansource{\texttt{meanForce} \quad PBGSO/\allowbreak{}ParticleMethodLp.lean, lines 159–159}
\begin{leancode}
\leanline{0}{def\ meanForce\ (s\ :\ A\ ×\ C)\ :\ A\ :=\ particleMeanForce\ D.F\ s.1\ (D.cloud\ s.2)}
\end{leancode}
\leansource{\texttt{samplerForce} \quad PBGSO/\allowbreak{}ParticleMethodLp.lean, lines 162–163}
\begin{leancode}
\leanline{0}{def\ samplerForce\ (s\ :\ A\ ×\ C)\ :\ C\ :=}
\leanline{2}{D.cloud.symm\ (particleSamplerForce\ D.G\ s.1\ (D.cloud\ s.2))}
\end{leancode}
\leanlabel{How the Lean proof follows the argument}
Both maps take one argument \texttt{s : A \ensuremath{\times} C}; \texttt{s.\allowbreak{}1} is the parameter and \texttt{s.\allowbreak{}2} the cloud coordinate. \texttt{mean\allowbreak{}Force} is \texttt{particle\allowbreak{}Mean\allowbreak{}Force D.\allowbreak{}F s.\allowbreak{}1 (D.\allowbreak{}cloud s.\allowbreak{}2)}; the map \texttt{particle\allowbreak{}Mean\allowbreak{}Force}, defined in the companion file on particle forces, has the body \texttt{(N : \ensuremath{\mathbb R})\ensuremath{{}^{-1}} \ensuremath{\bullet} ∑ i, F \ensuremath{\theta} (x i)}, where \texttt{\ensuremath{\bullet}} is scalar multiplication. \texttt{sampler\allowbreak{}Force} is \texttt{D.\allowbreak{}cloud.\allowbreak{}symm (particle\allowbreak{}Sampler\allowbreak{}Force D.\allowbreak{}G s.\allowbreak{}1 (D.\allowbreak{}cloud s.\allowbreak{}2))}; there \texttt{particle\allowbreak{}Sampler\allowbreak{}Force} has the body \texttt{fun i \ensuremath{\mapsto} G \ensuremath{\theta} (x i)}, and \texttt{D.\allowbreak{}cloud.\allowbreak{}symm} is \(\mathcal E^{-1}\). Mathlib equips the product \texttt{A \ensuremath{\times} C} with the maximum of the two norms, as in the definition.
\end{leanbox}
\subsection{\texorpdfstring{The two Lipschitz constants}{The two Lipschitz constants}}


\textbf{Definition.} Put


\[L_F=\ell_F\bigl(1+\|\mathcal E\|\bigr),\qquad L_x=\|\mathcal E^{-1}\|\,\ell_x\bigl(1+\|\mathcal E\|\bigr).\]


Both are nonnegative. They are the Lipschitz constants of \(\bar F\) and \(\mathcal F_x\) proved in the next two steps. For the standard Euclidean encoding of \(\mathcal C=\mathbb R^{Nd_x}\) (the last step of this section), the operator norms are \(\|\mathcal E\|=1\) and \(\|\mathcal E^{-1}\|=\sqrt N\) when \(d_x\ge1\), so that


\[L_F=2\ell_F,\qquad L_x=2\sqrt N\,\ell_x,\]


and \(\mathcal E=0\), \(\mathcal E^{-1}=0\) when \(d_x=0\), so that \(L_F=\ell_F\) and \(L_x=0\). These operator norms are computed in the results on the cloud encoding constants: for \(x\in\mathbb R^{Nd_x}\) one has \(\|x\|^2=\sum_{i}\|(\mathcal Ex)_i\|^2\), hence \(\|(\mathcal Ex)_i\|\le\|x\|\) for every \(i\) and \(\|\mathcal E^{-1}y\|^2=\sum_i\|y_i\|^2\le N\|y\|_\infty^2\); a single unit coordinate and a constant unit cloud show that both bounds are attained.

\begin{leanbox}{Lean: \texttt{meanLip}, \texttt{samplerForceLip}}
\leansource{\texttt{meanLip} \quad PBGSO/\allowbreak{}ParticleMethodLp.lean, lines 167–167}
\begin{leancode}
\leanline{0}{def\ meanLip\ :\ ℝ≥0\ :=\ D.parameterLip\ *\ (1\ +\ ‖D.cloud.toContinuousLinearMap‖₊)}
\end{leancode}
\leansource{\texttt{samplerForceLip} \quad PBGSO/\allowbreak{}ParticleMethodLp.lean, lines 168–170}
\begin{leancode}
\leanline{0}{def\ samplerForceLip\ :\ ℝ≥0\ :=}
\leanline{2}{‖D.cloud.symm.toContinuousLinearMap‖₊\ *\ D.samplerLip\ *}
\leanline{4}{(1\ +\ ‖D.cloud.toContinuousLinearMap‖₊)}
\end{leancode}
\leanlabel{How the Lean proof follows the argument}
\texttt{mean\allowbreak{}Lip} and \texttt{sampler\allowbreak{}Force\allowbreak{}Lip} are elements of \texttt{\ensuremath{\mathbb R}\ensuremath{\ge}0}. They are written with \texttt{\ensuremath{\Vert}D.\allowbreak{}cloud.\allowbreak{}to\allowbreak{}Continuous\allowbreak{}Linear\allowbreak{}Map\ensuremath{\Vert}\ensuremath{{}_{+}}}, the operator norm of \(\mathcal E\) regarded as a nonnegative real (\texttt{\ensuremath{\Vert}·\ensuremath{\Vert}\ensuremath{{}_{+}}}), and similarly for \texttt{D.\allowbreak{}cloud.\allowbreak{}symm}. The specialisations to the Euclidean encoding are proved in the file on cloud encoding constants as \texttt{mean\allowbreak{}Lip\_\allowbreak{}euclidean}, \texttt{sampler\allowbreak{}Force\allowbreak{}Lip\_\allowbreak{}euclidean} (for nonempty coordinate index, that is \(d_x\ge1\)) and \texttt{mean\allowbreak{}Lip\_\allowbreak{}euclidean\_\allowbreak{}zero}, \texttt{sampler\allowbreak{}Force\allowbreak{}Lip\_\allowbreak{}euclidean\_\allowbreak{}zero} (for \(d_x=0\)).
\end{leanbox}
\subsection{\texorpdfstring{The empirical mean force is globally Lipschitz}{The empirical mean force is globally Lipschitz}}


\textbf{Claim.} Let the force data satisfy N1 (in particular \(N\ge1\)). Then for all \((\theta,c),(\theta',c')\in\Theta\times\mathcal C\),


\[\bigl\|\bar F(\theta,c)-\bar F(\theta',c')\bigr\|\le L_F\max\{\|\theta-\theta'\|,\|c-c'\|\}.\]


\textbf{Proof.} \emph{Averaging.} Write \(x=\mathcal Ec\) and \(y=\mathcal Ec'\). By linearity of the finite sum,


\[\bar F(\theta,c)-\bar F(\theta',c')=\frac1N\sum_{i=1}^N\bigl(F(\theta,x_i)-F(\theta',y_i)\bigr).\]


Since \(N\ge1\), the factor \(1/N\) is positive, and the triangle inequality for finite sums gives


\[\bigl\|\bar F(\theta,c)-\bar F(\theta',c')\bigr\|\le\frac1N\sum_{i=1}^N\bigl\|F(\theta,x_i)-F(\theta',y_i)\bigr\|.\]


For each \(i\), (1) and \(\|x_i-y_i\|\le\max_j\|x_j-y_j\|=\|x-y\|_\infty\) give, because \(\ell_F\ge0\),


\[\bigl\|F(\theta,x_i)-F(\theta',y_i)\bigr\|\le\ell_F\bigl(\|\theta-\theta'\|+\|x-y\|_\infty\bigr).\]


The right-hand side does not depend on \(i\); summing \(N\) copies and multiplying by \(1/N\) leaves it unchanged. Hence


\[\bigl\|\bar F(\theta,c)-\bar F(\theta',c')\bigr\|\le\ell_F\bigl(\|\theta-\theta'\|+\|\mathcal Ec-\mathcal Ec'\|_\infty\bigr).\qquad(3)\]


\emph{Change of coordinates.} By linearity \(\mathcal Ec-\mathcal Ec'=\mathcal E(c-c')\), and by the definition of the operator norm \(\|\mathcal E(c-c')\|_\infty\le\|\mathcal E\|\,\|c-c'\|\). Put \(\delta=\max\{\|\theta-\theta'\|,\|c-c'\|\}\), the distance of the two points in \(\Theta\times\mathcal C\). Then \(\|\theta-\theta'\|\le\delta\) and \(\|c-c'\|\le\delta\), so


\[\|\theta-\theta'\|+\|\mathcal Ec-\mathcal Ec'\|_\infty\le\delta+\|\mathcal E\|\,\delta=(1+\|\mathcal E\|)\,\delta .\]


Multiplying by \(\ell_F\ge0\) and combining with (3) gives \(\|\bar F(\theta,c)-\bar F(\theta',c')\|\le\ell_F(1+\|\mathcal E\|)\delta=L_F\delta\). \(\square\)




The assumption \(N\ge1\) is used only to write the empirical average as an average of \(N\) terms; the bound does not depend on \(N\) except through \(\|\mathcal E\|\). The factor \(1+\|\mathcal E\|\) appears because (1) controls the sum of the two distances, whereas \(\Theta\times\mathcal C\) is normed by their maximum.

\begin{leanbox}{Lean: \texttt{meanForce\_\allowbreak{}lipschitz}}
\leansource{\texttt{meanForce\_\allowbreak{}lipschitz} \quad PBGSO/\allowbreak{}ParticleMethodLp.lean, lines 174–192}
\begin{leancode}
\leanline{0}{theorem\ meanForce\_lipschitz\ :\ LipschitzWith\ D.meanLip\ D.meanForce}
\end{leancode}
\leanlabel{How the Lean proof follows the argument}
The statement \texttt{Lipschitz\allowbreak{}With D.\allowbreak{}mean\allowbreak{}Lip D.\allowbreak{}mean\allowbreak{}Force} uses Mathlib's predicate \texttt{Lipschitz\allowbreak{}With K f} for \texttt{K : \ensuremath{\mathbb R}\ensuremath{\ge}0}: for all arguments the (extended) distance of the images is at most \(K\) times the distance of the arguments. In a normed space this is the inequality of the claim with \(K=L_F\).



\begin{enumerate}


\item \texttt{apply Lipschitz\allowbreak{}With.\allowbreak{}of\_\allowbreak{}dist\_\allowbreak{}le\_\allowbreak{}mul} reduces the goal to the inequality between distances for arbitrary \texttt{s t}, and \texttt{simp only [dist\_\allowbreak{}eq\_\allowbreak{}norm]} writes distances as norms of differences.

\item \texttt{hm} is inequality (3). It is the lemma \texttt{particle\allowbreak{}Mean\allowbreak{}Force\_\allowbreak{}lipschitz} from the companion file on particle forces, applied with \texttt{D.\allowbreak{}nonempty} (\(N\ge1\)), the constant \texttt{D.\allowbreak{}parameter\allowbreak{}Lip}, its nonnegativity \texttt{D.\allowbreak{}parameter\allowbreak{}Lip.\allowbreak{}coe\_\allowbreak{}nonneg} and assumption (1) \texttt{D.\allowbreak{}parameter\_\allowbreak{}bound}; that lemma states \(\|N^{-1}\sum_iF(\theta,x_i)-N^{-1}\sum_iF(\theta',y_i)\|\le\ell_F(\|\theta-\theta'\|+\|x-y\|_\infty)\) and its proof is the averaging argument above.

\item \texttt{he} is the bound \(\|\mathcal Ec-\mathcal Ec'\|_\infty\le\|\mathcal E\|\|c-c'\|\), the Lipschitz property of a bounded linear map (\texttt{D.\allowbreak{}cloud.\allowbreak{}to\allowbreak{}Continuous\allowbreak{}Linear\allowbreak{}Map.\allowbreak{}lipschitz\allowbreak{}With.\allowbreak{}norm\_\allowbreak{}sub\_\allowbreak{}le}). \texttt{h\allowbreak{}A} and \texttt{h\allowbreak{}C} are \(\|\theta-\theta'\|\le\delta\) and \(\|c-c'\|\le\delta\), from \texttt{norm\_\allowbreak{}fst\_\allowbreak{}le} and \texttt{norm\_\allowbreak{}snd\_\allowbreak{}le} (each component of a pair has norm at most the maximum norm of the pair); \texttt{henc} multiplies \texttt{h\allowbreak{}C} by \(\|\mathcal E\|\ge0\).

\item The \texttt{calc} block starts from \texttt{hm}; \texttt{gcongr} replaces \(\|\theta-\theta'\|\) by \(\delta\) and \(\|\mathcal Ec-\mathcal Ec'\|_\infty\) by \(\|\mathcal E\|\delta\) (the latter via \texttt{he.\allowbreak{}trans henc}), which is the change of coordinates. The final equality unfolds \texttt{mean\allowbreak{}Lip}, converts the product in \texttt{\ensuremath{\mathbb R}\ensuremath{\ge}0} into real numbers with \texttt{NNReal.\allowbreak{}coe\_\allowbreak{}mul} and \texttt{NNReal.\allowbreak{}coe\_\allowbreak{}add}, and closes with \texttt{ring}.

\end{enumerate}
\end{leanbox}
\subsection{\texorpdfstring{The stacked sampler force is globally Lipschitz}{The stacked sampler force is globally Lipschitz}}


\textbf{Claim.} Under N1, for all \((\theta,c),(\theta',c')\in\Theta\times\mathcal C\),


\[\bigl\|\mathcal F_x(\theta,c)-\mathcal F_x(\theta',c')\bigr\|\le L_x\max\{\|\theta-\theta'\|,\|c-c'\|\}.\]


\textbf{Proof.} Write \(x=\mathcal Ec\), \(y=\mathcal Ec'\) and \(\Phi(\theta,x)=(F_x(\theta,x_i))_{i=1}^N\in\mathcal X^N\), so that \(\mathcal F_x(\theta,c)=\mathcal E^{-1}\Phi(\theta,x)\). For each coordinate, (2) and \(\|x_i-y_i\|\le\|x-y\|_\infty\) give


\[\bigl\|F_x(\theta,x_i)-F_x(\theta',y_i)\bigr\|\le\ell_x\bigl(\|\theta-\theta'\|+\|x-y\|_\infty\bigr).\]


A vector of \(\mathcal X^N\) has norm at most a number \(r\ge0\) if and only if each of its coordinates has norm at most \(r\); the right-hand side above is nonnegative, so


\[\bigl\|\Phi(\theta,x)-\Phi(\theta',y)\bigr\|_\infty\le\ell_x\bigl(\|\theta-\theta'\|+\|x-y\|_\infty\bigr).\qquad(4)\]


No averaging and no factor depending on \(N\) occurs in (4). By linearity of \(\mathcal E^{-1}\) and the definition of its operator norm,


\[\bigl\|\mathcal F_x(\theta,c)-\mathcal F_x(\theta',c')\bigr\|=\bigl\|\mathcal E^{-1}\bigl(\Phi(\theta,x)-\Phi(\theta',y)\bigr)\bigr\|\le\|\mathcal E^{-1}\|\,\bigl\|\Phi(\theta,x)-\Phi(\theta',y)\bigr\|_\infty .\]


As in the previous step, with \(\delta=\max\{\|\theta-\theta'\|,\|c-c'\|\}\) we have \(\|\theta-\theta'\|+\|x-y\|_\infty\le(1+\|\mathcal E\|)\delta\). Inserting this in (4) and multiplying by \(\|\mathcal E^{-1}\|\ge0\) gives the bound \(\|\mathcal E^{-1}\|\ell_x(1+\|\mathcal E\|)\delta=L_x\delta\). \(\square\)




For the Euclidean encoding the factor \(\|\mathcal E^{-1}\|=\sqrt N\) is the cost of measuring the cloud in the Euclidean norm of \(\mathbb R^{Nd_x}\) while the force is controlled particle by particle.

\begin{leanbox}{Lean: \texttt{samplerForce\_\allowbreak{}lipschitz}}
\leansource{\texttt{samplerForce\_\allowbreak{}lipschitz} \quad PBGSO/\allowbreak{}ParticleMethodLp.lean, lines 195–216}
\begin{leancode}
\leanline{0}{theorem\ samplerForce\_lipschitz\ :\ LipschitzWith\ D.samplerForceLip\ D.samplerForce}
\end{leancode}
\leanlabel{How the Lean proof follows the argument}
\texttt{sampler\allowbreak{}Force\_\allowbreak{}lipschitz} states \texttt{Lipschitz\allowbreak{}With D.\allowbreak{}sampler\allowbreak{}Force\allowbreak{}Lip D.\allowbreak{}sampler\allowbreak{}Force} and follows the same pattern as the previous step.



\begin{enumerate}


\item \texttt{hg} is (4): the lemma \texttt{particle\allowbreak{}Sampler\allowbreak{}Force\_\allowbreak{}lipschitz} from the companion file on particle forces, applied to assumption (2) \texttt{D.\allowbreak{}sampler\_\allowbreak{}bound}. Its proof bounds every coordinate and uses the characterisation of the supremum norm described above.

\item \texttt{he} is the bound for \(\mathcal E\) and \texttt{hi} the Lipschitz bound of the bounded linear map \(\mathcal E^{-1}\) (\texttt{D.\allowbreak{}cloud.\allowbreak{}symm.\allowbreak{}to\allowbreak{}Continuous\allowbreak{}Linear\allowbreak{}Map.\allowbreak{}lipschitz\allowbreak{}With.\allowbreak{}norm\_\allowbreak{}sub\_\allowbreak{}le}), applied to the two stacked force vectors. \texttt{h\allowbreak{}A}, \texttt{h\allowbreak{}C}, \texttt{henc} are as before.

\item The \texttt{calc} block applies \texttt{hi}, then (inside \texttt{gcongr}) \texttt{hg} together with \texttt{he.\allowbreak{}trans henc}, and finally unfolds \texttt{sampler\allowbreak{}Force\allowbreak{}Lip} and closes with \texttt{ring}.

\end{enumerate}
\end{leanbox}
\subsection{\texorpdfstring{Pairing two square-integrable random variables}{Pairing two square-integrable random variables}}


Let \((\Omega,\mu)\) be a measure space and \(E\) a real normed space. We write \(L^2(\Omega;E)\) for the space of equivalence classes, modulo equality \(\mu\)-almost everywhere, of strongly measurable maps \(Y:\Omega\to E\) with \(\int\|Y\|^2d\mu<\infty\), normed by \(\|Y\|_{L^2}=(\int\|Y\|^2d\mu)^{1/2}\).




\textbf{Claim.} For \(a\in L^2(\Omega;\Theta)\) and \(X\in L^2(\Omega;\mathcal C)\) there is an element \((a,X)\in L^2(\Omega;\Theta\times\mathcal C)\) equal almost everywhere to \(\omega\mapsto(a(\omega),X(\omega))\). For all \(a,b\in L^2(\Omega;\Theta)\) and \(X,Y\in L^2(\Omega;\mathcal C)\),


\[\bigl\|(a,X)-(b,Y)\bigr\|_{L^2}\le\|a-b\|_{L^2}+\|X-Y\|_{L^2}.\]


No independence or other relation between \(a\) and \(X\) is required.




\textbf{Proof.} Let \(\iota_1(\theta)=(\theta,0)\) and \(\iota_2(c)=(0,c)\). These are bounded linear maps into \(\Theta\times\mathcal C\) with \(\|\iota_1(\theta)\|=\max\{\|\theta\|,0\}=\|\theta\|\) and \(\|\iota_2(c)\|=\|c\|\). Composition with a bounded linear map \(T\) sends \(L^2\) to \(L^2\), since \(\|T(Y(\omega))\|\le\|T\|\|Y(\omega)\|\) pointwise. Define \((a,X)=\iota_1\circ a+\iota_2\circ X\); almost everywhere it equals \(\omega\mapsto(a(\omega),0)+(0,X(\omega))=(a(\omega),X(\omega))\). By linearity of composition, \((a,X)-(b,Y)=\iota_1\circ(a-b)+\iota_2\circ(X-Y)\). Minkowski's inequality in \(L^2\) gives


\[\bigl\|(a,X)-(b,Y)\bigr\|_{L^2}\le\|\iota_1\circ(a-b)\|_{L^2}+\|\iota_2\circ(X-Y)\|_{L^2}.\]


Pointwise \(\|\iota_1(a(\omega)-b(\omega))\|=\|a(\omega)-b(\omega)\|\), and the \(L^2\) norm is monotone under almost-everywhere domination of the pointwise norms; hence \(\|\iota_1\circ(a-b)\|_{L^2}\le\|a-b\|_{L^2}\), and similarly for \(\iota_2\). \(\square\)

\begin{leanbox}{Lean: \texttt{pairLp}, \texttt{pairLp\_\allowbreak{}coe}, \texttt{pairLp\_\allowbreak{}norm\_\allowbreak{}sub\_\allowbreak{}le}}
\leansource{\texttt{pairLp} \quad PBGSO/\allowbreak{}ParticleMethodLp.lean, lines 27–29}
\begin{leancode}
\leanline{0}{def\ pairLp\ (a\ :\ Lp\ A\ 2\ Q)\ (x\ :\ Lp\ C\ 2\ Q)\ :\ Lp\ (A\ ×\ C)\ 2\ Q\ :=}
\leanline{2}{((ContinuousLinearMap.inl\ ℝ\ A\ C).compLpL\ 2\ Q)\ a\ +}
\leanline{4}{((ContinuousLinearMap.inr\ ℝ\ A\ C).compLpL\ 2\ Q)\ x}
\end{leancode}
\leansource{\texttt{pairLp\_\allowbreak{}coe} \quad PBGSO/\allowbreak{}ParticleMethodLp.lean, lines 32–43}
\begin{leancode}
\leanline{0}{theorem\ pairLp\_coe\ (a\ :\ Lp\ A\ 2\ Q)\ (x\ :\ Lp\ C\ 2\ Q)\ :}
\leanline{4}{pairLp\ Q\ a\ x\ =ᵐ[Q]\ fun\ ω\ ↦\ (a\ ω,\ x\ ω)}
\end{leancode}
\leansource{\texttt{pairLp\_\allowbreak{}norm\_\allowbreak{}sub\_\allowbreak{}le} \quad PBGSO/\allowbreak{}ParticleMethodLp.lean, lines 61–68}
\begin{leancode}
\leanline{0}{theorem\ pairLp\_norm\_sub\_le\ (a\ b\ :\ Lp\ A\ 2\ Q)\ (x\ y\ :\ Lp\ C\ 2\ Q)\ :}
\leanline{4}{‖pairLp\ Q\ a\ x\ -\ pairLp\ Q\ b\ y‖\ ≤\ ‖a\ -\ b‖\ +\ ‖x\ -\ y‖}
\end{leancode}
\leanlabel{How the Lean proof follows the argument}
\texttt{Lp A 2 Q} is Mathlib's space \(L^2(\Omega;\Theta)\) for the measure \texttt{Q}; an element is coerced to a function \(\omega\mapsto a(\omega)\) (a representative). \texttt{pair\allowbreak{}Lp Q a x} is defined as \texttt{((Continuous\allowbreak{}Linear\allowbreak{}Map.\allowbreak{}inl \ensuremath{\mathbb R} A C).\allowbreak{}comp\allowbreak{}Lp\allowbreak{}L 2 Q) a + ((Continuous\allowbreak{}Linear\allowbreak{}Map.\allowbreak{}inr \ensuremath{\mathbb R} A C).\allowbreak{}comp\allowbreak{}Lp\allowbreak{}L 2 Q) x}: \texttt{inl} and \texttt{inr} are \(\iota_1,\iota_2\), and \texttt{comp\allowbreak{}Lp\allowbreak{}L} turns a bounded linear map into the bounded linear map \(Y\mapsto T\circ Y\) on \texttt{Lp}.



\begin{enumerate}


\item \texttt{pair\allowbreak{}Lp\_\allowbreak{}coe} states \texttt{pair\allowbreak{}Lp Q a x =ᵐ[Q] fun \ensuremath{\omega} \ensuremath{\mapsto} (a \ensuremath{\omega}, x \ensuremath{\omega})}, where \texttt{=ᵐ[Q]} is equality \texttt{Q}-almost everywhere. \texttt{filter\_\allowbreak{}upwards} intersects the almost-everywhere identities \texttt{Lp.\allowbreak{}coe\allowbreak{}Fn\_\allowbreak{}add} (a representative of a sum is the sum of representatives) and \texttt{coe\allowbreak{}Fn\_\allowbreak{}comp\allowbreak{}Lp} (a representative of \(T\circ Y\) is \(T\) applied to a representative of \(Y\)).

\item In \texttt{pair\allowbreak{}Lp\_\allowbreak{}norm\_\allowbreak{}sub\_\allowbreak{}le}, \texttt{he} is the linearity identity, proved by \texttt{map\_\allowbreak{}sub} and \texttt{module}. Then \texttt{norm\_\allowbreak{}add\_\allowbreak{}le} is Minkowski's inequality in the normed space \texttt{Lp}, and the private lemmas \texttt{inl\allowbreak{}Lp\_\allowbreak{}norm\_\allowbreak{}le}, \texttt{inr\allowbreak{}Lp\_\allowbreak{}norm\_\allowbreak{}le} give the two domination bounds; each is proved by \texttt{Lp.\allowbreak{}norm\_\allowbreak{}le\_\allowbreak{}norm\_\allowbreak{}of\_\allowbreak{}ae\_\allowbreak{}le}, which states that \(\|f(\omega)\|\le\|g(\omega)\|\) almost everywhere implies \(\|f\|_{L^2}\le\|g\|_{L^2}\).

\end{enumerate}
\end{leanbox}
\subsection{\texorpdfstring{The forces at random arguments}{The forces at random arguments}}


From now on \(\mu=\mathbb P\) is a probability measure on \(\Omega\).




\textbf{Claim.} For \(\theta\in L^2(\Omega;\Theta)\) and \(X\in L^2(\Omega;\mathcal C)\) there are elements \(\bar F(\theta,X)\in L^2(\Omega;\Theta)\) and \(\mathcal F_x(\theta,X)\in L^2(\Omega;\mathcal C)\) such that, for \(\mathbb P\)-almost every \(\omega\),


\[\bar F(\theta,X)(\omega)=\frac1N\sum_{i=1}^NF\bigl(\theta(\omega),(\mathcal EX(\omega))_i\bigr),\qquad \mathcal F_x(\theta,X)(\omega)=\mathcal F_x\bigl(\theta(\omega),X(\omega)\bigr),\]


and, equivalently for the second identity, \(\bigl(\mathcal E\,\mathcal F_x(\theta,X)(\omega)\bigr)_i=F_x\bigl(\theta(\omega),(\mathcal EX(\omega))_i\bigr)\) for every \(i\), almost surely.




\textbf{Proof.} Let \(g:E\to K\) be a Lipschitz map between normed spaces with constant \(L\) and let \(Y\in L^2(\Omega;E)\). Then \(g\circ Y\) is strongly measurable (a continuous function of a strongly measurable map), and \(\|g(Y(\omega))\|\le\|g(0)\|+L\|Y(\omega)\|\). By Minkowski's inequality and \(\mathbb P(\Omega)=1\), \(\|g\circ Y\|_{L^2}\le\|g(0)\|+L\|Y\|_{L^2}<\infty\), so \(g\circ Y\in L^2(\Omega;K)\). Moreover \(\|g(Y(\omega))-g(Y'(\omega))\|\le L\|Y(\omega)-Y'(\omega)\|\) pointwise, hence


\[\|g\circ Y-g\circ Y'\|_{L^2}\le L\,\|Y-Y'\|_{L^2}.\qquad(5)\]


Apply this with \(g=\bar F\) and \(g=\mathcal F_x\), which are Lipschitz by the two previous steps, and with \(Y=(\theta,X)\), the pair of the previous step. Since \((\theta,X)(\omega)=(\theta(\omega),X(\omega))\) almost surely, the two displayed identities follow from the definitions of \(\bar F\) and \(\mathcal F_x\); the coordinate identity follows by applying \(\mathcal E\) and using \(\mathcal E\mathcal E^{-1}=\mathrm{id}\). \(\square\)




The lifted mean force is a random vector: at each sample it is the empirical average over the sampled cloud, with the same sampled parameter for every particle. It is not an expectation over the law of the cloud.

\begin{leanbox}{Lean: \texttt{meanForceLp}, \texttt{samplerForceLp}, \texttt{meanForceLp\_\allowbreak{}coe} and 2 more}
\leansource{\texttt{meanForceLp} \quad PBGSO/\allowbreak{}ParticleMethodLp.lean, lines 223–224}
\begin{leancode}
\leanline{0}{def\ meanForceLp\ (a\ :\ Lp\ A\ 2\ Q)\ (x\ :\ Lp\ C\ 2\ Q)\ :\ Lp\ A\ 2\ Q\ :=}
\leanline{2}{PBGSO.liftMap\ Q\ D.meanForce\ D.meanForce\_lipschitz\ (pairLp\ Q\ a\ x)}
\end{leancode}
\leansource{\texttt{samplerForceLp} \quad PBGSO/\allowbreak{}ParticleMethodLp.lean, lines 227–228}
\begin{leancode}
\leanline{0}{def\ samplerForceLp\ (a\ :\ Lp\ A\ 2\ Q)\ (x\ :\ Lp\ C\ 2\ Q)\ :\ Lp\ C\ 2\ Q\ :=}
\leanline{2}{PBGSO.liftMap\ Q\ D.samplerForce\ D.samplerForce\_lipschitz\ (pairLp\ Q\ a\ x)}
\end{leancode}
\leansource{\texttt{meanForceLp\_\allowbreak{}coe} \quad PBGSO/\allowbreak{}ParticleMethodLp.lean, lines 232–239}
\begin{leancode}
\leanline{0}{theorem\ meanForceLp\_coe\ (a\ :\ Lp\ A\ 2\ Q)\ (x\ :\ Lp\ C\ 2\ Q)\ :}
\leanline{4}{D.meanForceLp\ Q\ a\ x\ =ᵐ[Q]\ fun\ ω\ ↦}
\leanline{6}{(N\ :\ ℝ)⁻¹\ •\ ∑\ i,\ D.F\ (a\ ω)\ (D.cloud\ (x\ ω)\ i)}
\end{leancode}
\leansource{\texttt{samplerForceLp\_\allowbreak{}coe} \quad PBGSO/\allowbreak{}ParticleMethodLp.lean, lines 243–250}
\begin{leancode}
\leanline{0}{theorem\ samplerForceLp\_coe\ (a\ :\ Lp\ A\ 2\ Q)\ (x\ :\ Lp\ C\ 2\ Q)\ :}
\leanline{4}{∀ᵐ\ ω\ ∂Q,\ ∀\ i\ :\ Fin\ N,}
\leanline{6}{D.cloud\ (D.samplerForceLp\ Q\ a\ x\ ω)\ i\ =\ D.G\ (a\ ω)\ (D.cloud\ (x\ ω)\ i)}
\end{leancode}
\leansource{\texttt{samplerForceLp\_\allowbreak{}coe\_\allowbreak{}cloud} \quad PBGSO/\allowbreak{}ParticleMethodLp.lean, lines 253–258}
\begin{leancode}
\leanline{0}{theorem\ samplerForceLp\_coe\_cloud\ (a\ :\ Lp\ A\ 2\ Q)\ (x\ :\ Lp\ C\ 2\ Q)\ :}
\leanline{4}{D.samplerForceLp\ Q\ a\ x\ =ᵐ[Q]\ fun\ ω\ ↦\ D.samplerForce\ (a\ ω,\ x\ ω)}
\end{leancode}
\leanlabel{How the Lean proof follows the argument}
\texttt{mean\allowbreak{}Force\allowbreak{}Lp Q a x} is \texttt{PBGSO.\allowbreak{}lift\allowbreak{}Map Q D.\allowbreak{}mean\allowbreak{}Force D.\allowbreak{}mean\allowbreak{}Force\_\allowbreak{}lipschitz (pair\allowbreak{}Lp Q a x)}, and \texttt{sampler\allowbreak{}Force\allowbreak{}Lp} is defined in the same way. The operation \texttt{lift\allowbreak{}Map} (from the file on general processes) realises \(g\circ Y\) for a Lipschitz map \texttt{g}: it is defined as the composition with the centred map \(g-g(0)\), which vanishes at the origin, plus the constant random variable \(g(0)\); its lemma \texttt{lift\allowbreak{}Map\_\allowbreak{}coe} states that it equals \(\omega\mapsto g(Y(\omega))\) almost everywhere, and \texttt{lift\allowbreak{}Map\_\allowbreak{}lipschitz} gives (5). The section variable \texttt{[Is\allowbreak{}Probability\allowbreak{}Measure Q]} is the hypothesis \(\mathbb P(\Omega)=1\).



\begin{enumerate}


\item \texttt{mean\allowbreak{}Force\allowbreak{}Lp\_\allowbreak{}coe} combines \texttt{lift\allowbreak{}Map\_\allowbreak{}coe} with \texttt{pair\allowbreak{}Lp\_\allowbreak{}coe} by \texttt{filter\_\allowbreak{}upwards} and rewrites; the final \texttt{rfl} unfolds \texttt{mean\allowbreak{}Force} and \texttt{particle\allowbreak{}Mean\allowbreak{}Force}.

\item \texttt{sampler\allowbreak{}Force\allowbreak{}Lp\_\allowbreak{}coe} is the coordinate form: \texttt{\ensuremath{\forall}ᵐ \ensuremath{\omega} \ensuremath{\partial}Q, \ensuremath{\forall} i : Fin N, D.\allowbreak{}cloud (D.\allowbreak{}sampler\allowbreak{}Force\allowbreak{}Lp Q a x \ensuremath{\omega}) i = D.\allowbreak{}G (a \ensuremath{\omega}) (D.\allowbreak{}cloud (x \ensuremath{\omega}) i)}, where \texttt{\ensuremath{\forall}ᵐ \ensuremath{\omega} \ensuremath{\partial}Q} means "for \texttt{Q}-almost every \(\omega\)". After the same rewriting, \texttt{simp} uses \(\mathcal E\mathcal E^{-1}=\mathrm{id}\).

\item \texttt{sampler\allowbreak{}Force\allowbreak{}Lp\_\allowbreak{}coe\_\allowbreak{}cloud} is the identity \(\mathcal F_x(\theta,X)(\omega)=\mathcal F_x(\theta(\omega),X(\omega))\).

\end{enumerate}
\end{leanbox}
\subsection{\texorpdfstring{The lifted forces satisfy N1 in square mean}{The lifted forces satisfy N1 in square mean}}


\textbf{Claim.} Under N1, for all \(a,b\in L^2(\Omega;\Theta)\) and \(X,Y\in L^2(\Omega;\mathcal C)\),


\[\begin{aligned}\bigl\|\bar F(a,X)-\bar F(b,Y)\bigr\|_{L^2}&\le L_F\bigl(\|a-b\|_{L^2}+\|X-Y\|_{L^2}\bigr),\\ \bigl\|\mathcal F_x(a,X)-\mathcal F_x(b,Y)\bigr\|_{L^2}&\le L_x\bigl(\|a-b\|_{L^2}+\|X-Y\|_{L^2}\bigr).\end{aligned}\]


\textbf{Proof.} By (5) applied to \(g=\bar F\) with constant \(L_F\), \(\|\bar F(a,X)-\bar F(b,Y)\|_{L^2}\le L_F\|(a,X)-(b,Y)\|_{L^2}\). The pairing estimate gives \(\|(a,X)-(b,Y)\|_{L^2}\le\|a-b\|_{L^2}+\|X-Y\|_{L^2}\); multiplying by \(L_F\ge0\) gives the first inequality. The second is identical with \(g=\mathcal F_x\) and \(L_x\). \(\square\)




These two inequalities have exactly the form (1)–(2), now for the random parameter and the random cloud with \(L^2\) norms. They are the force hypotheses of the one-step error estimates for both methods, and they give the global Lipschitz constants of the kinetic and higher-order drifts.

\begin{leanbox}{Lean: \texttt{meanForceLp\_\allowbreak{}bound}, \texttt{samplerForceLp\_\allowbreak{}bound}}
\leansource{\texttt{meanForceLp\_\allowbreak{}bound} \quad PBGSO/\allowbreak{}ParticleMethodLp.lean, lines 262–267}
\begin{leancode}
\leanline{0}{theorem\ meanForceLp\_bound\ (a\ b\ :\ Lp\ A\ 2\ Q)\ (x\ y\ :\ Lp\ C\ 2\ Q)\ :}
\leanline{4}{‖D.meanForceLp\ Q\ a\ x\ -\ D.meanForceLp\ Q\ b\ y‖\ ≤}
\leanline{6}{(D.meanLip\ :\ ℝ)\ *\ (‖a\ -\ b‖\ +\ ‖x\ -\ y‖)}
\end{leancode}
\leansource{\texttt{samplerForceLp\_\allowbreak{}bound} \quad PBGSO/\allowbreak{}ParticleMethodLp.lean, lines 271–276}
\begin{leancode}
\leanline{0}{theorem\ samplerForceLp\_bound\ (a\ b\ :\ Lp\ A\ 2\ Q)\ (x\ y\ :\ Lp\ C\ 2\ Q)\ :}
\leanline{4}{‖D.samplerForceLp\ Q\ a\ x\ -\ D.samplerForceLp\ Q\ b\ y‖\ ≤}
\leanline{6}{(D.samplerForceLip\ :\ ℝ)\ *\ (‖a\ -\ b‖\ +\ ‖x\ -\ y‖)}
\end{leancode}
\leanlabel{How the Lean proof follows the argument}
Each proof is one line: \texttt{(PBGSO.\allowbreak{}lift\allowbreak{}Map\_\allowbreak{}lipschitz Q D.\allowbreak{}mean\allowbreak{}Force D.\allowbreak{}mean\allowbreak{}Force\_\allowbreak{}lipschitz).\allowbreak{}norm\_\allowbreak{}sub\_\allowbreak{}le} is (5) for the lifted map, applied to the two pairs; \texttt{.\allowbreak{}trans} chains it with \texttt{pair\allowbreak{}Lp\_\allowbreak{}norm\_\allowbreak{}sub\_\allowbreak{}le} multiplied by the nonnegative constant (\texttt{mul\_\allowbreak{}le\_\allowbreak{}mul\_\allowbreak{}of\_\allowbreak{}nonneg\_\allowbreak{}left} with \texttt{D.\allowbreak{}mean\allowbreak{}Lip.\allowbreak{}coe\_\allowbreak{}nonneg}). The sampler version uses \texttt{D.\allowbreak{}sampler\allowbreak{}Force\_\allowbreak{}lipschitz} and \texttt{D.\allowbreak{}sampler\allowbreak{}Force\allowbreak{}Lip.\allowbreak{}coe\_\allowbreak{}nonneg}.
\end{leanbox}
\subsection{\texorpdfstring{One step of either method acts sample by sample}{One step of either method acts sample by sample}}


Write \(\mathbf F(\theta,c)=(\bar F(\theta,c),\mathcal F_x(\theta,c))\in\Theta\times\mathcal C\) for the full coupled force. A kinetic phase is a pair \(z=(p,q)\) of a momentum \(p=(m,V)\) and a position \(q=(\theta,X)\), both in \(\Theta\times\mathcal C\). For \(h,\eta,\sigma\in\mathbb R\) and innovations \(\xi,\xi'\in\Theta\times\mathcal C\), one OBABO step is the composition of five maps:


\[\begin{aligned}\mathrm O_\xi&:(p,q)\mapsto(\eta p+\sigma\xi,\,q),\\ \mathrm B&:(p,q)\mapsto\bigl(p-\tfrac h2\mathbf F(q),\,q\bigr),\\ \mathrm A&:(p,q)\mapsto(p,\,q+hp),\end{aligned}\qquad\text{one step}=\mathrm O_{\xi'}\circ\mathrm B\circ\mathrm A\circ\mathrm B\circ\mathrm O_{\xi}.\]


With \(\eta=e^{-\gamma h/2}\), \(\sigma=\sqrt{1-e^{-\gamma h}}\) and innovations whose parameter component is zero, this is the paper's M-MFPO update: the O stage is the exact Ornstein–Uhlenbeck flow on the particle velocities together with the damping \(m\mapsto\eta m\) of the parameter momentum, B is the half kick by both forces, and A is the joint drift of parameter and positions. One HOL-MYPGD step with coefficients \(\mu_{ij}=\mu_{ij}(\alpha,\gamma,h)\) and innovation triple \((\xi^1,\xi^2,\xi^3)\) maps \((\theta,X,U,V)\) to \((\theta',X',U',V')\), where


\[\begin{aligned}\theta'&=\theta-h\bar F(\theta,X),\qquad \Delta=h\,\mathcal F_x(\theta',X),\\ X'&=X-\tfrac h2\Delta+\mu_{12}U+\mu_{13}V+\xi^1,\\ U'&=-\Delta+\mu_{22}U+\mu_{23}V+\xi^2,\\ V'&=\mu_{31}\Delta+\mu_{32}U+\mu_{33}V+\xi^3 .\end{aligned}\]


The parameter is updated first and the sampler force is evaluated at the new parameter, as in the paper.




\textbf{Claim.} Let all inputs \(z\), \(\xi\), \(\xi'\) (respectively \((\theta,X,U,V)\) and \((\xi^1,\xi^2,\xi^3)\)) be square-integrable random variables on \(\Omega\), and compute the step in \(L^2\) with the lifted forces of the previous steps. Then, for \(\mathbb P\)-almost every \(\omega\), the value of the result at \(\omega\) equals the same step applied to the sampled values \(z(\omega),\xi(\omega),\xi'(\omega)\) (respectively the sampled state and innovations) with the forces \(\mathbf F\), \(\bar F\), \(\mathcal F_x\) of the second step.




\textbf{Proof.} Each stage is built from sums, scalar multiples and the lifted forces. For sums and scalar multiples in \(L^2\), a representative of the result is the corresponding combination of representatives, almost surely; for the lifted forces this is the previous claim. Hence each of the five OBABO stages (respectively each of the defining formulas of the HOL step, starting with \(\theta'\) and \(\Delta\)) holds samplewise outside a null set. A finite union of null sets is null, so all stages hold simultaneously for almost every \(\omega\), and composing them gives the claim. \(\square\)




Consequently the iterates studied in Theorems New A and New B, which are computed in \(L^2\), are almost surely the iterates of the algorithm run on each sample.

\begin{leanbox}{Lean: \texttt{rawCoupledForce}, \texttt{coupledForceLp\_\allowbreak{}coe}, \texttt{obaboStep\_\allowbreak{}Lp\_\allowbreak{}coe} and 1 more}
\leansource{\texttt{rawCoupledForce} \quad PBGSO/\allowbreak{}ParticleMethodLp.lean, lines 279–279}
\begin{leancode}
\leanline{0}{def\ rawCoupledForce\ (q\ :\ A\ ×\ C)\ :\ A\ ×\ C\ :=\ (D.meanForce\ q,\ D.samplerForce\ q)}
\end{leancode}
\leansource{\texttt{coupledForceLp\_\allowbreak{}coe} \quad PBGSO/\allowbreak{}ParticleMethodLp.lean, lines 283–289}
\begin{leancode}
\leanline{0}{theorem\ coupledForceLp\_coe\ (q\ :\ Lp\ A\ 2\ Q\ ×\ Lp\ C\ 2\ Q)\ :}
\leanline{4}{∀ᵐ\ ω\ ∂Q,}
\leanline{6}{jointValue\ Q\ (coupledForce\ (D.meanForceLp\ Q)\ (D.samplerForceLp\ Q)\ q)\ ω\ =}
\leanline{8}{D.rawCoupledForce\ (jointValue\ Q\ q\ ω)}
\end{leancode}
\leansource{\texttt{obaboStep\_\allowbreak{}Lp\_\allowbreak{}coe} \quad PBGSO/\allowbreak{}ParticleMethodLp.lean, lines 308–326}
\begin{leancode}
\leanline{0}{theorem\ obaboStep\_Lp\_coe\ (h\ η\ σ\ :\ ℝ)}
\leanline{4}{(z\ :\ (Lp\ A\ 2\ Q\ ×\ Lp\ C\ 2\ Q)\ ×\ (Lp\ A\ 2\ Q\ ×\ Lp\ C\ 2\ Q))}
\leanline{4}{(ξ\ ξ'\ :\ Lp\ A\ 2\ Q\ ×\ Lp\ C\ 2\ Q)\ :}
\leanline{4}{∀ᵐ\ ω\ ∂Q,}
\leanline{6}{phaseValue\ Q\ (KineticLangevin.obaboStep}
\leanline{8}{(coupledForce\ (D.meanForceLp\ Q)\ (D.samplerForceLp\ Q))\ h\ η\ σ\ z\ ξ\ ξ')\ ω\ =}
\leanline{6}{KineticLangevin.obaboStep\ D.rawCoupledForce\ h\ η\ σ}
\leanline{8}{(phaseValue\ Q\ z\ ω)\ (jointValue\ Q\ ξ\ ω)\ (jointValue\ Q\ ξ'\ ω)}
\end{leancode}
\leansource{\texttt{higherOrderStep\_\allowbreak{}Lp\_\allowbreak{}coe} \quad PBGSO/\allowbreak{}ParticleMethodLp.lean, lines 331–389}
\begin{leancode}
\leanline{0}{theorem\ higherOrderStep\_Lp\_coe\ (α\ γ\ h\ :\ ℝ)}
\leanline{4}{(z\ :\ HOLState\ (Lp\ A\ 2\ Q)\ (Lp\ C\ 2\ Q))\ (ξ\ :\ Lp\ C\ 2\ Q\ ×\ Lp\ C\ 2\ Q\ ×\ Lp\ C\ 2\ Q)\ :}
\leanline{4}{∀ᵐ\ ω\ ∂Q,}
\leanline{6}{holValue\ Q\ (higherOrderStep\ (D.meanForceLp\ Q)\ (D.samplerForceLp\ Q)\ α\ γ\ h\ z\ ξ)\ ω\ =}
\leanline{8}{higherOrderStep\ (fun\ a\ x\ ↦\ D.meanForce\ (a,\ x))}
\leanline{10}{(fun\ a\ x\ ↦\ D.samplerForce\ (a,\ x))\ α\ γ\ h\ (holValue\ Q\ z\ ω)}
\leanline{10}{(ξ.1\ ω,\ ξ.2.1\ ω,\ ξ.2.2\ ω)}
\end{leancode}
\leanlabel{How the Lean proof follows the argument}
\texttt{raw\allowbreak{}Coupled\allowbreak{}Force} is \(\mathbf F\) on \texttt{A \ensuremath{\times} C}. \texttt{joint\allowbreak{}Value}, \texttt{phase\allowbreak{}Value} and \texttt{hol\allowbreak{}Value} evaluate pairs, phases and four-component states of \texttt{Lp} elements at a sample. The map \texttt{Kinetic\allowbreak{}Langevin.\allowbreak{}obabo\allowbreak{}Step} from the kinetic Langevin library applies, in this order, \texttt{ornstein\allowbreak{}Step \ensuremath{\delta} \ensuremath{\sigma} \ensuremath{\xi}}, \texttt{kick\allowbreak{}Step force (h /\allowbreak{} 2)}, \texttt{drift\allowbreak{}Step h}, \texttt{kick\allowbreak{}Step force (h /\allowbreak{} 2)} and \texttt{ornstein\allowbreak{}Step \ensuremath{\delta} \ensuremath{\sigma} \ensuremath{\xi}'}; \texttt{higher\allowbreak{}Order\allowbreak{}Step} is the HOL map with the printed coefficients \texttt{hol\allowbreak{}Coefficients \ensuremath{\alpha} \ensuremath{\gamma} h}.



\begin{enumerate}


\item \texttt{coupled\allowbreak{}Force\allowbreak{}Lp\_\allowbreak{}coe} combines \texttt{mean\allowbreak{}Force\allowbreak{}Lp\_\allowbreak{}coe} and \texttt{sampler\allowbreak{}Force\allowbreak{}Lp\_\allowbreak{}coe\_\allowbreak{}cloud} with \texttt{Prod.\allowbreak{}ext}.

\item The stage lemmas \texttt{ornstein\allowbreak{}Step\_\allowbreak{}Lp\_\allowbreak{}coe}, \texttt{kick\allowbreak{}Step\_\allowbreak{}Lp\_\allowbreak{}coe} and \texttt{drift\allowbreak{}Step\_\allowbreak{}Lp\_\allowbreak{}coe} are each proved from \texttt{joint\allowbreak{}Value\_\allowbreak{}add}, \texttt{joint\allowbreak{}Value\_\allowbreak{}sub}, \texttt{joint\allowbreak{}Value\_\allowbreak{}smul} (samplewise sums and scalar multiples) and \texttt{coupled\allowbreak{}Force\allowbreak{}Lp\_\allowbreak{}coe}.

\item \texttt{obabo\allowbreak{}Step\_\allowbreak{}Lp\_\allowbreak{}coe} names the four intermediate phases \texttt{z0}, …, \texttt{z3}, collects the five stage identities with \texttt{filter\_\allowbreak{}upwards} and rewrites them in reverse order (\texttt{rw [h4, h3, h2, h1, h0]}).

\item \texttt{higher\allowbreak{}Order\allowbreak{}Step\_\allowbreak{}Lp\_\allowbreak{}coe} proves the samplewise forms \texttt{h\ensuremath{\theta}} of \(\theta'\) and \texttt{h\allowbreak{}\ensuremath{\Delta}} of \(\Delta\) (from \texttt{mean\allowbreak{}Force\allowbreak{}Lp\_\allowbreak{}coe} and \texttt{sampler\allowbreak{}Force\allowbreak{}Lp\_\allowbreak{}coe\_\allowbreak{}cloud}), and \texttt{hx}, \texttt{hu}, \texttt{hv} of the three sampler rows from \texttt{linear5\allowbreak{}Lp\_\allowbreak{}coe}, the samplewise form of a linear combination of five \texttt{Lp} elements; after \texttt{filter\_\allowbreak{}upwards}, \texttt{dsimp} unfolds the update and \texttt{module} closes the remaining linear identity.

\end{enumerate}
\end{leanbox}
\subsection{\texorpdfstring{The standard Euclidean cloud encoding}{The standard Euclidean cloud encoding}}


\textbf{Definition.} Let \(d_x\in\mathbb N\) and let \(\mathcal C=\mathbb R^{N\times d_x}\) with the Euclidean norm \(\|x\|^2=\sum_{i=1}^N\sum_{j=1}^{d_x}x_{ij}^2\), and \(\mathcal X=\mathbb R^{d_x}\) with the Euclidean norm. The standard encoding \(\mathcal E_N:\mathbb R^{N\times d_x}\to(\mathbb R^{d_x})^N\) reads the rows:


\[(\mathcal E_Nx)_i=(x_{i1},\dots,x_{id_x}),\qquad (\mathcal E_N^{-1}y)_{ij}=(y_i)_j .\]


It is linear and bijective, and continuous in both directions because all spaces are finite-dimensional. The two displayed formulas hold by definition.




With this encoding the Euclidean norm of a cloud is \(\|x\|^2=\sum_i\|(\mathcal E_Nx)_i\|^2\), the sum of the squared particle norms. Therefore \(\|(\mathcal E_Nx)_i\|\le\|x\|\) for every \(i\), so \(\|\mathcal E_N\|\le1\), and \(\|\mathcal E_N^{-1}y\|^2=\sum_i\|y_i\|^2\le N\|y\|_\infty^2\), so \(\|\mathcal E_N^{-1}\|\le\sqrt N\). For \(N\ge1\) and \(d_x\ge1\) both bounds are attained (by a unit coordinate vector and by the cloud whose particles all equal one unit vector), which gives the constants \(L_F=2\ell_F\), \(L_x=2\sqrt N\ell_x\) stated above. The theorems New A and New B are proved for this encoding; independent standard Gaussian coordinates of \(\mathbb R^{N\times d_x}\) then give independent isotropic particle noises.

\begin{leanbox}{Lean: \texttt{euclideanParticleEncoding}, \texttt{euclideanParticleEncoding\_\allowbreak{}apply}, \texttt{euclideanParticleEncoding\_\allowbreak{}symm\_\allowbreak{}apply}}
\leansource{\texttt{euclideanParticleEncoding} \quad PBGSO/\allowbreak{}ParticleMethodLp.lean, lines 395–400}
\begin{leancode}
\leanline{0}{def\ euclideanParticleEncoding\ (N\ :\ ℕ)\ (I\ :\ Type*)\ [Fintype\ I]\ :}
\leanline{4}{EuclideanSpace\ ℝ\ (Fin\ N\ ×\ I)\ ≃L[ℝ]\ (Fin\ N\ →\ EuclideanSpace\ ℝ\ I)\ :=}
\leanline{2}{((WithLp.linearEquiv\ 2\ ℝ\ (Fin\ N\ ×\ I\ →\ ℝ)).trans}
\leanline{4}{((LinearEquiv.curry\ ℝ\ ℝ\ (Fin\ N)\ I).trans}
\leanline{6}{(LinearEquiv.piCongrRight\ (fun\ \_\ :\ Fin\ N\ ↦}
\leanline{8}{(WithLp.linearEquiv\ 2\ ℝ\ (I\ →\ ℝ)).symm)))).toContinuousLinearEquiv}
\end{leancode}
\leansource{\texttt{euclideanParticleEncoding\_\allowbreak{}apply} \quad PBGSO/\allowbreak{}ParticleMethodLp.lean, lines 403–405}
\begin{leancode}
\leanline{0}{@[simp]\ theorem\ euclideanParticleEncoding\_apply\ (N\ :\ ℕ)\ (I\ :\ Type*)\ [Fintype\ I]}
\leanline{4}{(x\ :\ EuclideanSpace\ ℝ\ (Fin\ N\ ×\ I))\ (i\ :\ Fin\ N)\ (j\ :\ I)\ :}
\leanline{4}{euclideanParticleEncoding\ N\ I\ x\ i\ j\ =\ x\ (i,\ j)}
\end{leancode}
\leansource{\texttt{euclideanParticleEncoding\_\allowbreak{}symm\_\allowbreak{}apply} \quad PBGSO/\allowbreak{}ParticleMethodLp.lean, lines 408–410}
\begin{leancode}
\leanline{0}{@[simp]\ theorem\ euclideanParticleEncoding\_symm\_apply\ (N\ :\ ℕ)\ (I\ :\ Type*)\ [Fintype\ I]}
\leanline{4}{(x\ :\ Fin\ N\ →\ EuclideanSpace\ ℝ\ I)\ (i\ :\ Fin\ N)\ (j\ :\ I)\ :}
\leanline{4}{(euclideanParticleEncoding\ N\ I).symm\ x\ (i,\ j)\ =\ x\ i\ j}
\end{leancode}
\leanlabel{How the Lean proof follows the argument}
\texttt{euclidean\allowbreak{}Particle\allowbreak{}Encoding N I} has type \texttt{Euclidean\allowbreak{}Space \ensuremath{\mathbb R} (Fin N \ensuremath{\times} I) ≃L[\ensuremath{\mathbb R}] (Fin N \ensuremath{\to} Euclidean\allowbreak{}Space \ensuremath{\mathbb R} I)}; \texttt{Euclidean\allowbreak{}Space \ensuremath{\mathbb R} \ensuremath{\iota}} is \(\mathbb R^\iota\) with the Euclidean norm, and the coordinate index \texttt{I} (a finite type) has \(d_x\) elements. The map is assembled from linear equivalences: \texttt{With\allowbreak{}Lp.\allowbreak{}linear\allowbreak{}Equiv} forgets the Euclidean norm, \texttt{Linear\allowbreak{}Equiv.\allowbreak{}curry} turns a function of pairs \((i,j)\) into a function of \(i\) with values functions of \(j\), and \texttt{Linear\allowbreak{}Equiv.\allowbreak{}pi\allowbreak{}Congr\allowbreak{}Right} restores the Euclidean structure on each particle; \texttt{to\allowbreak{}Continuous\allowbreak{}Linear\allowbreak{}Equiv} uses finite dimensionality to make it a continuous linear equivalence. \texttt{euclidean\allowbreak{}Particle\allowbreak{}Encoding\_\allowbreak{}apply} and \texttt{euclidean\allowbreak{}Particle\allowbreak{}Encoding\_\allowbreak{}symm\_\allowbreak{}apply} are the two evaluation formulas, proved by \texttt{rfl}. The norm identities and operator norms are proved in other files: \texttt{cloud\_\allowbreak{}norm\_\allowbreak{}sq} (metric comparisons) and \texttt{cloud\allowbreak{}Encoding\_\allowbreak{}op\allowbreak{}Norm\_\allowbreak{}eq\_\allowbreak{}one}, \texttt{cloud\allowbreak{}Encoding\_\allowbreak{}inverse\_\allowbreak{}op\allowbreak{}Norm\_\allowbreak{}eq\_\allowbreak{}sqrt} (cloud encoding constants).
\end{leanbox}
\section{\texorpdfstring{The signed Jacobian of the discrete splitting}{The signed Jacobian of the discrete splitting}}


For a differentiable force in finite dimension, the derivatives of the kick and drift are triangular shear matrices. Their diagonal blocks are identities, so the determinant is exactly one, including its sign. The chain rule preserves this value for the full B–A–B map.

\subsection{\texorpdfstring{Triangular derivatives and composition}{Triangular derivatives and composition}}


\textbf{Claim.} Let \(Q\) be a finite-dimensional complete real Hilbert space with its Borel, second-countable structure, and let \(F:Q\to Q\) be differentiable everywhere. For every \(h\in\mathbb R\) and \(z\in Q\times Q\), the B–A–B map has \(\det D(K_{h/2}\circ A_h\circ K_{h/2})(z)=1\).



\textbf{Proof.} With velocity-first coordinates,

\[DK_h(p,q)=\begin{pmatrix}I&-hDF(q)\\0&I\end{pmatrix},\qquad DA_h(p,q)=\begin{pmatrix}I&0\\hI&I\end{pmatrix}.\]

These expressions follow by differentiating each coordinate: the kick differentiates \(p-hF(q)\) and leaves \(q\) fixed, while the drift leaves \(p\) fixed and differentiates \(q+hp\). In a product basis the matrices are block triangular. Their determinants are the products of the diagonal-block determinants, hence one, also in zero dimension. The derivative chain rule expresses the derivative of the three-stage map as the product of these three derivatives evaluated at their successive inputs. Multiplicativity of determinant gives \(1\cdot1\cdot1=1\). The force need not be conservative or twice differentiable. \(\square\)

\begin{leanbox}{Lean: \texttt{interwovenBAB\_\allowbreak{}fderiv\_\allowbreak{}det\_\allowbreak{}one}}
\leansource{\texttt{interwovenBAB\_\allowbreak{}fderiv\_\allowbreak{}det\_\allowbreak{}one} \quad PBGSO/\allowbreak{}DiscreteVolume.lean, lines 130–136}
\begin{leancode}
\leanline{0}{theorem\ interwovenBAB\_fderiv\_det\_one}
\leanline{4}{\{Q\ :\ Type*\}\ [NormedAddCommGroup\ Q]\ [InnerProductSpace\ ℝ\ Q]}
\leanline{4}{[CompleteSpace\ Q]\ [MeasurableSpace\ Q]\ [BorelSpace\ Q]\ [SecondCountableTopology\ Q]}
\leanline{4}{[FiniteDimensional\ ℝ\ Q]\ (F\ :\ Q\ →\ Q)\ (hF\ :\ Differentiable\ ℝ\ F)}
\leanline{4}{(h\ :\ ℝ)\ (z\ :\ Q\ ×\ Q)\ :}
\leanline{4}{(fderiv\ ℝ\ (interwovenBAB\ F\ h)\ z).det\ =\ 1}
\end{leancode}
\leanlabel{How the Lean proof follows the argument}
\texttt{interwoven\allowbreak{}BAB\_\allowbreak{}fderiv\_\allowbreak{}det\_\allowbreak{}one} is the public complete-map theorem and invokes \texttt{Discrete\allowbreak{}Volume.\allowbreak{}verlet\allowbreak{}Step\_\allowbreak{}fderiv\_\allowbreak{}det\_\allowbreak{}one}. Inside the namespace, \texttt{upper\allowbreak{}Shear\_\allowbreak{}det} and \texttt{lower\allowbreak{}Shear\_\allowbreak{}det} represent the linear maps in a product basis and use the block-triangular determinant formulas. \texttt{kick\allowbreak{}Step\_\allowbreak{}has\allowbreak{}FDeriv\allowbreak{}At} and \texttt{drift\allowbreak{}Step\_\allowbreak{}has\allowbreak{}FDeriv\allowbreak{}At} compute the actual derivatives. \texttt{fderiv\_\allowbreak{}det\_\allowbreak{}comp} applies the derivative and determinant composition laws; the Verlet proof uses it twice and rewrites all three factors to one.
\end{leanbox}
\section{\texorpdfstring{Transferring frozen stationarity between Brownian sources}{Transferring frozen stationarity between Brownian sources}}


A frozen sampler invariant law can be realised on a different Brownian probability space from the coupled optimiser–sampler solution. The transfer uses actual strong existence on initial-law times Wiener-path space and pathwise uniqueness. An invariant output law is derived after transport; it is not assumed for the solution being studied.

\subsection{\texorpdfstring{The stationary complete-state product law}{The stationary complete-state product law}}


\textbf{Claim.} Fix \(T\ge0\), a probability law \(\nu_\star\), and optimiser state \(s_\star\). Suppose a frozen sampler on some filtered Brownian source has law \(\nu_\star\) at every \(0\le t\le T\). Its observation is integrable with mean zero, and \(\Psi(t,s_\star,0)=0\) throughout the interval. Assume a coupled strong solution exists on independent initial-law times Wiener-path space, with initial state \((s_\star,Z_0)\), \(Z_0\sim\nu_\star\). Assume pathwise uniqueness on the frozen source and on the original solution's source after continuous Brownian modification and null augmentation. If an original solution starts with optimiser \(s_\star\) and sampler law \(\nu_\star\), then its full joint law at each \(t\in[0,T]\) is \(\delta_{s_\star}\otimes\nu_\star\).



\textbf{Proof.} The frozen sampler together with the constant optimiser satisfies the coupled equation, because its mean observation and optimiser velocity vanish. The map sending a sample to its initial state and continuous Brownian path preserves the initial-law times Wiener probability measure. Pull back the canonical path-space solution through this map. The pulled-back solution and the frozen one have the same driver and initial state; pathwise uniqueness gives their equality. Therefore the canonical solution has constant optimiser and sampler law \(\nu_\star\). Repeat this argument for the original solution using its initial/\allowbreak{}path map and uniqueness. Its optimiser is \(s_\star\), its sampler law is \(\nu_\star\), and hence the joint pushforward is \(\delta_{s_\star}\otimes\nu_\star\). The last identity follows by pushing the sampler law through \(z\mapsto(s_\star,z)\). \(\square\)

\begin{leanbox}{Lean: \texttt{stationary\_\allowbreak{}product\_\allowbreak{}law\_\allowbreak{}of\_\allowbreak{}weak\_\allowbreak{}frozen}}
\leansource{\texttt{stationary\_\allowbreak{}product\_\allowbreak{}law\_\allowbreak{}of\_\allowbreak{}weak\_\allowbreak{}frozen} \quad PBGSO/\allowbreak{}RawPathSolutionTransfer.lean, lines 120–140}
\begin{leancode}
\leanline{0}{theorem\ stationary\_product\_law\_of\_weak\_frozen\ [MeasurableSpace\ S]}
\leanline{4}{\{sStar\ :\ S\}\ \{νStar\ :\ Measure\ Z\}\ [IsProbabilityMeasure\ νStar]}
\leanline{4}{(F\ :\ FrozenSamplerData\ Qw\ Dw\ M\ T\ sStar\ νStar)}
\leanline{4}{(hint\ :\ Integrable\ (fun\ z\ =>\ M.observation\ (M.parameter\ sStar)\ (M.position\ z))\ νStar)}
\leanline{4}{(hg\ :\ M.meanGradient\ νStar\ sStar\ =\ 0)}
\leanline{4}{(hfix\ :\ ∀\ t\ ∈\ Icc\ 0\ (T\ :\ ℝ),\ M.optimizer\ t\ sStar\ 0\ =\ 0)}
\leanline{4}{(huniquew\ :\ PathwiseUniqueness\ Qw\ (nullAugmentedDriver\ Qw\ (continuousDriver\ Qw\ Dw))\ M\ T)}
\leanline{4}{(hexists\ :\ ∃\ V\ :\ Data\ (BrownianPathDriver.probability\ (I\ :=\ I)\ νStar)}
\leanline{8}{(BrownianPathDriver.driver\ νStar)\ M\ T,}
\leanline{6}{V.optimizer\ 0\ =\ sStar\ ∧\ V.sampler.state\ 0\ =ᵐ[}
\leanline{8}{BrownianPathDriver.probability\ (I\ :=\ I)\ νStar]\ Prod.fst)}
\leanline{4}{(U\ :\ Data\ Q\ D\ M\ T)}
\leanline{4}{(hunique\ :\ PathwiseUniqueness\ Q\ (nullAugmentedDriver\ Q\ (continuousDriver\ Q\ D))\ M\ T)}
\leanline{4}{(hs\ :\ U.optimizer\ 0\ =\ sStar)\ (hν\ :\ Q.map\ (U.sampler.state\ 0)\ =\ νStar)}
\leanline{4}{(t\ :\ ℝ)\ (ht\ :\ t\ ∈\ Icc\ 0\ (T\ :\ ℝ))\ :}
\leanline{4}{Q.map\ (fun\ ω\ =>\ (U.optimizer\ t,\ U.sampler.state\ t\ ω))\ =\ (Measure.dirac\ sStar).prod\ νStar}
\end{leancode}
\leanlabel{How the Lean proof follows the argument}
\texttt{stationary\_\allowbreak{}product\_\allowbreak{}law\_\allowbreak{}of\_\allowbreak{}weak\_\allowbreak{}frozen} first obtains optimiser and law identities from \texttt{stationary\_\allowbreak{}of\_\allowbreak{}weak\_\allowbreak{}frozen}. That result constructs \texttt{F.\allowbreak{}constant\allowbreak{}Data}, transports it and the original solution to the same path-space solution using \texttt{law\_\allowbreak{}eq\_\allowbreak{}pathspace\_\allowbreak{}solution}, and combines the resulting equalities. The final proof rewrites the optimiser, uses \texttt{Measure.\allowbreak{}dirac\_\allowbreak{}prod}, then \texttt{Measure.\allowbreak{}map\_\allowbreak{}map}. The quantifier \texttt{ht} is the closed finite horizon, including zero and the terminal time.
\end{leanbox}
\section{\texorpdfstring{Lebesgue preservation by the continuous frictionless flow}{Lebesgue preservation by the continuous frictionless flow}}


Determinant one gives a measure-preservation statement once each time slice is globally invertible. Autonomous ODE uniqueness provides the inverse \(\Phi_{-t}\). The finite-dimensional change-of-variables theorem then proves preservation of Lebesgue measure on the complete phase space.

\subsection{\texorpdfstring{Uniqueness, inverse time and change of variables}{Uniqueness, inverse time and change of variables}}


\textbf{Claim.} For a \(C^1\) force and actual global frictionless phase trajectories as in the preceding result, every map \(x\mapsto\Phi(t,x)\) preserves Lebesgue measure, for each real \(t\).



\textbf{Proof.} For fixed \(s,x\), the two curves \(u\mapsto\Phi(s+u,x)\) and \(u\mapsto\Phi(u,\Phi(s,x))\) solve the same autonomous ODE and agree at zero. Uniqueness on a compact interval containing zero and the desired time gives \(\Phi(s+t,x)=\Phi(t,\Phi(s,x))\). Taking \(s=-t\) shows that \(\Phi(-t,\cdot)\) is a two-sided inverse. Therefore each time slice is bijective. It is \(C^1\) by the preceding smooth-dependence result, and its determinant equals one everywhere.



The change-of-variables identity for an injective differentiable map says that the pushforward of the measure \(|\det D\Phi_t(x)|dx\) is Lebesgue measure restricted to its image. Here the density is identically one and the image is the whole space by surjectivity. Thus \((\Phi_t)_\#dx=dx\). Measurability follows from continuity, completing the measure-preservation assertion. No common potential for the force is used. \(\square\)

\begin{leanbox}{Lean: \texttt{flow\_\allowbreak{}add}, \texttt{flow\_\allowbreak{}bijective}, \texttt{phase\_\allowbreak{}flow\_\allowbreak{}measurePreserving\_\allowbreak{}source}}
\leansource{\texttt{flow\_\allowbreak{}add} \quad PBGSO/\allowbreak{}ContinuousFlowMeasure.lean, lines 22–44}
\begin{leancode}
\leanline{0}{theorem\ flow\_add\ (Φ\ :\ ℝ\ ×\ E\ →\ E)\ (f\ :\ E\ →\ E)\ (hf\ :\ ContDiff\ ℝ\ 1\ f)}
\leanline{4}{(hODE\ :\ ∀\ t\ x,\ HasDerivAt\ (fun\ s\ =>\ Φ\ (s,\ x))\ (f\ (Φ\ (t,\ x)))\ t)}
\leanline{4}{(hinit\ :\ ∀\ x,\ Φ\ (0,\ x)\ =\ x)\ (t\ s\ :\ ℝ)\ (x\ :\ E)\ :}
\leanline{4}{Φ\ (t\ +\ s,\ x)\ =\ Φ\ (s,\ Φ\ (t,\ x))}
\end{leancode}
\leansource{\texttt{flow\_\allowbreak{}bijective} \quad PBGSO/\allowbreak{}ContinuousFlowMeasure.lean, lines 48–59}
\begin{leancode}
\leanline{0}{theorem\ flow\_bijective\ (Φ\ :\ ℝ\ ×\ E\ →\ E)\ (f\ :\ E\ →\ E)\ (hf\ :\ ContDiff\ ℝ\ 1\ f)}
\leanline{4}{(hODE\ :\ ∀\ t\ x,\ HasDerivAt\ (fun\ s\ =>\ Φ\ (s,\ x))\ (f\ (Φ\ (t,\ x)))\ t)}
\leanline{4}{(hinit\ :\ ∀\ x,\ Φ\ (0,\ x)\ =\ x)\ (t\ :\ ℝ)\ :}
\leanline{4}{Function.Bijective\ (fun\ x\ =>\ Φ\ (t,\ x))}
\end{leancode}
\leansource{\texttt{phase\_\allowbreak{}flow\_\allowbreak{}measurePreserving\_\allowbreak{}source} \quad PBGSO/\allowbreak{}ContinuousFlowMeasure.lean, lines 90–100}
\begin{leancode}
\leanline{0}{theorem\ phase\_flow\_measurePreserving\_source}
\leanline{4}{(F\ :\ (ι\ →\ ℝ)\ →\ (ι\ →\ ℝ))\ (hF\ :\ ContDiff\ ℝ\ 1\ F)}
\leanline{4}{(Φ\ :\ ℝ\ ×\ ContinuousVolume.Phase\ ι\ →\ ContinuousVolume.Phase\ ι)}
\leanline{4}{(hODE\ :\ ∀\ t\ x,\ HasDerivAt\ (fun\ s\ =>\ Φ\ (s,\ x))}
\leanline{6}{(ContinuousVolume.phaseField\ F\ (Φ\ (t,\ x)))\ t)}
\leanline{4}{(hinit\ :\ ∀\ x,\ Φ\ (0,\ x)\ =\ x)\ (t\ :\ ℝ)\ :}
\leanline{4}{MeasurePreserving\ (fun\ x\ =>\ Φ\ (t,\ x))\ volume\ volume}
\end{leancode}
\leanlabel{How the Lean proof follows the argument}
\texttt{flow\_\allowbreak{}add} applies \texttt{trajectories\_\allowbreak{}unique\_\allowbreak{}Icc} to the shifted and restarted curves. \texttt{flow\_\allowbreak{}bijective} obtains both injectivity and surjectivity from negative-time composition. \texttt{phase\_\allowbreak{}flow\_\allowbreak{}measure\allowbreak{}Preserving\_\allowbreak{}source} reduces to \texttt{divergence\allowbreak{}Free\_\allowbreak{}flow\_\allowbreak{}measure\allowbreak{}Preserving}, which derives joint regularity, bijectivity and determinant one before invoking \texttt{map\_\allowbreak{}with\allowbreak{}Density\_\allowbreak{}abs\_\allowbreak{}det\_\allowbreak{}fderiv\_\allowbreak{}eq\_\allowbreak{}add\allowbreak{}Haar}. Its simplification uses determinant one and surjectivity to turn the weighted image measure into unrestricted volume.
\end{leanbox}
\section{\texorpdfstring{Absorption of affine local errors by an exactly contracting flow}{Absorption of affine local errors by an exactly contracting flow}}


This section contains the discrete comparison argument behind Theorems New A and New B. Two sequences in a normed space \(H\) are compared: the numerical chain \(X_{n+1}=A_n(X_n)\) and the exact chain \(Y_{n+1}=\Phi_n(Y_n)\), where \(\Phi_n\) is the exact flow of the continuous dynamics over one step and \(A_n\) is one step of the numerical method. Distances are measured through a bounded linear map \(C:H\to K\) into another normed space. Three hypotheses are made:



\begin{itemize}


\item the exact maps contract: \(\|C(\Phi_nx-\Phi_ny)\|\le e^{-\rho h}\|C(x-y)\|\);

\item the local error is affine in the state: \(\|C(A_nx-\Phi_nx)\|\le h\sqrt h\,(a+b\|Cx\|)\);

\item the exact chain is bounded: \(\|CY_n\|\le R\) for all \(n\).

\end{itemize}



The numerical maps are not assumed Lipschitz or contracting, and no bound on the numerical chain is assumed. The local error at the numerical state is controlled through \(\|CX_n\|\le\|C(X_n-Y_n)\|+R\), and the resulting growth term \(h\sqrt h\,b\) is absorbed by the contraction when \(b\sqrt h\le\rho/4\) and \(\rho h\le1\), since then \(e^{-\rho h}+b\,h\sqrt h\le1-\rho h/4\). The conclusion, for chains started at the same point, is


\[\|X_n-Y_n\|\le\frac{4(a+bR)}{\rho\,m}\sqrt h\qquad\text{for every }n\in\mathbb N,\]


when \(m\|z\|\le\|Cz\|\). The bound is uniform in the number of iterations; the factor \(\sqrt h\) comes from dividing the local error \(h\sqrt h\) by the contraction gain \(\rho h/4\) per step.

\subsection{\texorpdfstring{Affine recurrences with a contraction factor}{Affine recurrences with a contraction factor}}


\textbf{Claim.} Let \((u_n)_{n\ge0}\) be real numbers and \(r,\beta\in\mathbb R\) with \(0\le r<1\). If \(u_{n+1}\le ru_n+\beta\) for all \(n\), then for all \(n\)


\[u_n\le r^nu_0+\frac{\beta}{1-r}\bigl(1-r^n\bigr).\]


\textbf{Proof.} Since \(r<1\), \(1-r\ne0\), and


\[r\cdot\frac{\beta}{1-r}+\beta=\frac{r\beta+(1-r)\beta}{1-r}=\frac{\beta}{1-r}.\qquad(1)\]


We argue by induction on \(n\). For \(n=0\) the right-hand side is \(u_0+\frac\beta{1-r}\cdot0=u_0\). Assume the bound for \(n\). Since \(r\ge0\), multiplying the induction hypothesis by \(r\) preserves the inequality, so


\[u_{n+1}\le ru_n+\beta\le r^{n+1}u_0+\frac{\beta}{1-r}\bigl(r-r^{n+1}\bigr)+\beta=r^{n+1}u_0+\frac{\beta}{1-r}\bigl(1-r^{n+1}\bigr),\]


where the last equality is (1). \(\square\)




No sign condition on \(\beta\) is needed.

\begin{leanbox}{Lean: \texttt{affine\_\allowbreak{}recurrence\_\allowbreak{}bound}}
\leansource{\texttt{affine\_\allowbreak{}recurrence\_\allowbreak{}bound} \quad PBGSO/\allowbreak{}UniformDiscrete.lean, lines 14–25}
\begin{leancode}
\leanline{0}{private\ theorem\ affine\_recurrence\_bound\ (u\ :\ ℕ\ →\ ℝ)\ (r\ β\ :\ ℝ)\ (hr\ :\ 0\ ≤\ r)}
\leanline{4}{(hr1\ :\ r\ <\ 1)\ (hstep\ :\ ∀\ n,\ u\ (n+1)\ ≤\ r*u\ n+β)\ (n\ :\ ℕ)\ :}
\leanline{4}{u\ n\ ≤\ r\textasciicircum{}n*u\ 0+β/(1-r)*(1-r\textasciicircum{}n)}
\end{leancode}
\leanlabel{How the Lean proof follows the argument}
\texttt{affine\_\allowbreak{}recurrence\_\allowbreak{}bound} is a private lemma of this file. \texttt{hden} is \(1-r\ne0\) and \texttt{h\allowbreak{}B} is identity (1), proved by \texttt{field\_\allowbreak{}simp} and \texttt{ring}. The induction uses \texttt{gcongr} to multiply the induction hypothesis by \(r\ge0\) and \texttt{nlinarith [h\allowbreak{}B]} with \texttt{pow\_\allowbreak{}succ} for the final equality.
\end{leanbox}
\subsection{\texorpdfstring{One-step recursion for the error of the numerical chain}{One-step recursion for the error of the numerical chain}}


\textbf{Claim.} Let \(H,K\) be real normed spaces, \(C:H\to K\) bounded linear, \(A_n,\Phi_n:H\to H\) maps and \(X_n,Y_n\in H\) with \(X_{n+1}=A_n(X_n)\), \(Y_{n+1}=\Phi_n(Y_n)\) for all \(n\). Let \(q,\varepsilon,a,b,R\in\mathbb R\) with \(\varepsilon\ge0\), \(b\ge0\), and assume for all \(n\) and \(x,y\in H\)


\[\begin{gathered}\|C(\Phi_nx-\Phi_ny)\|\le q\|C(x-y)\|,\qquad\|C(A_nx-\Phi_nx)\|\le\varepsilon\bigl(a+b\|Cx\|\bigr),\\ \|CY_n\|\le R .\end{gathered}\]


Then \(e_n=\|C(X_n-Y_n)\|\) satisfies \(e_{n+1}\le(q+\varepsilon b)\,e_n+\varepsilon(a+bR)\) for every \(n\).




\textbf{Proof.} First, \(CX_n=C(X_n-Y_n)+CY_n\) by linearity, so the triangle inequality and the reference bound give


\[\|CX_n\|\le e_n+R.\qquad(2)\]


Next, \(X_{n+1}-Y_{n+1}=\bigl(A_n(X_n)-\Phi_n(X_n)\bigr)+\bigl(\Phi_n(X_n)-\Phi_n(Y_n)\bigr)\). Applying \(C\), the triangle inequality, the local error bound at the numerical point \(x=X_n\) and the contraction gives


\[e_{n+1}\le\varepsilon\bigl(a+b\|CX_n\|\bigr)+q\,e_n\le\varepsilon\bigl(a+b(e_n+R)\bigr)+q\,e_n=(q+\varepsilon b)e_n+\varepsilon(a+bR),\]


where the second inequality uses (2) multiplied by \(\varepsilon b\ge0\). \(\square\)




The numerical maps \(A_n\) enter only through the local error at the actual numerical state; no Lipschitz property of \(A_n\) is used. The price is the factor \(\varepsilon b\) added to the contraction factor.

\begin{leanbox}{Lean: \texttt{chain\_\allowbreak{}error\_\allowbreak{}step}}
\leansource{\texttt{chain\_\allowbreak{}error\_\allowbreak{}step} \quad PBGSO/\allowbreak{}UniformDiscrete.lean, lines 33–55}
\begin{leancode}
\leanline{0}{theorem\ chain\_error\_step\ (C\ :\ H\ →L[ℝ]\ K)\ (A\ Φ\ :\ ℕ\ →\ H\ →\ H)\ (X\ Y\ :\ ℕ\ →\ H)}
\leanline{4}{(q\ ε\ a\ b\ R\ :\ ℝ)\ (hε\ :\ 0\ ≤\ ε)\ (hb\ :\ 0\ ≤\ b)}
\leanline{4}{(hX\ :\ ∀\ n,\ X\ (n+1)=A\ n\ (X\ n))\ (hY\ :\ ∀\ n,\ Y\ (n+1)=Φ\ n\ (Y\ n))}
\leanline{4}{(hcontract\ :\ ∀\ n\ x\ y,\ ‖C\ (Φ\ n\ x-Φ\ n\ y)‖\ ≤\ q*‖C\ (x-y)‖)}
\leanline{4}{(hdefect\ :\ ∀\ n\ x,\ ‖C\ (A\ n\ x-Φ\ n\ x)‖\ ≤\ ε*(a+b*‖C\ x‖))}
\leanline{4}{(hreference\ :\ ∀\ n,\ ‖C\ (Y\ n)‖\ ≤\ R)\ (n\ :\ ℕ)\ :}
\leanline{4}{‖C\ (X\ (n+1)-Y\ (n+1))‖\ ≤}
\leanline{6}{(q+ε*b)*‖C\ (X\ n-Y\ n)‖+ε*(a+b*R)}
\end{leancode}
\leanlabel{How the Lean proof follows the argument}
The maps are indexed by the step number (\texttt{A Φ : \ensuremath{\mathbb N} \ensuremath{\to} H \ensuremath{\to} H}), so that each step may use its own noise. \texttt{hradius} is (2); \texttt{he} is the decomposition of \(X_{n+1}-Y_{n+1}\), obtained from \texttt{h\allowbreak{}X}, \texttt{h\allowbreak{}Y} and linearity (\texttt{simp}). The \texttt{calc} block applies \texttt{norm\_\allowbreak{}add\_\allowbreak{}le}, then \texttt{hdefect n (X n)} and \texttt{hcontract n (X n) (Y n)}, then \texttt{gcongr} with \texttt{hradius}, and closes the identity with \texttt{ring}.
\end{leanbox}
\subsection{\texorpdfstring{An elementary exponential inequality}{An elementary exponential inequality}}


\textbf{Claim.} For \(0\le x\le1\), \(e^{-x}\le1-x/2\).




\textbf{Proof.} From \(e^x\ge1+x\) and \(e^xe^{-x}=1\) we get \(1\ge(1+x)e^{-x}\), so \(e^{-x}\le1/(1+x)\) because \(1+x>0\). Moreover \(1/(1+x)\le1-x/2\) is equivalent to \(1\le(1+x)(1-x/2)=1+x/2-x^2/2\), that is, to \(x(1-x)\ge0\), which holds for \(0\le x\le1\). \(\square\)

\begin{leanbox}{Lean: \texttt{exp\_\allowbreak{}neg\_\allowbreak{}le\_\allowbreak{}one\_\allowbreak{}sub\_\allowbreak{}half}}
\leansource{\texttt{exp\_\allowbreak{}neg\_\allowbreak{}le\_\allowbreak{}one\_\allowbreak{}sub\_\allowbreak{}half} \quad PBGSO/\allowbreak{}UniformDiscrete.lean, lines 98–107}
\begin{leancode}
\leanline{0}{theorem\ exp\_neg\_le\_one\_sub\_half\ (x\ :\ ℝ)\ (hx\ :\ 0\ ≤\ x)\ (hx1\ :\ x\ ≤\ 1)\ :}
\leanline{4}{Real.exp\ (-x)\ ≤\ 1-x/2}
\end{leancode}
\leanlabel{How the Lean proof follows the argument}
\texttt{he} multiplies \texttt{Real.\allowbreak{}add\_\allowbreak{}one\_\allowbreak{}le\_\allowbreak{}exp x} by \(e^{-x}>0\), and \texttt{heq} is \(e^xe^{-x}=1\) (\texttt{Real.\allowbreak{}exp\_\allowbreak{}add}). \texttt{hi} is \(e^{-x}\le1/(1+x)\), obtained with \texttt{le\_\allowbreak{}div\_\allowbreak{}iff\ensuremath{{}_{0}}}; the last comparison is \texttt{div\_\allowbreak{}le\_\allowbreak{}iff\ensuremath{{}_{0}}} followed by \texttt{nlinarith} with the product \(x(1-x)\ge0\).
\end{leanbox}
\subsection{\texorpdfstring{Absorption of the growth term for small steps}{Absorption of the growth term for small steps}}


\textbf{Claim.} Let \(\rho>0\), \(h>0\) and \(b\in\mathbb R\) with \(\rho h\le1\) and \(b\sqrt h\le\rho/4\). Then


\[e^{-\rho h}+b\,h\sqrt h\le1-\frac{\rho h}{4}.\]


\textbf{Proof.} The previous step with \(x=\rho h\in[0,1]\) gives \(e^{-\rho h}\le1-\rho h/2\). Multiplying \(b\sqrt h\le\rho/4\) by \(h>0\) gives \(b\,h\sqrt h\le\rho h/4\). Adding, \(e^{-\rho h}+b\,h\sqrt h\le1-\rho h/2+\rho h/4=1-\rho h/4\). \(\square\)




Half of the contraction gain \(\rho h/2\) per step is spent on the growth of the local error with the numerical state.

\begin{leanbox}{Lean: \texttt{one\_\allowbreak{}step\_\allowbreak{}absorption}}
\leansource{\texttt{one\_\allowbreak{}step\_\allowbreak{}absorption} \quad PBGSO/\allowbreak{}UniformDiscrete.lean, lines 111–117}
\begin{leancode}
\leanline{0}{theorem\ one\_step\_absorption\ (ρ\ h\ b\ :\ ℝ)\ (hρ\ :\ 0\ <\ ρ)\ (hh\ :\ 0\ <\ h)}
\leanline{4}{(hρh\ :\ ρ*h≤1)\ (hsmall\ :\ b*Real.sqrt\ h≤ρ/4)\ :}
\leanline{4}{Real.exp\ (-ρ*h)+b*(h*Real.sqrt\ h)\ ≤\ 1-ρ*h/4}
\end{leancode}
\leanlabel{How the Lean proof follows the argument}
\texttt{he} is \texttt{exp\_\allowbreak{}neg\_\allowbreak{}le\_\allowbreak{}one\_\allowbreak{}sub\_\allowbreak{}half} at \texttt{\ensuremath{\rho}*h} (rewritten with \texttt{neg\_\allowbreak{}mul}), \texttt{hs} multiplies \texttt{hsmall} by \(h\ge0\), and \texttt{nlinarith} adds them.
\end{leanbox}
\subsection{\texorpdfstring{A square-root error bound for every iteration}{A square-root error bound for every iteration}}


\textbf{Claim.} In the setting of the recursion step, let \(\rho>0\), \(h>0\), \(b\ge0\), \(\rho h\le1\), \(b\sqrt h\le\rho/4\), and assume for all \(n\) and \(x,y\in H\)


\[\begin{gathered}\|C(\Phi_nx-\Phi_ny)\|\le e^{-\rho h}\|C(x-y)\|,\qquad\|C(A_nx-\Phi_nx)\|\le h\sqrt h\,\bigl(a+b\|Cx\|\bigr),\\ \|CY_n\|\le R .\end{gathered}\]


Then for every \(n\in\mathbb N\), with \(p=1-\rho h/4\),


\[\|C(X_n-Y_n)\|\le p^n\|C(X_0-Y_0)\|+\frac{4(a+bR)}{\rho}\sqrt h\,\bigl(1-p^n\bigr).\]


\textbf{Proof.} Since \(0<\rho h\le1\), we have \(0\le p<1\). The recursion step with \(q=e^{-\rho h}\) and \(\varepsilon=h\sqrt h\ge0\) gives \(e_{k+1}\le(e^{-\rho h}+h\sqrt h\,b)e_k+h\sqrt h(a+bR)\). By the absorption step and \(e_k\ge0\), \((e^{-\rho h}+h\sqrt h\,b)e_k\le p\,e_k\), so


\[e_{k+1}\le p\,e_k+h\sqrt h\,(a+bR)\qquad(k\in\mathbb N).\]


The affine recurrence bound with \(r=p\) and \(\beta=h\sqrt h(a+bR)\) gives \(e_n\le p^ne_0+\frac{\beta}{1-p}(1-p^n)\), and


\[\frac{\beta}{1-p}=\frac{h\sqrt h\,(a+bR)}{\rho h/4}=\frac{4(a+bR)}{\rho}\sqrt h. \qquad\square\]\begin{leanbox}{Lean: \texttt{one\_\allowbreak{}step\_\allowbreak{}uniform\_\allowbreak{}error}}
\leansource{\texttt{one\_\allowbreak{}step\_\allowbreak{}uniform\_\allowbreak{}error} \quad PBGSO/\allowbreak{}UniformDiscrete.lean, lines 125–150}
\begin{leancode}
\leanline{0}{theorem\ one\_step\_uniform\_error\ (C\ :\ H\ →L[ℝ]\ K)\ (A\ Φ\ :\ ℕ\ →\ H\ →\ H)\ (X\ Y\ :\ ℕ\ →\ H)}
\leanline{4}{(ρ\ h\ a\ b\ R\ :\ ℝ)\ (hρ\ :\ 0\ <\ ρ)\ (hh\ :\ 0\ <\ h)\ (hb\ :\ 0\ ≤\ b)}
\leanline{4}{(hρh\ :\ ρ*h≤1)\ (hsmall\ :\ b*Real.sqrt\ h≤ρ/4)}
\leanline{4}{(hX\ :\ ∀\ n,\ X\ (n+1)=A\ n\ (X\ n))\ (hY\ :\ ∀\ n,\ Y\ (n+1)=Φ\ n\ (Y\ n))}
\leanline{4}{(hcontract\ :\ ∀\ n\ x\ y,\ ‖C\ (Φ\ n\ x-Φ\ n\ y)‖\ ≤\ Real.exp\ (-ρ*h)*‖C\ (x-y)‖)}
\leanline{4}{(hdefect\ :\ ∀\ n\ x,\ ‖C\ (A\ n\ x-Φ\ n\ x)‖\ ≤\ (h*Real.sqrt\ h)*(a+b*‖C\ x‖))}
\leanline{4}{(hreference\ :\ ∀\ n,\ ‖C\ (Y\ n)‖\ ≤\ R)\ (n\ :\ ℕ)\ :}
\leanline{4}{‖C\ (X\ n-Y\ n)‖\ ≤\ (1-ρ*h/4)\textasciicircum{}n*‖C\ (X\ 0-Y\ 0)‖+}
\leanline{6}{(4*(a+b*R)/ρ)*Real.sqrt\ h*(1-(1-ρ*h/4)\textasciicircum{}n)}
\end{leancode}
\leanlabel{How the Lean proof follows the argument}
\texttt{p} is \(1-\rho h/4\), with \texttt{hp} and \texttt{hp1} the bounds \(0\le p<1\). \texttt{hrec} is the displayed recursion: \texttt{chain\_\allowbreak{}error\_\allowbreak{}step} with \(q=e^{-\rho h}\), \(\varepsilon=h\sqrt h\), combined with \texttt{one\_\allowbreak{}step\_\allowbreak{}absorption} multiplied by \(e_k\ge0\). \texttt{he} is \texttt{affine\_\allowbreak{}recurrence\_\allowbreak{}bound}, and \texttt{hc} is the evaluation of \(\beta/(1-p)\), proved by \texttt{field\_\allowbreak{}simp} and \texttt{ring}.
\end{leanbox}
\subsection{\texorpdfstring{Synchronous initialisation}{Synchronous initialisation}}


\textbf{Claim.} Under the hypotheses of the previous step, assume in addition \(a\ge0\), \(R\ge0\) and \(X_0=Y_0\). Then for every \(n\in\mathbb N\)


\[\|C(X_n-Y_n)\|\le\frac{4(a+bR)}{\rho}\sqrt h .\]


\textbf{Proof.} With \(X_0=Y_0\) the first term of the previous bound vanishes. Since \(0\le p\le1\), \(0\le p^n\le1\) and \(1-p^n\le1\). The coefficient \(\frac{4(a+bR)}{\rho}\sqrt h\) is nonnegative because \(a,b,R\ge0\) and \(\rho>0\); multiplying \(1-p^n\le1\) by it gives the claim. \(\square\)

\begin{leanbox}{Lean: \texttt{one\_\allowbreak{}step\_\allowbreak{}same\_\allowbreak{}initial\_\allowbreak{}uniform}}
\leansource{\texttt{one\_\allowbreak{}step\_\allowbreak{}same\_\allowbreak{}initial\_\allowbreak{}uniform} \quad PBGSO/\allowbreak{}UniformDiscrete.lean, lines 153–166}
\begin{leancode}
\leanline{0}{theorem\ one\_step\_same\_initial\_uniform\ (C\ :\ H\ →L[ℝ]\ K)\ (A\ Φ\ :\ ℕ\ →\ H\ →\ H)\ (X\ Y\ :\ ℕ\ →\ H)}
\leanline{4}{(ρ\ h\ a\ b\ R\ :\ ℝ)\ (hρ\ :\ 0\ <\ ρ)\ (hh\ :\ 0\ <\ h)\ (ha\ :\ 0\ ≤\ a)\ (hb\ :\ 0\ ≤\ b)}
\leanline{4}{(hR\ :\ 0\ ≤\ R)\ (hρh\ :\ ρ*h≤1)\ (hsmall\ :\ b*Real.sqrt\ h≤ρ/4)}
\leanline{4}{(hX\ :\ ∀\ n,\ X\ (n+1)=A\ n\ (X\ n))\ (hY\ :\ ∀\ n,\ Y\ (n+1)=Φ\ n\ (Y\ n))}
\leanline{4}{(hcontract\ :\ ∀\ n\ x\ y,\ ‖C\ (Φ\ n\ x-Φ\ n\ y)‖\ ≤\ Real.exp\ (-ρ*h)*‖C\ (x-y)‖)}
\leanline{4}{(hdefect\ :\ ∀\ n\ x,\ ‖C\ (A\ n\ x-Φ\ n\ x)‖\ ≤\ (h*Real.sqrt\ h)*(a+b*‖C\ x‖))}
\leanline{4}{(hreference\ :\ ∀\ n,\ ‖C\ (Y\ n)‖\ ≤\ R)\ (hinit\ :\ X\ 0=Y\ 0)\ (n\ :\ ℕ)\ :}
\leanline{4}{‖C\ (X\ n-Y\ n)‖\ ≤\ (4*(a+b*R)/ρ)*Real.sqrt\ h}
\end{leancode}
\leanlabel{How the Lean proof follows the argument}
After \texttt{one\_\allowbreak{}step\_\allowbreak{}uniform\_\allowbreak{}error}, the rewrite \texttt{rw [hinit,sub\_\allowbreak{}self,map\_\allowbreak{}zero,norm\_\allowbreak{}zero,mul\_\allowbreak{}zero,zero\_\allowbreak{}add]} removes the initial term; \texttt{hp} is \(p^n\ge0\), \texttt{h\allowbreak{}K} is the nonnegativity of the coefficient, and \texttt{mul\_\allowbreak{}le\_\allowbreak{}of\_\allowbreak{}le\_\allowbreak{}one\_\allowbreak{}right} finishes.
\end{leanbox}
\subsection{\texorpdfstring{The error in the original norm}{The error in the original norm}}


\textbf{Claim.} Let \(H,K\), \(C\), \(A_n\), \(\Phi_n\), \(X_n\), \(Y_n\) be as in the recursion step, and let \(\rho,h,a,b,R,m\) be real numbers with \(\rho>0\), \(h>0\), \(a,b,R\ge0\), \(m>0\), \(\rho h\le1\) and \(b\sqrt h\le\rho/4\). Assume



\begin{enumerate}


\item \(m\|z\|\le\|Cz\|\) for all \(z\in H\) (coercivity of the metric);

\item \(\|C(\Phi_nx-\Phi_ny)\|\le e^{-\rho h}\|C(x-y)\|\) for all \(n,x,y\) (exact contraction);

\item \(\|C(A_nx-\Phi_nx)\|\le h\sqrt h\,(a+b\|Cx\|)\) for all \(n,x\) (affine local error);

\item \(\|CY_n\|\le R\) for all \(n\) (reference radius);

\item \(X_{n+1}=A_n(X_n)\), \(Y_{n+1}=\Phi_n(Y_n)\) and \(X_0=Y_0\).

\end{enumerate}



Then for every \(n\in\mathbb N\)


\[\|X_n-Y_n\|\le\frac{1}{m}\cdot\frac{4(a+bR)}{\rho}\sqrt h .\]


\textbf{Proof.} By coercivity and the previous step, \(m\|X_n-Y_n\|\le\|C(X_n-Y_n)\|\le\frac{4(a+bR)}{\rho}\sqrt h\). Divide by \(m>0\). \(\square\)




This is the form used for both methods: \(H\) is the space of random particle states with the maximum of the four \(L^2\) norms, \(C\) is the contracting metric constructed from A3, \(A_n\) is the numerical step with the \(n\)-th innovations and \(\Phi_n\) is the exact flow over \([nh,(n+1)h]\) driven by the same Brownian path.

\begin{leanbox}{Lean: \texttt{one\_\allowbreak{}step\_\allowbreak{}same\_\allowbreak{}initial\_\allowbreak{}uniform\_\allowbreak{}norm}}
\leansource{\texttt{one\_\allowbreak{}step\_\allowbreak{}same\_\allowbreak{}initial\_\allowbreak{}uniform\_\allowbreak{}norm} \quad PBGSO/\allowbreak{}UniformDiscrete.lean, lines 170–182}
\begin{leancode}
\leanline{0}{theorem\ one\_step\_same\_initial\_uniform\_norm\ (C\ :\ H\ →L[ℝ]\ K)\ (A\ Φ\ :\ ℕ\ →\ H\ →\ H)\ (X\ Y\ :\ ℕ\ →\ H)}
\leanline{4}{(ρ\ h\ a\ b\ R\ m\ :\ ℝ)\ (hρ\ :\ 0\ <\ ρ)\ (hh\ :\ 0\ <\ h)\ (ha\ :\ 0\ ≤\ a)\ (hb\ :\ 0\ ≤\ b)}
\leanline{4}{(hR\ :\ 0\ ≤\ R)\ (hm\ :\ 0\ <\ m)\ (hcoerc\ :\ ∀\ z,\ m*‖z‖≤‖C\ z‖)}
\leanline{4}{(hρh\ :\ ρ*h≤1)\ (hsmall\ :\ b*Real.sqrt\ h≤ρ/4)}
\leanline{4}{(hX\ :\ ∀\ n,\ X\ (n+1)=A\ n\ (X\ n))\ (hY\ :\ ∀\ n,\ Y\ (n+1)=Φ\ n\ (Y\ n))}
\leanline{4}{(hcontract\ :\ ∀\ n\ x\ y,\ ‖C\ (Φ\ n\ x-Φ\ n\ y)‖\ ≤\ Real.exp\ (-ρ*h)*‖C\ (x-y)‖)}
\leanline{4}{(hdefect\ :\ ∀\ n\ x,\ ‖C\ (A\ n\ x-Φ\ n\ x)‖\ ≤\ (h*Real.sqrt\ h)*(a+b*‖C\ x‖))}
\leanline{4}{(hreference\ :\ ∀\ n,\ ‖C\ (Y\ n)‖\ ≤\ R)\ (hinit\ :\ X\ 0=Y\ 0)\ (n\ :\ ℕ)\ :}
\leanline{4}{‖X\ n-Y\ n‖\ ≤\ ((4*(a+b*R)/ρ)*Real.sqrt\ h)/m}
\end{leancode}
\leanlabel{How the Lean proof follows the argument}
The conclusion \texttt{\ensuremath{\Vert}X n-\allowbreak{}Y n\ensuremath{\Vert} \ensuremath{\le} ((4*(a+b*R)/\allowbreak{}\ensuremath{\rho})*Real.\allowbreak{}sqrt h)/\allowbreak{}m} is proved with \texttt{le\_\allowbreak{}div\_\allowbreak{}iff\ensuremath{{}_{0}} hm}, which turns it into \(m\|X_n-Y_n\|\le\cdots\), followed by \texttt{nlinarith} with \texttt{one\_\allowbreak{}step\_\allowbreak{}same\_\allowbreak{}initial\_\allowbreak{}uniform} and \texttt{hcoerc (X n-\allowbreak{}Y n)}.
\end{leanbox}
\section{\texorpdfstring{Finite second moments need not give finite signed work}{Finite second moments need not give finite signed work}}


The scalar drift \(b(x)=-x-x^3\) is strongly dissipative, yet its signed work against the point zero contains a fourth power. This distinguishes a conditional work inequality from a universal finite-integrability assumption on \(\mathcal P_2\).

\subsection{\texorpdfstring{A finite-second-moment obstruction}{A finite-second-moment obstruction}}


\textbf{Claim.} There is a probability measure \(\mu\) on \(\mathbb R\) with \(\int x^2d\mu<\infty\) but \(\int x^4d\mu=\infty\). For \(b(x)=-x-x^3\), one has \((x-y)(b(x)-b(y))\le-(x-y)^2\) for every pair of reals, but \(x(b(x)-b(0))\) is not integrable for such a law.



\textbf{Proof.} If every probability law integrable against \(1+x^2\) also integrated \(x^4\), the universal-integrability growth theorem would imply a finite \(C\ge0\) with \(x^4\le C(1+x^2)\) for every real \(x\). That theorem is proved by a countable atomic measure concentrating successively on points where the ratio is unbounded. At \(x=\sqrt{C+1}\), the alleged inequality becomes \((C+1)^2\le C(C+2)\), or \(1\le0\). Hence such a law exists. Integrability of \(1+x^2\) implies integrability of \(x^2\).



Since \(x^3-y^3=(x-y)(x^2+xy+y^2)\) and \(x^2+xy+y^2\ge0\),

\[(x-y)(b(x)-b(y))=-(x-y)^2\bigl(1+x^2+xy+y^2\bigr)\le-(x-y)^2.\]

But \(x(b(x)-b(0))=-x^2-x^4\). If this were integrable, subtracting the integrable quadratic term from its negative would make \(x^4\) integrable, a contradiction. The example concerns the finite-work domain; it does not disprove a contraction estimate on laws for which the required work argument is justified. \(\square\)

\begin{leanbox}{Lean: \texttt{exists\_\allowbreak{}quadratic\_\allowbreak{}not\_\allowbreak{}quartic}, \texttt{drift}, \texttt{drift\_\allowbreak{}strongly\_\allowbreak{}dissipative} and 1 more}
\leansource{\texttt{exists\_\allowbreak{}quadratic\_\allowbreak{}not\_\allowbreak{}quartic} \quad PBGSO/\allowbreak{}FiniteWorkScopeAudit.lean, lines 18–34}
\begin{leancode}
\leanline{0}{theorem\ exists\_quadratic\_not\_quartic\ :}
\leanline{4}{∃\ μ\ :\ Measure\ ℝ,\ IsProbabilityMeasure\ μ\ ∧}
\leanline{6}{Integrable\ (fun\ x\ :\ ℝ\ =>\ 1\ +\ x\ \textasciicircum{}\ 2)\ μ\ ∧\ ¬\ Integrable\ (fun\ x\ :\ ℝ\ =>\ x\ \textasciicircum{}\ 4)\ μ}
\end{leancode}
\leansource{\texttt{drift} \quad PBGSO/\allowbreak{}FiniteWorkScopeAudit.lean, lines 36–36}
\begin{leancode}
\leanline{0}{def\ drift\ (x\ :\ ℝ)\ :\ ℝ\ :=\ -x-x\textasciicircum{}3}
\end{leancode}
\leansource{\texttt{drift\_\allowbreak{}strongly\_\allowbreak{}dissipative} \quad PBGSO/\allowbreak{}FiniteWorkScopeAudit.lean, lines 40–45}
\begin{leancode}
\leanline{0}{theorem\ drift\_strongly\_dissipative\ (x\ y\ :\ ℝ)\ :}
\leanline{4}{(x-y)*(drift\ x-drift\ y)\ ≤\ -(x-y)\textasciicircum{}2}
\end{leancode}
\leansource{\texttt{exists\_\allowbreak{}p2\_\allowbreak{}nonintegrable\_\allowbreak{}work} \quad PBGSO/\allowbreak{}FiniteWorkScopeAudit.lean, lines 49–66}
\begin{leancode}
\leanline{0}{theorem\ exists\_p2\_nonintegrable\_work\ :}
\leanline{4}{∃\ μ\ :\ Measure\ ℝ,\ IsProbabilityMeasure\ μ\ ∧\ MemLp\ (fun\ x\ :\ ℝ\ =>\ x)\ 2\ μ\ ∧}
\leanline{6}{¬\ Integrable\ (fun\ x\ :\ ℝ\ =>\ x*(drift\ x-drift\ 0))\ μ}
\end{leancode}
\leanlabel{How the Lean proof follows the argument}
\texttt{exists\_\allowbreak{}quadratic\_\allowbreak{}not\_\allowbreak{}quartic} argues by contradiction using \texttt{Universal\allowbreak{}Moment\allowbreak{}Growth.\allowbreak{}growth\_\allowbreak{}of\_\allowbreak{}probability\_\allowbreak{}integrability} and tests the resulting bound at the square root of \texttt{C+1}. \texttt{drift\_\allowbreak{}strongly\_\allowbreak{}dissipative} expands the polynomial and uses the nonnegative factor \texttt{x\textasciicircum{}2+x*y+y\textasciicircum{}2}. \texttt{exists\_\allowbreak{}p2\_\allowbreak{}nonintegrable\_\allowbreak{}work} obtains \texttt{Mem\allowbreak{}Lp} from integrability of the squared norm. Its final contradiction rewrites the quartic function as the negative work minus the quadratic term and invokes closure of integrable functions under subtraction.
\end{leanbox}
\section{\texorpdfstring{Local radii and time-uniform moments of integral equations with a dissipative drift}{Local radii and time-uniform moments of integral equations with a dissipative drift}}


This section studies continuous solutions \(Y:\mathbb R\to E\) of an integral equation with additive input,


\[Y(t)-Y(0)=\int_0^tf(Y(r))\,dr+B(t),\]


in a real Banach space \(E\), where \(f:E\to E\) is globally Lipschitz with constant \(L\ge0\) and \(B\) is a continuous input with \(B(0)=0\). In the applications \(E\) is a space of square-integrable random particle states, \(f\) is the kinetic or higher-order drift, and \(B\) is the scaled Brownian motion, whose increments satisfy \(\|B(a+s)-B(a)\|_{L^2}=\beta\sqrt s\).




On a time interval of length at most one, Grönwall's inequality bounds the state by the explicit radius \(r_L(\|f(0)\|,\Xi,\|Y(0)\|)=e^{L+1}(\|Y(0)\|+\|f(0)\|+\Xi)\), which is affine in the initial norm; consequences are bounds on the drift and a square-root modulus of continuity \(\|Y(t)-Y(0)\|\le c\sqrt t\). If moreover a bounded linear map \(C\) into a Hilbert space makes \(f\) dissipative, \(\langle C(x-y),C(f(x)-f(y))\rangle\le-\rho\|C(x-y)\|^2\) with \(\rho>0\), then two solutions driven by the same input contract at rate \(\rho\) in the metric \(\|C\cdot\|\). Comparing the solution on each unit time interval with the solution started at zero and driven by the same input increments gives the uniform moment bound


\[\|CY(t)\|\le\|CY(0)\|+\frac{S}{1-e^{-\rho}}+S\quad(t\ge0),\qquad S=\|C\|\,e^{L+1}\bigl(\|f(0)\|+\Xi\bigr).\]


The bound holds for all times with constants independent of \(t\); it is the reference radius in the proofs of Theorems New A and New B. No moment bound on the solution is assumed.

\subsection{\texorpdfstring{The one-unit radius}{The one-unit radius}}


\textbf{Definition.} For \(L\ge0\) and real numbers \(f_0,\Xi,y_0\) put


\[r_L(f_0,\Xi,y_0)=e^{L+1}\bigl(y_0+f_0+\Xi\bigr).\]


If \(f_0,\Xi,y_0\ge0\), then \(r_L(f_0,\Xi,y_0)\ge0\), since it is a product of the positive number \(e^{L+1}\) and a sum of nonnegative numbers. In the applications \(f_0=\|f(0)\|\) is the size of the drift at the origin, \(\Xi\) bounds the input on a unit interval and \(y_0\) is the initial norm. The radius is affine in \(y_0\) with slope \(e^{L+1}\).

\begin{leanbox}{Lean: \texttt{radius}, \texttt{radius\_\allowbreak{}nonneg}}
\leansource{\texttt{radius} \quad PBGSO/\allowbreak{}UniformIntegralMoments.lean, lines 19–20}
\begin{leancode}
\leanline{0}{def\ radius\ (L\ :\ ℝ≥0)\ (fzero\ noise\ initial\ :\ ℝ)\ :\ ℝ\ :=}
\leanline{2}{Real.exp\ ((L\ :\ ℝ)\ +\ 1)\ *\ (initial\ +\ fzero\ +\ noise)}
\end{leancode}
\leansource{\texttt{radius\_\allowbreak{}nonneg} \quad PBGSO/\allowbreak{}UniformIntegralMoments.lean, lines 22–24}
\begin{leancode}
\leanline{0}{theorem\ radius\_nonneg\ (L\ :\ ℝ≥0)\ \{fzero\ noise\ initial\ :\ ℝ\}}
\leanline{4}{(hf\ :\ 0\ ≤\ fzero)\ (hnoise\ :\ 0\ ≤\ noise)\ (hi\ :\ 0\ ≤\ initial)\ :}
\leanline{4}{0\ ≤\ radius\ L\ fzero\ noise\ initial}
\end{leancode}
\leanlabel{How the Lean proof follows the argument}
\texttt{radius (L : \ensuremath{\mathbb R}\ensuremath{\ge}0) (fzero noise initial : \ensuremath{\mathbb R})} is \texttt{Real.\allowbreak{}exp ((L : \ensuremath{\mathbb R}) + 1) * (initial + fzero + noise)}; the order of the arguments is \((L,f_0,\Xi,y_0)\). \texttt{radius\_\allowbreak{}nonneg} assumes \texttt{hf : 0 \ensuremath{\le} fzero}, \texttt{hnoise : 0 \ensuremath{\le} noise}, \texttt{hi : 0 \ensuremath{\le} initial} and is proved by \texttt{positivity}.
\end{leanbox}
\subsection{\texorpdfstring{A Grönwall bound on intervals of length at most one}{A Grönwall bound on intervals of length at most one}}


\textbf{Claim.} Let \(E\) be a real Banach space, \(f:E\to E\) Lipschitz with constant \(L\ge0\), and \(Y,B:\mathbb R\to E\) continuous. Let \(T\le1\), \(\Xi\ge0\), and assume that for all \(t\in[0,T]\)


\[Y(t)-Y(0)=\int_0^tf(Y(r))\,dr+B(t),\qquad\|B(t)\|\le\Xi .\]


Then \(\|Y(t)\|\le r_L(\|f(0)\|,\Xi,\|Y(0)\|)=e^{L+1}(\|Y(0)\|+\|f(0)\|+\Xi)\) for every \(t\in[0,T]\).




\textbf{Proof.} Fix \(t\in[0,T]\); then \(0\le T\), so \(0\in[0,T]\). Put \(K=L+1\) and \(c=\|f(0)\|+\Xi\).




\emph{The input vanishes at zero.} The equation at time \(0\) reads \(0=0+B(0)\), so \(B(0)=0\).




\emph{A differentiable quantity.} Let \(D(s)=Y(s)-B(s)\). By the equation, \(D(s)=Y(0)+\int_0^sf(Y(r))\,dr\) for \(s\in[0,T]\). The integrand \(r\mapsto f(Y(r))\) is continuous, because \(f\) is Lipschitz and \(Y\) is continuous. By the fundamental theorem of calculus, \(D\) has right derivative \(f(Y(s))\) at every \(s\in[0,T)\), and \(D\) is continuous.




\emph{Linear growth of the derivative.} For \(s\in[0,T)\), the Lipschitz bound and the triangle inequality give \(\|f(Y(s))\|\le\|f(Y(s))-f(0)\|+\|f(0)\|\le L\|Y(s)\|+\|f(0)\|\), and \(\|Y(s)\|\le\|D(s)\|+\|B(s)\|\le\|D(s)\|+\Xi\). Hence, using \(L\le K\), \(1\le K\) and the nonnegativity of all norms,


\[\|f(Y(s))\|\le L\|D(s)\|+L\Xi+\|f(0)\|\le K\|D(s)\|+Kc .\]


\emph{Grönwall's inequality.} We use it in the following form: if \(D:[0,T]\to E\) is continuous, has right derivatives \(D'(s)\) on \([0,T)\) with \(\|D'(s)\|\le K\|D(s)\|+\varepsilon\), and \(\|D(0)\|\le\delta\), then for \(K\ne0\) and \(s\in[0,T]\)


\[\|D(s)\|\le\delta e^{Ks}+\frac{\varepsilon}{K}\bigl(e^{Ks}-1\bigr).\]


With \(\delta=\|D(0)\|=\|Y(0)\|\) (because \(B(0)=0\)) and \(\varepsilon=Kc\) this gives \(\|D(t)\|\le\|Y(0)\|e^{Kt}+c(e^{Kt}-1)\).




\emph{Conclusion.} Since \(0\le t\le T\le1\) and \(K>0\), monotonicity of the exponential gives \(e^{Kt}\le e^K\). Therefore


\[\begin{aligned}\|Y(t)\|&\le\|D(t)\|+\Xi\\&\le(\|Y(0)\|+c)e^{Kt}-c+\Xi\\&=(\|Y(0)\|+c)e^{Kt}-\|f(0)\|\\&\le e^{K}\bigl(\|Y(0)\|+\|f(0)\|+\Xi\bigr).\qquad\square\end{aligned}\]


The hypothesis \(T\le1\) is used only in the last line; it makes the radius independent of the length of the interval.

\begin{leanbox}{Lean: \texttt{local\_\allowbreak{}radius}}
\leansource{\texttt{local\_\allowbreak{}radius} \quad PBGSO/\allowbreak{}UniformIntegralMoments.lean, lines 29–92}
\begin{leancode}
\leanline{0}{theorem\ local\_radius\ (f\ :\ E\ →\ E)\ (L\ :\ ℝ≥0)\ (hf\ :\ LipschitzWith\ L\ f)}
\leanline{4}{(Y\ B\ :\ ℝ\ →\ E)\ (hY\ :\ Continuous\ Y)\ (hB\ :\ Continuous\ B)}
\leanline{4}{(T\ Ξ\ :\ ℝ)\ (hT\ :\ T\ ≤\ 1)\ (hΞ\ :\ 0\ ≤\ Ξ)}
\leanline{4}{(heq\ :\ ∀\ t\ ∈\ Icc\ 0\ T,\ Y\ t\ -\ Y\ 0\ =\ (∫\ r\ in\ 0..t,\ f\ (Y\ r))\ +\ B\ t)}
\leanline{4}{(hinput\ :\ ∀\ t\ ∈\ Icc\ 0\ T,\ ‖B\ t‖\ ≤\ Ξ)}
\leanline{4}{(t\ :\ ℝ)\ (ht\ :\ t\ ∈\ Icc\ 0\ T)\ :}
\leanline{4}{‖Y\ t‖\ ≤\ radius\ L\ ‖f\ 0‖\ Ξ\ ‖Y\ 0‖}
\end{leancode}
\leanlabel{How the Lean proof follows the argument}
The integral \(\int_0^t\) is the interval integral \texttt{\ensuremath{\int} r in 0.\allowbreak{}.\allowbreak{}t, f (Y r)}, and \texttt{Icc 0 T} is \([0,T]\). In the proof \texttt{K} and \texttt{c} are as above.



\begin{enumerate}


\item \texttt{hidentity} is \(D(r)=Y(0)+\int_0^rf(Y)\). \texttt{hder} is the right derivative: \texttt{interval\allowbreak{}Integral.\allowbreak{}integral\_\allowbreak{}has\allowbreak{}Deriv\allowbreak{}At\_\allowbreak{}right} differentiates the integral at its upper limit (using continuity of \texttt{hforce}), and \texttt{congr\_\allowbreak{}of\_\allowbreak{}eventually\allowbreak{}Eq} transfers the derivative to \(D\) because the two functions agree on a right neighbourhood inside \([0,T]\) (\texttt{hmem}, \texttt{he}). \texttt{Has\allowbreak{}Deriv\allowbreak{}Within\allowbreak{}At g g' (Ici r) r} is Mathlib's right derivative at \(r\).

\item \texttt{hbzero} is \(B(0)=0\). \texttt{hbound} is the linear growth of the derivative; inside it, \texttt{hfn} and \texttt{hyn} are the two triangle inequalities and \texttt{nlinarith} combines them.

\item \texttt{hv} is Grönwall's inequality, the Mathlib lemma \texttt{norm\_\allowbreak{}le\_\allowbreak{}gronwall\allowbreak{}Bound\_\allowbreak{}of\_\allowbreak{}norm\_\allowbreak{}deriv\_\allowbreak{}right\_\allowbreak{}le}; its bound is \texttt{gronwall\allowbreak{}Bound \ensuremath{\delta} K \ensuremath{\varepsilon} x}, which for \(K\ne0\) equals \(\delta e^{Kx}+(\varepsilon/K)(e^{Kx}-1)\) (rewritten by \texttt{ite\_\allowbreak{}eq\_\allowbreak{}right}), and \texttt{hcdiv} simplifies \(Kc/K=c\).

\item \texttt{he} is \(e^{Kt}\le e^K\) (\texttt{Real.\allowbreak{}exp\_\allowbreak{}le\_\allowbreak{}exp}), \texttt{hprod} multiplies it by \(\|Y(0)\|+c\ge0\), \texttt{hyt} is \(\|Y(t)\|\le\|D(t)\|+\Xi\), and the final \texttt{nlinarith} is the last display.

\end{enumerate}
\end{leanbox}
\subsection{\texorpdfstring{Linear growth of a Lipschitz map}{Linear growth of a Lipschitz map}}


\textbf{Claim.} Let \(f:E\to E\) be Lipschitz with constant \(L\ge0\) on a normed space \(E\), and let \(x\in E\), \(P\in\mathbb R\) with \(\|x\|\le P\). Then \(\|f(x)\|\le\|f(0)\|+LP\).




\textbf{Proof.} By the triangle inequality and the Lipschitz bound, \(\|f(x)\|\le\|f(x)-f(0)\|+\|f(0)\|\le L\|x\|+\|f(0)\|\), and \(L\|x\|\le LP\) because \(L\ge0\). \(\square\)

\begin{leanbox}{Lean: \texttt{drift\_\allowbreak{}bound}}
\leansource{\texttt{drift\_\allowbreak{}bound} \quad PBGSO/\allowbreak{}UniformIntegralMoments.lean, lines 96–102}
\begin{leancode}
\leanline{0}{theorem\ drift\_bound\ (f\ :\ E\ →\ E)\ (L\ :\ ℝ≥0)\ (hf\ :\ LipschitzWith\ L\ f)}
\leanline{4}{(x\ :\ E)\ (P\ :\ ℝ)\ (hx\ :\ ‖x‖\ ≤\ P)\ :}
\leanline{4}{‖f\ x‖\ ≤\ ‖f\ 0‖\ +\ (L\ :\ ℝ)*P}
\end{leancode}
\leanlabel{How the Lean proof follows the argument}
\texttt{hb} combines \texttt{norm\_\allowbreak{}le\_\allowbreak{}norm\_\allowbreak{}sub\_\allowbreak{}add} (\(\|a\|\le\|a-b\|+\|b\|\)) with \texttt{hf.\allowbreak{}norm\_\allowbreak{}sub\_\allowbreak{}le x 0}; \texttt{nlinarith} adds the product of \texttt{hx} with \(L\ge0\).
\end{leanbox}
\subsection{\texorpdfstring{State, drift and increment bounds on a short interval}{State, drift and increment bounds on a short interval}}


\textbf{Claim.} Let \(E\), \(f\), \(L\), \(Y\), \(B\) be as in the Grönwall step, let \(0\le T\le1\), \(\Xi\ge0\), and assume that for all \(t\in[0,T]\)


\[Y(t)-Y(0)=\int_0^tf(Y(r))\,dr+B(t),\qquad\|B(t)\|\le\Xi\sqrt t .\]


Put \(P=r_L(\|f(0)\|,\Xi,\|Y(0)\|)\). Then for every \(t\in[0,T]\)


\[\|Y(t)\|\le P,\qquad\|f(Y(t))\|\le\|f(0)\|+LP,\qquad\|Y(t)-Y(0)\|\le\bigl(\|f(0)\|+LP+\Xi\bigr)\sqrt t .\]


\textbf{Proof.} For \(t\in[0,T]\) we have \(\sqrt t\le1\) because \(t\le1\), so \(\|B(t)\|\le\Xi\sqrt t\le\Xi\). The Grönwall step therefore gives \(\|Y(t)\|\le P\), and the previous step gives \(\|f(Y(t))\|\le\|f(0)\|+LP=:G\). For the increment, the equation at time \(t\) gives \(Y(t)-Y(0)=\int_0^tf(Y(r))\,dr+B(t)\). The norm of an interval integral is at most the length of the interval times a bound of the integrand, so \(\|\int_0^tf(Y(r))\,dr\|\le Gt\), and hence \(\|Y(t)-Y(0)\|\le Gt+\Xi\sqrt t\). Finally \(t\le\sqrt t\) for \(0\le t\le1\) (square both sides: \(t^2\le t\)), and \(G\ge0\), so \(Gt\le G\sqrt t\) and \(\|Y(t)-Y(0)\|\le(G+\Xi)\sqrt t\). \(\square\)




All three constants are affine in \(\|Y(0)\|\) and do not depend on the starting time of the interval. This is why the local errors in Theorems New A and New B grow at most linearly in the norm of the numerical state.

\begin{leanbox}{Lean: \texttt{local\_\allowbreak{}regularity}}
\leansource{\texttt{local\_\allowbreak{}regularity} \quad PBGSO/\allowbreak{}UniformIntegralMoments.lean, lines 106–136}
\begin{leancode}
\leanline{0}{theorem\ local\_regularity\ (f\ :\ E\ →\ E)\ (L\ :\ ℝ≥0)\ (hf\ :\ LipschitzWith\ L\ f)}
\leanline{4}{(Y\ B\ :\ ℝ\ →\ E)\ (hY\ :\ Continuous\ Y)\ (hB\ :\ Continuous\ B)}
\leanline{4}{(T\ Ξ\ :\ ℝ)\ (hT0\ :\ 0\ ≤\ T)\ (hT\ :\ T\ ≤\ 1)\ (hΞ\ :\ 0\ ≤\ Ξ)}
\leanline{4}{(heq\ :\ ∀\ t\ ∈\ Icc\ 0\ T,\ Y\ t\ -\ Y\ 0\ =\ (∫\ r\ in\ 0..t,\ f\ (Y\ r))\ +\ B\ t)}
\leanline{4}{(hinput\ :\ ∀\ t\ ∈\ Icc\ 0\ T,\ ‖B\ t‖\ ≤\ Ξ\ *\ Real.sqrt\ t)\ :}
\leanline{4}{∀\ t\ ∈\ Icc\ 0\ T,}
\leanline{6}{‖Y\ t‖\ ≤\ radius\ L\ ‖f\ 0‖\ Ξ\ ‖Y\ 0‖\ ∧}
\leanline{6}{‖f\ (Y\ t)‖\ ≤\ ‖f\ 0‖+(L\ :\ ℝ)*radius\ L\ ‖f\ 0‖\ Ξ\ ‖Y\ 0‖\ ∧}
\leanline{6}{‖Y\ t-Y\ 0‖\ ≤\ (‖f\ 0‖+(L\ :\ ℝ)*radius\ L\ ‖f\ 0‖\ Ξ\ ‖Y\ 0‖+Ξ)*Real.sqrt\ t}
\end{leancode}
\leanlabel{How the Lean proof follows the argument}
\texttt{P} and \texttt{D} are the radius and the drift bound \(G\). \texttt{hunit} derives \(\|B(t)\|\le\Xi\) using \texttt{Real.\allowbreak{}sqrt\_\allowbreak{}le\_\allowbreak{}one}. \texttt{hstate} is \texttt{local\_\allowbreak{}radius}, \texttt{hdrift} is \texttt{drift\_\allowbreak{}bound}, and \texttt{hbzero} is \(B(0)=0\). The increment bound \texttt{hm} is the lemma \texttt{integral\_\allowbreak{}segment\_\allowbreak{}sqrt\_\allowbreak{}modulus} from the file on integral moduli, applied with starting time \(a=0\) and length \(h=t\). That lemma assumes \(0\le a\), \(0\le h\le1\), \(a+h\le T\), a drift bound \(G\ge0\) on \([0,T]\) and an increment bound \(\|B(a+h)-B(a)\|\le\Xi\sqrt h\), and concludes \(\|Y(a+h)-Y(a)\|\le(G+\Xi)\sqrt h\); its proof is the integral estimate and \(h\le\sqrt h\) above.
\end{leanbox}
\subsection{\texorpdfstring{Contraction of two solutions with a common input}{Contraction of two solutions with a common input}}


\textbf{Claim.} Let \(E\) be a real Banach space, \(H\) a real inner product space, \(f:E\to E\) Lipschitz with constant \(L\ge0\), \(C:E\to H\) bounded linear and \(\rho\in\mathbb R\) such that


\[\bigl\langle C(x-y),C(f(x)-f(y))\bigr\rangle\le-\rho\|C(x-y)\|^2\qquad(x,y\in E).\qquad(1)\]


Let \(Y,Z,B:\mathbb R\to E\) with \(Y,Z\) continuous, \(T\in\mathbb R\), and assume \(Y(t)-Y(0)=\int_0^tf(Y)+B(t)\) and \(Z(t)-Z(0)=\int_0^tf(Z)+B(t)\) for \(t\in[0,T]\), with the same input \(B\). Then for \(t\in[0,T]\)


\[\|C(Y(t)-Z(t))\|\le e^{-\rho t}\,\|C(Y(0)-Z(0))\| .\]


\textbf{Proof.} We use the following result on additive-input equations, proved separately: if \(g:E\to E\) is Lipschitz, \(\rho\|C(x-y)\|^2\le\langle C(x-y),C(g(x)-g(y))\rangle\) for all \(x,y\), and \(Y,Z\) are continuous with \(Y(t)-Y(0)=-\int_0^tg(Y)+B(t)\), \(Z(t)-Z(0)=-\int_0^tg(Z)+B(t)\) on \([0,T]\), then \(\|C(Y(t)-Z(t))\|\le e^{-\rho t}\|C(Y(0)-Z(0))\|\) on \([0,T]\). Its proof: subtracting the equations, the common input cancels and \(Y-Z\) has right derivative \(-(g(Y)-g(Z))\); the energy \(\mathcal E(t)=\|C(Y(t)-Z(t))\|^2\) then has right derivative \(-2\langle C(Y-Z),C(g(Y)-g(Z))\rangle\le-2\rho\,\mathcal E(t)\), Grönwall's inequality gives \(\mathcal E(t)\le\mathcal E(0)e^{-2\rho t}\), and one takes square roots.




Apply this result with \(g=-f\), which is Lipschitz with the same constant. The monotonicity hypothesis holds: by linearity \(C(g(x)-g(y))=-C(f(x)-f(y))\), so \(\langle C(x-y),C(g(x)-g(y))\rangle=-\langle C(x-y),C(f(x)-f(y))\rangle\ge\rho\|C(x-y)\|^2\) by (1). The equations for \(Y\) and \(Z\) are of the required form because \(-\int_0^t(-f(Y))=\int_0^tf(Y)\). \(\square\)




No sign of \(\rho\) is assumed here. The noise is not required to be Brownian or adapted: the input cancels exactly because both solutions are driven by the same \(B\).

\begin{leanbox}{Lean: \texttt{contraction}}
\leansource{\texttt{contraction} \quad PBGSO/\allowbreak{}UniformIntegralMoments.lean, lines 143–161}
\begin{leancode}
\leanline{0}{theorem\ contraction\ (f\ :\ E\ →\ E)\ (L\ :\ ℝ≥0)\ (hf\ :\ LipschitzWith\ L\ f)}
\leanline{4}{(C\ :\ E\ →L[ℝ]\ H)\ (ρ\ :\ ℝ)}
\leanline{4}{(hmono\ :\ ∀\ x\ y,\ {\leanfallback ⟪}C\ (x-y),C\ (f\ x-f\ y){\leanfallback ⟫}\_ℝ\ ≤\ -ρ*‖C\ (x-y)‖\textasciicircum{}2)}
\leanline{4}{(Y\ Z\ B\ :\ ℝ\ →\ E)\ (hY\ :\ Continuous\ Y)\ (hZ\ :\ Continuous\ Z)\ (T\ :\ ℝ)}
\leanline{4}{(heY\ :\ ∀\ t\ ∈\ Icc\ 0\ T,\ Y\ t-Y\ 0\ =\ (∫\ r\ in\ 0..t,\ f\ (Y\ r))+B\ t)}
\leanline{4}{(heZ\ :\ ∀\ t\ ∈\ Icc\ 0\ T,\ Z\ t-Z\ 0\ =\ (∫\ r\ in\ 0..t,\ f\ (Z\ r))+B\ t)}
\leanline{4}{(t\ :\ ℝ)\ (ht\ :\ t\ ∈\ Icc\ 0\ T)\ :}
\leanline{4}{‖C\ (Y\ t-Z\ t)‖\ ≤\ Real.exp\ (-ρ*t)*‖C\ (Y\ 0-Z\ 0)‖}
\end{leancode}
\leanlabel{How the Lean proof follows the argument}
\texttt{contraction} applies \texttt{transformed\_\allowbreak{}additive\_\allowbreak{}integral\_\allowbreak{}contraction} from the file on additive contraction to \texttt{fun x => -\allowbreak{}f x} with the Lipschitz property \texttt{hf.\allowbreak{}neg}. The first goal is the monotonicity hypothesis: \texttt{he} is \(-f(x)-(-f(y))=-(f(x)-f(y))\), and \texttt{map\_\allowbreak{}neg}, \texttt{inner\_\allowbreak{}neg\_\allowbreak{}right} and \texttt{linarith} turn (1) into it. The other two goals are the integral equations, obtained from \texttt{he\allowbreak{}Y}, \texttt{he\allowbreak{}Z} with \texttt{interval\allowbreak{}Integral.\allowbreak{}integral\_\allowbreak{}neg} and \texttt{neg\_\allowbreak{}neg}.
\end{leanbox}
\subsection{\texorpdfstring{Moments uniform over all time}{Moments uniform over all time}}


\textbf{Claim.} Let \(E\), \(H\), \(f\), \(L\), \(C\) be as in the previous step, let \(\rho>0\) and assume (1). Let \(Y,B:\mathbb R\to E\) be continuous with


\[Y(t)-Y(0)=\int_0^tf(Y(r))\,dr+B(t)\qquad\text{for all }t\ge0,\]


and let \(\Xi\ge0\) be such that \(\|B(a+s)-B(a)\|\le\Xi\sqrt s\) for all \(a\ge0\) and \(s\in[0,1]\). Put \(S=\|C\|\,r_L(\|f(0)\|,\Xi,0)=\|C\|e^{L+1}(\|f(0)\|+\Xi)\). Then for every \(t\ge0\)


\[\|CY(t)\|\le\|CY(0)\|+\frac{S}{1-e^{-\rho}}+S .\]


\textbf{Proof.} Put \(q=e^{-\rho}\). Since \(\rho>0\), \(0\le q<1\) and \(1-q>0\). Also \(S\ge0\).




\emph{Step 1: one segment.} We show that for \(a\ge0\) and \(s\in[0,1]\)


\[\|CY(a+s)\|\le e^{-\rho s}\|CY(a)\|+S.\qquad(2)\]


Let \(\tilde B(r)=B(a+r)-B(a)\); it is continuous and \(\tilde B(0)=0\). By the existence theorem for Lipschitz integral equations with continuous additive input (Picard iteration, proved in the unadjusted Langevin library), there is a continuous \(Z\) with \(Z(0)=0\) and \(Z(r)=\int_0^rf(Z)+\tilde B(r)\) for \(r\in[0,1]\). The shifted trajectory \(r\mapsto Y(a+r)\) satisfies the same equation with the same input and initial value \(Y(a)\): subtracting the equations of \(Y\) at the times \(a\) and \(a+r\) gives \(Y(a+r)-Y(a)=\int_a^{a+r}f(Y)+\tilde B(r)=\int_0^rf(Y(a+u))\,du+\tilde B(r)\). The input of \(Z\) satisfies \(\|\tilde B(r)\|\le\Xi\sqrt r\) for \(r\in[0,1]\), so the short-interval bound (with \(T=1\) and initial norm \(0\)) gives \(\|Z(s)\|\le r_L(\|f(0)\|,\Xi,0)\), hence \(\|CZ(s)\|\le\|C\|\|Z(s)\|\le S\). The contraction step for the two solutions \(Y(a+\cdot)\) and \(Z\) gives \(\|C(Y(a+s)-Z(s))\|\le e^{-\rho s}\|C(Y(a)-0)\|\). By the triangle inequality,


\[\|CY(a+s)\|\le\|C(Y(a+s)-Z(s))\|+\|CZ(s)\|\le e^{-\rho s}\|CY(a)\|+S .\]


\emph{Step 2: integer times.} We show \(\|CY(n)\|\le\|CY(0)\|+S/(1-q)\) for \(n\in\mathbb N\) by induction. For \(n=0\) this holds because \(S/(1-q)\ge0\). If it holds for \(n\), then (2) with \(a=n\), \(s=1\) gives


\[\|CY(n+1)\|\le q\|CY(n)\|+S\le q\|CY(0)\|+\frac{qS}{1-q}+S\le\|CY(0)\|+\frac{S}{1-q},\]


where we used \(q\ge0\), then \(q\|CY(0)\|\le\|CY(0)\|\) (as \(q<1\)) and \(qS/(1-q)+S=S/(1-q)\).




\emph{Step 3: all times.} For \(t\ge0\) let \(n=\lfloor t\rfloor\) and \(s=t-n\in[0,1)\). By (2) with \(a=n\), and \(e^{-\rho s}\le1\) because \(\rho s\ge0\),


\[\|CY(t)\|\le e^{-\rho s}\|CY(n)\|+S\le\|CY(n)\|+S\le\|CY(0)\|+\frac{S}{1-q}+S. \qquad\square\]


The comparison with the solution \(Z\) started at zero is what removes any dependence on time: on each unit interval the state is within \(S\) of a contracted copy of its value at the start of the interval, and the contractions sum to the geometric series \(\sum_kq^k=1/(1-q)\).

\begin{leanbox}{Lean: \texttt{uniform\_\allowbreak{}moment}}
\leansource{\texttt{uniform\_\allowbreak{}moment} \quad PBGSO/\allowbreak{}UniformIntegralMoments.lean, lines 166–238}
\begin{leancode}
\leanline{0}{theorem\ uniform\_moment\ (f\ :\ E\ →\ E)\ (L\ :\ ℝ≥0)\ (hf\ :\ LipschitzWith\ L\ f)}
\leanline{4}{(C\ :\ E\ →L[ℝ]\ H)\ (ρ\ :\ ℝ)\ (hρ\ :\ 0\ <\ ρ)}
\leanline{4}{(hmono\ :\ ∀\ x\ y,\ {\leanfallback ⟪}C\ (x-y),C\ (f\ x-f\ y){\leanfallback ⟫}\_ℝ\ ≤\ -ρ*‖C\ (x-y)‖\textasciicircum{}2)}
\leanline{4}{(Y\ B\ :\ ℝ\ →\ E)\ (hY\ :\ Continuous\ Y)\ (hB\ :\ Continuous\ B)}
\leanline{4}{(heq\ :\ ∀\ t,\ 0\ ≤\ t\ →\ Y\ t-Y\ 0\ =\ (∫\ r\ in\ 0..t,\ f\ (Y\ r))+B\ t)}
\leanline{4}{(Ξ\ :\ ℝ)\ (hΞ\ :\ 0\ ≤\ Ξ)}
\leanline{4}{(hincrement\ :\ ∀\ a\ s,\ 0\ ≤\ a\ →\ s\ ∈\ Icc\ 0\ 1\ →}
\leanline{6}{‖B\ (a+s)-B\ a‖\ ≤\ Ξ*Real.sqrt\ s)}
\leanline{4}{(t\ :\ ℝ)\ (ht\ :\ 0\ ≤\ t)\ :}
\leanline{4}{‖C\ (Y\ t)‖\ ≤\ ‖C\ (Y\ 0)‖\ +}
\leanline{6}{(‖C‖*radius\ L\ ‖f\ 0‖\ Ξ\ 0)/(1-Real.exp\ (-ρ))\ +}
\leanline{6}{‖C‖*radius\ L\ ‖f\ 0‖\ Ξ\ 0}
\end{leancode}
\leanlabel{How the Lean proof follows the argument}
In the proof \texttt{S} and \texttt{q} are as above; \texttt{hq1} is \(q<1\) (\texttt{Real.\allowbreak{}exp\_\allowbreak{}lt\_\allowbreak{}one\_\allowbreak{}iff}) and \texttt{hgap} is \(1-q>0\).



\begin{enumerate}


\item \texttt{hsegment} is (2). \texttt{Bshift} is \(\tilde B\). \texttt{ULA.\allowbreak{}exists\_\allowbreak{}langevin\_\allowbreak{}integral} (Picard existence for \(Z(r)-x=-\int_0^rg(Z)+\tilde B(r)\) with Lipschitz \texttt{g}, here \texttt{fun x => -\allowbreak{}f x} and initial value \texttt{0}) produces \texttt{Z}; \texttt{he\allowbreak{}Z} rewrites its equation with \texttt{interval\allowbreak{}Integral.\allowbreak{}integral\_\allowbreak{}neg}. \texttt{he\allowbreak{}Y} is the shifted equation, from \texttt{Method\allowbreak{}Segments.\allowbreak{}integral\_\allowbreak{}shift}. \texttt{hloc} is the state bound of \texttt{local\_\allowbreak{}regularity} on \([0,1]\) for \(Z\) (with \texttt{hincrement} as input bound), \texttt{hc} is \texttt{contraction}, and \texttt{hcz} is \(\|CZ(s)\|\le\|C\|\|Z(s)\|\le S\) from \texttt{C.\allowbreak{}le\_\allowbreak{}op\allowbreak{}Norm}.

\item \texttt{hn} is Step 2, proved by \texttt{induction n}; in the \texttt{succ} case \texttt{hs} is (2) at \(a=n\), \(s=1\), \texttt{he} is \((1-q)\cdot S/(1-q)=S\), and \texttt{nlinarith} performs the displayed estimate.

\item Step 3 uses \texttt{n := ⌊t⌋\ensuremath{{}_{+}}}, the natural-number floor, with \texttt{Nat.\allowbreak{}floor\_\allowbreak{}le} and \texttt{Nat.\allowbreak{}lt\_\allowbreak{}floor\_\allowbreak{}add\_\allowbreak{}one} to show \(s=t-n\in[0,1]\), \texttt{hexp} for \(e^{-\rho s}\le1\) (\texttt{Real.\allowbreak{}exp\_\allowbreak{}le\_\allowbreak{}one\_\allowbreak{}iff}), and \texttt{linarith}.

\end{enumerate}
\end{leanbox}
\section{\texorpdfstring{A checked obstruction to the exact linear-flow interpretation}{A checked obstruction to the exact linear-flow interpretation}}


The printed coefficients define the numerical method, but an exact linear flow must also satisfy the differential equations of its fundamental matrix. At unit parameters, one necessary coefficient identity fails by a nonzero exact quantity.

\subsection{\texorpdfstring{Derivative mismatch at unit time}{Derivative mismatch at unit time}}


\textbf{Claim.} At \(\alpha=\gamma=1\), one has \(\mu_{22}'(t)=-1+e^{-t}\), and \(\mu_{22}'(1)-\mu_{32}(1)=1-3e^{-1}\ne0\). Consequently the printed coefficients cannot coincide at all positive times with the exact fundamental matrix of \(u'=v,\ v'=-u-v\).



\textbf{Proof.} Substituting unit parameters into the printed formula gives \(\mu_{22}(t)=2-t-e^{-t}\). Differentiating yields \(-1+e^{-t}\). Substituting the unit parameters and \(t=1\) in the other printed coefficient and subtracting gives exactly \(1-3e^{-1}\). Since \(e<3\) and \(e^{-1}>0\), multiplication gives \(1<3e^{-1}\), so the defect is negative and in particular nonzero. In an exact fundamental matrix, the first equation \(u'=v\) forces the derivative of the response of \(u\) to its initial \(u\) to equal the response of \(v\) to that same initial datum, namely \(\mu_{22}'=\mu_{32}\). Their mismatch contradicts that necessary identity. It does not alter the defined numerical map, and the local-error analysis above uses the actual printed coefficients instead. \(\square\)

\begin{leanbox}{Lean: \texttt{hol\_\allowbreak{}mu22\_\allowbreak{}unit\_\allowbreak{}hasDeriv}, \texttt{hol\_\allowbreak{}linear\_\allowbreak{}coefficient\_\allowbreak{}defect}, \texttt{hol\_\allowbreak{}linear\_\allowbreak{}coefficient\_\allowbreak{}not\_\allowbreak{}exact}}
\leansource{\texttt{hol\_\allowbreak{}mu22\_\allowbreak{}unit\_\allowbreak{}hasDeriv} \quad PBGSO/\allowbreak{}HOLCoefficientAudit.lean, lines 17–27}
\begin{leancode}
\leanline{0}{theorem\ hol\_mu22\_unit\_hasDeriv\ (t\ :\ ℝ)\ :}
\leanline{4}{HasDerivAt\ (fun\ h\ :\ ℝ\ =>\ (holCoefficients\ 1\ 1\ h).μ22)}
\leanline{6}{(-1\ +\ Real.exp\ (-t))\ t}
\end{leancode}
\leansource{\texttt{hol\_\allowbreak{}linear\_\allowbreak{}coefficient\_\allowbreak{}defect} \quad PBGSO/\allowbreak{}HOLCoefficientAudit.lean, lines 30–35}
\begin{leancode}
\leanline{0}{theorem\ hol\_linear\_coefficient\_defect\ :}
\leanline{4}{deriv\ (fun\ h\ :\ ℝ\ =>\ (holCoefficients\ 1\ 1\ h).μ22)\ 1\ -}
\leanline{6}{(holCoefficients\ 1\ 1\ 1).μ32\ =\ 1\ -\ 3\ *\ Real.exp\ (-1)}
\end{leancode}
\leansource{\texttt{hol\_\allowbreak{}linear\_\allowbreak{}coefficient\_\allowbreak{}not\_\allowbreak{}exact} \quad PBGSO/\allowbreak{}HOLCoefficientAudit.lean, lines 39–48}
\begin{leancode}
\leanline{0}{theorem\ hol\_linear\_coefficient\_not\_exact\ :}
\leanline{4}{deriv\ (fun\ h\ :\ ℝ\ =>\ (holCoefficients\ 1\ 1\ h).μ22)\ 1\ ≠}
\leanline{6}{(holCoefficients\ 1\ 1\ 1).μ32}
\end{leancode}
\leanlabel{How the Lean proof follows the argument}
\texttt{hol\_\allowbreak{}mu22\_\allowbreak{}unit\_\allowbreak{}has\allowbreak{}Deriv} differentiates the exponential with the chain rule and differentiates the two affine terms. \texttt{hol\_\allowbreak{}linear\_\allowbreak{}coefficient\_\allowbreak{}defect} rewrites by this derivative and simplifies the two printed coefficient expressions. \texttt{hol\_\allowbreak{}linear\_\allowbreak{}coefficient\_\allowbreak{}not\_\allowbreak{}exact} combines the exact defect with \texttt{Real.\allowbreak{}exp\_\allowbreak{}one\_\allowbreak{}lt\_\allowbreak{}three}, positivity of the inverse exponential, and \texttt{Real.\allowbreak{}exp\_\allowbreak{}add}. The contradiction is exact real arithmetic, not floating-point evaluation.
\end{leanbox}
\section{\texorpdfstring{Law contraction on arbitrary Brownian probability spaces}{Law contraction on arbitrary Brownian probability spaces}}


For probability laws \(\mu,\nu\) with finite second moment, let \(W_2(\mu,\nu)\) be the infimum of root mean squared displacement over their couplings. The optimiser–law distance is \(d((s,\mu),(s',\nu))=(\|s-s'\|^2+W_2(\mu,\nu)^2)^{1/2}\). The model may be realised on two different filtered Brownian probability spaces.

\subsection{\texorpdfstring{Identifying canonical laws before contracting}{Identifying canonical laws before contracting}}


\textbf{Claim.} Let two actual original strong solutions on \([0,T]\), \(T\ge0\), have initial optimiser states \(s_0,s'_0\) and probability sampler laws \(\mu,\nu\in\mathcal P_2\). Assume actual strong solutions also exist on each independent initial-law times canonical Brownian source and on each initial-law times Wiener-path source, with these initial conditions. Assume same-driver pathwise uniqueness on both original sources and both canonical sources after null augmentation and continuous modification. Suppose the model satisfies joint dissipativity with rate \(\rho\) in a quadratic geometry of condition number \(\kappa\).



In addition, assume the canonical solutions have jointly measurable drift and diffusion realizations, progressive drift, Borel diffusion sections, a bounded parameter path on the finite horizon, and a nonnegative linear-growth bound \(\|b(P(s),z)\|+\|(\sigma_i(P(s),z))_i\|\le C(1+\|P(s)\|+\|z\|)\). These are primitive coefficient regularity assumptions; their moment and energy consequences are proved.



Then the laws of the two original solutions satisfy

\[d((s_T,\mu_T),(s'_T,\nu_T))\le\sqrt\kappa\,e^{-\rho T}d((s_0,\mu),(s'_0,\nu)).\]

\textbf{Proof.} The canonical initial sampler is the first projection, whose law is the prescribed initial law because the independent Brownian factor has total mass one. For each original solution, the initial/\allowbreak{}path probability map transports its canonical path-space solution back to the original source. Pathwise uniqueness gives equality of the original and canonical optimiser trajectories and sampler marginal laws.



Apply the established canonical coupling contraction. Its construction starts with any initial coupling, uses common Brownian coordinates, derives the quadratic energy inequality and applies exponential comparison. Comparing the weighted and Euclidean energies contributes \(u/\ell=\kappa\); taking square roots contributes \(\sqrt\kappa\). Infimising over initial couplings produces \(W_2\), and therefore the displayed distance. Finally substitute the original/\allowbreak{}canonical law and optimiser identities. This last substitution introduces no assumption that the original sources were synchronously coupled. \(\square\)

\begin{leanbox}{Lean: \texttt{law\_\allowbreak{}distance\_\allowbreak{}decay}}
\leansource{\texttt{law\_\allowbreak{}distance\_\allowbreak{}decay} \quad PBGSO/\allowbreak{}RawSamplerGrowthArbitraryLaw.lean, lines 43–81}
\begin{leancode}
\leanline{0}{theorem\ law\_distance\_decay}
\leanline{4}{(U\ :\ RawLocalStationarity.Data\ Q\ D\ M\ T)\ (V\ :\ RawLocalStationarity.Data\ R\ E\ M\ T)}
\leanline{4}{(μ\ ν\ :\ Measure\ Z)\ [IsProbabilityMeasure\ μ]\ [IsProbabilityMeasure\ ν]}
\leanline{4}{(hμ\ :\ MemLp\ (fun\ z\ :\ Z\ =>\ z)\ 2\ μ)\ (hν\ :\ MemLp\ (fun\ z\ :\ Z\ =>\ z)\ 2\ ν)}
\leanline{4}{(hU\ :\ Q.map\ (U.sampler.state\ 0)\ =\ μ)\ (hV\ :\ R.map\ (V.sampler.state\ 0)\ =\ ν)}
\leanline{4}{(A\ :\ CanonicalSolution\ (M\ :=\ M)\ (T\ :=\ T)\ μ\ (U.optimizer\ 0))}
\leanline{4}{(B\ :\ CanonicalSolution\ (M\ :=\ M)\ (T\ :=\ T)\ ν\ (V.optimizer\ 0))}
\leanline{4}{(huniqueU\ :\ PathwiseUniqueness\ Q\ (nullAugmentedDriver\ Q\ (continuousDriver\ Q\ D))\ M\ T)}
\leanline{4}{(huniqueV\ :\ PathwiseUniqueness\ R\ (nullAugmentedDriver\ R\ (continuousDriver\ R\ E))\ M\ T)}
\leanline{4}{(huniqueA\ :\ PathwiseUniqueness\ (μ.prod\ (brownianProbability\ (I\ :=\ I)))}
\leanline{6}{(nullAugmentedDriver\ \_\ (continuousDriver\ \_\ (CanonicalBrownianDriver.driver\ (I\ :=\ I)\ μ)))\ M\ T)}
\leanline{4}{(huniqueB\ :\ PathwiseUniqueness\ (ν.prod\ (brownianProbability\ (I\ :=\ I)))}
\leanline{6}{(nullAugmentedDriver\ \_\ (continuousDriver\ \_\ (CanonicalBrownianDriver.driver\ (I\ :=\ I)\ ν)))\ M\ T)}
\leanline{4}{(hpathμ\ :\ ∃\ W\ :\ RawLocalStationarity.Data\ (BrownianPathDriver.probability\ (I\ :=\ I)\ μ)}
\leanline{8}{(BrownianPathDriver.driver\ μ)\ M\ T,}
\leanline{6}{W.optimizer\ 0\ =\ U.optimizer\ 0\ ∧}
\leanline{6}{W.sampler.state\ 0\ =ᵐ[BrownianPathDriver.probability\ (I\ :=\ I)\ μ]\ Prod.fst)}
\leanline{4}{(hpathν\ :\ ∃\ W\ :\ RawLocalStationarity.Data\ (BrownianPathDriver.probability\ (I\ :=\ I)\ ν)}
\leanline{8}{(BrownianPathDriver.driver\ ν)\ M\ T,}
\leanline{6}{W.optimizer\ 0\ =\ V.optimizer\ 0\ ∧}
\leanline{6}{W.sampler.state\ 0\ =ᵐ[BrownianPathDriver.probability\ (I\ :=\ I)\ ν]\ Prod.fst)}
\leanline{4}{(G\ :\ TwistedGeometry\ S\ Z)\ (ρ\ :\ ℝ)\ (hmono\ :\ TwistedJointMonotonicity\ M\ G\ ρ)\ :}
\leanline{4}{optimizerLawDistance\ (U.optimizer\ T)\ (V.optimizer\ T)}
\leanline{6}{(Q.map\ (U.sampler.state\ T))\ (R.map\ (V.sampler.state\ T))\ ≤}
\leanline{6}{Real.sqrt\ G.conditionNumber\ *\ Real.exp\ (-ρ\ *\ (T\ :\ ℝ))\ *}
\leanline{8}{optimizerLawDistance\ (U.optimizer\ 0)\ (V.optimizer\ 0)\ μ\ ν}
\end{leancode}
\leanlabel{How the Lean proof follows the argument}
\texttt{law\_\allowbreak{}distance\_\allowbreak{}decay} first proves \texttt{h\allowbreak{}A0} and \texttt{h\allowbreak{}B0} by pushforward of the first product projection. \texttt{law\_\allowbreak{}unique\_\allowbreak{}two\_\allowbreak{}sources} gives \texttt{h\allowbreak{}UA} and \texttt{h\allowbreak{}VB}, matching each original solution to its canonical one. It then calls \texttt{Raw\allowbreak{}Sampler\allowbreak{}Growth\allowbreak{}Law.\allowbreak{}law\_\allowbreak{}distance\_\allowbreak{}decay} with \texttt{A.\allowbreak{}regularity} and \texttt{B.\allowbreak{}regularity}. The final rewrites replace canonical initial states and terminal laws by the original data. \texttt{Mem\allowbreak{}Lp} at exponent two is square integrability, and \texttt{Q.\allowbreak{}map} denotes the distribution of a random variable under the probability measure \texttt{Q}.
\end{leanbox}
\section{\texorpdfstring{Law contraction on arbitrary Brownian probability spaces}{Law contraction on arbitrary Brownian probability spaces}}


For probability laws \(\mu,\nu\) with finite second moment, let \(W_2(\mu,\nu)\) be the infimum of root mean squared displacement over their couplings. The optimiser–law distance is \(d((s,\mu),(s',\nu))=(\|s-s'\|^2+W_2(\mu,\nu)^2)^{1/2}\). The model may be realised on two different filtered Brownian probability spaces.

\subsection{\texorpdfstring{Identifying canonical laws before contracting}{Identifying canonical laws before contracting}}


\textbf{Claim.} Let two actual original strong solutions on \([0,T]\), \(T\ge0\), have initial optimiser states \(s_0,s'_0\) and probability sampler laws \(\mu,\nu\in\mathcal P_2\). Assume actual strong solutions also exist on each independent initial-law times canonical Brownian source and on each initial-law times Wiener-path source, with these initial conditions. Assume same-driver pathwise uniqueness on both original sources and both canonical sources after null augmentation and continuous modification. Suppose the model satisfies joint dissipativity with rate \(\rho\) in a quadratic geometry of condition number \(\kappa\).



In addition, the complete optimiser, observation and sampler spaces are finite-dimensional where stipulated, \(\rho>0\), all sampler diffusion columns are constant across both states and parameters, the primitive drift, observation and optimiser maps are Borel measurable, and signed sampler coupling work is integrable for every probability coupling with finite second moment. This finite-work domain is an explicit assumption.



Then the laws of the two original solutions satisfy

\[d((s_T,\mu_T),(s'_T,\nu_T))\le\sqrt\kappa\,e^{-\rho T}d((s_0,\mu),(s'_0,\nu)).\]

\textbf{Proof.} The canonical initial sampler is the first projection, whose law is the prescribed initial law because the independent Brownian factor has total mass one. For each original solution, the initial/\allowbreak{}path probability map transports its canonical path-space solution back to the original source. Pathwise uniqueness gives equality of the original and canonical optimiser trajectories and sampler marginal laws.



Apply the established canonical coupling contraction. Its construction starts with any initial coupling, uses common Brownian coordinates, derives the quadratic energy inequality and applies exponential comparison. Comparing the weighted and Euclidean energies contributes \(u/\ell=\kappa\); taking square roots contributes \(\sqrt\kappa\). Infimising over initial couplings produces \(W_2\), and therefore the displayed distance. Finally substitute the original/\allowbreak{}canonical law and optimiser identities. This last substitution introduces no assumption that the original sources were synchronously coupled. \(\square\)

\begin{leanbox}{Lean: \texttt{law\_\allowbreak{}distance\_\allowbreak{}decay}}
\leansource{\texttt{law\_\allowbreak{}distance\_\allowbreak{}decay} \quad PBGSO/\allowbreak{}NaturalFiniteWorkArbitraryLaw.lean, lines 43–88}
\begin{leancode}
\leanline{0}{theorem\ law\_distance\_decay}
\leanline{4}{(U\ :\ RawLocalStationarity.Data\ Q\ D\ M\ T)\ (V\ :\ RawLocalStationarity.Data\ R\ E\ M\ T)}
\leanline{4}{(μ\ ν\ :\ Measure\ Z)\ [IsProbabilityMeasure\ μ]\ [IsProbabilityMeasure\ ν]}
\leanline{4}{(hμ\ :\ MemLp\ (fun\ z\ :\ Z\ =>\ z)\ 2\ μ)\ (hν\ :\ MemLp\ (fun\ z\ :\ Z\ =>\ z)\ 2\ ν)}
\leanline{4}{(hU\ :\ Q.map\ (U.sampler.state\ 0)\ =\ μ)\ (hV\ :\ R.map\ (V.sampler.state\ 0)\ =\ ν)}
\leanline{4}{(A\ :\ CanonicalSolution\ (M\ :=\ M)\ (T\ :=\ T)\ μ\ (U.optimizer\ 0))}
\leanline{4}{(B\ :\ CanonicalSolution\ (M\ :=\ M)\ (T\ :=\ T)\ ν\ (V.optimizer\ 0))}
\leanline{4}{(huniqueU\ :\ PathwiseUniqueness\ Q\ (nullAugmentedDriver\ Q\ (continuousDriver\ Q\ D))\ M\ T)}
\leanline{4}{(huniqueV\ :\ PathwiseUniqueness\ R\ (nullAugmentedDriver\ R\ (continuousDriver\ R\ E))\ M\ T)}
\leanline{4}{(huniqueA\ :\ PathwiseUniqueness\ (μ.prod\ (brownianProbability\ (I\ :=\ I)))}
\leanline{6}{(nullAugmentedDriver\ \_\ (continuousDriver\ \_\ (CanonicalBrownianDriver.driver\ (I\ :=\ I)\ μ)))\ M\ T)}
\leanline{4}{(huniqueB\ :\ PathwiseUniqueness\ (ν.prod\ (brownianProbability\ (I\ :=\ I)))}
\leanline{6}{(nullAugmentedDriver\ \_\ (continuousDriver\ \_\ (CanonicalBrownianDriver.driver\ (I\ :=\ I)\ ν)))\ M\ T)}
\leanline{4}{(hpathμ\ :\ ∃\ W\ :\ RawLocalStationarity.Data\ (BrownianPathDriver.probability\ (I\ :=\ I)\ μ)}
\leanline{8}{(BrownianPathDriver.driver\ μ)\ M\ T,}
\leanline{6}{W.optimizer\ 0\ =\ U.optimizer\ 0\ ∧}
\leanline{6}{W.sampler.state\ 0\ =ᵐ[BrownianPathDriver.probability\ (I\ :=\ I)\ μ]\ Prod.fst)}
\leanline{4}{(hpathν\ :\ ∃\ W\ :\ RawLocalStationarity.Data\ (BrownianPathDriver.probability\ (I\ :=\ I)\ ν)}
\leanline{8}{(BrownianPathDriver.driver\ ν)\ M\ T,}
\leanline{6}{W.optimizer\ 0\ =\ V.optimizer\ 0\ ∧}
\leanline{6}{W.sampler.state\ 0\ =ᵐ[BrownianPathDriver.probability\ (I\ :=\ I)\ ν]\ Prod.fst)}
\leanline{4}{(G\ :\ TwistedGeometry\ S\ Z)\ (ρ\ :\ ℝ)\ (hρ\ :\ 0\ <\ ρ)}
\leanline{4}{(hmono\ :\ TwistedJointMonotonicity\ M\ G\ ρ)}
\leanline{4}{(hfull\ :\ DiracOptimizerGrowth.FiniteCouplingWork\ M\ G)}
\leanline{4}{(hσ\ :\ ∀\ s\ s'\ z\ w\ i,\ M.diffusion\ (M.parameter\ s)\ z\ i\ =\ M.diffusion\ (M.parameter\ s')\ w\ i)}
\leanline{4}{(hb\ :\ ∀\ θ,\ Measurable\ (M.drift\ θ))}
\leanline{4}{(hObs\ :\ Measurable\ (fun\ p\ :\ S\ ×\ Z\ =>\ M.observation\ (M.parameter\ p.1)\ (M.position\ p.2)))}
\leanline{4}{(hDrift\ :\ Measurable\ (fun\ p\ :\ S\ ×\ Z\ =>\ M.drift\ (M.parameter\ p.1)\ p.2))}
\leanline{4}{(hOptimizer\ :\ Measurable\ (fun\ p\ :\ ℝ\ ×\ (S\ ×\ Θ)\ =>\ M.optimizer\ p.1\ p.2.1\ p.2.2))\ :}
\leanline{4}{optimizerLawDistance\ (U.optimizer\ T)\ (V.optimizer\ T)}
\leanline{6}{(Q.map\ (U.sampler.state\ T))\ (R.map\ (V.sampler.state\ T))\ ≤}
\leanline{6}{Real.sqrt\ G.conditionNumber\ *\ Real.exp\ (-ρ\ *\ (T\ :\ ℝ))\ *}
\leanline{8}{optimizerLawDistance\ (U.optimizer\ 0)\ (V.optimizer\ 0)\ μ\ ν}
\end{leancode}
\leanlabel{How the Lean proof follows the argument}
\texttt{law\_\allowbreak{}distance\_\allowbreak{}decay} proves \texttt{h\allowbreak{}A0} and \texttt{h\allowbreak{}B0} from first-projection product laws, then obtains the source comparison identities \texttt{h\allowbreak{}UA} and \texttt{h\allowbreak{}VB} using \texttt{law\_\allowbreak{}unique\_\allowbreak{}two\_\allowbreak{}sources}. \texttt{Natural\allowbreak{}Finite\allowbreak{}Work\allowbreak{}Law.\allowbreak{}law\_\allowbreak{}distance\_\allowbreak{}decay} receives the explicit finite-work, constant-noise and Borel hypotheses \texttt{hfull}, \texttt{h\ensuremath{\sigma}}, \texttt{hb}, \texttt{h\allowbreak{}Obs}, \texttt{h\allowbreak{}Drift} and \texttt{h\allowbreak{}Optimizer}. The final rewrites transfer that canonical result back. The finite-work premise remains visible, rather than being inferred from the conditional monotonicity inequality.
\end{leanbox}
\section{\texorpdfstring{Independent reference particles and the interaction error}{Independent reference particles and the interaction error}}


Take \(N>0\) independent copies \(\bar Z^i_t\) of one actual mean-field sampler trajectory, with its deterministic complete optimiser \(s_t\). The interacting empirical solution uses the corresponding Brownian coordinate in each particle equation. Its complete optimiser \(S^N_t\) is shared by the particles. The empirical comparison distance squared is \(\|S^N_t-s_t\|^2+W_2(\widehat\nu^N_t,\widehat{\bar\nu}^N_t)^2\).

\subsection{\texorpdfstring{A forced-energy bound for the actual empirical equation}{A forced-energy bound for the actual empirical equation}}


\textbf{Claim.} Suppose actual mean-field and empirical strong solutions exist on \([0,T]\). Assume Borel observation, drift and optimiser maps, joint dissipativity with \(\rho>0\), a quadratic metric of condition number \(\kappa\), and optimiser response Lipschitz constant \(L_\Psi\ge0\) in its mean-observation argument. The observation along every mean-field law is square integrable and has centred second moment at most \(V\ge0\). Initialise the empirical solution almost surely with the same optimiser and reference cloud. Then

\[\left(\mathbb E\big[\|S^N_T-s_T\|^2+W_2(\widehat\nu^N_T,\widehat{\bar\nu}^N_T)^2\big]\right)^{1/2}\le\frac{L_\Psi}{\rho}\sqrt{\frac{\kappa V}{N}}(1-e^{-\rho T}).\]

\textbf{Proof.} Construct the reference cloud on the \(N\)-fold product of the original probability space. By the product pushforward identity its coordinates are independent and each has the original mean-field law. Its empirical observation fluctuation is a sum of independent centred vectors. The mixed inner products have expectation zero; expanding the squared norm gives centred variance at most \(V/N\).



Pair each empirical particle with its own reference particle. Localisation of the actual coupled Brownian equations justifies the weighted energy derivative. Joint dissipativity bounds the matching drift part by \(-2\rho E\). The remaining optimiser response is at most \(L_\Psi\) times the empirical observation fluctuation; Cauchy–Schwarz bounds its energy contribution by \(2c\sqrt E\), with the metric factors giving the final constant \(L_\Psi\sqrt{\kappa V/N}\) after conversion to Euclidean distance. The initial paired error is zero. The forced square-root comparison yields \((c/\rho)(1-e^{-\rho T})\). The labelled particle pairs give a coupling of their empirical measures, so the empirical Wasserstein cost is no larger than their mean squared labelled displacement. Thus the same bound controls the stated distance. \(\square\)

\begin{leanbox}{Lean: \texttt{originalPair}, \texttt{original\_\allowbreak{}finite\_\allowbreak{}particle\_\allowbreak{}interaction\_\allowbreak{}bound}}
\leansource{\texttt{originalPair} \quad PBGSO/\allowbreak{}RawConstructedParticleBounds.lean, lines 31–32}
\begin{leancode}
\leanline{0}{def\ originalPair\ (r\ :\ ℝ)\ (ω\ :\ Fin\ N\ →\ Ω)\ :\ ParticlePair\ N\ S\ Z\ :=}
\leanline{2}{(oneEquiv\ N\ (Y.state\ r\ ω),(U.optimizer\ r,RawIndependentCopies.cloud\ Q\ D\ N\ U\ r\ ω))}
\end{leancode}
\leansource{\texttt{original\_\allowbreak{}finite\_\allowbreak{}particle\_\allowbreak{}interaction\_\allowbreak{}bound} \quad PBGSO/\allowbreak{}RawConstructedParticleBounds.lean, lines 109–132}
\begin{leancode}
\leanline{0}{theorem\ original\_finite\_particle\_interaction\_bound}
\leanline{4}{(hBorelF\ :\ Measurable\ (fun\ p\ :\ S\ ×\ Z\ =>\ M.observation\ (M.parameter\ p.1)\ (M.position\ p.2)))}
\leanline{4}{(hBorelb\ :\ Measurable\ (fun\ p\ :\ S\ ×\ Z\ =>\ M.drift\ (M.parameter\ p.1)\ p.2))}
\leanline{4}{(hBorelΨ\ :\ Measurable\ (fun\ p\ :\ ℝ\ ×\ (S\ ×\ Θ)\ =>\ M.optimizer\ p.1\ p.2.1\ p.2.2))}
\leanline{4}{(G\ :\ TwistedGeometry\ S\ Z)\ (ρ\ :\ ℝ)\ (hρ\ :\ 0\ <\ ρ)}
\leanline{4}{(hmono\ :\ TwistedJointMonotonicity\ M\ G\ ρ)\ (hN\ :\ 0\ <\ N)}
\leanline{4}{(LΨ\ :\ ℝ≥0)\ (V\ :\ ℝ)\ (hV\ :\ 0\ ≤\ V)}
\leanline{4}{(hΨ\ :\ ∀\ t\ ∈\ Icc\ 0\ (T\ :\ ℝ),\ ∀\ s\ g\ g',}
\leanline{6}{‖M.optimizer\ t\ s\ g-M.optimizer\ t\ s\ g'‖\ ≤\ (LΨ\ :\ ℝ)*‖g-g'‖)}
\leanline{4}{(hF\ :\ ∀\ t\ ∈\ Icc\ 0\ (T\ :\ ℝ),\ MemLp}
\leanline{6}{(fun\ z\ =>\ M.observation\ (M.parameter\ (U.optimizer\ t))\ (M.position\ z))\ 2}
\leanline{6}{(Q.map\ (U.sampler.state\ t)))}
\leanline{4}{(hVar\ :\ ∀\ t\ ∈\ Icc\ 0\ (T\ :\ ℝ),}
\leanline{6}{(∫\ z,\ ‖M.observation\ (M.parameter\ (U.optimizer\ t))\ (M.position\ z)-}
\leanline{8}{M.meanGradient\ (Q.map\ (U.sampler.state\ t))\ (U.optimizer\ t)‖\textasciicircum{}2}
\leanline{10}{∂Q.map\ (U.sampler.state\ t))\ ≤\ V)}
\leanline{4}{(hinit\ :\ ∀ᵐ\ ω\ ∂IndependentCopies.seeds\ (J\ :=\ Fin\ N)\ Q,}
\leanline{6}{oneEquiv\ N\ (Y.state\ 0\ ω)\ =\ (U.optimizer\ 0,RawIndependentCopies.cloud\ Q\ D\ N\ U\ 0\ ω))\ :}
\leanline{4}{Real.sqrt\ (∫\ ω,\ particleEmpiricalDistanceSq\ N\ (originalPair\ Q\ D\ M\ N\ U\ Y\ T\ ω)}
\leanline{6}{∂IndependentCopies.seeds\ (J\ :=\ Fin\ N)\ Q)\ ≤}
\leanline{6}{((LΨ\ :\ ℝ)/ρ)*Real.sqrt\ (G.conditionNumber*V/(N\ :\ ℝ))*(1-Real.exp\ (-ρ*T))}
\end{leancode}
\leanlabel{How the Lean proof follows the argument}
\texttt{original\allowbreak{}Pair} is the original interacting trajectory paired with the constructed independent copies. \texttt{original\_\allowbreak{}finite\_\allowbreak{}particle\_\allowbreak{}interaction\_\allowbreak{}bound} builds \texttt{Raw\allowbreak{}Constructed\allowbreak{}Particle\allowbreak{}Copy.\allowbreak{}build} from the three Borel hypotheses and calls \texttt{finite\_\allowbreak{}particle\_\allowbreak{}interaction\_\allowbreak{}bound}. That proof establishes the actual reference cloud product law with \texttt{reference\_\allowbreak{}cloud\_\allowbreak{}law}, optimiser identity with \texttt{reference\_\allowbreak{}optimizer}, and zero initial discrepancy with \texttt{initial\_\allowbreak{}coupling}. It applies the proved paired bound, then uses \texttt{pair\_\allowbreak{}state} and \texttt{integral\_\allowbreak{}congr\_\allowbreak{}ae} to replace the constructed pair by the original native state. These identities are conclusions of the construction.
\end{leanbox}
\section{\texorpdfstring{Parameter error by stationary and particle comparison}{Parameter error by stationary and particle comparison}}


The parameter projection is assumed \(1\)-Lipschitz in the complete optimiser norm. A stationary complete optimiser–sampler pair \((s_\star,\nu_\star)\) is compared first with the mean-field solution and then with the empirical solution. The bound contains an exponentially decaying initial-law term and a finite-particle term of order \(N^{-1/2}\).

\subsection{\texorpdfstring{The native parameter estimate with constant diffusion}{The native parameter estimate with constant diffusion}}


\textbf{Claim.} Assume actual original mean-field and empirical strong equations on \([0,T]\), \(T\ge0\), with \(N>0\), probability initial law \(\mu\in\mathcal P_2\), and synchronised initial empirical/\allowbreak{}reference cloud. Assume the Borel coefficient maps, joint dissipativity at \(\rho>0\), \(1\)-Lipschitz parameter projection, constant sampler diffusion, optimiser response bound \(L_\Psi\ge0\), and square-integrable observations with centred variance at most \(V\) throughout the horizon. Suppose a canonical strong solution with initial law \(\mu\) exists, with same-driver uniqueness on the original and canonical sources and a strong Wiener-path-source solution for that initial law. Suppose an actual canonical comparison solution has constant optimiser \(s_\star\), constant sampler law \(\nu_\star\in\mathcal P_2\), and integrable observation at that state. Then the parameter discrepancy is square integrable and

\[\begin{aligned}\|P(S^N_T)-P(s_\star)\|_{L^2}\le {}&\sqrt\kappa\,e^{-\rho T}d((s_0,\mu),(s_\star,\nu_\star))\\&+\frac{L_\Psi}{\rho}\sqrt{\frac{\kappa V}{N}}(1-e^{-\rho T}).\end{aligned}\]

\textbf{Proof.} The variance bound at time zero implies \(V\ge0\), since it bounds an integral of a squared norm. Apply the native paired interaction estimate: the triangle inequality in \(L^2\) and the projection Lipschitz bound give the finite-particle term plus \(\|s_T-s_\star\|\), and establish square integrability. The constant-noise law comparison uses the actual response and variance assumptions and the stationary comparison solution to give \(d((s_T,\nu_T),(s_\star,\nu_\star))\le\sqrt\kappa e^{-\rho T}d((s_0,\mu),(s_\star,\nu_\star))\). The optimiser norm is at most this distance, because its other squared summand is nonnegative. Substituting and commuting the two nonnegative error terms proves the claim. This particular comparison uses the response/\allowbreak{}variance hypotheses; it does not assume universal finite signed coupling work. \(\square\)

\begin{leanbox}{Lean: \texttt{original\_\allowbreak{}parameter\_\allowbreak{}bound\_\allowbreak{}constant}}
\leansource{\texttt{original\_\allowbreak{}parameter\_\allowbreak{}bound\_\allowbreak{}constant} \quad PBGSO/\allowbreak{}RawConstructedParticleParameter.lean, lines 98–153}
\begin{leancode}
\leanline{0}{theorem\ original\_parameter\_bound\_constant}
\leanline{4}{(U\ :\ RawLocalStationarity.Data\ Q\ W\ M\ T)}
\leanline{4}{(Y\ :\ RawEmpiricalDynamics.Solution\ M\ N\ (IndependentCopies.seeds\ (J\ :=\ Fin\ N)\ Q)}
\leanline{6}{(RawIndependentCopies.driver\ Q\ W\ N)\ T)}
\leanline{4}{(μ\ νStar\ :\ Measure\ Z)\ [IsProbabilityMeasure\ μ]\ [IsProbabilityMeasure\ νStar]}
\leanline{4}{(hμ\ :\ MemLp\ (fun\ z\ :\ Z\ =>\ z)\ 2\ μ)\ (hν\ :\ MemLp\ (fun\ z\ :\ Z\ =>\ z)\ 2\ νStar)}
\leanline{4}{(hUlaw\ :\ Q.map\ (U.sampler.state\ 0)\ =\ μ)}
\leanline{4}{(A\ :\ RawLocalStationarity.Data\ (μ.prod\ (brownianProbability\ (I\ :=\ I)))}
\leanline{6}{(CanonicalBrownianDriver.driver\ (I\ :=\ I)\ μ)\ M\ T)}
\leanline{4}{(hA0\ :\ A.optimizer\ 0\ =\ U.optimizer\ 0)}
\leanline{4}{(hAz\ :\ A.sampler.state\ 0\ =ᵐ[μ.prod\ (brownianProbability\ (I\ :=\ I))]\ Prod.fst)}
\leanline{4}{(V\ :\ RawLocalStationarity.Data\ (νStar.prod\ (brownianProbability\ (I\ :=\ I)))}
\leanline{6}{(CanonicalBrownianDriver.driver\ (I\ :=\ I)\ νStar)\ M\ T)}
\leanline{4}{(hVz\ :\ V.sampler.state\ 0\ =ᵐ[νStar.prod\ (brownianProbability\ (I\ :=\ I))]\ Prod.fst)}
\leanline{4}{(huniqueU\ :\ PathwiseUniqueness\ Q\ (nullAugmentedDriver\ Q\ (continuousDriver\ Q\ W))\ M\ T)}
\leanline{4}{(huniqueA\ :\ PathwiseUniqueness\ (μ.prod\ (brownianProbability\ (I\ :=\ I)))}
\leanline{6}{(nullAugmentedDriver\ \_\ (continuousDriver\ \_\ (CanonicalBrownianDriver.driver\ (I\ :=\ I)\ μ)))\ M\ T)}
\leanline{4}{(hpath\ :\ ∃\ B\ :\ RawLocalStationarity.Data\ (BrownianPathDriver.probability\ (I\ :=\ I)\ μ)}
\leanline{8}{(BrownianPathDriver.driver\ μ)\ M\ T,}
\leanline{6}{B.optimizer\ 0\ =\ U.optimizer\ 0\ ∧}
\leanline{6}{B.sampler.state\ 0\ =ᵐ[BrownianPathDriver.probability\ (I\ :=\ I)\ μ]\ Prod.fst)}
\leanline{4}{(G\ :\ TwistedGeometry\ S\ Z)\ (ρ\ :\ ℝ)\ (hρ\ :\ 0\ <\ ρ)}
\leanline{4}{(hmono\ :\ TwistedJointMonotonicity\ M\ G\ ρ)\ (hN\ :\ 0\ <\ N)}
\leanline{4}{(hProjection\ :\ LipschitzWith\ 1\ M.parameter)}
\leanline{4}{(hσ\ :\ ∀\ s\ s'\ z\ w\ i,\ M.diffusion\ (M.parameter\ s)\ z\ i\ =\ M.diffusion\ (M.parameter\ s')\ w\ i)}
\leanline{4}{(hb\ :\ ∀\ θ,\ Measurable\ (M.drift\ θ))}
\leanline{4}{(hF\ :\ Measurable\ (fun\ p\ :\ S\ ×\ Z\ =>\ M.observation\ (M.parameter\ p.1)\ (M.position\ p.2)))}
\leanline{4}{(hDrift\ :\ Measurable\ (fun\ p\ :\ S\ ×\ Z\ =>\ M.drift\ (M.parameter\ p.1)\ p.2))}
\leanline{4}{(hOptimizer\ :\ Measurable\ (fun\ p\ :\ ℝ\ ×\ (S\ ×\ Θ)\ =>\ M.optimizer\ p.1\ p.2.1\ p.2.2))}
\leanline{4}{(LΨ\ :\ ℝ≥0)\ (VF\ :\ ℝ)}
\leanline{4}{(hΨ\ :\ ∀\ r\ ∈\ Icc\ 0\ (T\ :\ ℝ),\ ∀\ s\ g\ g',}
\leanline{6}{‖M.optimizer\ r\ s\ g-M.optimizer\ r\ s\ g'‖\ ≤\ (LΨ\ :\ ℝ)*‖g-g'‖)}
\leanline{4}{(hU2\ :\ ∀\ r\ ∈\ Icc\ 0\ (T\ :\ ℝ),\ MemLp}
\leanline{6}{(fun\ z\ =>\ M.observation\ (M.parameter\ (U.optimizer\ r))\ (M.position\ z))\ 2\ (Q.map\ (U.sampler.state\ r)))}
\leanline{4}{(hVar\ :\ ∀\ r\ ∈\ Icc\ 0\ (T\ :\ ℝ),}
\leanline{6}{(∫\ z,\ ‖M.observation\ (M.parameter\ (U.optimizer\ r))\ (M.position\ z)-}
\leanline{8}{M.meanGradient\ (Q.map\ (U.sampler.state\ r))\ (U.optimizer\ r)‖\textasciicircum{}2\ ∂Q.map\ (U.sampler.state\ r))\ ≤\ VF)}
\leanline{4}{(sStar\ :\ S)}
\leanline{4}{(hopt\ :\ ∀\ r\ ∈\ Icc\ 0\ (T\ :\ ℝ),\ V.optimizer\ r\ =\ sStar)}
\leanline{4}{(hlaw\ :\ ∀\ r\ ∈\ Icc\ 0\ (T\ :\ ℝ),\ (νStar.prod\ (brownianProbability\ (I\ :=\ I))).map\ (V.sampler.state\ r)\ =\ νStar)}
\leanline{4}{(hobs\ :\ Integrable\ (fun\ z\ =>\ M.observation\ (M.parameter\ sStar)\ (M.position\ z))\ νStar)}
\leanline{4}{(hinit\ :\ ∀ᵐ\ ω\ ∂IndependentCopies.seeds\ (J\ :=\ Fin\ N)\ Q,}
\leanline{6}{oneEquiv\ N\ (Y.state\ 0\ ω)\ =\ (U.optimizer\ 0,RawIndependentCopies.cloud\ Q\ W\ N\ U\ 0\ ω))\ :}
\leanline{4}{MemLp\ (fun\ ω\ =>\ M.parameter\ (oneEquiv\ N\ (Y.state\ T\ ω)).1-M.parameter\ sStar)}
\leanline{6}{2\ (IndependentCopies.seeds\ (J\ :=\ Fin\ N)\ Q)\ ∧}
\leanline{4}{Real.sqrt\ (∫\ ω,\ ‖M.parameter\ (oneEquiv\ N\ (Y.state\ T\ ω)).1-M.parameter\ sStar‖\textasciicircum{}2}
\leanline{6}{∂IndependentCopies.seeds\ (J\ :=\ Fin\ N)\ Q)\ ≤}
\leanline{6}{Real.sqrt\ G.conditionNumber*Real.exp\ (-ρ*(T\ :\ ℝ))*optimizerLawDistance\ (U.optimizer\ 0)\ sStar\ μ\ νStar+}
\leanline{6}{((LΨ\ :\ ℝ)/ρ)*Real.sqrt\ (G.conditionNumber*VF/(N\ :\ ℝ))*(1-Real.exp\ (-ρ*T))}
\end{leancode}
\leanlabel{How the Lean proof follows the argument}
\texttt{original\_\allowbreak{}parameter\_\allowbreak{}bound\_\allowbreak{}constant} derives \texttt{h\allowbreak{}VF} using \texttt{integral\_\allowbreak{}nonneg} and the initial variance bound. \texttt{native\_\allowbreak{}parameter\_\allowbreak{}bias\_\allowbreak{}bound} provides both \texttt{Mem\allowbreak{}Lp} and the interaction-plus-relaxation estimate. \texttt{Raw\allowbreak{}Natural\allowbreak{}Constant\allowbreak{}Arbitrary\allowbreak{}Law.\allowbreak{}law\_\allowbreak{}decay} supplies the actual stationary comparison with the stated canonical/\allowbreak{}path-source hypotheses. \texttt{optimizer\_\allowbreak{}norm\_\allowbreak{}le\_\allowbreak{}law\allowbreak{}Distance} discards the nonnegative transport component. The last line adds these inequalities and reorders the two terms. The projection and original native empirical state are retained in the conclusion.
\end{leanbox}
\section{\texorpdfstring{Contracting metrics for the two particle methods from the joint monotonicity assumption}{Contracting metrics for the two particle methods from the joint monotonicity assumption}}


This section constructs, for each of the two numerical methods, a Hilbert metric on the random particle states in which the drift of the continuous particle system is dissipative with the rate \(\rho\) of assumption A3. The random state of M-MFPO is \(y=((m,V),(\theta,X))\): the parameter momentum \(m\) and parameter \(\theta\) are square-integrable random vectors in \(\Theta\), and the particle velocities \(V\) and positions \(X\) are square-integrable random cloud vectors in \(\mathcal C\). The random state of HOL-MYPGD is \(y=(\theta,X,U,V)\) with three random cloud vectors. In both cases the state space is normed by the maximum of the four \(L^2\) norms, written \(\|y\|\).




Assumption A3 is a statement about laws: for any two optimiser states and any coupling of two sampler laws, the weighted work of the drifts is at most \(-\rho\) times the weighted squared distance. Applied to the empirical laws of two clouds, it becomes a dissipativity inequality for the empirical drift on the finite-dimensional state space \(S\times Z^N\) (optimiser phase and \(N\) particle phases). Factoring the weight, integrating over the sample space and transporting by the linear regrouping of coordinates gives a bounded linear map \(T\) from the random states into \(L^2(\Omega;\mathbb R^n)\) with


\[m\|y\|\le\|Ty\|\le U\|y\|,\qquad\bigl\langle T(y-y'),T(b(y)-b(y'))\bigr\rangle\le-\rho\|T(y-y')\|^2,\]


where \(b\) is the kinetic or higher-order drift and \(m,U>0\) depend on the A3 eigenvalue bounds, on \(N\) and on the norms of the regrouping map. This metric is used for the exact-flow contraction and the uniform moment bounds of Theorems New A and New B. The section also proves the samplewise identity between the drift computed on the random variables and the empirical drift of the paper's model, which is what makes A3 applicable.

\subsection{\texorpdfstring{The particle drift of compatible primitive maps is the empirical drift}{The particle drift of compatible primitive maps is the empirical drift}}


We recall the paper's model. It consists of real Hilbert spaces \(S\) (optimiser states) and \(Z\) (sampler states), a parameter map \(\pi:S\to\Theta\), a position map \(\mathrm{pos}:Z\to\mathcal X\), an observation \(F:\Theta\times\mathcal X\to\Theta\), an optimiser response \(\psi_t(s,g)\in S\) for \(t\ge0\), \(s\in S\), \(g\in\Theta\), a sampler drift \(b(\theta,z)\in Z\) and diffusion columns \(\sigma_j(\theta,z)\in Z\), \(j\) in a finite index set. For \(N\ge1\) the \emph{empirical state space} is \(H_N=S\times Z^N\) with the Hilbert norm \(\|(s,z)\|^2=\|s\|^2+\sum_{i=1}^N\|z_i\|^2\), and the \emph{empirical drift} is


\[b_N\bigl(t,(s,z)\bigr)=\Bigl(\psi_t\Bigl(s,\frac1N\sum_{i=1}^NF\bigl(\pi(s),\mathrm{pos}(z_i)\bigr)\Bigr),\;\bigl(b(\pi(s),z_i)\bigr)_{i=1}^N\Bigr).\]


Primitive maps for the same dynamics are an observation \(o:S\times Z\to\Theta\), a response \(r:S\times\Theta\to S\) and a sampler field \(\beta:S\times Z\to Z\); their particle field is \(\Phi_N(s,z)=\bigl(r(s,N^{-1}\sum_io(s,z_i)),(\beta(s,z_i))_i\bigr)\). The model is \emph{compatible} with the primitive maps if \(F(\pi(s),\mathrm{pos}(z))=o(s,z)\), \(\psi_t(s,g)=r(s,g)\) for \(t\ge0\), \(b(\pi(s),z)=\beta(s,z)\), and the diffusion columns do not depend on the state.




\textbf{Claim.} If the model is compatible with the primitive maps, then \(\Phi_N(s,z)=b_N(t,(s,z))\) for every \(t\ge0\) and every \((s,z)\in S\times Z^N\).




\textbf{Proof.} Substitute the compatibility identities into the definition of \(b_N\): the average \(N^{-1}\sum_iF(\pi(s),\mathrm{pos}(z_i))\) equals \(N^{-1}\sum_io(s,z_i)\), the response \(\psi_t\) equals \(r\) because \(t\ge0\), and each particle drift \(b(\pi(s),z_i)\) equals \(\beta(s,z_i)\). The Hilbert space \(H_N\) has the same vectors as \(S\times Z^N\), only another norm, so the identity of vectors holds in either description. \(\square\)

\begin{leanbox}{Lean: \texttt{physical\_\allowbreak{}drift\_\allowbreak{}of\_\allowbreak{}compatibility}}
\leansource{\texttt{physical\_\allowbreak{}drift\_\allowbreak{}of\_\allowbreak{}compatibility} \quad PBGSO/\allowbreak{}ImplementationMetric.lean, lines 28–36}
\begin{leancode}
\leanline{0}{theorem\ physical\_drift\_of\_compatibility\ (N\ :\ ℕ)\ (L\ :\ LipschitzDynamics\ S\ Z\ Θ)}
\leanline{4}{(M\ :\ Model\ S\ Θ\ Z\ X\ I)\ (H\ :\ AdditiveModelCompatibility\ L\ M)}
\leanline{4}{(t\ :\ ℝ)\ (ht\ :\ 0≤t)\ (p\ :\ S\ ×\ (Fin\ N\ →\ Z))\ :}
\leanline{4}{(oneEquiv\ N).symm\ (L.particleField\ N\ p)=}
\leanline{6}{RawEmpiricalDynamics.drift\ M\ N\ t\ ((oneEquiv\ N).symm\ p)}
\end{leancode}
\leanlabel{How the Lean proof follows the argument}
The Hilbert space \(H_N\) is \texttt{One\allowbreak{}Hilbert N S Z}, and \texttt{one\allowbreak{}Equiv N} is the identification of \(H_N\) with \texttt{S \ensuremath{\times} (Fin N \ensuremath{\to} Z)}, a continuous linear equivalence that changes only the norm. \texttt{L : Lipschitz\allowbreak{}Dynamics S Z Θ} collects the primitive maps (with their Lipschitz constants), \texttt{M : Model S Θ Z X I} is the paper's model, and \texttt{H : Additive\allowbreak{}Model\allowbreak{}Compatibility L M} is the compatibility. \texttt{Raw\allowbreak{}Empirical\allowbreak{}Dynamics.\allowbreak{}drift M N t} is \(b_N(t,\cdot)\) on \texttt{One\allowbreak{}Hilbert N S Z} and \texttt{L.\allowbreak{}particle\allowbreak{}Field N} is \(\Phi_N\).




The proof applies injectivity of \texttt{one\allowbreak{}Equiv N}, rewrites with \texttt{Continuous\allowbreak{}Linear\allowbreak{}Equiv.\allowbreak{}apply\_\allowbreak{}symm\_\allowbreak{}apply} and the coordinate formula \texttt{Raw\allowbreak{}Empirical\allowbreak{}Dynamics.\allowbreak{}drift\_\allowbreak{}coordinates}, and concludes with \texttt{additive\_\allowbreak{}particle\_\allowbreak{}field}, which is the substitution of the compatibility identities.
\end{leanbox}
\subsection{\texorpdfstring{The kinetic and higher-order drifts in random coordinates}{The kinetic and higher-order drifts in random coordinates}}


We use the force data of N1 with forces \(F,F_x\), the lifted forces \(\bar F,\mathcal F_x\) on random arguments, and \(\mathbf F(\theta,X)=(\bar F(\theta,X),\mathcal F_x(\theta,X))\). The paper's two models are as follows.



\begin{itemize}


\item \emph{Kinetic model} (M-MFPO): \(S=\Theta\times\Theta\ni(\theta,m)\), \(Z=\mathcal X\times\mathcal X\ni(x,v)\), \(\pi(\theta,m)=\theta\), \(\mathrm{pos}(x,v)=x\), \(\psi_t((\theta,m),g)=(m,-g-\gamma m)\), \(b(\theta,(x,v))=(v,-F_x(\theta,x)-\gamma v)\), with constant diffusion columns.

\item \emph{Higher-order model} (HOL-MYPGD): \(S=\Theta\), \(Z=\mathcal X\times\mathcal X\times\mathcal X\ni(x,u,v)\), \(\pi=\mathrm{id}\), \(\mathrm{pos}(x,u,v)=x\), \(\psi_t(\theta,g)=-g\), \(b(\theta,(x,u,v))=(u,-F_x(\theta,x)+\gamma v,-\gamma u-\alpha v)\), with constant diffusion columns acting on the last coordinate only.

\end{itemize}



On random states the corresponding drifts are


\[\begin{aligned}b_M\bigl((m,V),(\theta,X)\bigr)&=\Bigl(-\mathbf F(\theta,X)-\gamma(m,V),\;(m,V)\Bigr),\\ b_H(\theta,X,U,V)&=\bigl(-\bar F(\theta,X),\;U,\;-\mathcal F_x(\theta,X)+\gamma V,\;-\gamma U-\alpha V\bigr).\end{aligned}\]


Let \(P_M\) (respectively \(P_H\)) be the linear map that regroups a random state into a random element of \(H_N\): at a sample it sends \(((m,V),(\theta,X))\) to the optimiser phase \((\theta,m)\) and the particle phases \(\bigl((\mathcal EX)_i,(\mathcal EV)_i\bigr)_{i=1}^N\) (respectively \((\theta,X,U,V)\) to \(\theta\) and \(\bigl((\mathcal EX)_i,(\mathcal EU)_i,(\mathcal EV)_i\bigr)_{i=1}^N\)). Both are continuous linear bijections onto \(L^2(\Omega;H_N)\) with continuous inverses.




\textbf{Claim.} For every random kinetic state \(y\), almost surely \(P_M(b_M(y))(\omega)=b_N\bigl(0,P_M(y)(\omega)\bigr)\), where \(b_N\) is the empirical drift of the kinetic model (with any constant diffusion columns). The same holds for \(P_H\), \(b_H\) and the higher-order model, for all \(\alpha,\gamma\in\mathbb R\).




\textbf{Proof.} Consider the kinetic case at a sample \(\omega\) where the samplewise identities of the lifted forces hold (a set of full probability). The parameter phase of \(P_M(b_M(y))(\omega)\) is \(\bigl(m,\,-\bar F(\theta,X)-\gamma m\bigr)\) evaluated at \(\omega\), and \(\bar F(\theta,X)(\omega)=N^{-1}\sum_iF(\theta(\omega),x_i)\) with \(x_i=(\mathcal EX(\omega))_i\); this is \(\psi_0\bigl((\theta,m),N^{-1}\sum_iF(\theta,x_i)\bigr)\). The phase of particle \(i\) is \(\bigl(v_i,\,-(\mathcal E\mathcal F_x(\theta,X)(\omega))_i-\gamma v_i\bigr)\) with \(v_i=(\mathcal EV(\omega))_i\), and \((\mathcal E\mathcal F_x(\theta,X)(\omega))_i=F_x(\theta(\omega),x_i)\); this is \(b(\theta,(x_i,v_i))\). Here linearity of \(\mathcal E\) was used to move it through \(-\gamma V\). These are exactly the components of \(b_N(0,P_M(y)(\omega))\). Formally, the primitive kinetic maps (observation \(F(\theta,x)\), response \((m,-g-\gamma m)\), sampler field \((v,-F_x(\theta,x)-\gamma v)\)) are compatible with the kinetic model, so the previous step identifies their particle field with \(b_N(0,\cdot)\). The higher-order case is the same computation with the response \(-g\) and the field \((u,-F_x(\theta,x)+\gamma v,-\gamma u-\alpha v)\). \(\square\)

\begin{leanbox}{Lean: \texttt{momentum\_\allowbreak{}drift\_\allowbreak{}ae}, \texttt{higherOrder\_\allowbreak{}drift\_\allowbreak{}ae}}
\leansource{\texttt{momentum\_\allowbreak{}drift\_\allowbreak{}ae} \quad PBGSO/\allowbreak{}ImplementationMetric.lean, lines 50–70}
\begin{leancode}
\leanline{0}{theorem\ momentum\_drift\_ae\ (γ\ :\ ℝ)\ (σ\ :\ I\ →\ HypocoercivePhase\ X)}
\leanline{4}{(z\ :\ (Lp\ A\ 2\ Q\ ×\ Lp\ C\ 2\ Q)\ ×\ (Lp\ A\ 2\ Q\ ×\ Lp\ C\ 2\ Q))\ :}
\leanline{4}{momentumPhysicalLpEquiv\ Q\ D}
\leanline{6}{(-coupledForce\ (D.meanForceLp\ Q)\ (D.samplerForceLp\ Q)\ z.2-γ\ •\ z.1,z.1)\ =ᵐ[Q]}
\leanline{6}{fun\ ω\ =>\ RawEmpiricalDynamics.drift\ (heavyBallModel\ γ\ D.F\ D.G\ σ)\ N\ 0}
\leanline{8}{(momentumPhysicalLpEquiv\ Q\ D\ z\ ω)}
\end{leancode}
\leansource{\texttt{higherOrder\_\allowbreak{}drift\_\allowbreak{}ae} \quad PBGSO/\allowbreak{}ImplementationMetric.lean, lines 73–92}
\begin{leancode}
\leanline{0}{theorem\ higherOrder\_drift\_ae\ (α\ γ\ :\ ℝ)\ (σ\ :\ I\ →\ X)}
\leanline{4}{(z\ :\ HOLState\ (Lp\ A\ 2\ Q)\ (Lp\ C\ 2\ Q))\ :}
\leanline{4}{higherOrderPhysicalLpEquiv\ Q\ D\ (holDrift\ (D.meanForceLp\ Q)\ (D.samplerForceLp\ Q)\ α\ γ\ z)\ =ᵐ[Q]}
\leanline{6}{fun\ ω\ =>\ RawEmpiricalDynamics.drift\ (higherOrderModel\ D.F\ D.G\ α\ γ\ σ)\ N\ 0}
\leanline{8}{(higherOrderPhysicalLpEquiv\ Q\ D\ z\ ω)}
\end{leancode}
\leanlabel{How the Lean proof follows the argument}
\texttt{momentum\allowbreak{}Physical\allowbreak{}Lp\allowbreak{}Equiv Q D} and \texttt{higher\allowbreak{}Order\allowbreak{}Physical\allowbreak{}Lp\allowbreak{}Equiv Q D} are \(P_M\) and \(P_H\); they are the regrouping maps \texttt{momentum\allowbreak{}Pack\allowbreak{}Lp}, \texttt{higher\allowbreak{}Order\allowbreak{}Pack\allowbreak{}Lp} composed with the identification \texttt{one\allowbreak{}Equiv}. The kinetic drift is written \texttt{(-\allowbreak{}coupled\allowbreak{}Force (D.\allowbreak{}mean\allowbreak{}Force\allowbreak{}Lp Q) (D.\allowbreak{}sampler\allowbreak{}Force\allowbreak{}Lp Q) z.\allowbreak{}2-\allowbreak{}\ensuremath{\gamma} \ensuremath{\bullet} z.\allowbreak{}1,z.\allowbreak{}1)} with \texttt{z.\allowbreak{}1} the momentum pair and \texttt{z.\allowbreak{}2} the position pair; the higher-order drift is \texttt{hol\allowbreak{}Drift (D.\allowbreak{}mean\allowbreak{}Force\allowbreak{}Lp Q) (D.\allowbreak{}sampler\allowbreak{}Force\allowbreak{}Lp Q) \ensuremath{\alpha} \ensuremath{\gamma} z}. The models are \texttt{heavy\allowbreak{}Ball\allowbreak{}Model \ensuremath{\gamma} D.\allowbreak{}F D.\allowbreak{}G \ensuremath{\sigma}} and \texttt{higher\allowbreak{}Order\allowbreak{}Model D.\allowbreak{}F D.\allowbreak{}G \ensuremath{\alpha} \ensuremath{\gamma} \ensuremath{\sigma}}.



\begin{enumerate}


\item \texttt{hcompat} is the compatibility of the primitive maps \texttt{momentum\allowbreak{}Force\allowbreak{}Dynamics D \ensuremath{\gamma}} (respectively \texttt{higher\allowbreak{}Order\allowbreak{}Force\allowbreak{}Dynamics D \ensuremath{\alpha} \ensuremath{\gamma}}) with the model, obtained from \texttt{momentum\_\allowbreak{}compatibility} (respectively \texttt{higher\allowbreak{}Order\_\allowbreak{}compatibility}), each of whose four fields holds by \texttt{rfl}.

\item \texttt{momentum\allowbreak{}Pack\allowbreak{}Lp\_\allowbreak{}drift} (respectively \texttt{higher\allowbreak{}Order\allowbreak{}Pack\allowbreak{}Lp\_\allowbreak{}drift}) rewrites the regrouped drift as \texttt{ULA.\allowbreak{}lp\allowbreak{}Drift} of the primitive particle field, that is, the particle field applied samplewise to the regrouped state.

\item \texttt{filter\_\allowbreak{}upwards} combines \texttt{coe\allowbreak{}Fn\_\allowbreak{}comp\allowbreak{}Lp\allowbreak{}L} (twice, for the identification) and \texttt{ULA.\allowbreak{}lp\allowbreak{}Drift\_\allowbreak{}coe}; the remaining pointwise identity is \texttt{physical\_\allowbreak{}drift\_\allowbreak{}of\_\allowbreak{}compatibility} at \(t=0\).

\end{enumerate}
\end{leanbox}
\subsection{\texorpdfstring{A contracting metric for the kinetic particle system}{A contracting metric for the kinetic particle system}}


A geometry for A3 consists of bounded symmetric linear operators \(W_S\) on \(S\) and \(W_Z\) on \(Z\) and numbers \(0<\lambda_-\le\lambda_+\) such that \(\lambda_-\|s\|^2\le\langle s,W_Ss\rangle\le\lambda_+\|s\|^2\) and \(\lambda_-\|z\|^2\le\langle z,W_Zz\rangle\le\lambda_+\|z\|^2\); write \(\|w\|_W^2=\langle w,Ww\rangle\). Assumption A3 with rate \(\rho\) states: for every \(t\ge0\), all \(s,s'\in S\), all probability laws \(\nu,\nu'\) on \(Z\) with finite second moments and every coupling \(\Lambda\) of \(\nu\) and \(\nu'\) for which the integrals below exist,


\[\begin{aligned}&\Bigl\langle s-s',W_S\bigl(\psi_t(s,G_\nu(s))-\psi_t(s',G_{\nu'}(s'))\bigr)\Bigr\rangle\\ &\quad+\int\Bigl[\bigl\langle z-z',W_Z\bigl(b(\pi s,z)-b(\pi s',z')\bigr)\bigr\rangle+\tfrac12\sum_j\bigl\|\sigma_j(\pi s,z)-\sigma_j(\pi s',z')\bigr\|_{W_Z}^2\Bigr]\Lambda(dz,dz')\\ &\le-\rho\Bigl(\|s-s'\|_{W_S}^2+\int\|z-z'\|_{W_Z}^2\,d\Lambda\Bigr),\end{aligned}\]


where \(G_\nu(s)=\int F(\pi s,\mathrm{pos}(z))\,\nu(dz)\). For \(N\ge1\) put \(\lambda_N^-=\lambda_-\min\{1,N^{-1}\}\) and \(\lambda_N^+=\lambda_+\max\{1,N^{-1}\}\).




We use the following result on the empirical metric, proved separately. \emph{Transported metric.} Let \(S\) and \(Z\) be finite-dimensional, let A3 hold with rate \(\rho\), let \(N\ge1\), let \(E\) be a normed space, let \(P:E\to L^2(\Omega;H_N)\) be a continuous linear bijection with continuous inverse, and let \(f:E\to E\) satisfy \(P(f(x))(\omega)=b_N(0,P(x)(\omega))\) almost surely, for every \(x\in E\). Then there are \(n\in\mathbb N\) and a bounded linear \(T:E\to L^2(\Omega;\mathbb R^n)\) such that, for all \(x,y\in E\),


\[\begin{gathered}\frac{\sqrt{\lambda_N^-}}{\|P^{-1}\|+1}\|x\|\le\|Tx\|\le\sqrt{\lambda_N^+}\,\|P\|\,\|x\|,\\ \bigl\langle T(x-y),T(f(x)-f(y))\bigr\rangle\le-\rho\|T(x-y)\|^2 .\end{gathered}\qquad(1)\]


Its proof: the weight \(\langle(s,z),W_N(s',z')\rangle=\langle s,W_Ss'\rangle+N^{-1}\sum_i\langle z_i,W_Zz'_i\rangle\) on \(H_N\) satisfies \(\lambda_N^-\|w\|^2\le\langle w,W_Nw\rangle\le\lambda_N^+\|w\|^2\); A3 applied to the empirical laws \(N^{-1}\sum_i\delta_{z_i}\), \(N^{-1}\sum_i\delta_{z'_i}\) and the coupling \(N^{-1}\sum_i\delta_{(z_i,z'_i)}\) gives \(\langle w-w',W_N(b_N(t,w)-b_N(t,w'))\rangle\le-\rho\langle w-w',W_N(w-w')\rangle\) after the nonnegative diffusion term is discarded; a finite-dimensional symmetric nonnegative operator factorises as \(\|Cw\|^2=\langle w,W_Nw\rangle\) with \(C:H_N\to\mathbb R^n\) linear; integrating the pointwise inequalities over \(\Omega\) and composing with \(P\) gives (1), where \(\|x\|\le\|P^{-1}\|\|Px\|\le(\|P^{-1}\|+1)\|Px\|\).




\textbf{Claim.} Let \(\Theta\) and \(\mathcal X\) be finite-dimensional real inner product spaces, let the force data satisfy N1, let \(\gamma\in\mathbb R\) and let the kinetic model with constant diffusion columns satisfy A3 with rate \(\rho\) for a geometry with constants \(\lambda_\pm\). Then there are \(n\in\mathbb N\) and a bounded linear map \(T\) from the random kinetic states into \(L^2(\Omega;\mathbb R^n)\) such that for all random kinetic states \(y,y'\)


\[\begin{gathered}\frac{\sqrt{\lambda_N^-}}{\|P_M^{-1}\|+1}\|y\|\le\|Ty\|\le\sqrt{\lambda_N^+}\,\|P_M\|\,\|y\|,\\ \bigl\langle T(y-y'),T(b_M(y)-b_M(y'))\bigr\rangle\le-\rho\|T(y-y')\|^2 .\end{gathered}\]


\textbf{Proof.} Apply the transported metric result with \(E\) the space of random kinetic states, \(P=P_M\) and \(f=b_M\). The spaces \(S=\Theta\times\Theta\) and \(Z=\mathcal X\times\mathcal X\) are finite-dimensional, and the hypothesis \(P_M(b_M(y))=b_N(0,P_M(y))\) almost surely is the previous step. \(\square\)




No sign condition on \(\gamma\) or \(\rho\) is needed here; positivity of \(\rho\) enters only when the dissipativity is turned into contraction. With the standard Euclidean encoding one has \(\|P_M\|\le2\) and \(\|P_M^{-1}\|\le1\), because \(\|P_My\|_{L^2}^2\) is the sum of the four squared \(L^2\) norms of the components; the constants then become \(m_N=\sqrt{\lambda_N^-}/2\) and \(U_N=2\sqrt{\lambda_N^+}\).

\begin{leanbox}{Lean: \texttt{momentum\_\allowbreak{}metric}}
\leansource{\texttt{momentum\_\allowbreak{}metric} \quad PBGSO/\allowbreak{}ImplementationMetric.lean, lines 96–111}
\begin{leancode}
\leanline{0}{theorem\ momentum\_metric\ [FiniteDimensional\ ℝ\ A]\ [FiniteDimensional\ ℝ\ X]}
\leanline{4}{(γ\ :\ ℝ)\ (σ\ :\ I\ →\ HypocoercivePhase\ X)}
\leanline{4}{(G\ :\ TwistedGeometry\ (HypocoercivePhase\ A)\ (HypocoercivePhase\ X))\ (ρ\ :\ ℝ)}
\leanline{4}{(hmono\ :\ TwistedJointMonotonicity\ (heavyBallModel\ γ\ D.F\ D.G\ σ)\ G\ ρ)\ :}
\leanline{4}{∃\ n\ :\ ℕ,\ ∃\ T\ :\ ((Lp\ A\ 2\ Q\ ×\ Lp\ C\ 2\ Q)\ ×\ (Lp\ A\ 2\ Q\ ×\ Lp\ C\ 2\ Q))\ →L[ℝ]}
\leanline{8}{Lp\ (EuclideanSpace\ ℝ\ (Fin\ n))\ 2\ Q,}
\leanline{6}{(∀\ x,\ (Real.sqrt\ (EmpiricalMetric.lower\ N\ G)/}
\leanline{8}{(‖(momentumPhysicalLpEquiv\ Q\ D).symm.toContinuousLinearMap‖+1))*‖x‖≤‖T\ x‖)\ ∧}
\leanline{6}{(∀\ x,\ ‖T\ x‖≤(Real.sqrt\ (EmpiricalMetric.upper\ N\ G)*}
\leanline{8}{‖(momentumPhysicalLpEquiv\ Q\ D).toContinuousLinearMap‖)*‖x‖)\ ∧}
\leanline{6}{(∀\ x\ y,\ {\leanfallback ⟪}T\ (x-y),T\ ((-coupledForce\ (D.meanForceLp\ Q)\ (D.samplerForceLp\ Q)\ x.2-γ\ •\ x.1,x.1)-}
\leanline{8}{(-coupledForce\ (D.meanForceLp\ Q)\ (D.samplerForceLp\ Q)\ y.2-γ\ •\ y.1,y.1)){\leanfallback ⟫}\_ℝ\ ≤\ -ρ*‖T\ (x-y)‖\textasciicircum{}2)}
\end{leancode}
\leanlabel{How the Lean proof follows the argument}
The hypotheses are the section variables (probability measure \texttt{Q}, force data \texttt{D}) together with \texttt{[Finite\allowbreak{}Dimensional \ensuremath{\mathbb R} A]}, \texttt{[Finite\allowbreak{}Dimensional \ensuremath{\mathbb R} X]}, the friction \texttt{\ensuremath{\gamma}}, constant columns \texttt{\ensuremath{\sigma} : I \ensuremath{\to} Hypocoercive\allowbreak{}Phase X}, a geometry \texttt{G : Twisted\allowbreak{}Geometry (Hypocoercive\allowbreak{}Phase A) (Hypocoercive\allowbreak{}Phase X)} and \texttt{hmono : Twisted\allowbreak{}Joint\allowbreak{}Monotonicity (heavy\allowbreak{}Ball\allowbreak{}Model \ensuremath{\gamma} D.\allowbreak{}F D.\allowbreak{}G \ensuremath{\sigma}) G \ensuremath{\rho}}, which is A3. \texttt{Hypocoercive\allowbreak{}Phase X} is \(\mathcal X\times\mathcal X\) with the Euclidean product norm. The structure \texttt{Twisted\allowbreak{}Geometry} carries \(W_S\), \(W_Z\), \(\lambda_\pm\) and the symmetry and spectral bounds; \texttt{Empirical\allowbreak{}Metric.\allowbreak{}lower N G} and \texttt{Empirical\allowbreak{}Metric.\allowbreak{}upper N G} are \(\lambda_N^-\) and \(\lambda_N^+\).




The conclusion is an existential statement: \texttt{\ensuremath{\exists} n : \ensuremath{\mathbb N}, \ensuremath{\exists} T : … \ensuremath{\to}L[\ensuremath{\mathbb R}] Lp (Euclidean\allowbreak{}Space \ensuremath{\mathbb R} (Fin n)) 2 Q}, followed by the two norm comparisons and the dissipativity, written with the real inner product \texttt{{\leanfallback ⟪}·,·{\leanfallback ⟫}\_\allowbreak{}\ensuremath{\mathbb R}} of \texttt{Lp}. The operator norms \texttt{\ensuremath{\Vert}(momentum\allowbreak{}Physical\allowbreak{}Lp\allowbreak{}Equiv Q D).\allowbreak{}symm.\allowbreak{}to\allowbreak{}Continuous\allowbreak{}Linear\allowbreak{}Map\ensuremath{\Vert}} and \texttt{\ensuremath{\Vert}(momentum\allowbreak{}Physical\allowbreak{}Lp\allowbreak{}Equiv Q D).\allowbreak{}to\allowbreak{}Continuous\allowbreak{}Linear\allowbreak{}Map\ensuremath{\Vert}} are \(\|P_M^{-1}\|\) and \(\|P_M\|\). The proof is a single application of \texttt{Empirical\allowbreak{}Metric.\allowbreak{}exists\_\allowbreak{}transported\_\allowbreak{}metric} (the transported metric result) with the equivalence \texttt{momentum\allowbreak{}Physical\allowbreak{}Lp\allowbreak{}Equiv Q D}, the model, \texttt{hmono}, \texttt{D.\allowbreak{}nonempty}, the time \texttt{0} with \texttt{le\_\allowbreak{}rfl}, the drift as a function of the state, and the drift identity \texttt{momentum\_\allowbreak{}drift\_\allowbreak{}ae Q D \ensuremath{\gamma} \ensuremath{\sigma}}. The explicit constants for the Euclidean encoding are derived from this theorem in the file on metric comparisons (\texttt{momentum\_\allowbreak{}metric\_\allowbreak{}explicit}).
\end{leanbox}
\subsection{\texorpdfstring{A contracting metric for the higher-order particle system}{A contracting metric for the higher-order particle system}}


\textbf{Claim.} Let \(\Theta\) and \(\mathcal X\) be finite-dimensional real inner product spaces, let the force data satisfy N1, let \(\alpha,\gamma\in\mathbb R\), and let the higher-order model with constant diffusion columns on the last coordinate satisfy A3 with rate \(\rho\) for a geometry on \(S=\Theta\), \(Z=\mathcal X^3\) with constants \(\lambda_\pm\). Then there are \(n\in\mathbb N\) and a bounded linear map \(T\) from the random higher-order states \((\theta,X,U,V)\) into \(L^2(\Omega;\mathbb R^n)\) such that for all random states \(y,y'\)


\[\begin{gathered}\frac{\sqrt{\lambda_N^-}}{\|P_H^{-1}\|+1}\|y\|\le\|Ty\|\le\sqrt{\lambda_N^+}\,\|P_H\|\,\|y\|,\\ \bigl\langle T(y-y'),T(b_H(y)-b_H(y'))\bigr\rangle\le-\rho\|T(y-y')\|^2 .\end{gathered}\]


\textbf{Proof.} Apply the transported metric result of the previous step with \(P=P_H\) and \(f=b_H\); the drift identity is the second step of this section. \(\square\)




The weight acts on the parameter and on the three sampler coordinates of each particle, so the metric couples \(\theta\), \(X\), \(U\) and \(V\); this coupling is what allows dissipation although the drift of \(X\) is \(U\) and contains no damping. Under the standard Euclidean encoding \(\|P_H\|\le2\) and \(\|P_H^{-1}\|\le1\), giving again \(m_N=\sqrt{\lambda_N^-}/2\) and \(U_N=2\sqrt{\lambda_N^+}\).

\begin{leanbox}{Lean: \texttt{higherOrder\_\allowbreak{}metric}}
\leansource{\texttt{higherOrder\_\allowbreak{}metric} \quad PBGSO/\allowbreak{}ImplementationMetric.lean, lines 114–129}
\begin{leancode}
\leanline{0}{theorem\ higherOrder\_metric\ [FiniteDimensional\ ℝ\ A]\ [FiniteDimensional\ ℝ\ X]}
\leanline{4}{(α\ γ\ :\ ℝ)\ (σ\ :\ I\ →\ X)}
\leanline{4}{(G\ :\ TwistedGeometry\ A\ (HigherOrderSampler\ X))\ (ρ\ :\ ℝ)}
\leanline{4}{(hmono\ :\ TwistedJointMonotonicity\ (higherOrderModel\ D.F\ D.G\ α\ γ\ σ)\ G\ ρ)\ :}
\leanline{4}{∃\ n\ :\ ℕ,\ ∃\ T\ :\ HOLState\ (Lp\ A\ 2\ Q)\ (Lp\ C\ 2\ Q)\ →L[ℝ]}
\leanline{8}{Lp\ (EuclideanSpace\ ℝ\ (Fin\ n))\ 2\ Q,}
\leanline{6}{(∀\ x,\ (Real.sqrt\ (EmpiricalMetric.lower\ N\ G)/}
\leanline{8}{(‖(higherOrderPhysicalLpEquiv\ Q\ D).symm.toContinuousLinearMap‖+1))*‖x‖≤‖T\ x‖)\ ∧}
\leanline{6}{(∀\ x,\ ‖T\ x‖≤(Real.sqrt\ (EmpiricalMetric.upper\ N\ G)*}
\leanline{8}{‖(higherOrderPhysicalLpEquiv\ Q\ D).toContinuousLinearMap‖)*‖x‖)\ ∧}
\leanline{6}{(∀\ x\ y,\ {\leanfallback ⟪}T\ (x-y),T\ (holDrift\ (D.meanForceLp\ Q)\ (D.samplerForceLp\ Q)\ α\ γ\ x-}
\leanline{8}{holDrift\ (D.meanForceLp\ Q)\ (D.samplerForceLp\ Q)\ α\ γ\ y){\leanfallback ⟫}\_ℝ\ ≤\ -ρ*‖T\ (x-y)‖\textasciicircum{}2)}
\end{leancode}
\leanlabel{How the Lean proof follows the argument}
The statement is the same as for the kinetic system with \texttt{HOLState (Lp A 2 Q) (Lp C 2 Q)} as the state space, \texttt{G : Twisted\allowbreak{}Geometry A (Higher\allowbreak{}Order\allowbreak{}Sampler X)} (where \texttt{Higher\allowbreak{}Order\allowbreak{}Sampler X} is \(\mathcal X\times(\mathcal X\times\mathcal X)\) with Euclidean product norms), \texttt{hmono : Twisted\allowbreak{}Joint\allowbreak{}Monotonicity (higher\allowbreak{}Order\allowbreak{}Model D.\allowbreak{}F D.\allowbreak{}G \ensuremath{\alpha} \ensuremath{\gamma} \ensuremath{\sigma}) G \ensuremath{\rho}}, and the drift \texttt{hol\allowbreak{}Drift (D.\allowbreak{}mean\allowbreak{}Force\allowbreak{}Lp Q) (D.\allowbreak{}sampler\allowbreak{}Force\allowbreak{}Lp Q) \ensuremath{\alpha} \ensuremath{\gamma}}. The proof applies \texttt{Empirical\allowbreak{}Metric.\allowbreak{}exists\_\allowbreak{}transported\_\allowbreak{}metric} with \texttt{higher\allowbreak{}Order\allowbreak{}Physical\allowbreak{}Lp\allowbreak{}Equiv Q D} and \texttt{higher\allowbreak{}Order\_\allowbreak{}drift\_\allowbreak{}ae Q D \ensuremath{\alpha} \ensuremath{\gamma} \ensuremath{\sigma}}. The explicit version for the Euclidean encoding is \texttt{higher\allowbreak{}Order\_\allowbreak{}metric\_\allowbreak{}explicit} in the file on metric comparisons.
\end{leanbox}
\section{\texorpdfstring{The local error of the OBABO step against the kinetic integral equation}{The local error of the OBABO step against the kinetic integral equation}}


This section bounds the error made by one OBABO step of length \(h\in(0,1]\) relative to the exact kinetic dynamics over the same time interval, when both are driven by the same Brownian increments. The state is a pair \(z=(p,q)\in E\times E\) of a momentum and a position in a real Banach space \(E\), normed by \(\max\{\|p\|,\|q\|\}\); the force \(G:E\to E\) is Lipschitz with constant \(L\), the friction is \(\gamma\ge0\), and the exact dynamics is the integral equation


\[Y(t)-Y(0)=\int_0^tf(Y(r))\,dr+\bigl(b(t),0\bigr),\qquad f(p,q)=\bigl(-G(q)-\gamma p,\;p\bigr),\]


with an additive input \(b\) in the momentum. In M-MFPO, \(E\) is the space of square-integrable random pairs (parameter, cloud), \(G\) is the coupled force \((\bar F,\mathcal F_x)\) and \(b(t)=\sqrt{2\gamma}\) times the Brownian motion of the particle velocities; the parameter momentum receives no noise.




The main result is that the sum of the momentum and position errors after one step is at most \((c_{0,M}+c_{1,M}\|Y(0)\|)\,h\sqrt h\), where \(c_{0,M}\) and \(c_{1,M}\) are explicit and depend only on \(L\), \(\gamma\), \(\|G(0)\|\), the bound \(\Xi\) of the two half-step innovations and the Brownian increment constant \(\beta\). All moments and moduli of the exact solution that enter are derived from the integral equation on \([0,h]\); in particular the bound is affine in the initial norm and does not depend on the starting time. This affine structure is what the absorption argument of Theorem New A requires.

\subsection{\texorpdfstring{The coefficients of the affine local error}{The coefficients of the affine local error}}


\textbf{Definition.} For \(L,\gamma\ge0\) and real \(f_0,\Xi,\beta\) put \(K=L+\gamma+1\) and


\[\begin{aligned}a&=\frac{\gamma^2+\gamma+L}{2},\qquad b=\frac{\gamma+1}{2}+\frac L4,\qquad \nu=\frac L2+\sqrt\gamma+\frac32\gamma\sqrt\gamma,\\ c_0&=b\,f_0+\nu\,\Xi+\frac23K\Bigl(f_0+Ke^{K+1}(f_0+\beta)+\beta\Bigr),\\ c_1&=a+b\,L+\frac23K^2e^{K+1}.\end{aligned}\]


With \(L=L_M\), \(f_0=b_\circ\), \(\Xi=\sqrt d\) and \(\beta=\sqrt\gamma\sqrt{2d}\) these are the coefficients \(c_{0,M}\) and \(c_{1,M}\) of Theorem New A; in that notation \(K=K_M\) and \(f_0+Ke^{K+1}(f_0+\beta)+\beta=T_M\). The coefficient \(c_1\) depends only on \(L\) and \(\gamma\).




The coefficients arise as follows. The OBABO local error (cited in the third step) is bounded by \(D\,h\sqrt h\) with


\[D(P,F,R)=aP+bF+\nu\Xi+\frac23KR,\qquad(1)\]


where \(P\) bounds the initial momentum, \(F\) bounds the initial force and \(R\) is a square-root modulus of the exact solution. If \(P=p_0+p_1u\), \(F=f_0+f_1u\) and \(R=r_0+r_1u\) are affine in a number \(u\), then \(D=(ap_0+bf_0+\nu\Xi+\tfrac23Kr_0)+(ap_1+bf_1+\tfrac23Kr_1)u\). The coefficients \(c_0,c_1\) are this decomposition for \(p_0=0\), \(p_1=1\), \(f_1=L\), \(r_0=f_0+Ke^{K+1}(f_0+\beta)+\beta\), \(r_1=Ke^{K+1}\), with \(u=\|Y(0)\|\).

\begin{leanbox}{Lean: \texttt{momentumUniformDefect}}
\leansource{\texttt{momentumUniformDefect} \quad PBGSO/\allowbreak{}MomentumUniformDefect.lean, lines 15–19}
\begin{leancode}
\leanline{0}{def\ momentumUniformDefect\ (L\ γ\ :\ ℝ≥0)\ (fzero\ Ξ\ β\ :\ ℝ)\ :\ ℝ\ ×\ ℝ\ :=}
\leanline{2}{let\ K\ :\ ℝ≥0\ :=\ L\ +\ γ\ +\ 1}
\leanline{2}{let\ c\ :=\ Real.exp\ ((K\ :\ ℝ)\ +\ 1)}
\leanline{2}{momentumDefectAffine\ L\ γ\ Ξ\ 0\ 1\ fzero\ L}
\leanline{4}{(fzero\ +\ (K\ :\ ℝ)\ *\ c\ *\ (fzero\ +\ β)\ +\ β)\ ((K\ :\ ℝ)\ *\ c)}
\end{leancode}
\leanlabel{How the Lean proof follows the argument}
\texttt{momentum\allowbreak{}Uniform\allowbreak{}Defect L \ensuremath{\gamma} fzero Ξ \ensuremath{\beta}} returns the pair \((c_0,c_1)\). It is defined as \texttt{momentum\allowbreak{}Defect\allowbreak{}Affine L \ensuremath{\gamma} Ξ 0 1 fzero L (fzero + (K : \ensuremath{\mathbb R}) * c * (fzero + \ensuremath{\beta}) + \ensuremath{\beta}) ((K : \ensuremath{\mathbb R}) * c)} with \texttt{K = L + \ensuremath{\gamma} + 1} and \texttt{c = Real.\allowbreak{}exp ((K : \ensuremath{\mathbb R}) + 1)}; \texttt{momentum\allowbreak{}Defect\allowbreak{}Affine}, defined in the file on the OBABO local estimates, is the affine decomposition of (1) with the coefficients \texttt{a}, \texttt{b}, \texttt{c} (there \texttt{c} denotes \(\tfrac23K\)) and the noise coefficient \(\nu\). The arguments \texttt{L \ensuremath{\gamma} : \ensuremath{\mathbb R}\ensuremath{\ge}0} are nonnegative by type. The closed formula is also certified by \texttt{momentum\allowbreak{}Uniform\allowbreak{}Defect\_\allowbreak{}expanded} in the file on constant expansions.
\end{leanbox}
\subsection{\texorpdfstring{Component form of the kinetic integral equation}{Component form of the kinetic integral equation}}


\textbf{Claim.} Let \(E\) be a real Banach space, \(G:E\to E\) Lipschitz, \(\gamma\ge0\), \(f(p,q)=(-G(q)-\gamma p,p)\), \(Y=(p,q):\mathbb R\to E\times E\) continuous, \(b:\mathbb R\to E\) and \(h\in\mathbb R\). If \(Y(h)-Y(0)=\int_0^hf(Y(r))\,dr+(b(h),0)\), then


\[p(h)=p(0)-\int_0^h\bigl(G(q(r))+\gamma p(r)\bigr)\,dr+b(h),\qquad q(h)=q(0)+\int_0^hp(r)\,dr .\]


\textbf{Proof.} The map \(f\) is Lipschitz (with constant \(L+\gamma+1\), see the next step), so \(r\mapsto f(Y(r))\) is continuous and hence integrable on the interval between \(0\) and \(h\). The coordinate projections \((p,q)\mapsto p\) and \((p,q)\mapsto q\) are bounded linear maps, and a bounded linear map commutes with the integral of an integrable function. Hence the first component of \(\int_0^hf(Y)\) is \(\int_0^h(-G(q)-\gamma p)=-\int_0^h(G(q)+\gamma p)\) and the second is \(\int_0^hp\). Taking the two components of the equation and rearranging gives the claim. \(\square\)

\begin{leanbox}{Lean: \texttt{uniformMomentum\_\allowbreak{}component\_\allowbreak{}equations}}
\leansource{\texttt{uniformMomentum\_\allowbreak{}component\_\allowbreak{}equations} \quad PBGSO/\allowbreak{}MomentumUniformDefect.lean, lines 23–61}
\begin{leancode}
\leanline{0}{theorem\ uniformMomentum\_component\_equations\ (G\ :\ E\ →\ E)\ (L\ γ\ :\ ℝ≥0)}
\leanline{4}{(hLip\ :\ LipschitzWith\ L\ G)\ (Y\ :\ ℝ\ →\ E\ ×\ E)\ (b\ :\ ℝ\ →\ E)}
\leanline{4}{(hY\ :\ Continuous\ Y)\ (h\ :\ ℝ)}
\leanline{4}{(heq\ :\ Y\ h\ -\ Y\ 0\ =}
\leanline{6}{(∫\ r\ in\ 0..h,\ uniformMomentumDrift\ G\ γ\ (Y\ r))\ +\ (b\ h,\ 0))\ :}
\leanline{4}{(Y\ h).1\ =\ (Y\ 0).1\ -}
\leanline{6}{(∫\ r\ in\ 0..h,\ G\ (Y\ r).2\ +\ (γ\ :\ ℝ)\ •\ (Y\ r).1)\ +\ b\ h\ ∧}
\leanline{4}{(Y\ h).2\ =\ (Y\ 0).2\ +\ (∫\ r\ in\ 0..h,\ (Y\ r).1)}
\end{leancode}
\leanlabel{How the Lean proof follows the argument}
The kinetic drift is \texttt{uniform\allowbreak{}Momentum\allowbreak{}Drift G \ensuremath{\gamma} z = (-\allowbreak{}G z.\allowbreak{}2 -\allowbreak{} \ensuremath{\gamma} \ensuremath{\bullet} z.\allowbreak{}1, z.\allowbreak{}1)}, and its Lipschitz property is \texttt{uniform\allowbreak{}Momentum\allowbreak{}Drift\_\allowbreak{}lipschitz}. \texttt{hi} is the integrability, from continuity (\texttt{hc.\allowbreak{}interval\allowbreak{}Integrable}). \texttt{hp}, \texttt{hq} are the two components of \texttt{heq} (\texttt{congr\allowbreak{}Arg Prod.\allowbreak{}fst}, \texttt{congr\allowbreak{}Arg Prod.\allowbreak{}snd}), and \texttt{hip}, \texttt{hiq} are the commutation of the projections \texttt{Continuous\allowbreak{}Linear\allowbreak{}Map.\allowbreak{}fst}, \texttt{Continuous\allowbreak{}Linear\allowbreak{}Map.\allowbreak{}snd} with the interval integral (\texttt{interval\allowbreak{}Integral\_\allowbreak{}comp\_\allowbreak{}comm}). \texttt{hm} rewrites the first component of the drift, and \texttt{interval\allowbreak{}Integral.\allowbreak{}integral\_\allowbreak{}neg} takes the sign out of the integral; \texttt{module} performs the final rearrangements.
\end{leanbox}
\subsection{\texorpdfstring{The OBABO local error is affine in the initial norm}{The OBABO local error is affine in the initial norm}}


For \(h>0\), \(\eta=e^{-\gamma h/2}\), \(\sigma=\sqrt{1-e^{-\gamma h}}\) and \(\xi,\xi'\in E\), the OBABO step \(\Psi(p,q;\xi,\xi')\) is the composition of the maps \(\mathrm O_\xi(p,q)=(\eta p+\sigma\xi,q)\), \(\mathrm B(p,q)=(p-\tfrac h2G(q),q)\), \(\mathrm A(p,q)=(p,q+hp)\), \(\mathrm B\) and \(\mathrm O_{\xi'}\), in this order. The phase error of two pairs is \(\mathrm e(z,w)=\|z_1-w_1\|+\|z_2-w_2\|\).




We use the following one-step estimate, proved in the file on the OBABO local estimates. \emph{OBABO local error.} Let \(E\) be a real Banach space, \(G:E\to E\) with \(\|G(x)-G(y)\|\le L\|x-y\|\), \(L,\gamma\ge0\), \(0<h\le1\), and \(P,F,\Xi,R\ge0\). Let \(p,q:\mathbb R\to E\) be integrable on \([0,h]\) together with \(G(q)+\gamma p\), with


\[p(h)=p(0)-\int_0^h\bigl(G(q)+\gamma p\bigr)+b,\qquad q(h)=q(0)+\int_0^hp,\]


\(\|p(t)-p(0)\|\le R\sqrt t\) and \(\|q(t)-q(0)\|\le R\sqrt t\) for \(t\in[0,h]\), \(\|p(0)\|\le P\), \(\|G(q(0))\|\le F\), \(\|\xi\|,\|\xi'\|\le\Xi\) and \(b=\sqrt{\gamma h}\,(\xi+\xi')\). Then \(\mathrm e\bigl(\Psi(p(0),q(0);\xi,\xi'),(p(h),q(h))\bigr)\le D(P,F,R)\,h\sqrt h\), with \(D\) as in (1). Its proof compares three maps: the OBABO step differs from the Euler step \((p_0-hG(q_0)-\gamma hp_0+\tilde b,\;q_0+hp_0)\) with the combined innovation \(\tilde b=\sigma(\eta\xi+\xi')\) by at most \((aP+bF+\tfrac L2\Xi)h^2+\sqrt\gamma\,\Xi\,h\sqrt h\), using \(0\le e^{-x}-1+x\le x^2/2\); replacing \(\tilde b\) by \(b\) costs \(\|\tilde b-b\|\le\tfrac32\gamma\sqrt\gamma\,\Xi\,h\sqrt h\); and the Euler step with input \(b\) differs from \((p(h),q(h))\) by at most \(\tfrac23(L+\gamma+1)R\,h\sqrt h\), because the integrands deviate from their initial values by at most \((L+\gamma)R\sqrt t\) and \(R\sqrt t\) and \(\int_0^h\sqrt t\,dt=\tfrac23h\sqrt h\). Finally \(h^2\le h\sqrt h\).




\textbf{Claim.} Let \(E\) be a real Banach space, \(G:E\to E\) Lipschitz with constant \(L\ge0\), \(\gamma\ge0\), \(f(p,q)=(-G(q)-\gamma p,p)\). Let \(Y=(p,q):\mathbb R\to E\times E\) and \(b:\mathbb R\to E\) be continuous, \(0<h\le1\), \(\Xi\ge0\), \(\beta\ge0\), and assume



\begin{enumerate}


\item \(Y(t)-Y(0)=\int_0^tf(Y(r))\,dr+(b(t),0)\) for \(t\in[0,h]\);

\item \(\|b(t)\|\le\beta\sqrt t\) for \(t\in[0,h]\);

\item \(\xi,\xi'\in E\) satisfy \(\|\xi\|\le\Xi\), \(\|\xi'\|\le\Xi\) and \(b(h)=\sqrt{\gamma h}\,(\xi+\xi')\).

\end{enumerate}



Then, with \((c_0,c_1)\) the coefficients of the first step for \(f_0=\|G(0)\|\),


\[\mathrm e\bigl(\Psi(Y(0);\xi,\xi'),\,Y(h)\bigr)\le\bigl(c_0+c_1\|Y(0)\|\bigr)\,h\sqrt h .\]


\textbf{Proof.} Put \(u=\|Y(0)\|\) and \(K=L+\gamma+1\).




\emph{Regularity of the exact solution.} The drift \(f\) is Lipschitz with constant \(K\): for \(z=(p,q)\), \(w=(p',q')\) the first component of \(f(z)-f(w)\) has norm at most \(\|G(q)-G(q')\|+\gamma\|p-p'\|\le(L+\gamma)\|z-w\|\), and the second is \(\|p-p'\|\le\|z-w\|\). Moreover \(\|f(0)\|=\|(-G(0),0)\|=\|G(0)\|\). Hypotheses 1 and 2 are those of the short-interval regularity result for the integral equation (with input \((b,0)\), whose norm is \(\|b(t)\|\le\beta\sqrt t\), and \(h\le1\)). It gives, for \(t\in[0,h]\),


\[\|Y(t)-Y(0)\|\le R\sqrt t,\qquad R=\|G(0)\|+K\,e^{K+1}\bigl(u+\|G(0)\|+\beta\bigr)+\beta .\]


Each component of a pair has norm at most the norm of the pair, so \(\|p(t)-p(0)\|\le R\sqrt t\) and \(\|q(t)-q(0)\|\le R\sqrt t\).




\emph{Initial bounds.} \(\|p(0)\|\le u\), and by the linear growth of a Lipschitz map, \(\|G(q(0))\|\le\|G(0)\|+L\|q(0)\|\le\|G(0)\|+Lu\).




\emph{Equations at time \(h\).} By the previous step, hypothesis 1 at \(t=h\) gives the two component equations required by the OBABO local error, with \(b=b(h)\). The integrability conditions hold because the integrands are continuous.




\emph{Application.} The OBABO local error with \(P=u\), \(F=\|G(0)\|+Lu\), the modulus \(R\) and hypothesis 3 gives \(\mathrm e\le D(u,\|G(0)\|+Lu,R)\,h\sqrt h\). Expanding (1),


\[\begin{aligned}D&=au+b\bigl(\|G(0)\|+Lu\bigr)+\nu\Xi+\frac23K\Bigl(\|G(0)\|+Ke^{K+1}(\|G(0)\|+\beta)+\beta\Bigr)+\frac23K^2e^{K+1}u\\ &=c_0+c_1u. \qquad\square\end{aligned}\]


The coupling hypothesis 3 is the only link between the numerical innovations and the exact dynamics. In M-MFPO, \(\xi\) and \(\xi'\) are the normalised Brownian increments over the two half steps, \(\xi=\sqrt{2/h}\,(W(h/2)-W(0))\) and \(\xi'=\sqrt{2/h}\,(W(h)-W(h/2))\), so that \(\sqrt{\gamma h}(\xi+\xi')=\sqrt{2\gamma}\,(W(h)-W(0))\).

\begin{leanbox}{Lean: \texttt{obabo\_\allowbreak{}affine\_\allowbreak{}local\_\allowbreak{}defect}}
\leansource{\texttt{obabo\_\allowbreak{}affine\_\allowbreak{}local\_\allowbreak{}defect} \quad PBGSO/\allowbreak{}MomentumUniformDefect.lean, lines 66–119}
\begin{leancode}
\leanline{0}{theorem\ obabo\_affine\_local\_defect\ (G\ :\ E\ →\ E)\ (L\ γ\ :\ ℝ≥0)}
\leanline{4}{(hLip\ :\ LipschitzWith\ L\ G)\ (Y\ :\ ℝ\ →\ E\ ×\ E)\ (b\ :\ ℝ\ →\ E)}
\leanline{4}{(hY\ :\ Continuous\ Y)\ (hb\ :\ Continuous\ b)}
\leanline{4}{(Ξ\ β\ h\ :\ ℝ)\ (hΞ\ :\ 0\ ≤\ Ξ)\ (hβ\ :\ 0\ ≤\ β)\ (hh\ :\ 0\ <\ h)\ (hh1\ :\ h\ ≤\ 1)}
\leanline{4}{(heq\ :\ ∀\ t\ ∈\ Icc\ 0\ h,\ Y\ t\ -\ Y\ 0\ =}
\leanline{6}{(∫\ r\ in\ 0..t,\ uniformMomentumDrift\ G\ γ\ (Y\ r))\ +\ (b\ t,\ 0))}
\leanline{4}{(hinput\ :\ ∀\ t\ ∈\ Icc\ 0\ h,\ ‖b\ t‖\ ≤\ β\ *\ Real.sqrt\ t)}
\leanline{4}{(ξ\ ξ'\ :\ E)\ (hξ\ :\ ‖ξ‖\ ≤\ Ξ)\ (hξ'\ :\ ‖ξ'‖\ ≤\ Ξ)}
\leanline{4}{(hbrown\ :\ b\ h\ =\ Real.sqrt\ ((γ\ :\ ℝ)\ *\ h)\ •\ (ξ\ +\ ξ'))\ :}
\leanline{4}{phaseError}
\leanline{6}{(KineticLangevin.obaboStep\ G\ h\ (Real.exp\ (-(γ\ :\ ℝ)\ *\ h\ /\ 2))}
\leanline{8}{(Real.sqrt\ (1\ -\ Real.exp\ (-(γ\ :\ ℝ)\ *\ h)))\ (Y\ 0)\ ξ\ ξ')}
\leanline{6}{(Y\ h)\ ≤}
\leanline{6}{((momentumUniformDefect\ L\ γ\ ‖G\ 0‖\ Ξ\ β).1\ +}
\leanline{8}{(momentumUniformDefect\ L\ γ\ ‖G\ 0‖\ Ξ\ β).2\ *\ ‖Y\ 0‖)\ *\ h\ *\ Real.sqrt\ h}
\end{leancode}
\leanlabel{How the Lean proof follows the argument}
The conclusion is written with \texttt{phase\allowbreak{}Error} (the sum of the two component errors) and \texttt{Kinetic\allowbreak{}Langevin.\allowbreak{}obabo\allowbreak{}Step G h (Real.\allowbreak{}exp (-\allowbreak{}(\ensuremath{\gamma} : \ensuremath{\mathbb R}) * h /\allowbreak{} 2)) (Real.\allowbreak{}sqrt (1 -\allowbreak{} Real.\allowbreak{}exp (-\allowbreak{}(\ensuremath{\gamma} : \ensuremath{\mathbb R}) * h))) (Y 0) \ensuremath{\xi} \ensuremath{\xi}'}. Hypothesis 1 is \texttt{heq}, hypothesis 2 is \texttt{hinput}, and hypothesis 3 is \texttt{h\ensuremath{\xi}}, \texttt{h\ensuremath{\xi}'}, \texttt{hbrown}.



\begin{enumerate}


\item \texttt{K} and \texttt{R} are as in the proof. \texttt{hreg} is \texttt{Uniform\allowbreak{}Integral.\allowbreak{}local\_\allowbreak{}regularity} for the drift \texttt{uniform\allowbreak{}Momentum\allowbreak{}Drift G \ensuremath{\gamma}} with Lipschitz constant \texttt{K} (from \texttt{uniform\allowbreak{}Momentum\allowbreak{}Drift\_\allowbreak{}lipschitz}), input \texttt{fun t \ensuremath{\mapsto} (b t, 0)} and input coefficient \texttt{\ensuremath{\beta}}; \texttt{uniform\allowbreak{}Momentum\allowbreak{}Drift\_\allowbreak{}zero\_\allowbreak{}norm} replaces \(\|f(0)\|\) by \(\|G(0)\|\).

\item \texttt{hmod} projects the increment bound onto the two components (\texttt{norm\_\allowbreak{}fst\_\allowbreak{}le}, \texttt{norm\_\allowbreak{}snd\_\allowbreak{}le}). \texttt{hf} is the initial force bound from \texttt{Uniform\allowbreak{}Integral.\allowbreak{}drift\_\allowbreak{}bound}. \texttt{he} is \texttt{uniform\allowbreak{}Momentum\_\allowbreak{}component\_\allowbreak{}equations} at time \texttt{h}.

\item \texttt{hlocal} is \texttt{obabo\_\allowbreak{}integral\_\allowbreak{}local\_\allowbreak{}sqrt\_\allowbreak{}bound}, the OBABO local error, with \(P=\|Y(0)\|\), \(F=\|G(0)\|+L\|Y(0)\|\) and the modulus \texttt{R}.

\item \texttt{hconstant} is the algebraic identity of the last display: after unfolding \texttt{R}, \texttt{Uniform\allowbreak{}Integral.\allowbreak{}radius}, \texttt{momentum\allowbreak{}Uniform\allowbreak{}Defect}, \texttt{momentum\allowbreak{}Defect\allowbreak{}Affine}, \texttt{momentum\allowbreak{}Local\allowbreak{}Defect\allowbreak{}Constant} and \texttt{momentum\allowbreak{}Local\allowbreak{}Consistency\allowbreak{}Constant}, it is closed by \texttt{ring}.

\end{enumerate}
\end{leanbox}
\subsection{\texorpdfstring{Nonnegativity of the coefficients}{Nonnegativity of the coefficients}}


\textbf{Claim.} If \(L,\gamma\ge0\) and \(f_0,\Xi,\beta\ge0\), then \(c_0\ge0\) and \(c_1\ge0\).




\textbf{Proof.} Every quantity in the definition of \(a,b,\nu,K\) is nonnegative, \(e^{K+1}>0\), and \(c_0\), \(c_1\) are sums of products of nonnegative numbers. \(\square\)




Nonnegativity of \(c_1\) is needed in the absorption argument, where \(c_1\) multiplies an upper bound for the norm of the numerical state.

\begin{leanbox}{Lean: \texttt{momentumUniformDefect\_\allowbreak{}nonneg}}
\leansource{\texttt{momentumUniformDefect\_\allowbreak{}nonneg} \quad PBGSO/\allowbreak{}MomentumUniformDefect.lean, lines 122–130}
\begin{leancode}
\leanline{0}{theorem\ momentumUniformDefect\_nonneg\ (L\ γ\ :\ ℝ≥0)\ (fzero\ Ξ\ β\ :\ ℝ)}
\leanline{4}{(hfzero\ :\ 0\ ≤\ fzero)\ (hΞ\ :\ 0\ ≤\ Ξ)\ (hβ\ :\ 0\ ≤\ β)\ :}
\leanline{4}{0\ ≤\ (momentumUniformDefect\ L\ γ\ fzero\ Ξ\ β).1\ ∧}
\leanline{6}{0\ ≤\ (momentumUniformDefect\ L\ γ\ fzero\ Ξ\ β).2}
\end{leancode}
\leanlabel{How the Lean proof follows the argument}
The proof unfolds \texttt{momentum\allowbreak{}Uniform\allowbreak{}Defect} and applies \texttt{momentum\allowbreak{}Defect\allowbreak{}Affine\_\allowbreak{}nonneg}, which requires nonnegativity of each affine input; the two remaining inputs \(f_0+Ke^{K+1}(f_0+\beta)+\beta\) and \(Ke^{K+1}\) are shown nonnegative by \texttt{positivity}.
\end{leanbox}
\subsection{\texorpdfstring{From the phase error to a linear metric}{From the phase error to a linear metric}}


\textbf{Claim.} Let \(E\) and \(H\) be real normed spaces and \(C:E\times E\to H\) a bounded linear map, where \(E\times E\) carries the norm \(\max\{\|z_1\|,\|z_2\|\}\). Then for all \(z,w\in E\times E\)


\[\|C(z-w)\|\le\|C\|\,\bigl(\|z_1-w_1\|+\|z_2-w_2\|\bigr).\]


\textbf{Proof.} By the definition of the operator norm, \(\|C(z-w)\|\le\|C\|\|z-w\|=\|C\|\max\{\|z_1-w_1\|,\|z_2-w_2\|\}\), and a maximum of two nonnegative numbers is at most their sum. Multiply by \(\|C\|\ge0\). \(\square\)




This converts the local error of the previous steps into the metric \(\|C\cdot\|\) in which the exact flow contracts.

\begin{leanbox}{Lean: \texttt{transformed\_\allowbreak{}norm\_\allowbreak{}sub\_\allowbreak{}le\_\allowbreak{}phaseError}}
\leansource{\texttt{transformed\_\allowbreak{}norm\_\allowbreak{}sub\_\allowbreak{}le\_\allowbreak{}phaseError} \quad PBGSO/\allowbreak{}MomentumUniformDefect.lean, lines 135–143}
\begin{leancode}
\leanline{0}{theorem\ transformed\_norm\_sub\_le\_phaseError}
\leanline{4}{\{H\ :\ Type*\}\ [NormedAddCommGroup\ H]\ [NormedSpace\ ℝ\ H]}
\leanline{4}{(C\ :\ (E\ ×\ E)\ →L[ℝ]\ H)\ (z\ w\ :\ E\ ×\ E)\ :}
\leanline{4}{‖C\ (z-w)‖\ ≤\ ‖C‖\ *\ phaseError\ z\ w}
\end{leancode}
\leanlabel{How the Lean proof follows the argument}
\texttt{C.\allowbreak{}le\_\allowbreak{}op\allowbreak{}Norm} is the operator-norm inequality; \texttt{change} exposes the product norm as \texttt{max \ensuremath{\Vert}z.\allowbreak{}1-\allowbreak{}w.\allowbreak{}1\ensuremath{\Vert} \ensuremath{\Vert}z.\allowbreak{}2-\allowbreak{}w.\allowbreak{}2\ensuremath{\Vert}}, and \texttt{max\_\allowbreak{}le} with \texttt{le\_\allowbreak{}add\_\allowbreak{}of\_\allowbreak{}nonneg\_\allowbreak{}right}, \texttt{le\_\allowbreak{}add\_\allowbreak{}of\_\allowbreak{}nonneg\_\allowbreak{}left} bounds the maximum by the sum. The lemma does not need completeness of \texttt{E} (the \texttt{omit} line removes that section assumption).
\end{leanbox}
\section{\texorpdfstring{Constructed exact restarts of the higher-order dynamics}{Constructed exact restarts of the higher-order dynamics}}


Write \(z=(\theta,x,u,v)\) for the four-coordinate higher-order state. Its drift is \(f(z)=(-\bar F(\theta,x),u,-G(\theta,x)+\gamma v,-\gamma u-\alpha v)\). The parameter space is complete, while the three sampler coordinates are square-integrable Euclidean random variables on a probability space carrying the actual Brownian coordinates. The forces have nonnegative Lipschitz bounds \(L_F,L_G\) jointly in the parameter and position.

\subsection{\texorpdfstring{Restart existence, agreement and contraction}{Restart existence, agreement and contraction}}


\textbf{Claim.} Suppose the stated complete-space and global force Lipschitz assumptions hold. For each initial state \(z\) and restart time \(a\), there is a continuous solution

\[r_z(t)=z+\int_0^tf(r_z(s))ds+B_a(t),\]

where \(B_a\) is the same shifted physical Brownian increment, acting only in the last sampler coordinate. The constructed solution agrees with any original solution restarted from its value at nonnegative time \(a\). If a bounded linear map \(C\) into a real Hilbert space satisfies \(\langle C(x-y),C(f(x)-f(y))\rangle\le-\rho\|C(x-y)\|^2\) for every \(x,y\), then for every \(t\ge0\)

\[\|C(r_z(t)-r_w(t))\|\le e^{-\rho t}\|C(z-w)\|.\]

The real rate \(\rho\) need not be positive for this particular statement.



\textbf{Proof.} The force bounds make the entire four-coordinate drift globally Lipschitz. The Brownian input is continuous in the square-integrable norm and starts at zero. Global Picard construction with this continuous additive input gives the integral solution. Subtracting two solutions with the same input leaves an ordinary integral of their drift difference, so Grönwall gives uniqueness. Subtracting the original integral equation at time \(a\) from that at \(a+t\) gives the restart equation, since Brownian increments telescope. Uniqueness proves restart agreement.



For two initial states with the same input the difference \(\delta=r_z-r_w\) is continuously differentiable, with \(\delta'=f(r_z)-f(r_w)\). The chain rule and the dissipativity assumption yield \((\|C\delta\|^2)'=2\langle C\delta,C\delta'\rangle\le-2\rho\|C\delta\|^2\). Exponential energy comparison followed by the nonnegative square root gives the result. No contraction of the numerical update is assumed. \(\square\)

\begin{leanbox}{Lean: \texttt{higherOrderRestart}, \texttt{higherOrderRestart\_\allowbreak{}equation}, \texttt{higherOrderRestart\_\allowbreak{}agrees} and 1 more}
\leansource{\texttt{higherOrderRestart} \quad PBGSO/\allowbreak{}HOLRestart.lean, lines 112–115}
\begin{leancode}
\leanline{0}{def\ higherOrderRestart}
\leanline{4}{(z\ :\ HOLState\ Θ\ (Lp\ (EuclideanSpace\ ℝ\ I)\ 2\ Q))\ (a\ :\ ℝ)\ :}
\leanline{4}{ℝ\ →\ HOLState\ Θ\ (Lp\ (EuclideanSpace\ ℝ\ I)\ 2\ Q)\ :=}
\leanline{2}{Classical.choose\ (exists\_higherOrder\_restart\ Q\ hf\ Fbar\ G\ LF\ LG\ α\ γ\ hLF\ hLG\ hF\ hG\ z\ a)}
\end{leancode}
\leansource{\texttt{higherOrderRestart\_\allowbreak{}equation} \quad PBGSO/\allowbreak{}HOLRestart.lean, lines 129–135}
\begin{leancode}
\leanline{0}{theorem\ higherOrderRestart\_equation\ (t\ :\ ℝ)\ (ht\ :\ 0\ ≤\ t)\ :}
\leanline{4}{higherOrderRestart\ Q\ hf\ Fbar\ G\ LF\ LG\ α\ γ\ hLF\ hLG\ hF\ hG\ z\ a\ t\ -\ z\ =}
\leanline{6}{(∫\ s\ in\ 0..t,\ holDrift\ Fbar\ G\ α\ γ}
\leanline{8}{(higherOrderRestart\ Q\ hf\ Fbar\ G\ LF\ LG\ α\ γ\ hLF\ hLG\ hF\ hG\ z\ a\ s))\ +}
\leanline{10}{higherOrderRestartInput\ Q\ hf\ α\ a\ t}
\end{leancode}
\leansource{\texttt{higherOrderRestart\_\allowbreak{}agrees} \quad PBGSO/\allowbreak{}HOLRestart.lean, lines 162–177}
\begin{leancode}
\leanline{0}{theorem\ higherOrderRestart\_agrees}
\leanline{4}{(Y\ :\ ℝ\ →\ HOLState\ Θ\ (Lp\ (EuclideanSpace\ ℝ\ I)\ 2\ Q))\ (hY\ :\ Continuous\ Y)}
\leanline{4}{(heq\ :\ ∀\ t,\ 0\ ≤\ t\ →\ Y\ t\ -\ Y\ 0\ =}
\leanline{6}{(∫\ s\ in\ 0..t,\ holDrift\ Fbar\ G\ α\ γ\ (Y\ s))\ +\ higherOrderRestartInput\ Q\ hf\ α\ 0\ t)}
\leanline{4}{(a\ t\ :\ ℝ)\ (ha\ :\ 0\ ≤\ a)\ (ht\ :\ 0\ ≤\ t)\ :}
\leanline{4}{higherOrderRestart\ Q\ hf\ Fbar\ G\ LF\ LG\ α\ γ\ hLF\ hLG\ hF\ hG\ (Y\ a)\ a\ t\ =\ Y\ (a\ +\ t)}
\end{leancode}
\leansource{\texttt{higherOrderRestart\_\allowbreak{}contracts} \quad PBGSO/\allowbreak{}HOLRestart.lean, lines 183–216}
\begin{leancode}
\leanline{0}{theorem\ higherOrderRestart\_contracts}
\leanline{4}{\{H\ :\ Type*\}\ [NormedAddCommGroup\ H]\ [InnerProductSpace\ ℝ\ H]}
\leanline{4}{(C\ :\ HOLState\ Θ\ (Lp\ (EuclideanSpace\ ℝ\ I)\ 2\ Q)\ →L[ℝ]\ H)\ (ρ\ :\ ℝ)}
\leanline{4}{(hmono\ :\ ∀\ x\ y,}
\leanline{6}{inner\ ℝ\ (C\ (x\ -\ y))\ (C\ (holDrift\ Fbar\ G\ α\ γ\ x\ -\ holDrift\ Fbar\ G\ α\ γ\ y))\ ≤}
\leanline{8}{-ρ\ *\ ‖C\ (x\ -\ y)‖\ \textasciicircum{}\ 2)}
\leanline{4}{(z\ w\ :\ HOLState\ Θ\ (Lp\ (EuclideanSpace\ ℝ\ I)\ 2\ Q))\ (t\ :\ ℝ)\ (ht\ :\ 0\ ≤\ t)\ :}
\leanline{4}{‖C\ (higherOrderRestart\ Q\ hf\ Fbar\ G\ LF\ LG\ α\ γ\ hLF\ hLG\ hF\ hG\ z\ a\ t\ -}
\leanline{8}{higherOrderRestart\ Q\ hf\ Fbar\ G\ LF\ LG\ α\ γ\ hLF\ hLG\ hF\ hG\ w\ a\ t)‖\ ≤}
\leanline{6}{Real.exp\ (-ρ\ *\ t)\ *\ ‖C\ (z\ -\ w)‖}
\end{leancode}
\leanlabel{How the Lean proof follows the argument}
\texttt{higher\allowbreak{}Order\allowbreak{}Restart} chooses the actual solution from \texttt{exists\_\allowbreak{}higher\allowbreak{}Order\_\allowbreak{}restart}; \texttt{higher\allowbreak{}Order\allowbreak{}Restart\_\allowbreak{}equation} extracts its integral equation. \texttt{higher\allowbreak{}Order\allowbreak{}Restart\_\allowbreak{}agrees} uses integral shifting and \texttt{higher\allowbreak{}Order\allowbreak{}Restart\allowbreak{}Input\_\allowbreak{}increment}, then the proved uniqueness result. \texttt{higher\allowbreak{}Order\allowbreak{}Restart\_\allowbreak{}contracts} changes the drift sign to match the shared integral-contraction theorem and applies \texttt{transformed\_\allowbreak{}additive\_\allowbreak{}integral\_\allowbreak{}contraction}. Continuity and both integral identities are supplied by the constructed restart properties. The final simplification substitutes their initial states.
\end{leanbox}
\section{\texorpdfstring{Theorem New A: a time-uniform error bound for the M-MFPO particle method}{Theorem New A: a time-uniform error bound for the M-MFPO particle method}}


This section proves Theorem New A. The M-MFPO particle system has parameter \(\theta\), parameter momentum \(m\), particle positions \(X^i\) and velocities \(V^i\), \(i=1,\dots,N\), and evolves by


\[\dot\theta=m,\quad\dot m=-\frac1N\sum_iF(\theta,X^i)-\gamma m,\quad dX^i=V^i\,dt,\quad dV^i=\bigl(-F_x(\theta,X^i)-\gamma V^i\bigr)dt+\sqrt{2\gamma}\,dB^i .\]


The numerical method is the OBABO splitting with step \(h\), Gaussian half-step innovations for the velocities and no innovation for the parameter momentum. Both are started from the same random state: a deterministic optimiser state \((\theta_0,m_0)\) and \(N\) independent particle phases with a common law \(\mu\) of finite second moment, independent of the Brownian motions.




Under N1 (global Lipschitz forces), A3 with rate \(\rho>0\), A5 with variance bound \(V\), N3 (a frozen stationary trajectory for the target law \(\nu_\star\)) and zero mean observation at the target \(\theta_\star\), the theorem states: for every step \(0<h\le h_{\star,M}\) and every iteration \(n\),


\[\|\hat\theta_n-\theta_\star\|_{L^2}\le D_M\sqrt h+\sqrt\kappa\,e^{-\rho nh}\,d_0+\frac{2+\gamma}{\rho}\sqrt{\frac{\kappa V}{N}}\bigl(1-e^{-\rho nh}\bigr),\]


where \(\kappa=\lambda_+/\lambda_-\), \(d_0\) is the distance of the initial optimiser state and initial law from the target, and \(D_M\), \(h_{\star,M}>0\) are given by explicit formulas in the data. The constants do not depend on \(h\), \(n\) or a time horizon; they may depend on \(N\). The first term is the discretisation error, proved for the whole random state; the other two terms are the continuous-time relaxation and the finite-particle error of the paper's parameter theorem. The section first defines every constant and then proves the bound.

\subsection{\texorpdfstring{The metric comparison constants}{The metric comparison constants}}


\textbf{Definition.} Let \(0<\lambda_-\le\lambda_+\) be the eigenvalue bounds of the A3 geometry and \(N\ge1\). Put \(\lambda_N^-=\lambda_-\min\{1,N^{-1}\}\), \(\lambda_N^+=\lambda_+\max\{1,N^{-1}\}\) and


\[m_N=\frac{\sqrt{\lambda_N^-}}{2},\qquad U_N=2\sqrt{\lambda_N^+}.\]


Since \(N\ge1\), \(m_N=\tfrac12\sqrt{\lambda_-/N}\) and \(U_N=2\sqrt{\lambda_+}\). These are the constants of the contracting metric \(T\) for the kinetic particle system under the standard Euclidean encoding: \(m_N\|y\|\le\|Ty\|\le U_N\|y\|\) for every random state \(y\), where \(\|y\|\) is the maximum of the four \(L^2\) norms of \(y=((m,V),(\theta,X))\). The factor \(2\) and \(\tfrac12\) come from the norm of the regrouping map into the empirical Hilbert state and of its inverse.

\begin{leanbox}{Lean: \texttt{metricCoercivity}, \texttt{metricUpper}}
\leansource{\texttt{metricCoercivity} \quad PBGSO/\allowbreak{}FullyUniformMomentum.lean, lines 45–45}
\begin{leancode}
\leanline{0}{def\ metricCoercivity\ :\ ℝ\ :=\ Real.sqrt\ (EmpiricalMetric.lower\ N\ Geometry)\ /\ 2}
\end{leancode}
\leansource{\texttt{metricUpper} \quad PBGSO/\allowbreak{}FullyUniformMomentum.lean, lines 48–48}
\begin{leancode}
\leanline{0}{def\ metricUpper\ :\ ℝ\ :=\ 2\ *\ Real.sqrt\ (EmpiricalMetric.upper\ N\ Geometry)}
\end{leancode}
\leanlabel{How the Lean proof follows the argument}
\texttt{metric\allowbreak{}Coercivity} is \texttt{Real.\allowbreak{}sqrt (Empirical\allowbreak{}Metric.\allowbreak{}lower N Geometry) /\allowbreak{} 2} and \texttt{metric\allowbreak{}Upper} is \texttt{2 * Real.\allowbreak{}sqrt (Empirical\allowbreak{}Metric.\allowbreak{}upper N Geometry)}, where \texttt{Geometry} is the A3 geometry on \texttt{Hypocoercive\allowbreak{}Phase A} and \texttt{Hypocoercive\allowbreak{}Phase (Euclidean\allowbreak{}Space \ensuremath{\mathbb R} J)}; the finite type \texttt{J} indexes the \(d_x\) coordinates of one particle.
\end{leanbox}
\subsection{\texorpdfstring{The noise radius and the local error coefficients}{The noise radius and the local error coefficients}}


\textbf{Definition.} Let \(d=Nd_x\). The probability space \(\Omega\) carries the initial phases of the \(N\) particles (independent with law \(\mu\)) and \(N\) independent standard Brownian motions in \(\mathbb R^{d_x}\), independent of the initial phases; under the Euclidean encoding the velocity noise is a standard Brownian motion \(W\) in \(\mathbb R^{d}\). Put


\[\beta_M=\sqrt\gamma\,\sqrt{2d},\qquad\Xi_M=\sqrt d,\qquad L_M=2\max\{L_F,L_x\},\qquad b_\circ=\|\mathbf F(0)\|_{L^2},\]


where \(\mathbf F=(\bar F,\mathcal F_x)\) is the coupled force on random arguments and \(L_F,L_x\) are its two Lipschitz constants. Let \((c_{0,M},c_{1,M})\) be the coefficients of the OBABO local error for \(L=L_M\), friction \(\gamma\), \(f_0=b_\circ\), \(\Xi=\Xi_M\), \(\beta=\beta_M\); explicitly, with \(K_M=L_M+\gamma+1\) and \(T_M=b_\circ+K_Me^{K_M+1}(b_\circ+\beta_M)+\beta_M\),


\[\begin{aligned}c_{0,M}&=\Bigl(\frac{\gamma+1}{2}+\frac{L_M}{4}\Bigr)b_\circ+\Bigl(\frac{L_M}{2}+\sqrt\gamma+\frac32\gamma\sqrt\gamma\Bigr)\sqrt d+\frac23K_MT_M,\\ c_{1,M}&=\frac{\gamma^2+\gamma+L_M}{2}+\Bigl(\frac{\gamma+1}{2}+\frac{L_M}{4}\Bigr)L_M+\frac23K_M^2e^{K_M+1}.\end{aligned}\]


The meaning of the three inputs is as follows. \(\beta_M\) is the increment constant of the physical input: \(\|\sqrt{2\gamma}(W(a+s)-W(a))\|_{L^2}=\sqrt{2\gamma}\sqrt{ds}=\beta_M\sqrt s\). \(\Xi_M\) is the \(L^2\) norm of a normalised half-step Gaussian innovation \((W((k+1)h/2)-W(kh/2))/\sqrt{h/2}\), whose \(d\) coordinates are independent standard Gaussians. \(L_M\) is a Lipschitz constant of \(\mathbf F\) for the maximum norm, because \(\|\mathbf F(q)-\mathbf F(q')\|\le\max\{L_F,L_x\}(\|\theta-\theta'\|+\|X-X'\|)\le2\max\{L_F,L_x\}\|q-q'\|\). For the Euclidean encoding, \(b_\circ=\max\{\|F(0,0)\|,\sqrt N\|F_x(0,0)\|\}\), since \(\bar F(0,0)=F(0,0)\) and \(\mathcal F_x(0,0)\) is the cloud with all particles equal to \(F_x(0,0)\).

\begin{leanbox}{Lean: \texttt{noiseRadius}, \texttt{localConstant0}, \texttt{localConstant1}}
\leansource{\texttt{noiseRadius} \quad PBGSO/\allowbreak{}FullyUniformMomentum.lean, lines 51–52}
\begin{leancode}
\leanline{0}{def\ noiseRadius\ (hγ\ :\ 0\ ≤\ γ)\ :\ ℝ\ :=}
\leanline{2}{momentumUniformBrownianRadius\ (I\ :=\ Fin\ N\ ×\ J)\ ⟨γ,\ hγ⟩}
\end{leancode}
\leansource{\texttt{localConstant0} \quad PBGSO/\allowbreak{}FullyUniformMomentum.lean, lines 55–59}
\begin{leancode}
\leanline{0}{def\ localConstant0\ (hγ\ :\ 0\ ≤\ γ)\ :\ ℝ\ :=}
\leanline{2}{let\ Q\ :=\ ParticleNoiseSpace.probability\ (J\ :=\ J)\ μ\ N}
\leanline{2}{(momentumUniformDefect\ (momentumUniformForceLip\ D)\ ⟨γ,\ hγ⟩}
\leanline{4}{‖coupledForce\ (D.meanForceLp\ Q)\ (D.samplerForceLp\ Q)\ 0‖}
\leanline{4}{(Real.sqrt\ (Fintype.card\ (Fin\ N\ ×\ J)\ :\ ℝ))\ (noiseRadius\ (N\ :=\ N)\ (J\ :=\ J)\ γ\ hγ)).1}
\end{leancode}
\leansource{\texttt{localConstant1} \quad PBGSO/\allowbreak{}FullyUniformMomentum.lean, lines 62–66}
\begin{leancode}
\leanline{0}{def\ localConstant1\ (hγ\ :\ 0\ ≤\ γ)\ :\ ℝ\ :=}
\leanline{2}{let\ Q\ :=\ ParticleNoiseSpace.probability\ (J\ :=\ J)\ μ\ N}
\leanline{2}{(momentumUniformDefect\ (momentumUniformForceLip\ D)\ ⟨γ,\ hγ⟩}
\leanline{4}{‖coupledForce\ (D.meanForceLp\ Q)\ (D.samplerForceLp\ Q)\ 0‖}
\leanline{4}{(Real.sqrt\ (Fintype.card\ (Fin\ N\ ×\ J)\ :\ ℝ))\ (noiseRadius\ (N\ :=\ N)\ (J\ :=\ J)\ γ\ hγ)).2}
\end{leancode}
\leanlabel{How the Lean proof follows the argument}
\texttt{noise\allowbreak{}Radius \ensuremath{\gamma} h\ensuremath{\gamma}} is \texttt{momentum\allowbreak{}Uniform\allowbreak{}Brownian\allowbreak{}Radius (I := Fin N \ensuremath{\times} J) \ensuremath{\langle}\ensuremath{\gamma}, h\ensuremath{\gamma}\ensuremath{\rangle}}, that is \(\sqrt\gamma\sqrt{2\,|\mathrm{Fin}\,N\times J|}\); the friction is packaged with its nonnegativity \texttt{h\ensuremath{\gamma} : 0 \ensuremath{\le} \ensuremath{\gamma}} as an element of \texttt{\ensuremath{\mathbb R}\ensuremath{\ge}0}. \texttt{local\allowbreak{}Constant0} and \texttt{local\allowbreak{}Constant1} are the two components of \texttt{momentum\allowbreak{}Uniform\allowbreak{}Defect (momentum\allowbreak{}Uniform\allowbreak{}Force\allowbreak{}Lip D) \ensuremath{\langle}\ensuremath{\gamma}, h\ensuremath{\gamma}\ensuremath{\rangle} \ensuremath{\Vert}coupled\allowbreak{}Force (D.\allowbreak{}mean\allowbreak{}Force\allowbreak{}Lp Q) (D.\allowbreak{}sampler\allowbreak{}Force\allowbreak{}Lp Q) 0\ensuremath{\Vert} (Real.\allowbreak{}sqrt (Fintype.\allowbreak{}card (Fin N \ensuremath{\times} J) : \ensuremath{\mathbb R})) (noise\allowbreak{}Radius …)}, where \texttt{Q := Particle\allowbreak{}Noise\allowbreak{}Space.\allowbreak{}probability (J := J) \ensuremath{\mu} N} is the product probability space described above and \texttt{momentum\allowbreak{}Uniform\allowbreak{}Force\allowbreak{}Lip D} is \texttt{2*max D.\allowbreak{}mean\allowbreak{}Lip D.\allowbreak{}sampler\allowbreak{}Force\allowbreak{}Lip}. The identity \(b_\circ=\max\{\|F(0,0)\|,\sqrt N\|F_x(0,0)\|\}\) is \texttt{coupled\allowbreak{}Force\_\allowbreak{}zero\_\allowbreak{}norm} in the file on origin constants.
\end{leanbox}
\subsection{\texorpdfstring{Radii of the exact flow}{Radii of the exact flow}}


\textbf{Definition.} Let \(y_0\) be the common initial random state of the particle system and of the numerical method, and \(\|y_0\|\) the maximum of its four \(L^2\) norms. Put


\[\begin{aligned}r_{0,M}&=e^{K_M+1}\bigl(b_\circ+\beta_M\bigr),\\ S_M&=U_N\,r_{0,M}=U_Ne^{K_M+1}(b_\circ+\beta_M),\\ R_M&=U_N\|y_0\|+\frac{S_M}{1-e^{-\rho}}+S_M.\end{aligned}\]


Here \(r_{0,M}\) is the one-unit radius \(r_{K_M}(b_\circ,\beta_M,0)\) of the kinetic integral equation started at zero: the drift \((p,q)\mapsto(-\mathbf F(q)-\gamma p,p)\) is Lipschitz with constant \(K_M\) and has norm \(b_\circ\) at the origin. The number \(R_M\) bounds the metric norm \(\|Ty(t)\|\) of the continuous particle system at all times \(t\ge0\): this is the uniform moment bound with \(\|Ty(0)\|\le U_N\|y_0\|\) and \(\|T\|\le U_N\). For \(\rho>0\) the denominator \(1-e^{-\rho}\) is positive.

\begin{leanbox}{Lean: \texttt{zeroStartRadius}, \texttt{zeroStartMoment}, \texttt{referenceRadius}}
\leansource{\texttt{zeroStartRadius} \quad PBGSO/\allowbreak{}FullyUniformMomentum.lean, lines 69–74}
\begin{leancode}
\leanline{0}{def\ zeroStartRadius\ (hγ\ :\ 0\ ≤\ γ)\ :\ ℝ\ :=}
\leanline{2}{let\ Q\ :=\ ParticleNoiseSpace.probability\ (J\ :=\ J)\ μ\ N}
\leanline{2}{let\ γnn\ :\ ℝ≥0\ :=\ ⟨γ,\ hγ⟩}
\leanline{2}{UniformIntegral.radius\ (momentumUniformForceLip\ D\ +\ γnn\ +\ 1)}
\leanline{4}{‖coupledForce\ (D.meanForceLp\ Q)\ (D.samplerForceLp\ Q)\ 0‖}
\leanline{4}{(noiseRadius\ (N\ :=\ N)\ (J\ :=\ J)\ γ\ hγ)\ 0}
\end{leancode}
\leansource{\texttt{zeroStartMoment} \quad PBGSO/\allowbreak{}FullyUniformMomentum.lean, lines 77–78}
\begin{leancode}
\leanline{0}{def\ zeroStartMoment\ (hγ\ :\ 0\ ≤\ γ)\ :\ ℝ\ :=}
\leanline{2}{metricUpper\ (N\ :=\ N)\ Geometry\ *\ zeroStartRadius\ D\ γ\ μ\ hγ}
\end{leancode}
\leansource{\texttt{referenceRadius} \quad PBGSO/\allowbreak{}FullyUniformMomentum.lean, lines 81–84}
\begin{leancode}
\leanline{0}{def\ referenceRadius\ (hγ\ :\ 0\ ≤\ γ)\ (ρ\ :\ ℝ)\ :\ ℝ\ :=}
\leanline{2}{let\ S\ :=\ zeroStartMoment\ D\ γ\ Geometry\ μ\ hγ}
\leanline{2}{metricUpper\ (N\ :=\ N)\ Geometry\ *\ ‖initial\ D\ γ\ Geometry\ μ\ hμ\ a\ Variance\ hVariance‖\ +}
\leanline{4}{S\ /\ (1\ -\ Real.exp\ (-ρ))\ +\ S}
\end{leancode}
\leanlabel{How the Lean proof follows the argument}
\texttt{zero\allowbreak{}Start\allowbreak{}Radius} is \texttt{Uniform\allowbreak{}Integral.\allowbreak{}radius (momentum\allowbreak{}Uniform\allowbreak{}Force\allowbreak{}Lip D + \ensuremath{\gamma}nn + 1) \ensuremath{\Vert}coupled\allowbreak{}Force … 0\ensuremath{\Vert} (noise\allowbreak{}Radius …) 0}, with \texttt{\ensuremath{\gamma}nn = \ensuremath{\langle}\ensuremath{\gamma}, h\ensuremath{\gamma}\ensuremath{\rangle}}; by the definition of \texttt{radius} this is \(e^{K_M+1}(0+b_\circ+\beta_M)\). \texttt{zero\allowbreak{}Start\allowbreak{}Moment} multiplies it by \texttt{metric\allowbreak{}Upper}. \texttt{reference\allowbreak{}Radius} uses \texttt{\ensuremath{\Vert}initial D \ensuremath{\gamma} Geometry \ensuremath{\mu} h\ensuremath{\mu} a Variance h\allowbreak{}Variance\ensuremath{\Vert}}, the norm of the raw initial state \texttt{initial} (defined in the file on the fully constructed kinetic bound as the raw coordinates of the constructed particle system at time zero).
\end{leanbox}
\subsection{\texorpdfstring{The discretisation constant and the step-size threshold}{The discretisation constant and the step-size threshold}}


\textbf{Definition.} Put


\[\begin{aligned}a_M&=U_Nc_{0,M},&b_M&=\frac{U_Nc_{1,M}}{m_N},\\D_M&=\frac{4\,(a_M+b_MR_M)}{\rho\,m_N},\\h_{\star,M}&=\min\Bigl\{1,\;\frac1\rho,\;\Bigl(\frac{\rho}{4(b_M+1)}\Bigr)^2\Bigr\}.\end{aligned}\]


\textbf{Claim.} Let \(\rho>0\) and \(b\ge0\), and let \(h_\star(\rho,b)=\min\{1,1/\rho,(\rho/(4(b+1)))^2\}\). Then \(h_\star(\rho,b)>0\), and every \(h\) with \(0<h\le h_\star(\rho,b)\) satisfies


\[h\le1,\qquad\rho h\le1,\qquad b\sqrt h\le\frac\rho4 .\]


\textbf{Proof.} The three numbers \(1\), \(1/\rho\) and \((\rho/(4(b+1)))^2\) are positive because \(\rho>0\) and \(b+1\ge1\), so their minimum is positive. If \(0<h\le h_\star\), then \(h\le1\) and \(h\le1/\rho\), hence \(\rho h\le1\). Also \(h\le(\rho/(4(b+1)))^2\) with \(\rho/(4(b+1))\ge0\), so \(\sqrt h\le\rho/(4(b+1))\) by monotonicity of the square root. Since \(b\ge0\) and \(b\le b+1\),


\[b\sqrt h\le\frac{b\,\rho}{4(b+1)}\le\frac{(b+1)\rho}{4(b+1)}=\frac\rho4. \qquad\square\]


Applied with \(b=b_M\ge0\), the threshold \(h_{\star,M}=h_\star(\rho,b_M)\) makes all three conditions of the absorption argument hold. The term \(b+1\) instead of \(b\) keeps the threshold finite when \(b_M=0\).

\begin{leanbox}{Lean: \texttt{radiusGrowth}, \texttt{discretizationConstant}, \texttt{maximumStep}}
\leansource{\texttt{radiusGrowth} \quad PBGSO/\allowbreak{}FullyUniformMomentum.lean, lines 87–88}
\begin{leancode}
\leanline{0}{def\ radiusGrowth\ (hγ\ :\ 0\ ≤\ γ)\ :\ ℝ\ :=}
\leanline{2}{metricUpper\ (N\ :=\ N)\ Geometry\ *\ localConstant1\ D\ γ\ μ\ hγ\ /\ metricCoercivity\ (N\ :=\ N)\ Geometry}
\end{leancode}
\leansource{\texttt{discretizationConstant} \quad PBGSO/\allowbreak{}FullyUniformMomentum.lean, lines 91–94}
\begin{leancode}
\leanline{0}{def\ discretizationConstant\ (hγ\ :\ 0\ ≤\ γ)\ (ρ\ :\ ℝ)\ :\ ℝ\ :=}
\leanline{2}{4\ *\ (metricUpper\ (N\ :=\ N)\ Geometry\ *\ localConstant0\ D\ γ\ μ\ hγ\ +}
\leanline{4}{radiusGrowth\ D\ γ\ Geometry\ μ\ hγ\ *\ referenceRadius\ D\ γ\ Geometry\ μ\ hμ\ a\ Variance\ hVariance\ hγ\ ρ)\ /}
\leanline{6}{(ρ\ *\ metricCoercivity\ (N\ :=\ N)\ Geometry)}
\end{leancode}
\leansource{\texttt{maximumStep} \quad PBGSO/\allowbreak{}FullyUniformMomentum.lean, lines 97–98}
\begin{leancode}
\leanline{0}{def\ maximumStep\ (hγ\ :\ 0\ ≤\ γ)\ (ρ\ :\ ℝ)\ :\ ℝ\ :=}
\leanline{2}{UniformStep.maximumStep\ ρ\ (radiusGrowth\ D\ γ\ Geometry\ μ\ hγ)}
\end{leancode}
\leanlabel{How the Lean proof follows the argument}
\texttt{radius\allowbreak{}Growth} is \texttt{metric\allowbreak{}Upper * local\allowbreak{}Constant1 /\allowbreak{} metric\allowbreak{}Coercivity} (\(b_M\)), \texttt{discretization\allowbreak{}Constant} is \(D_M\) with \texttt{reference\allowbreak{}Radius} as \(R_M\), and \texttt{maximum\allowbreak{}Step} is \texttt{Uniform\allowbreak{}Step.\allowbreak{}maximum\allowbreak{}Step \ensuremath{\rho} (radius\allowbreak{}Growth …)}. The function \texttt{Uniform\allowbreak{}Step.\allowbreak{}maximum\allowbreak{}Step \ensuremath{\rho} b} (file on the uniform step size) is \texttt{min 1 (min (1 /\allowbreak{} \ensuremath{\rho}) ((\ensuremath{\rho} /\allowbreak{} (4 * (b + 1))) \textasciicircum{} 2))}; the claim consists of its two lemmas \texttt{Uniform\allowbreak{}Step.\allowbreak{}maximum\allowbreak{}Step\_\allowbreak{}pos} (positivity, by \texttt{positivity} and \texttt{sq\_\allowbreak{}pos\_\allowbreak{}of\_\allowbreak{}pos}) and \texttt{Uniform\allowbreak{}Step.\allowbreak{}admissible} (the three conditions, using \texttt{le\_\allowbreak{}div\_\allowbreak{}iff\ensuremath{{}_{0}}}, \texttt{Real.\allowbreak{}sqrt\_\allowbreak{}le\_\allowbreak{}iff} and a \texttt{calc} chain that is the last display).
\end{leanbox}
\subsection{\texorpdfstring{The explicit guarantee}{The explicit guarantee}}


\textbf{Setting.} Let \(\Theta\) be a finite-dimensional real inner product space, \(\mathcal X=\mathbb R^{d_x}\), \(\mathcal C=\mathbb R^{Nd_x}\) with the standard Euclidean encoding, and let the force data satisfy N1 with \(N\ge1\). Let \(\gamma\ge0\). Let \(\mu\) be a probability law on \(\mathbb R^{d_x}\times\mathbb R^{d_x}\) with finite second moment (the initial phase of each particle) and \(a=(\theta_0,m_0)\in\Theta\times\Theta\) the initial optimiser state. Let \(V\in\mathbb R\) bound the observation variance \(\int\|F(\vartheta_r,x)-G_{\nu_r}(\vartheta_r)\|^2\nu_r(dx\,dv)\) of the mean-field solution \((\vartheta_r,\nu_r)\) started at \((a,\mu)\), for all \(r\ge0\) (A5). Let \(\theta_\star\in\Theta\), and let \(\nu_\star\) be a probability law with finite second moment for which there is a frozen trajectory (N3): a continuous square-integrable process \(Z^\star\), defined on \(\nu_\star\) times the Wiener measure, with \(Z^\star(0)\) the initial coordinate, \(Z^\star(t)-Z^\star(0)=\int_0^tb_{\theta_\star}(Z^\star(r))dr+\sqrt{2\gamma}(0,B(t))\) and \(\operatorname{Law}(Z^\star(t))=\nu_\star\) for \(t\ge0\). Assume A3 for the kinetic model with rate \(\rho>0\) and geometry constants \(\lambda_\pm\), the zero mean observation \(G_{\nu_\star}(\theta_\star)=0\), and \(V\ge0\). Put \(\kappa=\lambda_+/\lambda_-\) and


\[d_0=\Bigl(\|(\theta_0,m_0)-(\theta_\star,0)\|^2+W_2(\mu,\nu_\star)^2\Bigr)^{1/2},\]


where \(W_2(\mu,\nu)^2=\inf_\pi\int\|z-z'\|^2\pi(dz,dz')\) over all couplings \(\pi\) of \(\mu\) and \(\nu\). Let \(y(t)\) be the continuous interacting particle system in the random coordinates \(((m,V),(\theta,X))\) and \(\hat y_n\) the OBABO iterates with step \(h\) and the Gaussian half-step innovations of the same Brownian motion, started from \(\hat y_0=y(0)\); \(\hat\theta_n\) is the parameter component of \(\hat y_n\).




\textbf{Claim.} In this setting \(h_{\star,M}>0\), \(D_M\ge0\), and for every \(h\in(0,h_{\star,M}]\) and every \(n\in\mathbb N\)


\[\|\hat y_n-y(nh)\|\le D_M\sqrt h,\qquad\|\hat\theta_n-\theta_\star\|_{L^2}\le D_M\sqrt h+\sqrt\kappa\,e^{-\rho nh}d_0+\frac{2+\gamma}{\rho}\sqrt{\frac{\kappa V}{N}}\bigl(1-e^{-\rho nh}\bigr).\]


\textbf{Proof.} We use four results proved elsewhere.



\begin{enumerate}


\item \emph{Particle system.} The interacting particle system is constructed on \(\Omega\) as a continuous process; in the random coordinates it satisfies, for \(t\ge0\), \(y(t)-y(0)=\int_0^tb_M(y(r))dr+\bigl((0,\sqrt{2\gamma}\,W(t)),(0,0)\bigr)\) with \(b_M((m,V),(\theta,X))=(-\mathbf F(\theta,X)-\gamma(m,V),(m,V))\). Its initial state has the deterministic optimiser state \(a\) and independent particle phases with law \(\mu\).

\item \emph{Metric.} A3 gives a bounded linear \(T\) into a Hilbert space with \(m_N\|y\|\le\|Ty\|\le U_N\|y\|\) and \(\langle T(y-y'),T(b_M(y)-b_M(y'))\rangle\le-\rho\|T(y-y')\|^2\) for all random states (contracting metric for the kinetic system, with the Euclidean bounds on the regrouping map).

\item \emph{Uniform comparison.} Let \(C\) be a bounded linear map into a Hilbert space and \(m>0\) with \(m\|x\|\le\|Cx\|\) and \(\langle C(x-y),C(b_M(x)-b_M(y))\rangle\le-\rho\|C(x-y)\|^2\), let \(\rho>0\), let \(y\) be as in item 1, and let \(0<h\le1\), \(\rho h\le1\) and \((\|C\|c_{1,M}/m)\sqrt h\le\rho/4\). Then for every \(n\)

\end{enumerate}

\[\|\hat y_n-y(nh)\|\le\frac{4\bigl(\|C\|c_{0,M}+(\|C\|c_{1,M}/m)\,R_C\bigr)}{\rho\,m}\sqrt h,\qquad R_C=\|Cy(0)\|+\frac{\|C\|r_{0,M}}{1-e^{-\rho}}+\|C\|r_{0,M}.\]


This comparison is the composition of the results of the preceding sections: the exact flow over \([kh,(k+1)h]\) is constructed by Picard iteration with the shifted Brownian input and contracts by \(e^{-\rho h}\) in \(\|C\cdot\|\); it reproduces \(y\) on the grid by uniqueness; its local error against one OBABO step is at most \(h\sqrt h(c_{0,M}+c_{1,M}\|x\|)\le h\sqrt h(c_{0,M}+c_{1,M}\|Cx\|/m)\), with \(\Xi_M=\sqrt d\) the \(L^2\) norm of the half-step innovations and \(\beta_M\) the Brownian increment constant; the grid values satisfy \(\|Cy(kh)\|\le R_C\) by the uniform moment bound; and the absorption argument concludes.



\begin{enumerate}


\item \emph{Continuous-time parameter bound} (the paper's parameter theorem for the constructed system). Under A3 with \(\rho>0\), N3, the zero mean observation and A5 with \(V\ge0\), the parameter \(\theta^N_t\) of the particle system satisfies, for \(t\ge0\), \(\|\theta^N_t-\theta_\star\|_{L^2}\le\sqrt\kappa e^{-\rho t}d_0+\frac{L_\psi}{\rho}\sqrt{\kappa V/N}(1-e^{-\rho t})\), where \(L_\psi\) is the Lipschitz constant of the optimiser response \(((\theta,m),g)\mapsto(m,-g-\gamma m)\); here \(L_\psi=2+|\gamma|=2+\gamma\). The target \((\theta_\star,0)\) is fixed by the response: \((0,-0-\gamma\cdot0)=0\).

\end{enumerate}



\emph{Step 1: constants.} By item 2, \(\|T\|\le U_N\), and \(m_N>0\) because \(\lambda_N^->0\). Since \(c_{0,M},c_{1,M}\ge0\) and \(r_{0,M}\ge0\), we get \(\|T\|c_{0,M}\le a_M\), \(\|T\|c_{1,M}/m_N\le b_M\), \(\|T\|r_{0,M}\le S_M\), and \(\|Ty(0)\|\le U_N\|y(0)\|=U_N\|y_0\|\); hence \(R_T\le R_M\) and the constant of item 3 for \(C=T\), \(m=m_N\) is at most \(D_M\). In particular \(D_M\ge0\), and \(h_{\star,M}>0\) by the previous step since \(b_M\ge0\).




\emph{Step 2: the full state.} Let \(0<h\le h_{\star,M}\). The previous step gives \(h\le1\), \(\rho h\le1\) and \(b_M\sqrt h\le\rho/4\), hence \((\|T\|c_{1,M}/m_N)\sqrt h\le\rho/4\). Item 3 with \(C=T\), \(m=m_N\) and Step 1 give \(\|\hat y_n-y(nh)\|\le D_M\sqrt h\).




\emph{Step 3: the parameter.} \(\hat\theta_n-\theta(nh)\) is one of the four components of \(\hat y_n-y(nh)\), so \(\|\hat\theta_n-\theta(nh)\|_{L^2}\le D_M\sqrt h\). The parameter component of \(y(nh)\) is \(\theta^N_{nh}\), and item 4 at \(t=nh\ge0\) bounds \(\|\theta^N_{nh}-\theta_\star\|_{L^2}\). The triangle inequality in \(L^2(\Omega;\Theta)\), \(\|\hat\theta_n-\theta_\star\|\le\|\hat\theta_n-\theta(nh)\|+\|\theta(nh)-\theta_\star\|\), gives the second bound. \(\square\)

\begin{leanbox}{Lean: \texttt{explicit\_\allowbreak{}guarantee}}
\leansource{\texttt{explicit\_\allowbreak{}guarantee} \quad PBGSO/\allowbreak{}FullyUniformMomentum.lean, lines 101–235}
\begin{leancode}
\leanline{0}{theorem\ explicit\_guarantee}
\leanline{4}{(hcloud\ :\ D.cloud\ =\ euclideanParticleEncoding\ N\ J)\ (hγ\ :\ 0\ ≤\ γ)}
\leanline{4}{(star\ :\ A)\ (νStar\ :\ Measure\ (HypocoercivePhase\ (EuclideanSpace\ ℝ\ J)))\ [IsProbabilityMeasure\ νStar]}
\leanline{4}{(hνStar\ :\ MemLp\ (fun\ z\ :\ HypocoercivePhase\ (EuclideanSpace\ ℝ\ J)\ ↦\ z)\ 2\ νStar)}
\leanline{4}{(Frozen\ :\ AdditiveStationary.FrozenTrajectory\ (ULA.CanonicalBrownian.vectorProbability\ J)}
\leanline{6}{(momentumForceDynamics\ D\ γ)\ νStar\ hνStar\ (targetState\ star)\ (input\ γ))}
\leanline{4}{(ρ\ :\ ℝ)\ (hρ\ :\ 0\ <\ ρ)\ (hmono\ :\ TwistedJointMonotonicity\ (model\ D\ γ)\ Geometry\ ρ)}
\leanline{4}{(hg\ :\ (model\ D\ γ).meanGradient\ νStar\ (targetState\ star)\ =\ 0)}
\leanline{4}{(hV\ :\ 0\ ≤\ Variance)\ :}
\leanline{4}{0\ <\ maximumStep\ D\ γ\ Geometry\ μ\ hγ\ ρ\ ∧}
\leanline{6}{0\ ≤\ discretizationConstant\ D\ γ\ Geometry\ μ\ hμ\ a\ Variance\ hVariance\ hγ\ ρ\ ∧}
\leanline{6}{∀\ h\ :\ ℝ,\ 0\ <\ h\ →\ h\ ≤\ maximumStep\ D\ γ\ Geometry\ μ\ hγ\ ρ\ →\ ∀\ n\ :\ ℕ,}
\leanline{8}{(‖momentumParticleIterate\ (ParticleNoiseSpace.probability\ (J\ :=\ J)\ μ\ N)}
\leanline{10}{(particleBrownian\ (J\ :=\ J)\ μ\ N)\ (particleBrownian\_isPreBrownian\ (J\ :=\ J)\ μ\ N)}
\leanline{10}{D\ γ\ h\ (initial\ D\ γ\ Geometry\ μ\ hμ\ a\ Variance\ hVariance)\ n\ -}
\leanline{10}{momentumUnpackLp\ (ParticleNoiseSpace.probability\ (J\ :=\ J)\ μ\ N)\ D}
\leanline{12}{((particles\ D\ γ\ Geometry\ μ\ hμ\ a\ Variance\ hVariance).state\ (n*h))‖\ ≤\ discretizationConstant\ D\ γ\ Geometry\ μ\ hμ\ a\ Variance\ hVariance\ hγ\ ρ\ *\ Real.sqrt\ h)\ ∧}
\leanline{8}{‖(momentumParticleIterate\ (ParticleNoiseSpace.probability\ (J\ :=\ J)\ μ\ N)}
\leanline{10}{(particleBrownian\ (J\ :=\ J)\ μ\ N)\ (particleBrownian\_isPreBrownian\ (J\ :=\ J)\ μ\ N)}
\leanline{10}{D\ γ\ h\ (initial\ D\ γ\ Geometry\ μ\ hμ\ a\ Variance\ hVariance)\ n).2.1\ -}
\leanline{10}{Lp.const\ 2\ (ParticleNoiseSpace.probability\ (J\ :=\ J)\ μ\ N)\ star‖\ ≤}
\leanline{8}{discretizationConstant\ D\ γ\ Geometry\ μ\ hμ\ a\ Variance\ hVariance\ hγ\ ρ\ *\ Real.sqrt\ h\ +}
\leanline{8}{(Real.sqrt\ Geometry.conditionNumber\ *\ Real.exp\ (-ρ\ *\ (n*h))\ *}
\leanline{10}{optimizerLawDistance\ a\ (targetState\ star)\ μ\ νStar\ +}
\leanline{10}{((2+γ)/ρ)\ *\ Real.sqrt\ (Geometry.conditionNumber\ *\ Variance/(N\ :\ ℝ))\ *}
\leanline{12}{(1-Real.exp\ (-ρ\ *\ (n*h))))}
\end{leancode}
\leanlabel{How the Lean proof follows the argument}
The hypotheses are the section variables (\texttt{D}, \texttt{\ensuremath{\gamma}}, \texttt{Geometry}, \texttt{\ensuremath{\mu}} with \texttt{h\ensuremath{\mu} : Mem\allowbreak{}Lp (fun z \ensuremath{\mapsto} z) 2 \ensuremath{\mu}}, \texttt{a}, \texttt{Variance}, \texttt{h\allowbreak{}Variance}) and \texttt{hcloud : D.\allowbreak{}cloud = euclidean\allowbreak{}Particle\allowbreak{}Encoding N J}, \texttt{h\ensuremath{\gamma}}, \texttt{star} (\(\theta_\star\)), \texttt{\ensuremath{\nu}\allowbreak{}Star}, \texttt{h\ensuremath{\nu}\allowbreak{}Star}, \texttt{Frozen} (N3, a \texttt{Additive\allowbreak{}Stationary.\allowbreak{}Frozen\allowbreak{}Trajectory} for \texttt{target\allowbreak{}State star} \(=(\theta_\star,0)\)), \texttt{h\ensuremath{\rho}}, \texttt{hmono} (A3), \texttt{hg} (zero mean observation) and \texttt{h\allowbreak{}V}. \texttt{Mem\allowbreak{}Lp f 2 \ensuremath{\mu}} means that \(f\) is square-integrable for \(\mu\). The numerical iterate is \texttt{momentum\allowbreak{}Particle\allowbreak{}Iterate Q W h\allowbreak{}W D \ensuremath{\gamma} h (initial …) n}, the continuous system is \texttt{momentum\allowbreak{}Unpack\allowbreak{}Lp Q D ((particles …).\allowbreak{}state (n*h))}, and \texttt{Lp.\allowbreak{}const 2 Q star} is the constant random variable \(\theta_\star\); \texttt{.\allowbreak{}2.\allowbreak{}1} selects the parameter. \texttt{optimizer\allowbreak{}Law\allowbreak{}Distance} is \(d_0\) and \texttt{Geometry.\allowbreak{}condition\allowbreak{}Number} is \(\kappa\).



\begin{enumerate}


\item Item 1: \texttt{Y} is the constructed process \texttt{(particles …).\allowbreak{}state} and \texttt{z} its raw coordinates; \texttt{heq} is its integral equation (with \texttt{Particle\allowbreak{}Brownian\allowbreak{}Input.\allowbreak{}packed\_\allowbreak{}input\_\allowbreak{}eq}), and \texttt{hze} is the raw form obtained with \texttt{momentum\allowbreak{}Unpack\allowbreak{}Lp\_\allowbreak{}integral} and \texttt{momentum\allowbreak{}Unpack\allowbreak{}Lp\_\allowbreak{}brownian} (the latter uses \texttt{hcloud}).

\item Item 2: \texttt{Method\allowbreak{}Metric\allowbreak{}Comparisons.\allowbreak{}momentum\_\allowbreak{}metric\_\allowbreak{}explicit} yields \texttt{C} with \texttt{hcoerc}, \texttt{hupper}, \texttt{hdrift}; \texttt{hm} is \(m_N>0\) (from \texttt{Empirical\allowbreak{}Metric.\allowbreak{}lower\_\allowbreak{}pos}) and \texttt{h\allowbreak{}Cnorm} is \(\|T\|\le U_N\) (\texttt{C.\allowbreak{}op\allowbreak{}Norm\_\allowbreak{}le\_\allowbreak{}bound}).

\item Step 1: \texttt{raw\allowbreak{}AB} is \((c_{0,M},c_{1,M})\), \texttt{actual\allowbreak{}AB} the coefficients with \(\|T\|\), \texttt{AB} the pair \((a_M,b_M)\), \texttt{base\allowbreak{}R} is \(r_{0,M}\), \texttt{R} is \(R_M\) and \texttt{U} is \(D_M\). The comparisons are \texttt{h\allowbreak{}Acomp}, \texttt{h\allowbreak{}Bcomp}, \texttt{h\allowbreak{}Scomp}, \texttt{h\allowbreak{}Rcomp} and \texttt{h\allowbreak{}Ucomp}; nonnegativity uses \texttt{momentum\allowbreak{}Uniform\allowbreak{}Defect\_\allowbreak{}nonneg} and \texttt{Uniform\allowbreak{}Integral.\allowbreak{}radius\_\allowbreak{}nonneg}.

\item Step 2: after \texttt{refine} with \texttt{Uniform\allowbreak{}Step.\allowbreak{}maximum\allowbreak{}Step\_\allowbreak{}pos}, \texttt{Uniform\allowbreak{}Step.\allowbreak{}admissible} gives \texttt{hh1}, \texttt{h\ensuremath{\rho}h}, \texttt{hsmall}, and \texttt{hsmall\allowbreak{}Actual} transfers the last one to \(\|T\|\). \texttt{hnum} is item 3, the theorem \texttt{momentum\allowbreak{}Particle\_\allowbreak{}time\_\allowbreak{}uniform} (file on the particle OBABO bound), and \texttt{herror} replaces its constant by \texttt{U}.

\item Step 3: \texttt{hcoord} identifies the parameter coordinate (\texttt{momentum\allowbreak{}Unpack\allowbreak{}Lp\_\allowbreak{}parameter}); \texttt{hsource} is item 4, \texttt{Additive\allowbreak{}Constructed.\allowbreak{}parameter\_\allowbreak{}bound}, with the 1-Lipschitz parameter map \texttt{hparam} and the fixed point \texttt{target\_\allowbreak{}fixed}; \texttt{hresponse} computes the response constant \(2+|\gamma|=2+\gamma\); \texttt{hp} is the component bound and \texttt{numerical\_\allowbreak{}target\_\allowbreak{}triangle} the triangle inequality.

\end{enumerate}
\end{leanbox}
\subsection{\texorpdfstring{Theorem New A: the full-state discretisation bound}{Theorem New A: the full-state discretisation bound}}


\textbf{Claim.} In the setting of the previous step, for every \(h\) with \(0<h\le h_{\star,M}\) and every \(n\in\mathbb N\),


\[\begin{aligned}\|\hat y_n-y(nh)\|=\max\bigl\{&\|\hat m_n-m(nh)\|_{L^2},\|\hat V_n-V(nh)\|_{L^2},\\&\|\hat\theta_n-\theta(nh)\|_{L^2},\|\hat X_n-X(nh)\|_{L^2}\bigr\}\\&\le D_M\sqrt h.\end{aligned}\]


\textbf{Proof.} This is the first conclusion of the previous step. \(\square\)




The bound compares the numerical method with the continuous particle system driven by the same Brownian motion, for all iterations at once. The constant \(D_M\) contains the initial norm \(\|y_0\|\) through \(R_M\), the force constants through \(c_{0,M},c_{1,M}\), the dimension through \(\sqrt d\), and the A3 data through \(m_N\), \(U_N\) and \(\rho\).

\begin{leanbox}{Lean: \texttt{explicit\_\allowbreak{}discretization\_\allowbreak{}bound}}
\leansource{\texttt{explicit\_\allowbreak{}discretization\_\allowbreak{}bound} \quad PBGSO/\allowbreak{}FullyUniformMomentum.lean, lines 237–253}
\begin{leancode}
\leanline{0}{theorem\ explicit\_discretization\_bound}
\leanline{4}{(hcloud\ :\ D.cloud\ =\ euclideanParticleEncoding\ N\ J)\ (hγ\ :\ 0\ ≤\ γ)}
\leanline{4}{(star\ :\ A)\ (νStar\ :\ Measure\ (HypocoercivePhase\ (EuclideanSpace\ ℝ\ J)))\ [IsProbabilityMeasure\ νStar]}
\leanline{4}{(hνStar\ :\ MemLp\ (fun\ z\ :\ HypocoercivePhase\ (EuclideanSpace\ ℝ\ J)\ ↦\ z)\ 2\ νStar)}
\leanline{4}{(Frozen\ :\ AdditiveStationary.FrozenTrajectory\ (ULA.CanonicalBrownian.vectorProbability\ J)}
\leanline{6}{(momentumForceDynamics\ D\ γ)\ νStar\ hνStar\ (targetState\ star)\ (input\ γ))}
\leanline{4}{(ρ\ :\ ℝ)\ (hρ\ :\ 0\ <\ ρ)\ (hmono\ :\ TwistedJointMonotonicity\ (model\ D\ γ)\ Geometry\ ρ)}
\leanline{4}{(hg\ :\ (model\ D\ γ).meanGradient\ νStar\ (targetState\ star)\ =\ 0)}
\leanline{4}{(hV\ :\ 0\ ≤\ Variance)}
\leanline{4}{(h\ :\ ℝ)\ (hh\ :\ 0\ <\ h)\ (hhmax\ :\ h\ ≤\ maximumStep\ D\ γ\ Geometry\ μ\ hγ\ ρ)\ (n\ :\ ℕ)\ :}
\leanline{0}{‖momentumParticleIterate\ (ParticleNoiseSpace.probability\ (J\ :=\ J)\ μ\ N)}
\leanline{10}{(particleBrownian\ (J\ :=\ J)\ μ\ N)\ (particleBrownian\_isPreBrownian\ (J\ :=\ J)\ μ\ N)}
\leanline{10}{D\ γ\ h\ (initial\ D\ γ\ Geometry\ μ\ hμ\ a\ Variance\ hVariance)\ n\ -}
\leanline{10}{momentumUnpackLp\ (ParticleNoiseSpace.probability\ (J\ :=\ J)\ μ\ N)\ D}
\leanline{12}{((particles\ D\ γ\ Geometry\ μ\ hμ\ a\ Variance\ hVariance).state\ (n*h))‖\ ≤\ discretizationConstant\ D\ γ\ Geometry\ μ\ hμ\ a\ Variance\ hVariance\ hγ\ ρ\ *\ Real.sqrt\ h}
\end{leancode}
\leanlabel{How the Lean proof follows the argument}
The proof is \texttt{((explicit\_\allowbreak{}guarantee …).\allowbreak{}2.\allowbreak{}2 h hh hhmax n).\allowbreak{}1}: the third component of the conjunction, applied to the step and the iteration, and its first part.
\end{leanbox}
\subsection{\texorpdfstring{Theorem New A: the parameter bound}{Theorem New A: the parameter bound}}


\textbf{Claim (Theorem New A).} In the setting of the explicit guarantee (N1–N3, \(\gamma\ge0\), A3 with \(\rho>0\), A5 with \(V\ge0\), zero mean observation at \(\theta_\star\), the standard Euclidean encoding, and identical initial states of the numerical method and the particle system), for every \(h\) with \(0<h\le h_{\star,M}\) and every \(n\in\mathbb N\),


\[\|\hat\theta_n-\theta_\star\|_{L^2}\le D_M\sqrt h+\sqrt\kappa\,e^{-\rho nh}d_0+\frac{2+\gamma}{\rho}\sqrt{\frac{\kappa V}{N}}\bigl(1-e^{-\rho nh}\bigr).\]


\textbf{Proof.} This is the second conclusion of the explicit guarantee. \(\square\)




The three terms are the discretisation error, uniform in \(n\); the exponential relaxation of the continuous dynamics from the initial distance \(d_0\); and the finite-particle error of order \(N^{-1/2}\). Letting \(n\to\infty\) gives \(\limsup_n\|\hat\theta_n-\theta_\star\|_{L^2}\le D_M\sqrt h+\frac{2+\gamma}\rho\sqrt{\kappa V/N}\). The discretisation constant may depend on \(N\), so no statement uniform in \(N\) is made.

\begin{leanbox}{Lean: \texttt{explicit\_\allowbreak{}parameter\_\allowbreak{}bound}}
\leansource{\texttt{explicit\_\allowbreak{}parameter\_\allowbreak{}bound} \quad PBGSO/\allowbreak{}FullyUniformMomentum.lean, lines 256–276}
\begin{leancode}
\leanline{0}{theorem\ explicit\_parameter\_bound}
\leanline{4}{(hcloud\ :\ D.cloud\ =\ euclideanParticleEncoding\ N\ J)\ (hγ\ :\ 0\ ≤\ γ)}
\leanline{4}{(star\ :\ A)\ (νStar\ :\ Measure\ (HypocoercivePhase\ (EuclideanSpace\ ℝ\ J)))\ [IsProbabilityMeasure\ νStar]}
\leanline{4}{(hνStar\ :\ MemLp\ (fun\ z\ :\ HypocoercivePhase\ (EuclideanSpace\ ℝ\ J)\ ↦\ z)\ 2\ νStar)}
\leanline{4}{(Frozen\ :\ AdditiveStationary.FrozenTrajectory\ (ULA.CanonicalBrownian.vectorProbability\ J)}
\leanline{6}{(momentumForceDynamics\ D\ γ)\ νStar\ hνStar\ (targetState\ star)\ (input\ γ))}
\leanline{4}{(ρ\ :\ ℝ)\ (hρ\ :\ 0\ <\ ρ)\ (hmono\ :\ TwistedJointMonotonicity\ (model\ D\ γ)\ Geometry\ ρ)}
\leanline{4}{(hg\ :\ (model\ D\ γ).meanGradient\ νStar\ (targetState\ star)\ =\ 0)}
\leanline{4}{(hV\ :\ 0\ ≤\ Variance)}
\leanline{4}{(h\ :\ ℝ)\ (hh\ :\ 0\ <\ h)\ (hhmax\ :\ h\ ≤\ maximumStep\ D\ γ\ Geometry\ μ\ hγ\ ρ)\ (n\ :\ ℕ)\ :}
\leanline{8}{‖(momentumParticleIterate\ (ParticleNoiseSpace.probability\ (J\ :=\ J)\ μ\ N)}
\leanline{10}{(particleBrownian\ (J\ :=\ J)\ μ\ N)\ (particleBrownian\_isPreBrownian\ (J\ :=\ J)\ μ\ N)}
\leanline{10}{D\ γ\ h\ (initial\ D\ γ\ Geometry\ μ\ hμ\ a\ Variance\ hVariance)\ n).2.1\ -}
\leanline{10}{Lp.const\ 2\ (ParticleNoiseSpace.probability\ (J\ :=\ J)\ μ\ N)\ star‖\ ≤}
\leanline{8}{discretizationConstant\ D\ γ\ Geometry\ μ\ hμ\ a\ Variance\ hVariance\ hγ\ ρ\ *\ Real.sqrt\ h\ +}
\leanline{8}{(Real.sqrt\ Geometry.conditionNumber\ *\ Real.exp\ (-ρ\ *\ (n*h))\ *}
\leanline{10}{optimizerLawDistance\ a\ (targetState\ star)\ μ\ νStar\ +}
\leanline{10}{((2+γ)/ρ)\ *\ Real.sqrt\ (Geometry.conditionNumber\ *\ Variance/(N\ :\ ℝ))\ *}
\leanline{12}{(1-Real.exp\ (-ρ\ *\ (n*h))))}
\end{leancode}
\leanlabel{How the Lean proof follows the argument}
The proof is \texttt{((explicit\_\allowbreak{}guarantee …).\allowbreak{}2.\allowbreak{}2 h hh hhmax n).\allowbreak{}2}. The statement displays every hypothesis of the theorem; the constants \texttt{discretization\allowbreak{}Constant} and \texttt{maximum\allowbreak{}Step} are the closed expressions defined in this file.
\end{leanbox}
\subsection{\texorpdfstring{The existential form}{The existential form}}


\textbf{Claim.} In the same setting there exist \(h_{\max}>0\) and \(D\ge0\) such that both bounds of the explicit guarantee hold, with \(D\) in place of \(D_M\), for all \(h\in(0,h_{\max}]\) and all \(n\).




\textbf{Proof.} Take \(h_{\max}=h_{\star,M}\) and \(D=D_M\). \(\square\)

\begin{leanbox}{Lean: \texttt{numerical\_\allowbreak{}parameter\_\allowbreak{}bound}}
\leansource{\texttt{numerical\_\allowbreak{}parameter\_\allowbreak{}bound} \quad PBGSO/\allowbreak{}FullyUniformMomentum.lean, lines 279–306}
\begin{leancode}
\leanline{0}{theorem\ numerical\_parameter\_bound}
\leanline{4}{(hcloud\ :\ D.cloud\ =\ euclideanParticleEncoding\ N\ J)\ (hγ\ :\ 0\ ≤\ γ)}
\leanline{4}{(star\ :\ A)\ (νStar\ :\ Measure\ (HypocoercivePhase\ (EuclideanSpace\ ℝ\ J)))\ [IsProbabilityMeasure\ νStar]}
\leanline{4}{(hνStar\ :\ MemLp\ (fun\ z\ :\ HypocoercivePhase\ (EuclideanSpace\ ℝ\ J)\ ↦\ z)\ 2\ νStar)}
\leanline{4}{(Frozen\ :\ AdditiveStationary.FrozenTrajectory\ (ULA.CanonicalBrownian.vectorProbability\ J)}
\leanline{6}{(momentumForceDynamics\ D\ γ)\ νStar\ hνStar\ (targetState\ star)\ (input\ γ))}
\leanline{4}{(ρ\ :\ ℝ)\ (hρ\ :\ 0\ <\ ρ)\ (hmono\ :\ TwistedJointMonotonicity\ (model\ D\ γ)\ Geometry\ ρ)}
\leanline{4}{(hg\ :\ (model\ D\ γ).meanGradient\ νStar\ (targetState\ star)\ =\ 0)}
\leanline{4}{(hV\ :\ 0\ ≤\ Variance)\ :}
\leanline{4}{∃\ hmax\ U\ :\ ℝ,\ 0\ <\ hmax\ ∧\ 0\ ≤\ U\ ∧}
\leanline{6}{∀\ h\ :\ ℝ,\ 0\ <\ h\ →\ h\ ≤\ hmax\ →\ ∀\ n\ :\ ℕ,}
\leanline{8}{(‖momentumParticleIterate\ (ParticleNoiseSpace.probability\ (J\ :=\ J)\ μ\ N)}
\leanline{10}{(particleBrownian\ (J\ :=\ J)\ μ\ N)\ (particleBrownian\_isPreBrownian\ (J\ :=\ J)\ μ\ N)}
\leanline{10}{D\ γ\ h\ (initial\ D\ γ\ Geometry\ μ\ hμ\ a\ Variance\ hVariance)\ n\ -}
\leanline{10}{momentumUnpackLp\ (ParticleNoiseSpace.probability\ (J\ :=\ J)\ μ\ N)\ D}
\leanline{12}{((particles\ D\ γ\ Geometry\ μ\ hμ\ a\ Variance\ hVariance).state\ (n*h))‖\ ≤\ U*Real.sqrt\ h)\ ∧}
\leanline{8}{‖(momentumParticleIterate\ (ParticleNoiseSpace.probability\ (J\ :=\ J)\ μ\ N)}
\leanline{10}{(particleBrownian\ (J\ :=\ J)\ μ\ N)\ (particleBrownian\_isPreBrownian\ (J\ :=\ J)\ μ\ N)}
\leanline{10}{D\ γ\ h\ (initial\ D\ γ\ Geometry\ μ\ hμ\ a\ Variance\ hVariance)\ n).2.1\ -}
\leanline{10}{Lp.const\ 2\ (ParticleNoiseSpace.probability\ (J\ :=\ J)\ μ\ N)\ star‖\ ≤}
\leanline{8}{U\ *\ Real.sqrt\ h\ +}
\leanline{8}{(Real.sqrt\ Geometry.conditionNumber\ *\ Real.exp\ (-ρ\ *\ (n*h))\ *}
\leanline{10}{optimizerLawDistance\ a\ (targetState\ star)\ μ\ νStar\ +}
\leanline{10}{((2+γ)/ρ)\ *\ Real.sqrt\ (Geometry.conditionNumber\ *\ Variance/(N\ :\ ℝ))\ *}
\leanline{12}{(1-Real.exp\ (-ρ\ *\ (n*h))))}
\end{leancode}
\leanlabel{How the Lean proof follows the argument}
The existential witnesses are \texttt{maximum\allowbreak{}Step …} and \texttt{discretization\allowbreak{}Constant …}, and the proof term is \texttt{explicit\_\allowbreak{}guarantee}.
\end{leanbox}
\section{\texorpdfstring{Derived local error of the printed Gaussian update}{Derived local error of the printed Gaussian update}}


The exact restarted process and numerical step use the same Brownian path. The numerical three-component Gaussian innovation is built from the prescribed kernels; its fourth component supplies the exact endpoint Brownian increment. The error norm here is the sum of the four coordinate norms, which bounds the ordinary product norm.

\subsection{\texorpdfstring{An affine initial-radius local estimate}{An affine initial-radius local estimate}}


\textbf{Claim.} Under the complete-space and force bounds of the restart result, assume \(\alpha>0\), \(\gamma\ge0\), restart time \(a\ge0\), and \(0<h\le1\). Let \(L\) be the full drift Lipschitz constant, \(b_0=\|f(0)\|\), \(f_0=\|\bar F(0,0)\|\), \(g_0=\|G(0,0)\|\), and \(\Xi=\sqrt\alpha\sqrt{2d}\), with \(d\) Brownian coordinates. Let \(H\) be the explicit coefficient-based deterministic growth constant and \(J\) the explicit scalar Gaussian-kernel error constant. Put

\[\begin{aligned}c_0&=H+\sqrt d\,J+\tfrac83L\{b_0+Le^{L+1}(b_0+\Xi)+\Xi\},\\c_1&=H+\tfrac83L^2e^{L+1}.\end{aligned}\]

The constants are fully determined as follows. Let \(A_0=f_0+2L_F\), \(B_0=g_0+2L_G+L_GA_0\), and let \(c_{ij}\) be the seven explicit remainder constants from the coefficient result. Then \(H=(1/2+c_{31})B_0+c_{12}+c_{22}+c_{32}+c_{13}+c_{23}+c_{33}+L_GA_0\), \(J=\sqrt{2\alpha}(3\gamma/2+\alpha+c_{33})\), and \(L=4(L_F+L_G+1+2\gamma+\alpha)\). For every square-integrable initial state \(z\), the printed Gaussian step \(A_{a,h}z\) and actual exact restart satisfy \(\|A_{a,h}z-r_z(h)\|_{\rm sum}\le(c_0+c_1\|z\|)h\sqrt h\).



\textbf{Proof.} The actual input satisfies \(\|B_a(t)\|\le\Xi\sqrt t\), by the Brownian second moment, independently of the possibly random initial state. On \([0,h]\), the integral equation and global Lipschitz bound give radius \(P=e^{L+1}(\|z\|+b_0+\Xi)\). Its increment is bounded by \(M\sqrt t\), with \(M=b_0+LP+\Xi\); use \(t\le\sqrt t\) on \([0,1]\). Therefore \(\|f(r_z(t))-f(z)\|\le LM\sqrt t\).



At the endpoint the exact input equals the fourth innovation from the very same Gaussian kernel block. Compare both updates with their common Euler map. The printed coefficient remainders and force evaluation bound contribute \(H(1+\|z\|)h\sqrt h\). The computed kernel mismatch contributes \(\sqrt d\,Jh\sqrt h\). Integrating the drift modulus gives the remaining bound \((8/3)LMh\sqrt h\). Substitution of \(P\) into this sum and collection of its constant and \(\|z\|\) coefficients gives exactly \(c_0+c_1\|z\|\). All constants are independent of \(a\) and of the initial state. \(\square\)

\begin{leanbox}{Lean: \texttt{holUniformNoiseRadius}, \texttt{holUniformLocalConstant0}, \texttt{holUniformLocalConstant1} and 1 more}
\leansource{\texttt{holUniformNoiseRadius} \quad PBGSO/\allowbreak{}HOLRestartLocal.lean, lines 16–17}
\begin{leancode}
\leanline{0}{def\ holUniformNoiseRadius\ (I\ :\ Type*)\ [Fintype\ I]\ (α\ :\ ℝ)\ :\ ℝ\ :=}
\leanline{2}{Real.sqrt\ α\ *\ Real.sqrt\ (2\ *\ (Fintype.card\ I\ :\ ℝ))}
\end{leancode}
\leansource{\texttt{holUniformLocalConstant0} \quad PBGSO/\allowbreak{}HOLRestartLocal.lean, lines 20–26}
\begin{leancode}
\leanline{0}{def\ holUniformLocalConstant0\ (I\ :\ Type*)\ [Fintype\ I]}
\leanline{4}{(α\ γ\ LF\ LG\ fzero\ gzero\ bzero\ :\ ℝ)\ :\ ℝ\ :=}
\leanline{2}{let\ L\ :\ ℝ\ :=\ holDriftGlobalLip\ LF\ LG\ α\ γ}
\leanline{2}{let\ Ξ\ :=\ holUniformNoiseRadius\ I\ α}
\leanline{2}{holLocalGrowthConstant\ α\ γ\ LF\ LG\ fzero\ gzero\ +}
\leanline{4}{Real.sqrt\ (Fintype.card\ I\ :\ ℝ)\ *\ holScalarNoiseConstant\ α\ γ\ +}
\leanline{4}{(8\ /\ 3\ :\ ℝ)\ *\ L\ *\ (bzero\ +\ L\ *\ Real.exp\ (L\ +\ 1)\ *\ (bzero\ +\ Ξ)\ +\ Ξ)}
\end{leancode}
\leansource{\texttt{holUniformLocalConstant1} \quad PBGSO/\allowbreak{}HOLRestartLocal.lean, lines 29–32}
\begin{leancode}
\leanline{0}{def\ holUniformLocalConstant1\ (α\ γ\ LF\ LG\ fzero\ gzero\ :\ ℝ)\ :\ ℝ\ :=}
\leanline{2}{let\ L\ :\ ℝ\ :=\ holDriftGlobalLip\ LF\ LG\ α\ γ}
\leanline{2}{holLocalGrowthConstant\ α\ γ\ LF\ LG\ fzero\ gzero\ +}
\leanline{4}{(8\ /\ 3\ :\ ℝ)\ *\ L\ \textasciicircum{}\ 2\ *\ Real.exp\ (L\ +\ 1)}
\end{leancode}
\leansource{\texttt{higherOrderRestart\_\allowbreak{}local\_\allowbreak{}error} \quad PBGSO/\allowbreak{}HOLRestartLocal.lean, lines 60–121}
\begin{leancode}
\leanline{0}{theorem\ higherOrderRestart\_local\_error}
\leanline{4}{(hα\ :\ 0\ <\ α)\ (hγ\ :\ 0\ ≤\ γ)\ (a\ h\ :\ ℝ≥0)}
\leanline{4}{(hh\ :\ 0\ <\ (h\ :\ ℝ))\ (hh1\ :\ (h\ :\ ℝ)\ ≤\ 1)}
\leanline{4}{(z\ :\ HOLState\ Θ\ (Lp\ (EuclideanSpace\ ℝ\ I)\ 2\ Q))\ :}
\leanline{4}{holError\ (higherOrderStep\ Fbar\ G\ α\ γ\ h\ z}
\leanline{8}{(holPullbackVectorNoise\ I\ hf\ α\ γ\ a\ h\ 0,}
\leanline{9}{holPullbackVectorNoise\ I\ hf\ α\ γ\ a\ h\ 1,}
\leanline{9}{holPullbackVectorNoise\ I\ hf\ α\ γ\ a\ h\ 2))}
\leanline{6}{(higherOrderRestart\ Q\ hf\ Fbar\ G\ LF\ LG\ α\ γ\ hLF\ hLG\ hF\ hG\ z\ a\ h)\ ≤}
\leanline{6}{(holUniformLocalConstant0\ I\ α\ γ\ LF\ LG\ ‖Fbar\ 0\ 0‖\ ‖G\ 0\ 0‖}
\leanline{10}{‖holDrift\ Fbar\ G\ α\ γ\ 0‖\ +}
\leanline{8}{holUniformLocalConstant1\ α\ γ\ LF\ LG\ ‖Fbar\ 0\ 0‖\ ‖G\ 0\ 0‖\ *\ ‖z‖)\ *}
\leanline{10}{(h\ :\ ℝ)\ *\ Real.sqrt\ (h\ :\ ℝ)}
\end{leancode}
\leanlabel{How the Lean proof follows the argument}
The three definitions show the exact local coefficients; \texttt{hol\allowbreak{}Local\allowbreak{}Growth\allowbreak{}Constant} and \texttt{hol\allowbreak{}Scalar\allowbreak{}Noise\allowbreak{}Constant} are the already derived printed-method constants used in them. \texttt{higher\allowbreak{}Order\allowbreak{}Restart\_\allowbreak{}local\_\allowbreak{}error} constructs \texttt{r}, \texttt{B}, \texttt{P} and \texttt{M}. \texttt{Uniform\allowbreak{}Integral.\allowbreak{}local\_\allowbreak{}regularity} proves the radius and increments, then \texttt{hmod} multiplies the increment estimate by the drift Lipschitz constant. \texttt{hend} replaces the endpoint input by the fourth actual innovation using \texttt{higher\allowbreak{}Order\allowbreak{}Restart\allowbreak{}Input\_\allowbreak{}endpoint}. It applies \texttt{higher\allowbreak{}Order\_\allowbreak{}pullback\_\allowbreak{}integral\_\allowbreak{}local\_\allowbreak{}error}, unfolds the two constants, and closes by exact algebra. Thus no local-error or reference-modulus premise is supplied.
\end{leanbox}
\section{\texorpdfstring{Time-uniform discretisation by exact-flow absorption}{Time-uniform discretisation by exact-flow absorption}}


Let \(X_k\) be the numerical iterates and \(Y(kh)\) the continuous integral trajectory, initially equal. The exact restart from \(X_k\) separates the one-step comparison into a numerical local defect and a contracting exact-flow difference. Uniform moments of the exact trajectory control the state dependence of the defect.

\subsection{\texorpdfstring{A uniform bound for every iteration}{A uniform bound for every iteration}}


\textbf{Claim.} Assume the global force bounds and actual common Brownian input above, \(\alpha>0\), \(\gamma\ge0\), and a bounded linear map \(C\) satisfying \(m\|z\|\le\|Cz\|\) for \(m>0\) and drift dissipativity with \(\rho>0\). Let \(Y\) be any continuous trajectory satisfying the original integral equation for all nonnegative times. Define

\[S=\|C\|e^{L+1}(b_0+\Xi),\quad R=\|CY(0)\|+\frac S{1-e^{-\rho}}+S,\quad A=\|C\|c_0,\quad B=\|C\|c_1/m.\]

If \(0<h\le1\), \(\rho h\le1\), and \(B\sqrt h\le\rho/4\), then all numerical iterates started at \(Y(0)\) obey

\[\|X_n-Y(nh)\|\le\frac{4(A+BR)}{m\rho}\sqrt h\qquad(n\ge0).\]

\textbf{Proof.} On each unit interval, compare the exact trajectory with the zero-start exact solution driven by the same increment. The latter has transformed radius at most \(S\), and exact contraction gives at integer times \(r_{k+1}\le e^{-\rho}r_k+S\). Summing this geometric recurrence, then allowing the remainder interval of length at most one, gives \(\|CY(t)\|\le R\) uniformly in \(t\).



Write \(e_k=\|C(X_k-Y(kh))\|\). The local bound at \(X_k\), coercivity, and \(\|CX_k\|\le e_k+R\) give

\[e_{k+1}\le(e^{-\rho h}+Bh\sqrt h)e_k+(A+BR)h\sqrt h.\]

For \(0\le\rho h\le1\), the scalar exponential bound gives \(e^{-\rho h}\le1-\rho h/2\). The small-step condition therefore makes the recurrence coefficient at most \(1-\rho h/4\), a nonnegative contraction. Its geometric sum, with \(e_0=0\), is bounded by \(4(A+BR)\sqrt h/\rho\). Divide by \(m\) to obtain the ordinary norm. The result is uniform over every iteration, rather than just a fixed finite horizon. \(\square\)

\begin{leanbox}{Lean: \texttt{holUniformReferenceRadius}, \texttt{higherOrder\_\allowbreak{}time\_\allowbreak{}uniform\_\allowbreak{}of\_\allowbreak{}integral}}
\leansource{\texttt{holUniformReferenceRadius} \quad PBGSO/\allowbreak{}HOLUniform.lean, lines 19–24}
\begin{leancode}
\leanline{0}{def\ holUniformReferenceRadius}
\leanline{4}{\{E\ H\ :\ Type*\}\ [NormedAddCommGroup\ E]\ [NormedSpace\ ℝ\ E]}
\leanline{4}{[NormedAddCommGroup\ H]\ [NormedSpace\ ℝ\ H]}
\leanline{4}{(f\ :\ E\ →\ E)\ (L\ :\ ℝ≥0)\ (C\ :\ E\ →L[ℝ]\ H)\ (ρ\ Ξ\ :\ ℝ)\ (z\ :\ E)\ :\ ℝ\ :=}
\leanline{2}{let\ S\ :=\ ‖C‖\ *\ PBGSO.UniformIntegral.radius\ L\ ‖f\ 0‖\ Ξ\ 0}
\leanline{2}{‖C\ z‖\ +\ S\ /\ (1\ -\ Real.exp\ (-ρ))\ +\ S}
\end{leancode}
\leansource{\texttt{higherOrder\_\allowbreak{}time\_\allowbreak{}uniform\_\allowbreak{}of\_\allowbreak{}integral} \quad PBGSO/\allowbreak{}HOLUniform.lean, lines 43–124}
\begin{leancode}
\leanline{0}{theorem\ higherOrder\_time\_uniform\_of\_integral}
\leanline{4}{(hα\ :\ 0\ <\ α)\ (hγ\ :\ 0\ ≤\ γ)}
\leanline{4}{(C\ :\ HOLState\ Θ\ (Lp\ (EuclideanSpace\ ℝ\ I)\ 2\ Q)\ →L[ℝ]\ H)}
\leanline{4}{(m\ ρ\ :\ ℝ)\ (hm\ :\ 0\ <\ m)\ (hρ\ :\ 0\ <\ ρ)}
\leanline{4}{(hcoerc\ :\ ∀\ x,\ m\ *\ ‖x‖\ ≤\ ‖C\ x‖)}
\leanline{4}{(hmono\ :\ ∀\ x\ y,}
\leanline{6}{inner\ ℝ\ (C\ (x\ -\ y))\ (C\ (holDrift\ Fbar\ G\ α\ γ\ x\ -\ holDrift\ Fbar\ G\ α\ γ\ y))\ ≤}
\leanline{8}{-ρ\ *\ ‖C\ (x\ -\ y)‖\ \textasciicircum{}\ 2)}
\leanline{4}{(Y\ :\ ℝ\ →\ HOLState\ Θ\ (Lp\ (EuclideanSpace\ ℝ\ I)\ 2\ Q))\ (hY\ :\ Continuous\ Y)}
\leanline{4}{(heq\ :\ ∀\ t,\ 0\ ≤\ t\ →\ Y\ t\ -\ Y\ 0\ =}
\leanline{6}{(∫\ s\ in\ 0..t,\ holDrift\ Fbar\ G\ α\ γ\ (Y\ s))\ +\ higherOrderRestartInput\ Q\ hf\ α\ 0\ t)}
\leanline{4}{(h\ :\ ℝ≥0)\ (hh\ :\ 0\ <\ (h\ :\ ℝ))\ (hh1\ :\ (h\ :\ ℝ)\ ≤\ 1)}
\leanline{4}{(hρh\ :\ ρ\ *\ (h\ :\ ℝ)\ ≤\ 1)}
\leanline{4}{(hsmall\ :\ (‖C‖\ *\ holUniformLocalConstant1\ α\ γ\ LF\ LG\ ‖Fbar\ 0\ 0‖\ ‖G\ 0\ 0‖\ /\ m)\ *}
\leanline{6}{Real.sqrt\ (h\ :\ ℝ)\ ≤\ ρ\ /\ 4)\ (n\ :\ ℕ)\ :}
\leanline{4}{let\ c₀\ :=\ holUniformLocalConstant0\ I\ α\ γ\ LF\ LG\ ‖Fbar\ 0\ 0‖\ ‖G\ 0\ 0‖\ ‖holDrift\ Fbar\ G\ α\ γ\ 0‖}
\leanline{4}{let\ c₁\ :=\ holUniformLocalConstant1\ α\ γ\ LF\ LG\ ‖Fbar\ 0\ 0‖\ ‖G\ 0\ 0‖}
\leanline{4}{let\ R\ :=\ holUniformReferenceRadius\ (holDrift\ Fbar\ G\ α\ γ)\ (holDriftGlobalLip\ LF\ LG\ α\ γ)}
\leanline{6}{C\ ρ\ (holUniformNoiseRadius\ I\ α)\ (Y\ 0)}
\leanline{4}{‖trajectory\ (higherOrderStep\ Fbar\ G\ α\ γ\ h)}
\leanline{8}{(higherOrderPullbackGaussianNoise\ I\ hf\ α\ γ\ h)\ (Y\ 0)\ n\ -\ Y\ ((n\ :\ ℝ)\ *\ (h\ :\ ℝ))‖\ ≤}
\leanline{6}{((4\ *\ (‖C‖\ *\ c₀\ +\ (‖C‖\ *\ c₁\ /\ m)\ *\ R)\ /\ ρ)\ *\ Real.sqrt\ (h\ :\ ℝ))\ /\ m}
\end{leancode}
\leanlabel{How the Lean proof follows the argument}
\texttt{higher\allowbreak{}Order\_\allowbreak{}time\_\allowbreak{}uniform\_\allowbreak{}of\_\allowbreak{}integral} defines numerical \texttt{X}, sampled reference \texttt{Z}, step \texttt{A} and exact restart \texttt{Φ}. \texttt{hreference} invokes \texttt{Uniform\allowbreak{}Integral.\allowbreak{}uniform\_\allowbreak{}moment} with actual shifted-input bounds. \texttt{h\allowbreak{}Z} follows from \texttt{higher\allowbreak{}Order\allowbreak{}Restart\_\allowbreak{}agrees}, and \texttt{hcontract} from \texttt{higher\allowbreak{}Order\allowbreak{}Restart\_\allowbreak{}contracts}. \texttt{hdefect} combines \texttt{higher\allowbreak{}Order\allowbreak{}Restart\_\allowbreak{}local\_\allowbreak{}error}, the operator norm and coercivity. Finally \texttt{Uniform\allowbreak{}Discrete.\allowbreak{}one\_\allowbreak{}step\_\allowbreak{}same\_\allowbreak{}initial\_\allowbreak{}uniform\_\allowbreak{}norm} applies the established absorption recurrence. The norm in the conclusion is the actual product-space norm of the iterates.
\end{leanbox}
\section{\texorpdfstring{Explicit higher-order parameter error for all iterations}{Explicit higher-order parameter error for all iterations}}


The primitive particle force assumptions construct both the exact interacting trajectory and the actual printed Gaussian scheme. The metric is derived from joint source dissipativity. Its coercivity, local error constants and Brownian radius yield a positive, data-defined step interval and a time-uniform numerical error term.

\subsection{\texorpdfstring{Combining discretisation, relaxation and particle errors}{Combining discretisation, relaxation and particle errors}}


\textbf{Claim.} Let the parameter space be finite-dimensional and let \(N>0\). Assume the particle force data: independent globally Lipschitz parameter and sampler forces, the specified normalized Euclidean cloud encoding, and the stated square-integrable initial sampler law \(\mu\) and deterministic optimiser initial value \(a\). Assume the mean-field observation variance is bounded by \(V\ge0\) for all nonnegative time. Take \(\alpha>0\), \(\gamma\ge0\), and source joint dissipativity at rate \(\rho>0\) with quadratic condition number \(\kappa\). Let \((a_\star,\nu_\star)\) be a frozen stationary comparison, with \(\nu_\star\in\mathcal P_2\), zero mean observation and fixed optimiser. Set \(m=\sqrt{\ell_N}/2\), \(U=2\sqrt{u_N}\), where \(\ell_N,u_N\) are the constructed empirical metric bounds. With the data-only local constants \(c_0,c_1\), noise radius and zero-start radius above, define

\[\begin{aligned}S&=Ue^{L+1}(b_0+\Xi),& R&=U\|Y(0)\|+S/(1-e^{-\rho})+S,\\B&=Uc_1/m,& K&=4(Uc_0+BR)/(m\rho),\\h_{\max}&=\min\{1,1/\rho,(\rho/(4(B+1)))^2\}.\end{aligned}\]

For \(0<h\le h_{\max}\), the actual numerical parameter satisfies, at every integer \(n\ge0\),

\[\|\theta_n-a_\star\|_{L^2}\le K\sqrt h+\sqrt\kappa e^{-\rho nh}d((a,\mu),(a_\star,\nu_\star))+\frac1\rho\sqrt{\frac{\kappa V}{N}}(1-e^{-\rho nh}).\]

\textbf{Proof.} The exact empirical metric factor obeys \(m\|z\|\le\|Cz\|\) and \(\|Cz\|\le U\|z\|\). Source joint dissipativity supplies its drift pairing inequality. The actual packed particle integral equation becomes the raw four-coordinate equation under the prescribed Euclidean cloud encoding, and the physical Brownian input becomes exactly the restart input. Thus the preceding uniform discretisation result gives \(\|X_n-Y(nh)\|\le K\sqrt h\). The definition of \(h_{\max}\) implies all its step restrictions: \(h\le1\), \(\rho h\le1\), and \(B\sqrt h\le\rho/4\). Its denominators are positive even when \(B=0\), hence \(h_{\max}>0\).



The constructed exact particle comparison supplies the last two terms, using the stationary reference, initial law distance and variance bound. Projection onto the parameter coordinate has norm at most one, so the numerical-to-exact parameter discrepancy is at most \(K\sqrt h\). The triangle inequality in \(L^2\) adds this to the exact parameter bound. \(\square\)

\begin{leanbox}{Lean: \texttt{explicit\_\allowbreak{}parameter\_\allowbreak{}bound}}
\leansource{\texttt{explicit\_\allowbreak{}parameter\_\allowbreak{}bound} \quad PBGSO/\allowbreak{}FullyUniformHigherOrder.lean, lines 261–312}
\begin{leancode}
\leanline{0}{theorem\ explicit\_parameter\_bound\ [FiniteDimensional\ ℝ\ A]}
\leanline{4}{(hcloud\ :\ D.cloud\ =\ euclideanParticleEncoding\ N\ J)}
\leanline{4}{(hα\ :\ 0\ <\ α)\ (hγ\ :\ 0\ ≤\ γ)}
\leanline{4}{(νStar\ :\ Measure\ (HigherOrderSampler\ (EuclideanSpace\ ℝ\ J)))\ [IsProbabilityMeasure\ νStar]}
\leanline{4}{(hνStar\ :\ MemLp\ (fun\ z\ :\ HigherOrderSampler\ (EuclideanSpace\ ℝ\ J)\ ↦\ z)\ 2\ νStar)}
\leanline{4}{(aStar\ :\ A)}
\leanline{4}{(Frozen\ :\ AdditiveStationary.FrozenTrajectory\ (ULA.CanonicalBrownian.vectorProbability\ J)}
\leanline{6}{(higherOrderForceDynamics\ D\ α\ γ)\ νStar\ hνStar\ aStar\ (input\ α))}
\leanline{4}{(ρ\ :\ ℝ)\ (hρ\ :\ 0\ <\ ρ)\ (hmono\ :\ TwistedJointMonotonicity\ (sourceModel\ D\ α\ γ)\ Geo\ ρ)}
\leanline{4}{(hg\ :\ (sourceModel\ D\ α\ γ).meanGradient\ νStar\ aStar\ =\ 0)}
\leanline{4}{(hV\ :\ 0\ ≤\ Variance)}
\leanline{4}{(h\ :\ ℝ≥0)\ (n\ :\ ℕ)\ (hh\ :\ 0\ <\ (h\ :\ ℝ))}
\leanline{4}{(hhstep\ :\ (h\ :\ ℝ)\ ≤\ maximumStep\ D\ α\ γ\ Geo\ μ\ ρ)\ :}
\leanline{8}{‖(higherOrderPullbackParticleIterate\ (Fin\ N\ ×\ J)}
\leanline{10}{(ParticleNoiseSpace.probability\ (J\ :=\ J)\ μ\ N)}
\leanline{10}{(ParticleNoiseSpace.noise\_preserving\ (J\ :=\ J)\ μ\ N)\ D\ α\ γ\ h}
\leanline{10}{(initial\ D\ α\ γ\ Geo\ μ\ hμ\ a\ Variance\ hVariance)\ n).1\ -}
\leanline{10}{Lp.const\ 2\ (ParticleNoiseSpace.probability\ (J\ :=\ J)\ μ\ N)\ aStar‖\ ≤}
\leanline{8}{discretizationConstant\ D\ α\ γ\ Geo\ μ\ hμ\ a\ Variance\ hVariance\ ρ\ *\ Real.sqrt\ (h\ :\ ℝ)\ +}
\leanline{8}{(Real.sqrt\ Geo.conditionNumber\ *\ Real.exp\ (-ρ\ *\ (n\ *\ (h\ :\ ℝ)))\ *}
\leanline{10}{optimizerLawDistance\ a\ aStar\ μ\ νStar\ +}
\leanline{10}{(1\ /\ ρ)\ *\ Real.sqrt\ (Geo.conditionNumber\ *\ Variance\ /\ (N\ :\ ℝ))\ *}
\leanline{12}{(1\ -\ Real.exp\ (-ρ\ *\ (n\ *\ (h\ :\ ℝ)))))}
\end{leancode}
\leanlabel{How the Lean proof follows the argument}
\texttt{explicit\_\allowbreak{}parameter\_\allowbreak{}bound} obtains \texttt{hdisc} from \texttt{explicit\_\allowbreak{}discretization\_\allowbreak{}bound}. Its underlying guarantee constructs the metric with \texttt{Method\allowbreak{}Metric\allowbreak{}Comparisons.\allowbreak{}higher\allowbreak{}Order\_\allowbreak{}metric\_\allowbreak{}explicit}, proves the raw integral and Brownian identities, bounds the actual constants by their data-defined versions, and invokes \texttt{Uniform\allowbreak{}Step.\allowbreak{}admissible}. The final parameter proof calls \texttt{Additive\allowbreak{}Constructed.\allowbreak{}parameter\_\allowbreak{}bound}, identifies the actual parameter projection through \texttt{higher\allowbreak{}Order\allowbreak{}Unpack\allowbreak{}Lp\_\allowbreak{}parameter}, and bounds the numerical parameter discrepancy by \texttt{norm\_\allowbreak{}fst\_\allowbreak{}le}. \texttt{numerical\_\allowbreak{}target\_\allowbreak{}triangle} then adds discretisation and exact-particle errors. The exact response Lipschitz constant for this model is one, explaining the particle coefficient.
\end{leanbox}

\end{document}
